
%%%%%%%%%%%%%%%%%%%%%%%%%%%%%%%% SCRIPT FOR E-RARA AND MDZ FILES     %%%%%%%%%%%%%%%%%%%%%%%%%%%%%%%%%%%%%%%%%%%%%%%%
%%%%%%%%%%%%%%%%%%%%%%%%%%% fini le 30.04.2025 par F. GOY            %%%%%%%%%%%%%%%%%%%%%%%%%%%%%%%%%%%%%%%%%%%%%%%%
% !TeX TS-program = lualatex
\documentclass{article}
\usepackage[T1]{fontenc}
\usepackage{microtype}
\usepackage[pdfusetitle,hidelinks]{hyperref}

\usepackage{polyglossia}
\setmainlanguage{english}
\setotherlanguages{latin,greek}
\usepackage[series={},nocritical,noend,noeledsec,nofamiliar,noledgroup]{reledmac}
\usepackage{reledpar}

\usepackage{fontspec}
\setmainfont{TeX Gyre Termes}

\usepackage{sectsty}
\usepackage{xcolor}

\usepackage{fancyhdr}
\pagestyle{fancy}
\fancyhf{}
\fancyhead[LE,RO]{\nouppercase{\leftmark}}  
\cfoot{\thepage}
\renewcommand{\headrulewidth}{0.4pt}

% Redefine \section to remove numbering
\usepackage{titlesec}
\titleformat{\section}[block]{\normalfont\scshape\color{gray}}{}{0pt}{} % no number in heading
\titleformat{\subsection}[hang]{\normalfont}{}{0pt}{} % also remove subsection number
\titleformat{\subsubsection}[hang]{\normalfont\footnotesize\color{black}}{}{0pt}{}

% Modify how section marks are stored to exclude numbers
\makeatletter
\renewcommand{\sectionmark}[1]{%
	\markboth{#1}{}} % Only store the section title, without number
\renewcommand{\subsectionmark}[1]{%
	\markright{#1}} % Only store the subsection title, without number
\renewcommand{\numberline}[1]{} % Hide the section number in TOC
\makeatother

\begin{document}

\date{}
        \title{In priorem epistolam Pauli ad Timotheum commemtarii : [Thomas Cajetan], [1531]}
\maketitle
\tableofcontents
\clearpage
\begin{pages} 
\beginnumbering
        
\phantomsection
\addcontentsline{toc}{subsection}{\textit{IOHAN NES BVGEN HAGIVS PO meranus, Lectori- }}
\subsection*{\textit{IOHAN NES BVGEN HAGIVS PO meranus, Lectori- }}\pstart VLTI LIBRI nunc, Optime Lector, excu- duntur inuitis autoribus, licet uiuentibus, per quod quam male consulatur posteritati, quis ignorat, nisi sit ualde stu pidus? Inter quos libros, in primis plane habeo meum Hiob, qui iam his nundinis Francofordianis, Anno domint M. D. XXVI. ue- nit in manus meas, haud gratus bospes. Vidi enim non esse eum, quem in ecclesia nostra uulgo interpretatus fue- ram ante annum, sed quem seclusa auditorum turba, ante quatuor annos, cum quibusdam Monachis con- tuleram. Iudico plane ex eo plus emolumenti proue nire uenditori, quam lectori: quando etiam non huic, sed illi consulitur per talia, quae nolentibus nobis inuul- gantur. Nemo huic pesti cupit subuentum, ne illae quidem ciuitates, quae Euangelio Christi, Domini no stri fauere dicuntur: Soli Typographi sine legibus egunt. Verum, quae nunc uides in Epistolam Pauli ad  \pend
\vspace{0.5cm}\noindent
\vspace{0.2cm}\rule{1cm}{0.2pt}\\ 
\hspace{0.2cm}\textit{mg}
\footnotesize Typograpbi sine legibus uiuunt. 
\normalsize| 
\section*{COMMENT. POMER.  }\pstart Romanos, non grauatim (ne nos illae Harpyiae ante uerterent) permisimus Secerio nostro, uiro ut erudi - to, ita et fidelt, quamprimum inuulganda: quae ta- men a me non sunt scripta, sed a Doctore Ambrosio Maiobano, iam apud Vratislauienses, sacri Euange- lii ministro, me publice in schola non dictante, sed di- cente, excepta: quo fit, ut admirandum sit potius, quod haec omnia sic excipi potuerint, quam accusari- dum, si quid alicubi desyderaueris. Nam ut nobis non uacat, ita neque necessarium ducimus, ex singulis nostris praelectionibus, singulos libros mundo obtru- dere. Ego tamen ista omnia peruidi, in quibus sic ex- ceptis, si quid est, quod probaueris, Doctori Ambro- sio Maiobano, secundum Deum refer acceptum.  Vale.   \pend
\phantomsection
\addcontentsline{toc}{subsection}{\textit{STATVS EPI- }}
\subsection*{\textit{STATVS EPI- }}
\phantomsection
\addcontentsline{toc}{subsection}{\textit{STOLAE.  }}
\subsection*{\textit{STOLAE.  }}\pstart \huge\textbf{I}\normalsize N primis id obseruandum in hac Epistola, Paulum hoc agere, ut omnes homines peccatores, plancque ́ filios irae, et damnationis aeternae esse ostendat. Vbi hoc egit, declarat unde tustificemur, nempe ex sola fide per Christum. Et hic est status huius Epistolae, huc omnia referuntur, quaecunque in tota Epistola tra- ctantur. Pudeat igitur eos, qui alibi quaerunt iusti- tiam ueram, faciunt autem id omnes homines. Iudae os nemo ignorat in suis Ceremoniis quaerere iustui- am, per quod ipsimet inter se inustos esse fatentur.  \pend
\vspace{0.5cm}\noindent
\vspace{0.2cm}\rule{1cm}{0.2pt}\\ 
\hspace{0.2cm}\textit{mg}
\footnotesize Summa totius Epistolae. 
\normalsize| 
\section*{IN EPIST. AD ROM.  }
\marginpar{[ p.2 ]}
\marginpar{[ p.1 ]}
\marginpar{[ p.11 ]}\pstart Gentes item, impio Idolorum cultu idem quaerunt, accusantes et ipsi suam iniustitiam. Sed dum illam per opera quaerunt, quid aliud faciunt miseri, quam quod stercus stercore abluere conantur? Nequaquam enim potest arbor mala, fructus bonos ferre. Dein- de, nec Ifraël sectando iustitiam legis, peruenit ad iu- stitiam, Roma. 9.   \pend
\phantomsection
\addcontentsline{toc}{subsection}{\textit{SVMMA TOTIVS EPI- }}
\subsection*{\textit{SVMMA TOTIVS EPI- }}\pstart \huge\textbf{c}\normalsize Ertum est in nullo libro sacrae scripturae clarius et apertius describi legem et gratiam, quam in hac Epistola, sed quo ordine id faciat Apostolus, dicemus, ut uideamus curam et diligentiam hanc uni cam fuisse Paulo, ut compendium tale totius scriptu- rae formaret. In quo simul omnia capita scripturae, quibus nobis ad cognitionem nostri peccati, et miseri cordiae diuinae opus est, ordine ac perspicue complecte- retur. Qui id ignorant, difficilimam et obscurissimam (quod multis, qui docti haberi uolunt, accidisse scio) hanc Epistolam esse somniant. Sed capita uideamus.  In primo capite omnes Gentes licet naturae legem habeant, esse peccatores, probat.   \pend
\phantomsection
\addcontentsline{toc}{subsection}{\textit{stolae per singula capita.  }}
\subsection*{\textit{stolae per singula capita.  }}
\section*{COMMENT. POMER.  }
\marginpar{[ p.III ]}
\marginpar{[ p.IIII ]}
\marginpar{[ p.V ]}
\marginpar{[ p.VI ]}
\marginpar{[ p.VII ]}\pstart In tertio, infert eam propositione, quae est status, et cardo Epistolae, dicens: Ergo omnes homines sunt pec- catores. Atque hactenus de morbo dixit, mox subiicit remedium, dicens: Vnam et solam esse iustitiam, quae est per fidem in Christum . Proinde, ne quis hanc iustitiam esse nouum quiddam, aut confictum a quoq, dicit, ipsam testimonium habere ex lege et Prophetis.   \pend\pstart In quarto, confirmat multis argumentis ex scriptu- ris petitis, hanc iustitiam per Christum.  In quinto, qui sit fructus, et quae utilitas huius iu stitiae, ostendit, nempe pax et gaudium conscientia- rum, antea peccato miserrime adflictarum: quae qui- dem nemo adsequitur, nisi per hanc iustitiam. Dein- de confert gratiae et peccatt uim, dicens, plus posse gratiam per Christum, quam peccatum per Adam.  In sexto, ponit occupationem contra eos, qui ubi audiunt ex peccato abundare gratiam, statim infe- runt fortiter, ergo peccabimus: et ut maxime id ore non dicant, tamen in corde ita sentiunt. Pii uero gau- dent se libertatem conscientiae adsecutos, nec aliud ur- gent, ut impii, qui tamen primi uolunt esse inter Chri- stianos. Paulus respondet. Non est porro peccandum, quando aboletur peccatum: conscientiarum pax est ibi, quod ad Deum, et non ad peccatum attinet, Spiritus enim semper aduersus carnem pugnat.   \pend\pstart In septimo docet, libere quiduis faciendum fore, si- uetus Adam totus mortuus esset: atque id dicit contra  \pend
\section*{IN EPIST. AD ROM.  }
\marginpar{[ p.8 ]}
\marginpar{[ p.VIII ]}\pstart eos, qui libertatem carnis somniant, et putant tum licere ut peccent impune, quando gratia praedicatur.  Verum, ex ista responsione offenduntur piae mentes, ut tandem dicant, O Paule, quam nobis libertatem praedicas, cum peccatum adhuc semper adsit: hos in sequenti capite solatur.   \pend\pstart In octauo consolatur pios, adhuc iniuriis carnis et mundi adflictos, dicens, ut maxime adhuc ad - sit peccatum, tamen non potest damnare. Saeuiunt etiam quantumuis propter spiritum Christi in uobis: estis filii Dei, a qua spe, nulla creatura potest nos se- parare. Atque haec hactenus. Quid potuit commodi- us et dilucidius de iustificatione dici, quam quod pri- mo morbum, secundo remedium indicarit? Tertio uero, nequaquam esse diffidendum, sed firmiter in Christo haerendum, docuerit.   \pend
\phantomsection
\addcontentsline{toc}{subsection}{\textit{SVMMA TRIVM SEQVEN- tium. IX. X. XI.  }}
\subsection*{\textit{SVMMA TRIVM SEQVEN- tium. IX. X. XI.  }}\pstart Vi praedicta recte intelligunt, obuiis ulnis, ut dicitur, amplectuntur ea quae de praedestina-  tione dicuntur in sequentibus cap. Prae scire igitur oportet ista quae in prioribus cap. do- cuit Paulus. Nam cum ipse gratiam commendare uo- luit, non potuit omittere locum praedestinationis: hic is maxime ad gratiam pertinet: magnum quiddam est, quod certus sim salutem meam non in mea manu,  \pend
\vspace{0.5cm}\noindent
\vspace{0.2cm}\rule{1cm}{0.2pt}\\ 
\hspace{0.2cm}\textit{mg}
\footnotesize IX. X. XI 
\normalsize| 
\section*{COMMENT. POMER.  }
\marginpar{[ p.XII XIII XIIII XV ]}\pstart sed Det esse- Hic nihil sunt uires tuae, nihil Liberum Arbitrium. Nam quas uires habuisti, quando non eras? Fecit nos ille charos ante mundum, in dilecto fi lio suo, Ephes.1. Vbi tum fuerunt bona opera? Atque  isto ordine docuit nos Paulus, in scripturis sanctis uer fart. Non uoluit, ut a praedestinatione inciperemus, ut quidam stulti et carnalissimi sapientes faciunt, quos hic locus tantum perturbat, et dementat. Rursus piis est con solationi, sciunt certo se non posse perire. Nam impii illi quamprimum cogitant praedestinationem, eô caecitatis proru unt, ut dicant: Ergo quiduis faciamus, peccemus, quia a eonclusum est, nihil retulerit sic uel sic nos uixisse. Hi ostendunt se non esse Christianos, et digni sunt audire hor rendam responsionem Pauli: Quorum damnatio iusta est-  \pend
\phantomsection
\addcontentsline{toc}{subsection}{\textit{SVMMA RELIQVORVM }}
\subsection*{\textit{SVMMA RELIQVORVM }}\pstart \huge\textbf{P}\normalsize Ostremo in sequentibus Cap-docet, quae sint bo na opera, atque hunc ordinem uides in omnibus scriptis Apostolicis obseruatum . Non est igitur inci- piendum a bonis operibus. Verum primum opus est, ut iustus sis per fidem, sicque facies bona opera. Do- cet igitur, ut nos primum offeramus Deo secundum Verbum suum: seruiamus donis spiritus, aliis: ne blas phemetur propter nos Dei Verbum. In summa, uult, ut teipsum agnoscas, et per Christum ueram iustitiam consequaris, utque  iustificatus, et coram Deo et coram hoibus pie uiuas. Id quod diligentissime docet haec Epi- stola: praeter hoc, quid aliud cupies?0  \pend
\phantomsection
\addcontentsline{toc}{subsection}{\textit{capitum.  }}
\subsection*{\textit{capitum.  }}
\vspace{0.5cm}\noindent
\vspace{0.2cm}\rule{1cm}{0.2pt}\\ 
\hspace{0.2cm}\textit{mg}
\footnotesize Contra Lib Arbitrium. 
\normalsize| 
\section*{4 IN EPIST. AD ROM.  }
\marginpar{[ p.]}
\endnumbering\beginnumbering\section{CAP .I.}\pstart Oc nomen Paulo post conuersionem famliare erat, siquidem Gentibus familiarius quiddam est Paulus, quam Saulus.   \pend
\phantomsection
\addcontentsline{toc}{subsection}{\textit{\huge\textbf{H}\normalsize Paulus. rum }}
\subsection*{\textit{\huge\textbf{H}\normalsize Paulus. rum }}\pstart De ministerio suo loquitur, ipsi enim commissum erat munus Apostolicum. Neque uero hic tantum ser- uum intelliges, eum qui colit Christum, qualis est om- nis Christianus, sed eum qui legatione fungitur, ad praedicandum Verbum, quique ad id uocatus est a Deo, ut tantum sit legatus Verbi. Per hoc damnans eos, qui sua sponte se ingerunt, et currunt non missi, sua quaerentes. Exponit uero hoc cum dicit  \pend
\phantomsection
\addcontentsline{toc}{subsection}{\textit{Seruus, }}
\subsection*{\textit{Seruus, }}
\phantomsection
\addcontentsline{toc}{subsection}{\textit{Segregatus. }}
\subsection*{\textit{Segregatus. }}\pstart Non profecto uulgaris haec est uocatio. q. d. Non passim omnibus est commissum, sed quibusdam saltem, ut praedicent Euangelium. Ego (ut in Act. habe- Vocatio vera tur) inquit, seorsum in hoc a Deo electus sum, haec est mea seruitus, haec mea functio, et meus Apostola- tus, ut praedicem Verbum. Atque hic est titulus Pau- li, quo necessario utitur, per hunc enim certus nos fa- cit de suis scriptis, et praedicatione Verbi. Nam nisi conscientiae certae sint de Verbo, quod ipsum sit Ver- bum Dei, erit inefficax, non figetradices . Itaque non  \pend
\vspace{0.5cm}\noindent
\vspace{0.2cm}\rule{1cm}{0.2pt}\\ 
\hspace{0.2cm}\textit{mg}
\footnotesize Seruus quo- modo accipi- endum.  
\normalsize| 
\hspace{0.2cm}\textit{mg}
\footnotesize 
\normalsize| 
\section*{COMMENT. POMER.  }\pstart iactabundo haec scribit spiritu de se Paulus, sed neces sario. Proinde non loquitur hic de praedestinatione (ut in Galatis) ut quidam somniarunt.  \pend
\phantomsection
\addcontentsline{toc}{subsection}{\textit{( Quod ante promilerat.) }}
\subsection*{\textit{( Quod ante promilerat.) }}\pstart Haec parenthesi includenda sunt. Hic uides quod sit Euangelium, in quod sit segregatus Apostolus, ne habeant quod calumnientur Papistae, qui dicunt, se quoque praedicare Euangelium, putantes simplicem historiam de Christo, per quatuor Euangelistas de- scriptam, tantum esse Euangelium: respicientesque ́ in solam literam, non in solidam gratiae Dei promissi- onem per Christum. Itaque non loquitur de omni Euam gelio, id est, laeto nuncio, sed tantum de Euangelio Christi. Quid enim peccatori laetius dici possit, quod audiat? quam, Christus filius Dei tibi natus est, pro- te passus, te ab omni tyrannide Sathanae redemit.  Atque illud uere est Euangelium, non historia simplex de Christo: fruitio igitur est Euangelium: est pro- missio Dei de filio suo etc. Interim quoque hic aduer- te, semper Deum antea promisisse, quae magnifice et potenter operaturus erat in mundo: ut per hoc conscientiae infirmiores non diffidant Verbo, sed ma- gis illi credant, qui solus est ueritas.  Euangelium ue- rum quod.   \pend
\phantomsection
\addcontentsline{toc}{subsection}{\textit{Per Prophetas. }}
\subsection*{\textit{Per Prophetas. }}\pstart Indicat per quem promissus sit nobis Christus.  Per Prophetas inquit. Sed quorsum attinet, quod ad  \pend\pstart  \pend
\vspace{0.5cm}\noindent
\vspace{0.2cm}\rule{1cm}{0.2pt}\\ 
\hspace{0.2cm}\textit{mg}
\footnotesize 
\normalsize| 
\section*{IN EPIST. AD ROM.  }
\marginpar{[ p.5 ]}\pstart iecit, In scripturis sanctis? Respondeo, utrunque neces sarium fuisse. Nam cogitare potuissent infirmae con- scientiae, Paulus dicit, per Prophetas Christum pro- missum: sed quis scit, an uerum sit, an id dixerint Prophetae? Potuit fortasse haec Paulus, item caeteri Apostoli fingere: quis nos certos faciet? quis nouit? Solet enim semper Verbum contemni et fastidiri, quo- ties ad ministerium respicias. Sunt enim homines in- firmi, caro et sanguis, qui praedicant Euangelium, ut nos sumus. Ita et Christum uidemus contemptum: hinc etiam dictum esse, Petram scandali. Attamen istud uile Verbum, quod uiles personae coram mun- do praedicant, est potentia Dei, operans mirifica in cordibus. Ideoque  adiectum est, in scripturis, ne di- cas Pauliun haec effinxisse ex suo capite.  udicium mun- di de Verbo-  \pend
\phantomsection
\addcontentsline{toc}{subsection}{\textit{De filio suo. }}
\subsection*{\textit{De filio suo. }}\pstart His omnibus commendat suum officium: et ne quis contemnat suum Euangelium, dicit: Hoc est quod praedico Euangelium de filio Dei, cui qui credit, sal- uus est . Praeterea hic filius est Dei natura, non adopti- one, ut nos, per quem solum nos adoptamur et con- stituimur cum eo in communem haereditatem.   \pend
\phantomsection
\addcontentsline{toc}{subsection}{\textit{Qui genitus. }}
\subsection*{\textit{Qui genitus. }}\pstart Hic uides humanitatem in Christo filio Dei: per bunc hominem solum oportuit satisfieri pro peccatis nostris secundum scripturam, ne aliam imaginaremur satis factionem.  humanitas so- lii Dei.  \pend
\vspace{0.5cm}\noindent
\vspace{0.2cm}\rule{1cm}{0.2pt}\\ 
\hspace{0.2cm}\textit{mg}
\footnotesize 
\normalsize| 
\hspace{0.2cm}\textit{mg}
\footnotesize 
\normalsize| 
\section*{COMMENT . POMER.  }\pstart Hoc est illud semen benedictum, primum Abrahae promissum: secundo clarius, ex Iuda quoque uenturum: Promissum postremo certissime ex Dauide.  \pend\pstart Nemo sensit Christum aliam, praeter humanam, habuisse naturam. Non putabatur in hoc mortali cor- pore esse filius Dei:nihil aliud, quam uilitas, infirmi- tas, et abiectio quaedam esse uidebatur. Erat quidem Deus, sed latebat. In nobis non erat Deus, sed in se- tantum. Ideo addit  \pend
\phantomsection
\addcontentsline{toc}{subsection}{\textit{Ex semine. }}
\subsection*{\textit{Ex semine. }}
\phantomsection
\addcontentsline{toc}{subsection}{\textit{Secundum carnem. }}
\subsection*{\textit{Secundum carnem. }}\pstart Id est, potenter et efficaciter in cordibus nostris percrebuit ipsum esse filium Dei. Nisi enim poten- tia Det corda nostra uicisset, et regenerasset, nunquam credere potuissemus, hunc uilem hominem esse filium Dei. Iam igitur omnes, in quorum cordibus spiritus Dei operatur, fateri coguntur, uelint nolint, Christum esse filium Dei: Deum iustum, Christum esse iustiti- am: Deum misericordem, Christum misericordiam.  In summa, Christum esse quicquid Deus est.   \pend
\phantomsection
\addcontentsline{toc}{subsection}{\textit{Qui declaratus. }}
\subsection*{\textit{Qui declaratus. }}
\phantomsection
\addcontentsline{toc}{subsection}{\textit{Secundum spiritum. }}
\subsection*{\textit{Secundum spiritum. }}\pstart Id est, per quem sanctificantur omnia. Est enim Hebraice dictum, Sptritus sanctificationis, id est, qui solus sanctificat, quale et illud est. Filius pecca-  \pend
\vspace{0.5cm}\noindent
\vspace{0.2cm}\rule{1cm}{0.2pt}\\ 
\hspace{0.2cm}\textit{mg}
\footnotesize Iudicium hu- manae ratio- nis de Chri- sto.  
\normalsize| 
\section*{IN EPIST. AD ROM.  }
\marginpar{[ p.6 ]}\pstart ti, id est, qui autor est omnis peccati, qualis hactenus fuit Papa . Missus est itaque ille spiritus sanctificans in corda credentium, post Christi ascensionem. Anima enim cum adhuc esset in carne mortali, non id crede- re potuit. Et ut dicit Esaias. Non reputauimus eum.  Necessario igitur moriebatur, ut ueniret aliquando il le spiritus, qui declararet cum potentia in cordibus ho- minum iplum esse filium Dei.   \pend
\phantomsection
\addcontentsline{toc}{subsection}{\textit{Resurrexit emortuis. }}
\subsection*{\textit{Resurrexit emortuis. }}\pstart Sic declaratus et ostensus est manifeste coram omni- bus, in confusionem impiorum, quod esset non homo- tantu, sed et Deus. Hominem nemo ignorabat, quia ex se mine Dauidis erat genttus, sed Deum nemo esse Jensit, donec uirtute Det resuscitatus est. Quis enim hominum unquam e morte potuit suis uiribus resurge- re? Haec ideo dicuntur, ut confidamus firmiter per eum nos saluandos, qui Deus et homo est, et media- tor inter Deum et hominem.   \pend
\phantomsection
\addcontentsline{toc}{subsection}{\textit{Per quem accepimus. }}
\subsection*{\textit{Per quem accepimus. }}\pstart Iam redit ad id quod de sua uocatione dicere ceperat.   \pend
\phantomsection
\addcontentsline{toc}{subsection}{\textit{Gratiam. }}
\subsection*{\textit{Gratiam. }}\pstart Non possunt homines munus Apostolicum accipe- re, nisi prius gratiam accipiant, nisi Deus illis fa- ueat, et mittat suum spiritum, per quem praedicent efficaciter Verbum. Deinde addit: Ad hoc accepimu- nus Apostolicum, uel legationem Euangelii.   \pend
\vspace{0.5cm}\noindent
\vspace{0.2cm}\rule{1cm}{0.2pt}\\ 
\hspace{0.2cm}\textit{mg}
\footnotesize Christus solus me diator.  
\normalsize| 
\section*{COMMENT. POMER.  }\pstart Id est, ut omnes Gentes subdantur fidei, et cre dant Euangelio. Nullus hominum maneat in suis iu- stitiis et uiribus: decedant a fiducia sui, sintque ; in no uo regno fidet: neminem uult excipi.  \pend
\phantomsection
\addcontentsline{toc}{subsection}{\textit{Vt obediatur. }}
\subsection*{\textit{Vt obediatur. }}\pstart Hoc addit, ut sciamus, de qua fide dicat, finxe- runt enim Gentes et Iudaei hypocritae, sibi fidem quan- dam, fidebant suis ceremoniis et iustitiis carnis. Ve- rum illa coram Deo non ualet. Fides ergo uera est in- unigenitum Dei filium, ut est in Iohan. Nemo itaque  sperare potest salutem, ex ullius creaturae uirtute, aut ullis suis meritis. Non est aliud nomen sub coelo, in quo saluamur, ut ait Petrus in Actis:  \pend
\phantomsection
\addcontentsline{toc}{subsection}{\textit{Super ipsius nomine. }}
\subsection*{\textit{Super ipsius nomine. }}
\phantomsection
\addcontentsline{toc}{subsection}{\textit{Quorum de numero. }}
\subsection*{\textit{Quorum de numero. }}
\phantomsection
\addcontentsline{toc}{subsection}{\textit{Id est, eorum, qui sanctificandi sunt per fidem.  }}
\subsection*{\textit{Id est, eorum, qui sanctificandi sunt per fidem.  }}\pstart Quasi dicat, Vocaticii. Vocari, in scriptura bifa- riam accipitur. Nam aliud est uocari ad munus Aposto- licum, ut supra dixit Paulus de seipso, Vocatus ad mu- nus Apostolicum. Aliud item uocari ad Christum per Euangelii praedicationem. Ad Christum uocati sunt omnes Christiani, quia omnes illius nomen confiteri debent. Verum non omnes uocati sunt ad praedican- dum Euangelium. Est igitur praedicatio, ministerium  \pend
\phantomsection
\addcontentsline{toc}{subsection}{\textit{Vocati. }}
\subsection*{\textit{Vocati. }}
\vspace{0.5cm}\noindent
\vspace{0.2cm}\rule{1cm}{0.2pt}\\ 
\hspace{0.2cm}\textit{mg}
\footnotesize Fides uera.  
\normalsize| 
\hspace{0.2cm}\textit{mg}
\footnotesize Vocari in scri- ptura bifariam accipitur.  
\normalsize| 
\section*{IN EPIST. AD ROM.  }
\marginpar{[ p.4 ]}\pstart quod homines exequuntur. Per hoc enim inuitantur caetert, qui uocantur ad Christum, ad coenam illam magnam, id est, Verbum Dei. Ex iis facile intelli- ges, quomodo illud Christi intelligendum sit, Multi sunt uocati etc. Deinde et illud Pauli, Quos uoca- uit, illos et iustificauit. Vocaticii itaque sumus, non operarii:in uocatione, non in opere, est nostra salus, cum dicat Christus, Ego uos elegi, uos non elegistis me.   \pend
\phantomsection
\addcontentsline{toc}{subsection}{\textit{SVPERSCRIPTIO. . }}
\subsection*{\textit{SVPERSCRIPTIO. . }}
\phantomsection
\addcontentsline{toc}{subsection}{\textit{Omnibus }}
\subsection*{\textit{Omnibus }}\pstart Scilicet scribo dilectis. Hoc dicit, ne quis sibi me- ritum fingat.   \pend
\phantomsection
\addcontentsline{toc}{subsection}{\textit{Vocatis sanctis. }}
\subsection*{\textit{Vocatis sanctis. }}\pstart Quasi dicat. Non operantibus, qui se non fecerunt sanctos suis operibus, uerum sancti facti sunt per Deum Sancti ueri qui  \pend
\phantomsection
\addcontentsline{toc}{subsection}{\textit{Gratia uobis et pax. }}
\subsection*{\textit{Gratia uobis et pax. }}\pstart Haec duo ubiq imprecatur. Nam non aliis indige- mus praeter haec. Nam gratia fauor est, quo nobis be- ne uult Deus. Pax autem est tranquillitas conscientiae nostrae erga Deum: haec in Deo, illa in nobis est: habe- mus autem ista a Deo patre, non iudice, qui factus est no- bis per Christum pater.  Gratie et pa- cis discrimen.   \pend
\phantomsection
\addcontentsline{toc}{subsection}{\textit{Et domino. }}
\subsection*{\textit{Et domino. }}\pstart Hic uides quantum errarint, qui gratiam Dei suis meritis uoluerunt impetrare. Nam si gratia esset ex nscritis, non esset gratia, ut patet Cap. 4. A Deo Nulla gratia ex meritis no stris.  \pend
\vspace{0.5cm}\noindent
\vspace{0.2cm}\rule{1cm}{0.2pt}\\ 
\hspace{0.2cm}\textit{mg}
\footnotesize p 
\normalsize| 
\hspace{0.2cm}\textit{mg}
\footnotesize 
\normalsize| 
\hspace{0.2cm}\textit{mg}
\footnotesize 
\normalsize| 
\hspace{0.2cm}\textit{mg}
\footnotesize 
\normalsize| 
\section*{COMMENT. POMER.  }\pstart patre nostro, et domino lesu Christo, uenit ista gra tia et pax, non ab ulla alia creatura, ut Philip.4- Et pax Dei etc. Item Col. 3. Proinde Deo patris no men adscribit, quia id nomen est maiestatis, quo om- nis creatura terretur. Sed qui audit nomen patris, non terretur nomine Det, sed magis prouocatur, ut  1  maiestatis est additum, nam dominum uocat. Idque  ideo factum, quia uidemus ipsum abiectum et humilia tum propter nos. Neque illud nomem, in Christo, nos terret, quandoquidem scimus ipsum esse Dei patris nostri filium. Imo potius prouocamur, ut ipsum dili gamus. Dominus est, quia regnum habet, in quo nos defendat contra Sathanam et mundum. Ergo re cte gratiam et pacem precatur.   \pend\pstart Exordium est, a gratulatione ductum, quod in om- nibus Epistolis uides, nam hactenus suum officium com- mendauit, ut certi essemus, Verbum esse Det, non homi- nis, et missi a Deo, non Pseudoapostoli. Captat autem beneuolentiam a sua persona, quando dicit, se gratulari Romanis. Item a Romanis, quando dicit, fidem eorum ubique  praedicari. Ideoque et orare se pro eis ait, quod est os- ficium uert Episcopi.   \pend
\phantomsection
\addcontentsline{toc}{subsection}{\textit{Primum. }}
\subsection*{\textit{Primum. }}
\phantomsection
\addcontentsline{toc}{subsection}{\textit{Per lesum. }}
\subsection*{\textit{Per lesum. }}
\vspace{0.5cm}\noindent
\vspace{0.2cm}\rule{1cm}{0.2pt}\\ 
\hspace{0.2cm}\textit{mg}
\footnotesize Deus cur pa ter dicatur. 
\normalsize| 
\section*{IN EPIST. AD ROM.  }
\marginpar{[ p.]}\pstart Quoties Paulus sic loquitur, Christum mediatorem facit, quare non opus est, ex uno atque altero loco scri- pturae urgere pontificatum Christi, cum is in mille locis clarissime exprimatur. In talibus locis ubique Paulus ostendit nos nihil posse, nostra opera nihil esse. Vn- tam Sophistae id solum aduertissent, cur semper per Christum dicatur. Christi.  \pend
\phantomsection
\addcontentsline{toc}{subsection}{\textit{In toto mundo. }}
\subsection*{\textit{In toto mundo. }}\pstart Vulgata est in scrip. hyperbole, Latine diceretur ubique . q. d. quocunque uento, istic mihi de uestra fide- dicitur. Itaq hic iterum uides Episcopi ueri officium, non solum curantis, ut credant, sed ut quoque pro- grediantur: ut quod praedicando non potest efficere, praestet oratione.  \pend
\phantomsection
\addcontentsline{toc}{subsection}{\textit{Testis. }}
\subsection*{\textit{Testis. }}\pstart Hic manifeste iurat Paulus, dum Deum testem citat.  Nobis uero non licet iurare, neque per Deum, neque per- ullam creaturam nostra caussa. Item nec contendere in tudicio, iuxta dictum Matth. 5. Et certum est, Chri- stianos non pro se et suis contendere. Paulus uero et Christus, item Christiantiurant propter nos: quemadmo- dum et Deus in Veteri Testamento, iurauit per semet- ipsum (nam maiorem se non habet ) ut confirmentur conscientiae nostrae. Tanta est Dei misericordia, ut non solum promittat, sed et iuret nostra caussa, cum alte- rum satis esset. Hic Paulus irat nostra caussa, ut cre- damus uerbo illius. Iurare in scri- ptura quid sit  \pend
\vspace{0.5cm}\noindent
\vspace{0.2cm}\rule{1cm}{0.2pt}\\ 
\hspace{0.2cm}\textit{mg}
\footnotesize Pontificatus 
\normalsize| 
\hspace{0.2cm}\textit{mg}
\footnotesize 
\normalsize| 
\section*{COMMENT. POMER.  }
\phantomsection
\addcontentsline{toc}{subsection}{\textit{Spiritu meo. }}
\subsection*{\textit{Spiritu meo. }}\pstart Non elementis huius mundi, aut externis operi - bus, sed operibus quae sequuntur internum spiritum.  Nam sine illo, mera sunt hypocrisis. Spiritu Deum is colit, qui non sua quaerit, sed gloriam De', et salutem proximi.   \pend
\phantomsection
\addcontentsline{toc}{subsection}{\textit{In Euangelio. }}
\subsection*{\textit{In Euangelio. }}\pstart Scilicet praedicando. Quia antea dixit se seruum Christi, qui est Dominus. Praedicatio fieri non potest sine externis: tamen is est cultus Deiin spiritu, cum scilicet non adultero Verbum, mea quaerendo, sed tan- tum ueritatem illius annuncio. Quasi dicat Paulus.  Sic Deus cor meum finxit, ut nil aliud possim praedica- re, quam Euangelium.   \pend
\phantomsection
\addcontentsline{toc}{subsection}{\textit{Mentionem faciam. }}
\subsection*{\textit{Mentionem faciam. }}\pstart Id est, Coram Deo soleo etiam pro Romanis orare, ubi pro aliis oro. Oratio quam commendat scriptura, est desyderium cordis, non multiloquium. Desyderium pauperum exaudit Deus. Quare nemo uere orare potest, nisi urgeatur necessitate ad orandum. Quan do sentis tua mala, tum non potes non orare, etiam dormiens et bibens. Et hoc est quando Christus iubet sine intermissione orare. Quod nescio, quomodo ex- posuerint quidam, nam clauem Dauidis non habue- runt, de horis statutis. Et bis pueri facti sunt, dum uoluerunt esse sapientes, hoc est, simul non intelli-  \pend
\vspace{0.5cm}\noindent
\vspace{0.2cm}\rule{1cm}{0.2pt}\\ 
\hspace{0.2cm}\textit{mg}
\footnotesize Praedicatore. ueri qui. 
\normalsize| 
\hspace{0.2cm}\textit{mg}
\footnotesize Orare quan- do debeamus 
\normalsize| 
\section*{IN EPIST. AD ROM.  }
\marginpar{[ p.9 ]}\pstart gere, et simul negare uerba Christi. Huius oromnis et de- syderii exemplum uidere licet in captiuo: qui cum sciat se- laqueo debere perire, tamen perpetuo non cessat de- syderare, ut liberetur, etiam quando dormit, sibi li- bertatem somniat aliquam. Si itaque sic adfectus es, sen- per oras, uidens prae oculis imminens periculum. Sic Paulus semper orabat pro Romanis, quia semper illo- rum fidem praedicari audiebat.   \pend\pstart Contra Libus Arb. haec est sententia Pauli, qui non audebat proponere ut ueniret ad Romanos. Non cu- rat hic suam bonam intentionem Paulus, sed in uo- luntatem Dei respicit. Nulla enim intentio nostra est bona, msi rectificata per Verbum Dei.   \pend
\phantomsection
\addcontentsline{toc}{subsection}{\textit{Ad uos. }}
\subsection*{\textit{Ad uos. }}
\phantomsection
\addcontentsline{toc}{subsection}{\textit{Vt impertiar. }}
\subsection*{\textit{Vt impertiar. }}\pstart Id est, ut praedicem uobis Euangelium, quo con- firmemini, scilicet, in fide uestra, quam de uobis au- dio . Vides modestissime loqui Paulum, et tamen ni- mium a se dictum putat, quare addit correctionem Rhetoricam.  \pend
\phantomsection
\addcontentsline{toc}{subsection}{\textit{Communem capiam consolationem. }}
\subsection*{\textit{Communem capiam consolationem. }}\pstart Quasi dicat: In hoc ueniam ad uos, ut mea fides uos consoletur, et me uestra: uel ut simul gaudea- mus, audientes fidem nostram mutuam, uti solent Christiani. Verum aliud interim agebat Paulus in cor-  \pend
\vspace{0.5cm}\noindent
\vspace{0.2cm}\rule{1cm}{0.2pt}\\ 
\hspace{0.2cm}\textit{mg}
\footnotesize Contra Libus  Arbitrium. 
\normalsize| 
\section*{O MMENT. POMER.  }
\marginpar{[ p.tur. ]}\pstart de, quia timebat eos feductos per fiduciam operum, ut eos sic in ueram uiam reuocaret.   \pend
\phantomsection
\addcontentsline{toc}{subsection}{\textit{Proposueram. }}
\subsection*{\textit{Proposueram. }}\pstart Propositum nostrum nihil est, nisi uelit Deus, in cuius manu sunt omnia. Qua uero occasione tandem peruenerit Rhomam Paulus, Acta Apostolorum de clarant. Sed ita agit Deus, non Paulus.   \pend
\phantomsection
\addcontentsline{toc}{subsection}{\textit{Fructum. }}
\subsection*{\textit{Fructum. }}\pstart Scilicet Euangelii, idipsum praedicando inter uos.  \pend
\phantomsection
\addcontentsline{toc}{subsection}{\textit{Debitor sum. }}
\subsection*{\textit{Debitor sum. }}\pstart Quia omnium Gentium sum Apostolus, ideo et in uos ius praedicandi habeo. Graeci dicebantur sapien- tes et Philosophi, inter quos erant et Romani . Bar- bari, ettam Iudaet erant, et aliae nationes. Haec decla- rat, cum dicit: Eruditis, per quos Graecos: rudi- bus, per quos Barbaros, intelligit. Non minus rustici ad Euangelium pertinent, quam sapientes. Non minus abiecti, quam qui magni habentur. Siquidem donum fidei non literarum est, sed Dei.   \pend
\phantomsection
\addcontentsline{toc}{subsection}{\textit{Quantum in me est. }}
\subsection*{\textit{Quantum in me est. }}\pstart Id est, non quod uires meae possunt, sed quantum- mea permiserit uocatio, ut supra dixit: Volente Deo  \pend
\phantomsection
\addcontentsline{toc}{subsection}{\textit{Non me pudet. }}
\subsection*{\textit{Non me pudet. }}\pstart Quando hactenus praedicaui aliis Gentibus, et uo- bis praedicabo, scio enim, quam insignis res sit et praeciosa Euangelium, tametsi carni periculosa uidea-  \pend
\vspace{0.5cm}\noindent
\vspace{0.2cm}\rule{1cm}{0.2pt}\\ 
\hspace{0.2cm}\textit{mg}
\footnotesize Ad omnes na- tiones fuit mis- sus Paulus. 
\normalsize| 
\section*{IN EPIST. AD ROM.  10 }
\marginpar{[ p.]}\pstart tur. Est autem hic totus locus attentionis, quo uideas, quam rem grandem sit dicturus. Sed utinam hanc sen- tentiam hodie ratam haberent nostri, qui persecutio- nem patiuntur propter Euangelium. Siquidem tanta- res est, quae explicari non possit. Ideo exponit, dicens.   \pend
\phantomsection
\addcontentsline{toc}{subsection}{\textit{Potentia Dei est. }}
\subsection*{\textit{Potentia Dei est. }}\pstart Scilicet Euangelium. Tanta potentia Dei exerce- tur in cordibus nostris quando credimus, quanta Chri- stum suum suscitauit a mortuis. Quis igitur non ui- det a ratione nostra fidem, et Euangelium non posse comprehendi, cum tamen interim somniemus quan- dam carnis libertatem: quando audimus praedicari Euangelium, putamus id esse in nostra potestate, sed non est-Est ergo Euangelium instrumentum quoddam, quo Deus salutem nostram in cordibus nostris opera- tur: Siquidem Deus per nihil aliud salutem operatur, quam per Euangelium, hoc modo uetustatem no- stram interficiendo.  Euangelium in- strumentum est  \pend
\phantomsection
\addcontentsline{toc}{subsection}{\textit{Omni credenti . }}
\subsection*{\textit{Omni credenti . }}\pstart Id est, corde suscipienti Euangelium, nam corde creditur ad iustitiam, non ore, aut alio membro, uel operibus. Non est igitur aliunde expectanda salus, quam a Deo per Euangelium, quod mihi fit salus et potentia Dei, cum credo. Ergo filii maledictionis sunt et damnationis, qui non credant.  Nulla salus, nisi per Euan- gelium.   \pend
\phantomsection
\addcontentsline{toc}{subsection}{\textit{}}
\subsection*{\textit{}}
\vspace{0.5cm}\noindent
\vspace{0.2cm}\rule{1cm}{0.2pt}\\ 
\hspace{0.2cm}\textit{mg}
\footnotesize 
\normalsize| 
\hspace{0.2cm}\textit{mg}
\footnotesize 
\normalsize| 
\section*{COMMENT. P POMER.  }\pstart Exponit quid sit, credenti, ne quis se excuset, ue- ludaei, qui gloriabantur se semen Abrahae esse, quasi Christo nihil opus habeant. Et Gentes, qui suas car- nis iustitias et uirtutes iactabant. Ex hoc loco at- tentionis narrationem adgreditur, ex qua propositio- nem colligimus, Sola nos fide iustificari.   \pend\pstart Haec diligentissime sunt spectanda, et quo ordine omma connectat. qd. Euangelium aperit nobis ue- Euangelium ue- ram iustit iam ram aperit iu   mundo, nisi quibus reuelata erat, ut Abrahamo, et caeteris. De hac Dauid saepe: In tua iustitia libera me.  Hac nihil est dulcius piis: sed ea carent impii et iusti- tiarii. Verum cum talta nos olim in Psalmis legeremus, abhorrebat ab ea lectione animus noster. Ideoque ; cer- tum est, uocabulum iustitiae, nunquam eo tempore nos intellexisse, neque etiam eam iustitiam nos habuisse-  \pend
\phantomsection
\addcontentsline{toc}{subsection}{\textit{lustitia enim Dei. }}
\subsection*{\textit{lustitia enim Dei. }}\pstart Dicit, non fit. q. d. Erat quidem illa iustitia ab aeterno apud Deum, sed non cognoscebatur, nisi re uelata. Deinde ita patefit in cordibus nostris, ut in- telligas eam indies in nobis augeri, et non statim per- fecte illius nobis contingere cognitionem. In quo eti- am ueram fidei rationem et naturam est cernere. Est  \pend
\phantomsection
\addcontentsline{toc}{subsection}{\textit{Patefit }}
\subsection*{\textit{Patefit }}
\vspace{0.5cm}\noindent
\vspace{0.2cm}\rule{1cm}{0.2pt}\\ 
\hspace{0.2cm}\textit{mg}
\footnotesize stitiam.  
\normalsize| 
\hspace{0.2cm}\textit{mg}
\footnotesize Vera iustitia indies auge- tur. 
\normalsize| 
\section*{VI IN EPIST. AD ROM.  }
\marginpar{[ p.]}\pstart eutem ea, ut subinde per Euangelium illa in nobis crescat, quod Paulus in Corinth. dixit: A gloria in glortam, tanquam a domini spiritu. Item in Ephes. a Donec crescamus in uirum perfectum. Item Psalm.  De uirtute in uirtutem. Haec omnia sunt contra satu- ros et crassos Christianos, qui se putant omnia adse- cutos, cum tamen piorum sit semper sitire, ut Matth. 6.  Ideo subiungit testimonium Prophete.   \pend
\phantomsection
\addcontentsline{toc}{subsection}{\textit{Iustus exfide uiuet. }}
\subsection*{\textit{Iustus exfide uiuet. }}
\phantomsection
\addcontentsline{toc}{subsection}{\textit{Palam. }}
\subsection*{\textit{Palam. }}\pstart Non ex operibus, aut creaturis. Non est ergo fi- des quiddam ociosum in homine, sed ipsa uita, nunt quiescens. Viuet igitur, ut dicit Propheta, qui Chri- stum fide expectat. Perdetur is, qui in suis iustitiis confidit.  Iustus ex fi- de uiuct.   \pend\pstart Quae hactenus sunt dicta ad exordium pertine- bant: Nunc, quod et erat faciendum, primum le- ge omnes homines damnat, dicens, et Iudaeos et Gen- tes esse peccatores, atque hoc agit in tribus primis ca- pitibus. Nam necessario Paulus docturus, quae esset uera iustitia, legem, id est, iram praedicat, cunt dixe- rat: Euangelium est potentia Dei, et iustitia Dei per illud patefit. Quid aliud consequi oportebat, quam omnibus hominibus Deum irasci, quia omnes uacui sunt iustitia Det? Itaque Euangelium cum praedicatur, primum irae plenum est et indignationis, non potest Lex praedicat- iram.   \pend
\vspace{0.5cm}\noindent
\vspace{0.2cm}\rule{1cm}{0.2pt}\\ 
\hspace{0.2cm}\textit{mg}
\footnotesize 
\normalsize| 
\hspace{0.2cm}\textit{mg}
\footnotesize 
\normalsize| 
\section*{COMMENT. POMER.  }\pstart illud sustinere caro, quae semper suas iustitias extollit et magnifacit: tum fieri nequit, ut non maxime odiat illa Euangelium et uoluntatem Dei, qua decla- rat, nihil esse, imó maledictum, et abhominationem, quicquid conatur. Alians enim non diceret Scriptura: In te benedicentur omnes cognationes terrae. Hoc ubi Apostolus absoluit, nempe morbi explicationem, sta- tim remedium subiicit, scilicet per Christum solum esse- iustitiam. Sed diligenter uidere uerba Pauli oportet Le- ctorem, nam nihil is inutile adposuit. Nam si iustitia Dei patefit per Euangelium, ergo etiam ira. Nam sic- alterum alteri infertur. Iustitia enim Dei omnia damnat et iudicat in nobis. Sic itaque Paulus uim legis decla - rat, ut uideamus nostram paupertatem, atq in ea euan- gelizemur. Sic Christus Matth. 5. legem interpreta- tur, non fert legem, ut quidam inepte ibi ex Christo le- gislatorem faciunt. Hunc ergo modum praedicandi obserua- re debent Praedicatores, ut primum nobis declarent quid a nobis exigat lex, non ferant legem, quando iam olim in cordibus nostris scripta est. Lex, bonum cor a nobis requirit, atq id natura non habemus: ergo statim ubi praedicatur Euangelium, cogimur uidere, quales simus, nempe filii irae et perditionis.   \pend\pstart Per Euangelium iustitia Det arguit nostram iniusti- tiam. Haec autem peccati reuelatio contingit nobis tan  \pend
\phantomsection
\addcontentsline{toc}{subsection}{\textit{De coelo. }}
\subsection*{\textit{De coelo. }}
\vspace{0.5cm}\noindent
\vspace{0.2cm}\rule{1cm}{0.2pt}\\ 
\hspace{0.2cm}\textit{mg}
\footnotesize Modus praedi- candi.  
\normalsize| 
\section*{IN EPIST. AD ROM.  12 }
\marginpar{[ p.]}\pstart tum e coelo. Vnde et Christus dicit: Spiritussanctus arguet mundum de peccato. Hoc tudicium non ex hu- manis uiribus uenit: latro qui furtum fecit, non nouit radicem peccati, tametsi sciat furtum esse peceatum.  Solus Sptritussanctus nostram reuelat hypocrisim, non ratio humana.   \pend\pstart Veritas ea est, quod agnoscant quaedam de Deo.  Item sciunt quid faciendum sit, quid omittendum in mo- ribus, illam inquam ueritatem praetereunt, et non ue- Veritas quid.  nerantur Deum timore et fiducia in ipsum, sed suam iniustitiam, id est, quicquid est infidelitatis ipsorum, hoc loco ueritatis honorant, faciuntque ex suo peccato et ignominia Deum, et ex Deo peccatum. Hoc uitii om- nibus innatum est hominibus, nec liberantur ab eo, ni- si per Spiritumsanctum.  \pend
\phantomsection
\addcontentsline{toc}{subsection}{\textit{Qui ueritatem. }}
\subsection*{\textit{Qui ueritatem. }}
\phantomsection
\addcontentsline{toc}{subsection}{\textit{Quod de Deo. }}
\subsection*{\textit{Quod de Deo. }}\pstart Hic non accusantur a Scriptura quod non cognorint mysteria Trinitatis, item rationes chororum Angelicorum, sed quod illud quod de Deo potest cognosci, neglexe- rint. Videbant opera Dei, uices temporum, coelum dare pluuiam, Solem et Lunam splendorem suum. Item reliqua innumera beneficia, et quod ipse nos fecerit, non nos, ut inquit Psalm. Haec omniapraeteribant, ex qui- bus Deum potuissent cognoscere, ac deinde uenera- ri. Et exponit id Paulus in sequentibus.   \pend
\vspace{0.5cm}\noindent
\vspace{0.2cm}\rule{1cm}{0.2pt}\\ 
\hspace{0.2cm}\textit{mg}
\footnotesize 
\normalsize| 
\section*{COMMENT. POMER.  }\pstart Scilicet, per opera. Quae ubiq iubet Scriptura uos intelligere et contemplari. Vide totum Psalmum; Bonum est confidere etc.  \pend\pstart Id est, res illius Dei inuisibiles, nempe diuinitatem et illius potentiam, quas in Deo nunquam uidere pos- sumus, ut Iohan. 1. Deum nemo uidit unquam. In ope- ribus autem et creaturis uidetur.   \pend
\phantomsection
\addcontentsline{toc}{subsection}{\textit{Siquidem inuisibilia. }}
\subsection*{\textit{Siquidem inuisibilia. }}
\phantomsection
\addcontentsline{toc}{subsection}{\textit{Patefecit. }}
\subsection*{\textit{Patefecit. }}\pstart Id est, a condito mundo. Semper Deus per opera potuit intelligi, non plus fuit cognitus post legem la tam, quam ante legem. Sed uidere per opera potue- runt, glorificar Egent gloria Dei. Sed hoc ipsum etiam quod uidere potuerunt, contemptum est ab eis, per stultam sapien- tiam ipsorum: et obscuratum est insipiens cor eorum, ita ut dicat insipiens in corde suo, non est Deus, dum nemo intelligit aut requirit Deum, aut facit bonum-  \pend
\phantomsection
\addcontentsline{toc}{subsection}{\textit{A creatione mundi. }}
\subsection*{\textit{A creatione mundi. }}
\phantomsection
\addcontentsline{toc}{subsection}{\textit{Vt sint inexcusabiles }}
\subsection*{\textit{Vt sint inexcusabiles }}\pstart Scilicet illi, quibus Deus ista reuelauit, et ob ocu- los posuit: quia cum Deum cognorint, scilicet, ex- operibus, haec de lege naturae intelliguntur. Nam certum est ex Paulo, et aliis Apostolis, ut uides in Acti-Gentes Deum cognouisse. Illud autem non pro- babant ex sacris Scripturis Apostoli, quia eas risissent Gentes, nam illis non credebant.   \pend
\vspace{0.5cm}\noindent
\vspace{0.2cm}\rule{1cm}{0.2pt}\\ 
\hspace{0.2cm}\textit{mg}
\footnotesize Deus quomo- do inuisibilis.  
\normalsize| 
\hspace{0.2cm}\textit{mg}
\footnotesize Videtur Deus peropera.  
\normalsize| 
\section*{IN EPIST. AD ROM.  }
\marginpar{[ p.2 ]}\pstart Id est, non fisi sunt Deo.Iste est unicus cultus, et honor Dei, ut sciam unum tantum esse Deum (etiam si legem, aut Euangelium, non legerim, aut non au- dierim) quem glorificem, non illis aut aliis operibus, sed fide: quomodo? uit credam me omnia ex illius ma- nibus accepisse, et accepturum. Idque unicumre- quirit a nobis Deus, ut hoc solo simus grati pro bene- ficiis suis nobis collatis, quemadmodum uides Act. 14.  Fecit coelum et c. ubi de creatione loquitur. Item Act. 17.  insignis est locus de mundi creatione, quae erat prae- dicanda Gentibus, cum Scripturis non crederent.  Et profecto nisi Gentes Deum cognouissent, non di- scissent, Iouis omnia plena, ut in Arato et Vergi- lio uides.  \pend
\phantomsection
\addcontentsline{toc}{subsection}{\textit{Non ut Deum. }}
\subsection*{\textit{Non ut Deum. }}
\phantomsection
\addcontentsline{toc}{subsection}{\textit{Frustrati sunt. }}
\subsection*{\textit{Frustrati sunt. }}\pstart Ecce quam bonae sint nostrae cogitationes, ubi a Deo derelinquimur: quisque fingit sibi Deum, somniat quod sibi uidetur de illo. Sed illud cogitare, nihil aliud est, quam in densissimis tenebris, sursum ac deorsum currere, nihilque ; uidere quid agas: sic fit, ut negliga- mus ea, etiam quae natura nobis praemonstrantur. Haec excaecatio est prima poena peccati, et contemptus Dei . Altera est horrenda dictu, cum Deus nos no- bis relinquit in foedissimis operibus erga proximum.   \pend
\phantomsection
\addcontentsline{toc}{subsection}{\textit{Mutarunt gloriam. }}
\subsection*{\textit{Mutarunt gloriam. }}
\vspace{0.5cm}\noindent
\vspace{0.2cm}\rule{1cm}{0.2pt}\\ 
\hspace{0.2cm}\textit{mg}
\footnotesize Verus et unl- us cultus Dei 
\normalsize| 
\hspace{0.2cm}\textit{mg}
\footnotesize Cogitationes nostrae uane. 
\normalsize| 
\section*{LOMMENT. POMER. }\pstart Hic uides Idololatriam statim esse comitem excaeca- tionis: quomodo uariis imaginibus coluerint Gen- tes Deos, declarant historiae, item Prophetae.  \pend\pstart Deus iustus est iudex, ut ipse statim ubi effinxeri mus nobis in corde nostro Idolum, cui placeant no- stra stercoria opera, (sicut Monachi cum Blattis et Cappis, Gentes cum suis talpis et anseribus fecerunt,) puniat peceatum peccato, iusto iudicio. Nam cum in peccatum aliquod manifestum et crassum incidimus, nihil certius esse nobis debet, quam hoc peccatum fie- ri propter secretius aliud peccatum, quod non uides.  Terreat igitur nos nostra uita, in qua nihil aliud cer- nimus, quam crassissima flagitia: haec sint nobis indi cia aliorum maiorum peccatorum, quae non possumus uidere. Si igitur etiam peccatum illud manifeste cras- sum non agnouerimus, Vae nobis, certe actum est de nobis. St autem agnouerimus, et poenitentia fue- rimus ducti, salua sunt omnia, siquidem potens est Deus ex stercore aurum facere-  \pend
\phantomsection
\addcontentsline{toc}{subsection}{\textit{Quapropter. }}
\subsection*{\textit{Quapropter. }}
\phantomsection
\addcontentsline{toc}{subsection}{\textit{Cordiumsuorum. }}
\subsection*{\textit{Cordiumsuorum. }}\pstart Horrendum illud est, non dicit per Sathanam tra- didit, sed per illud quod iam in corde habent, quod gra uissimum est periculum. Contra quod et oramus: Et ne nos inducas in tentationem.   \pend
\phantomsection
\addcontentsline{toc}{subsection}{\textit{In immunditiam. }}
\subsection*{\textit{In immunditiam. }}\pstart Hanc post exponit quaesit.   \pend
\vspace{0.5cm}\noindent
\vspace{0.2cm}\rule{1cm}{0.2pt}\\ 
\hspace{0.2cm}\textit{mg}
\footnotesize Nostra uita, nil nisi pecca tum. 
\normalsize| 
\section*{10 IN EPIST. AD ROM.  }
\marginpar{[ p.]}
\marginpar{[ p.]}\pstart Id est, agnouerunt unum Deum, et beneficia eius in creaturis, quem solum coli oportebat, sed illi in Idola conuerterunt, et mendaces cultus: ut nos quoq- fecimus, qui pro iustitia fidei, iustitiam operum, et pro sana doctrina, doctrinam daemoniorum, accepi- mus. CONDITA. Id est, creaturas. VSVM. Scilicet coeundi. PRAEMIVM. Quia tradidit eos in concupiscentias. NON PROBAVERVNT. Idest, non uidebatur eis bonum. AGNOSCERENT. Quod li- cuisset ex operibus Det. IN REPROBAM.  Id est, contraria de Deo sentientem.   \pend
\phantomsection
\addcontentsline{toc}{subsection}{\textit{Veritatem. }}
\subsection*{\textit{Veritatem. }}
\phantomsection
\addcontentsline{toc}{subsection}{\textit{Omni uerlutia, }}
\subsection*{\textit{Omni uerlutia, }}\pstart Plura et pauciora hic potuissent enumerari, sed haec certis numeris non debent adstringi.  \pend
\phantomsection
\addcontentsline{toc}{subsection}{\textit{Qui cum Dei iustitiam. }}
\subsection*{\textit{Qui cum Dei iustitiam. }}
\phantomsection
\addcontentsline{toc}{subsection}{\textit{Assentiuntur. }}
\subsection*{\textit{Assentiuntur. }}\pstart Perstringit quarundam Gentium hypocrifim, qui tanen in speciem iusti uideri potuerint.   \pend
\phantomsection
\addcontentsline{toc}{subsection}{\textit{Finis primi capitis.  }}
\subsection*{\textit{Finis primi capitis.  }}
\section*{COMMENT. POMER.  }
\section*{CAP .II. }\pstart Icut superiori capite ostendit, omnes Gentes es- se peccatores, et crassissimos quidem, ita et hoc S capite Iudaeos peccatores declarat. Et ne quid haerea- mus in nominibus     nifesti peccatores in Scrip.ut qui sciunt se scortatores, publicanos, meretrices, aliisque obnoxios peccatis: quemadmodum in Euangelio publicanos et meretri- ces legimus: hi non possunt negare se esse peccato- res, neque negant unquam. Nomine uero Iudaeorum, ueniunt ii, quos Iustitiarios uocamus, a iustitiis carnis, qui olim ludaei et hypocritae uocabantur: ii inquam nolunt esse, aut dici peccatores, quanquam omnibus peccatis sint obnoxii. Et quemadmodum Iudaei om- nia contra impios in Scripturis dicta, in Gentes trans- ferebant, sic nostri hodie in speciem probi et sancti, in alios torquent scripturas. Et talium in primis est iu- dicare alios: Nam hic fructus ex eo prodit, quod pu- tant se tustos. Hi audent dicere coram Deo et homi- nibus: Non sum sicut caeteri hominum, id est, non sum scortator, non fur, non homicida: cum interim- cor ipsorum plenum sit infidelitate, et lingua blas-  \pend\pstart  \pend
\vspace{0.5cm}\noindent
\vspace{0.2cm}\rule{1cm}{0.2pt}\\ 
\hspace{0.2cm}\textit{mg}
\footnotesize Quid uocabu lo Gentuum et Iudaeorum sib uelit Paulus 
\normalsize| 
\hspace{0.2cm}\textit{mg}
\footnotesize Nostri Iudaei qui sint.  
\normalsize| 
\section*{IN EPIST. AD ROM.  }
\marginpar{[ p.13 ]}\pstart phema : Hoc est trabem habere in oculo, et ean non uidere: festucam uero in alio statim consydera- re . Seduide in qua damnatione sint tales. Paulus eis cor durum, et poenitere nescium attribuit: Nam difficulter conuertuntur ad gratiam: facilius crassi pec- catores, qui non ita tument suis iustitiis: ut in Phari- seo apud Lucam uides. Idcirco rectius sperandum de- iis, de quibus superiori cap-dixit Paulus, quam de illis. Nam ii Euangelium audiunt, contra illi, tantum sua iustitia fidunt. Hinc uides diuersum iudicium Dei et mundi: Nam mundus de iustitiariis sic iudicat.  Esto, quod sint peccatores, tamen non tam magni, ut palam impii, id est, scortatores, homicidae etc. Contra Deus, maximos eos peccatores iudicat. Itaque ̃ solus Spiritussanctus arguit illos, in speciem probos. Nolunt credere ex uerbo Euangelii esse salutem, ideoqque ̃ pereunt iusto iudicio Dei. Sic igitur ex hoc cap. clarum fiet, omnes homines esse peccatores.   \pend
\phantomsection
\addcontentsline{toc}{subsection}{\textit{Quapropter. }}
\subsection*{\textit{Quapropter. }}\pstart Haec non solum ad literam de Iudaeis accipienda sunt, uerum de omnibus hominibus, qui faciles sunt ad damnandum et iudicandum alios, tamen praeci- pue haec de Iudaeis dicuntur. Sed haec quomodo cohae- rent cum superioribus, cum inferantur ad illa? Opti- mo ordine. Nam ubi hominum impietatem descripsit, mox poenam eiusdem subiecit: deinde et qui seque-  \pend
\vspace{0.5cm}\noindent
\vspace{0.2cm}\rule{1cm}{0.2pt}\\ 
\hspace{0.2cm}\textit{mg}
\footnotesize Diuersum iu- dicium Dei et mundi. 
\normalsize| 
\section*{COMMENT. POMER.  }\pstart rentur fructus: et tandem concludit, Qui cum Dei iustitiam ex lege naturae, uel ex conscientia naturali nouerint, scilicet, furtum facere, occidere proximum, falsum testimonium, esse malum. At unde nouissent, nisi esset aliqua naturalis conscientia, quae dictaret istae esse mala, et addidit: Non solum etc. quo accusa- uit non solum eos qui in manifesto huiusmodi scelera perpetrarent, sed quibus talia placerent in corde, et iis in corde delectantur, et uolunt uocari boni uiri, et sunt in speciem tales, qualia sunt omnia corda homi- num, nisi gratia Spiritussancti liberentur: aut manife- ste haec faciunt, aut tantum in corde, id quod est coram Deo malefacere. Sed et non aliter homines iudicant, ut si quis proditorum, qui in animo statuerat prodere ciuitatem, ante proditionem deprehensus fuerit: nonne dicetur prodidisse? nam animo et uoluntate id fecisse conuincitur, is profecto damnabitur mortis sententia- Cum igitur etiam homines ita iudicent, quid de Deo di- cemus, qui scrutatur corda, qui iudex est et testis eo- rum operum, quae corde fiunt? Contra, si etiam homicidium feceris, non uoluntate, non adest opus cordis, itaque  id non imputatur apud Deum, neque apud homines: si iu- ste utique tudicare uoluerim, illud externum opus sic ma- nufactum, non est opus, quia corde non fecisti, sed opus internum cordis, etiam apud hoies tibi imputatur, ut iam de proditione dixi. Ergo omnes hoïes sunt tales petonres, ettansi inspeciem talia non uideantur facere: Vt Socrates Camillus, Cato, ubi spiritu Dei non illuminantur .   \pend
\vspace{0.5cm}\noindent
\vspace{0.2cm}\rule{1cm}{0.2pt}\\ 
\hspace{0.2cm}\textit{mg}
\footnotesize Deus iudicat animu et uo- luntatem no- stram.  
\normalsize| 
\section*{ IN EPIST. AD ROM.  }
\marginpar{[ p.]}\pstart Quam saepe Cicero egerit contra conscientiam, testantur etus scripta. Sic Philosophi alii non ulti sunt iniurias, sed a cupiditate uindictae non fuerunt liberi. Itaque in quit, qui adsentiuntur tantum, et non faciunt externe, isti quoquam; sunt obnoxii morti. Sunt autem in hoc nume- ro omnes homines: Id faciunt Philosophi, faciunt et Iudaei, na utrique in speciem boni uideri uolunt, et ali- os tudicare. Ideo dicit: O homo quisquis es, id est, si- ue ex ludaeis aut Gentibus, qui bonus in speciem adpa- res, hypocrita es, bonus diceris, iudicas alios, sed Turpe est doctori, cum culpa redarguit ipsum. Nisi putes te posse fallere Deum.   \pend
\phantomsection
\addcontentsline{toc}{subsection}{\textit{Eadem enim tu facis. }}
\subsection*{\textit{Eadem enim tu facis. }}\pstart Scilicet corde, nam adsentiris, ut supra dictum est. Sed non- potest perpetuo latere hypocrisis, quin se prodat ope- re externo, ut est nunc uidere in Monachis, quorum etiam uita in prouerbium abiit, et puert in plateis clamant.   \pend
\phantomsection
\addcontentsline{toc}{subsection}{\textit{Scimus autem. }}
\subsection*{\textit{Scimus autem. }}\pstart scilicet, uel naturali conscientia, uel lege scripta, ne- mo hoc ignorat, siue sit ex Iudaeis, siue ex Gentibus.   \pend
\phantomsection
\addcontentsline{toc}{subsection}{\textit{Secundum ueritatem. }}
\subsection*{\textit{Secundum ueritatem. }}\pstart Id est, quod Deus alio utatur iudicio, cor intuetur, licet non eadem facis coram homimbus, tamen coram Deo, ideoreus est. EADEM. Scilicet corde.   \pend
\phantomsection
\addcontentsline{toc}{subsection}{\textit{Suffugies. }}
\subsection*{\textit{Suffugies. }}\pstart Vt maxime suffugias iudicium humanum, et te ta- lem non agnoscant homines. Vide igitur quibus uerbis  \pend
\vspace{0.5cm}\noindent
\vspace{0.2cm}\rule{1cm}{0.2pt}\\ 
\hspace{0.2cm}\textit{mg}
\footnotesize Hypocrisis diu se non la- tet. 
\normalsize| 
\section*{COMMENT. POMER.  }\pstart describat sanctos illos homines, quid illis tribuat, im poenitentiam cordis, et qui non possint resipiscere-  \pend
\phantomsection
\addcontentsline{toc}{subsection}{\textit{Diuitias. }}
\subsection*{\textit{Diuitias. }}\pstart Id est, abundantiam bonitatis, id est, quod tam benigne nobiseum agat. Toleran. Id est, quod suffert peccata nostra. Longani. Quod non sta- tim punit.   \pend
\phantomsection
\addcontentsline{toc}{subsection}{\textit{Ignorans. }}
\subsection*{\textit{Ignorans. }}\pstart Ecce ignorantiam Dei illis adscribit, qui putant se Deo proximos, et alios abiectos: Iuxta illud Psal- 13. Dixit insipiens in corde suo, non est Deus. Item tribuit illis contemptum et superbiam, nam conten- nunt has diuitias, quas peccatores et meretrices susci- piunt et admirantur.   \pend
\phantomsection
\addcontentsline{toc}{subsection}{\textit{Poenitentiam. }}
\subsection*{\textit{Poenitentiam. }}\pstart Vt agnoscas te habere cor, quod non timeat Deum, non seruet legem, quae dicit: Dilige Deum et proxi- mum ex corde, sed adscribis eam externis obseruan- tiis, et gloriaris quasi serues.   \pend
\phantomsection
\addcontentsline{toc}{subsection}{\textit{Colligis tibiipsi iram. }}
\subsection*{\textit{Colligis tibiipsi iram. }}\pstart Id est, iudicium Dei facis tibi semper maius, tua ficta sanctitate et hypocrisi: cum interim alii pecca- tores suscipiant Euangelium, et a iudicio Dei libe- rentur.   \pend
\phantomsection
\addcontentsline{toc}{subsection}{\textit{Patefiet. }}
\subsection*{\textit{Patefiet. }}
\vspace{0.5cm}\noindent
\vspace{0.2cm}\rule{1cm}{0.2pt}\\ 
\hspace{0.2cm}\textit{mg}
\footnotesize Ignorant De- um, qui pu- tant se Deo proximos, et alios abie- dtos.  
\normalsize| 
\section*{17 IN EPIST. AD ROM.  }
\marginpar{[ p.]}\pstart Quod nunc absconditum est a mundo: qui supe- riori Cap. describuntur, a mundo iudicart possunt, quis enim tales non diceret malos? Sed de his iudicium est absconditum.   \pend\pstart Scripturarum ignari huiusmodi locos arripiunt, et uno uersiculo uolunt pugnare contra totam Scri- Pturam, et hic contra scopum totius Epistolae, non uidentes praecedentia et sequentia, dicuntque ́: Ecce hic locus confirmat, salutem esse ex bonis operibus, clamant, intonant, uociferantur: Ergo fides non iu- stificat, quia et in Psalmo scriptum est, Reddet uni- cuique secundum opera sua. O caecos homines, qui non uident fidelem facere bona opera, et hunc saluari pro- pter fidem: infidelem uero mala, et hunc damnari propter infidelitatem. Spiritus Det regnans in cordi- bus nostris facit bona opera, reliqua sine spiritu Dei sunt impia, Omnis enim plantatio, quam non plan- tauit pater coelestis, eradicabitur. Porro, dum su- mus in carne, illa opera tamen pura esse non possunt propter carnem, quare ne per illa quidem iustificari possemus. Itaque ex fide tantum est iustitia, non ex ope- ribus, quae testimonium solum sunt pii et fidelis cordis.   \pend
\phantomsection
\addcontentsline{toc}{subsection}{\textit{Qui redditurus. }}
\subsection*{\textit{Qui redditurus. }}
\phantomsection
\addcontentsline{toc}{subsection}{\textit{His quidem. }}
\subsection*{\textit{His quidem. }}\pstart Distributio est, aliquibus est redditurus uitam.  quibus?  \pend\pstart  \pend
\vspace{0.5cm}\noindent
\vspace{0.2cm}\rule{1cm}{0.2pt}\\ 
\hspace{0.2cm}\textit{mg}
\footnotesize Quicquid fit sine spiritu Dei, est impium. 
\normalsize| 
\section*{COMMENT. POMER.  }
\phantomsection
\addcontentsline{toc}{subsection}{\textit{Qui perseuerantes. }}
\subsection*{\textit{Qui perseuerantes. }}\pstart Quod nemo potest facere, nisi fidelis, hypocri- ta non potest impietatem diu caelare, et pertinet hoc uocabulum, perseuerantes, ad fidem.   \pend
\phantomsection
\addcontentsline{toc}{subsection}{\textit{Quaerunt. }}
\subsection*{\textit{Quaerunt. }}\pstart Scilicet a Deo, id quod quoque est fidei: alioqui quid carnalius dici posset? cum Scriptura hoc prohi- beat, et uult ut Dei gloriam quaeramus. HIS. Id est, aliis, mutat hic figurate orationem .   \pend
\phantomsection
\addcontentsline{toc}{subsection}{\textit{Contentiosi. }}
\subsection*{\textit{Contentiosi. }}\pstart Qui contendunt contra ueritatem, non solum co- ram hominibus, sed etiam in corde, nihil enim certi possunt apud se de ueritate constituere, fluctuant et impelluntur uelutt undae, cogitationes cordis eorum.   \pend
\phantomsection
\addcontentsline{toc}{subsection}{\textit{Qui ueritati. }}
\subsection*{\textit{Qui ueritati. }}\pstart Scilicet Dei, quae solius Det est. Nam omnis homo mendax. Sic enim semper apud se disceptant, forte uerum est, forte non, quaerunt quid obiectent ueritati.   \pend
\phantomsection
\addcontentsline{toc}{subsection}{\textit{Iniustitiae. }}
\subsection*{\textit{Iniustitiae. }}\pstart Prudens hic Paulus abstinet a uocabulis fidei et infidelitatis, at per iniustitiam, infidelitatem intelligit.   \pend
\phantomsection
\addcontentsline{toc}{subsection}{\textit{Indignatio. }}
\subsection*{\textit{Indignatio. }}\pstart Scilicet Dei, haec duo in Deo sunt, sed anxietas et afflictioin nobis. PERPETRANTIS.  Scilicet animo, et excorde.   \pend
\vspace{0.5cm}\noindent
\vspace{0.2cm}\rule{1cm}{0.2pt}\\ 
\hspace{0.2cm}\textit{mg}
\footnotesize Gloriam Dei debemus quae- rere.  
\normalsize| 
\section*{28 IN EPIST. AD ROM.  }
\marginpar{[ p.]}\pstart Non hic differentiam facit inter Iudaeos et Grae- cos, neque his priores, illos posteriores, ut exponit mox. GLORIA. Scilicet a Deo.   \pend
\phantomsection
\addcontentsline{toc}{subsection}{\textit{Iudaei. }}
\subsection*{\textit{Iudaei. }}\pstart Ex puro corde, alioqui non potest esse bonum opus, id quod uides per totam Epistolam. Nam quic quid non est exfide, peccatum est, Rom. 14.   \pend
\phantomsection
\addcontentsline{toc}{subsection}{\textit{Operanti. }}
\subsection*{\textit{Operanti. }}
\phantomsection
\addcontentsline{toc}{subsection}{\textit{Non est personarum. }}
\subsection*{\textit{Non est personarum. }}\pstart Exponit quod ante dixerat, Iudaeo primum simul et Graeco. Iudaeum non ideo saluat, quia Iudeus est: neque Graecum aut aliquem ex Gentibus ideo damnat, quia ex Gentibus est. Sed si Gentilis seruat legem, saluatur: si Iudaeus est transgressor legis, damnatur.  Nullus autem hominum unquam seruauit legem ex sic is uiribus, ut Petrus in Actibus, Neque nos, neq pa- tres nostri potuerunt portare onus legis. Hypocritae eum respiciant tantum in externa opera, decipiun- tur, putant se legem implere, cum nihil minus faciant.  Gentilis si ser- vat lege, sal- vatur.   \pend
\phantomsection
\addcontentsline{toc}{subsection}{\textit{Quicunque . }}
\subsection*{\textit{Quicunque . }}\pstart Hic aufert Iudaeis fiduciam legis scriptae, nam iactabant se habere legem scriptam, ut nos hodie de- Euangelio. At ex hoc nulla est iustificatio quod praedicatur Verbum, nisi corde arripiamus . Iudaet dicere potuerunt, Non fecit taliter omni nationi et c.   \pend
\vspace{0.5cm}\noindent
\vspace{0.2cm}\rule{1cm}{0.2pt}\\ 
\hspace{0.2cm}\textit{mg}
\footnotesize 
\normalsize| 
\section*{COMMENT. POMER.  }\pstart Qui annunciat Verbum suum Iacob etc. Sicque con- tendebant Gentes: et fuerunt Romae propter haec dis- cordes inter se Christiani ex Iudaeis et Gentibus: di- cebant illi, Ex nostro sanguine est Christus natus: hae contra, Vos Messiam occidistis, non recepistis eum qui missus erat uobis, ergo Deus nos Gentes recepit in gratiam. Itaque ut aequet Gentes Iudaeis, dicit, tam impii sunt Iudaei in lege scripta, quam sunt Gentes in lege naturali. Sicut illae naturalem uel legem conscien- tiae non seruarunt, sic nec isti legem scriptam. Duplex itaque est lex, naturalis et scripta. Naturalis, est ipsa con- scientia, quae est in omnibus hominibus, et Iudaeis et Gentibus, nisi quod eam homines non uideant, pro- pter eam excaecationem, de qua dixit superiori Cap.  Lex vero scripta est, quae data est per Mosen, in hoc ut innouaretur lex naturalis in hominibus. Siquidem eadem est cum illa, utraque docet colendum Deum, et non laedendum proximum. Erat lex naturalas oblite- rata per malitiam in cordibus hominum, ergo Deus uoluit eam rursus inculcare per Mosen. Dicit itaque , Sicut lex naturalis non facit sanctos, sic nec lex scri- pta, sed seruata. Concludit igitur: Idem iudicium est super Gentes et Iudaeos, utrique ; sunt impii, et con- temptores uoluntatis diuinae.   \pend\pstart Scilicet sine lege scripta, habent conscientiam quae eos iudicat et arguit.   \pend
\phantomsection
\addcontentsline{toc}{subsection}{\textit{Peribunt. }}
\subsection*{\textit{Peribunt. }}
\vspace{0.5cm}\noindent
\vspace{0.2cm}\rule{1cm}{0.2pt}\\ 
\hspace{0.2cm}\textit{mg}
\footnotesize Lae est du- plex.  
\normalsize| 
\section*{IN EPIST. AD ROM.  }
\marginpar{[ p.19 ]}\pstart Id est, Iudaei, quibus data est lex, iudicabuntur per eam legem: est enim Moses qui ipsos accusat.  Excaecabat eadem opinio Gentes et Iudaeos, Gentes dicebant, qui benefacit, bonus est, sic et iudaei. Sed interim non intelligebant, quae essent bona opera, quod uidere potuissent, si non fuissent excaecati, Nempe se- non amare legem, sed omnia facere quodam fastidio et odio cordis aut amore sui.   \pend
\phantomsection
\addcontentsline{toc}{subsection}{\textit{In lege. }}
\subsection*{\textit{In lege. }}
\phantomsection
\addcontentsline{toc}{subsection}{\textit{Qui audiunt. }}
\subsection*{\textit{Qui audiunt. }}\pstart De externo auditu intellige, praedicatur eis lex, legunt, sed frustra. Alians saepe, audire, in scripturis credere est, ut Audiifrael etc. FACTIS.  Id est, operibus.   \pend
\phantomsection
\addcontentsline{toc}{subsection}{\textit{Legem. }}
\subsection*{\textit{Legem. }}\pstart Id est, id quod lex requirit, nempe, ne concupiscant, id est, ne scortentur corde, bene uelint proximo. Vi- ua illa sunt facta Christianorum, siquidem sunt ex uir- tute Spiritussancti: facta hypocritarum sunt mortua, id est, in corde talia non sunt, qualia foris adparent.  Phariseus in Luca, non est occisor foris, uerum in cor- de omnes homines occidit: non sum, inquit, ut caete- ri etc. Huiusmodi sententiis abutuntur stulti interpre- tes pro sua libidine contra Scripturam.   \pend
\phantomsection
\addcontentsline{toc}{subsection}{\textit{Nam cum Gentes. }}
\subsection*{\textit{Nam cum Gentes. }}\pstart Quidam de Gentibus conversus ad Christum, haec  \pend
\vspace{0.5cm}\noindent
\vspace{0.2cm}\rule{1cm}{0.2pt}\\ 
\hspace{0.2cm}\textit{mg}
\footnotesize Iudei iudica- buntur per Mosen.  
\normalsize| 
\hspace{0.2cm}\textit{mg}
\footnotesize Opera hypo- critarum mor- tua.  
\normalsize| 
\section*{COMMENT. POMER.  }\pstart intelligi uolunt, sed simplicius est, si conditione dictum intelligas, utpote, Nam cum Gentes fecerint etc. id est, si facerent Gentes ex naturali lege, ea quae sunt legis scriptae: sed quando faciunt? nunquam: con- scientiam habent quodreuereri debeant Deum, non facere malum. At ubi faciunt mala uel bona, excu- sat eas uel accusat conscientia eorum: habent con- scientiam bonam uel malam, secundum quod faciunt.  Itaque haec lex naturalis, id est, conscientia, est in omnibus hominibus, ut possint uidere quae sint mala, quae bona, aut facere habent conscientiam in quibus dam se benefecisse, in quibusdam male.  \pend\pstart Scilicet scriptam. Sententia est, uos Iudaei putatis uos praestare Gentibus propter legem, an ignoratis apud Deum non esse contentionem de lege scripta, uel non scripta, sed impleta? Si aliquis Gentium im- pleat legem quam mandauit Deus, nonne melior erit Iudaeo non faciente. Itaque nihil prodest lex scripta, nisi sit impleta.   \pend
\phantomsection
\addcontentsline{toc}{subsection}{\textit{Quae legem. }}
\subsection*{\textit{Quae legem. }}\pstart Id est, naturali lege, habent enim conscientiam, que dictat talia facienda uel omittenda, quae lex scripta-  \pend
\phantomsection
\addcontentsline{toc}{subsection}{\textit{Natura. }}
\subsection*{\textit{Natura. }}
\phantomsection
\addcontentsline{toc}{subsection}{\textit{Dibi sunt lex, }}
\subsection*{\textit{Dibi sunt lex, }}\pstart Id est, per conscientiam suam, pro lege. AVDIVNT.  uel negligendo.  \pend
\vspace{0.5cm}\noindent
\vspace{0.2cm}\rule{1cm}{0.2pt}\\ 
\hspace{0.2cm}\textit{mg}
\footnotesize Conscientia docet, quae na Deo debea- mus. 
\normalsize| 
\hspace{0.2cm}\textit{mg}
\footnotesize Implere legem.  
\normalsize| 
\hspace{0.2cm}\textit{mg}
\footnotesize sua corda habent Scilicet implendo 
\normalsize| 
\section*{}
\marginpar{[ p.]}
\phantomsection
\addcontentsline{toc}{subsection}{\textit{}}
\subsection*{\textit{}}
\phantomsection
\addcontentsline{toc}{subsection}{\textit{}}
\subsection*{\textit{}}\pstart  \pend\pstart  \pend\pstart  \pend
\phantomsection
\addcontentsline{toc}{subsection}{\textit{}}
\subsection*{\textit{}}
\vspace{0.5cm}\noindent
\vspace{0.2cm}\rule{1cm}{0.2pt}\\ 
\hspace{0.2cm}\textit{mg}
\footnotesize 
\normalsize| 
\section*{IN EPIST. AD ROM.  }
\marginpar{[ p.21 ]}\pstart Id est, ea quae bona sunt. Vnde hoc nosti? quia institutus es in Mose et lege. CONFIDIS- VE. Non dico quod es. CAECORVM.  Id est, ignorantium.   \pend
\phantomsection
\addcontentsline{toc}{subsection}{\textit{Laudas eximia. }}
\subsection*{\textit{Laudas eximia. }}
\phantomsection
\addcontentsline{toc}{subsection}{\textit{Lumen in tenebris. }}
\subsection*{\textit{Lumen in tenebris. }}\pstart Id est, eorum qui errant, non nouerunt Deum, ueritatem etc. ERVDITOREM. Copia est in iis uerbis.  \pend\pstart Id est, qui habeat. PER LEGEM. i. scriptu- ram sanctam, in qua nihil non adsecutunte putas. Haec omnia secundum affectum illorum hominum in spe- ciem sanctorum locutus est. DOCES. Sci- licet ex lege, scripturis sanctis. TEIPSVM.  Id est, ea ipsa quae doces, non imples .   \pend
\phantomsection
\addcontentsline{toc}{subsection}{\textit{Habentem. }}
\subsection*{\textit{Habentem. }}
\phantomsection
\addcontentsline{toc}{subsection}{\textit{Furaris. }}
\subsection*{\textit{Furaris. }}\pstart Id est, in corde concupisces aliena, ut in. 7. Cap.  uides. Loquitur uulgaribus uocabulis de peccatis, abstinet interim a uocabulis, pietatis et impietatis, usque in suum locum. Gentes dicebant, Bonus uir est, qui non furatur, non facit malum etc. Iudaei, Bo- nus uir est, qui facit bona opera: Aliter uero Spiri- tussanctus, qui non curat mores istos philosophicos nec hypocrisin sanctulorum: ibi est peccatum, ubi mundus putat esse sanctitatem. ideoque hic dicit, pu-  \pend
\vspace{0.5cm}\noindent
\vspace{0.2cm}\rule{1cm}{0.2pt}\\ 
\hspace{0.2cm}\textit{mg}
\footnotesize Hypocrisis maximum pec- catum.  
\normalsize| 
\section*{OMMENT. POMER.  }\pstart tas te non esse furem, cum tamen sis fur-  \pend\pstart Scilicet corde, ut Christus Matth.5. imterpreta- tur, Qui uidet mulierem ad concupiscendum. Quis illorum est, qui hodie uoluerunt esse castissimi, qui non hic peccarit? Hic, si possunt nostri Monachi et Monachae, defendant suam castitatem, nisi ipse Chri- stus et Paulus sic interpretaretur, nemo posset id es- se adulterium credere, quod in corde committitur.   \pend
\phantomsection
\addcontentsline{toc}{subsection}{\textit{Committis. }}
\subsection*{\textit{Committis. }}\pstart Furaris Deo suum honorem, nam sacrilegium est furto auferre rem sacram. Est autem sacratissimae res honor Dei, quem Iusticiarii sibi arrogant, tribuem̈- tes iustitiam suis operibus. Quod potest esse maius sa- crilegium? Gloria et honor iustificationis soli Deo de- betur, non nostris uiribus aut meritis. Nam dixit in Propheta, Gloriam meam alteri non dabo. O im- pios Hypocritas, qui externa Idola, et Gentes idolo latras execrantur, interim tamen, quod maximum et sanctissimum est in Deo, sibi arrogant: Christus uo- luit Matth. 23. ut tantum unum magistrum agnosce- remus, Quid igitur stulti nobis sumimus?  \pend
\phantomsection
\addcontentsline{toc}{subsection}{\textit{Sacrilegium admittis. }}
\subsection*{\textit{Sacrilegium admittis. }}
\phantomsection
\addcontentsline{toc}{subsection}{\textit{Dehonestas. }}
\subsection*{\textit{Dehonestas. }}
\vspace{0.5cm}\noindent
\vspace{0.2cm}\rule{1cm}{0.2pt}\\ 
\hspace{0.2cm}\textit{mg}
\footnotesize Sacrilegiun quid.  
\normalsize| 
\section*{EPIST. AD ROM.  22 }
\marginpar{[ p.]}\pstart Id est, transgrederis legem. Tu dicis, Gentes sunt maledictae, quia non habent legem: At interim, tu hypocrita, non uides quod ipse sis transgressor legis.  Quicunque enim non erudiuntur spiritu, non possunt nisi crassa peccata cernere: bypocritarum peccata la- tent et al sconduntur in eorum ipsorum damnationem, non uident infidelitatem esse peccatum, et soli Deo fidendum, nesciunt: qui hoc non uident, non uident odium et inuidiam quo laborant, nec uident concu- piscentiam esse peccatum. Somniant interim sibi iu- sticiam ex suis operibus quae faciunt. Sunt igitur de- teriores Gentibus, contemptores sunt et irrisores no- minis Dei in cordibus suis, ut Psalm. Dixit insipiens etc. 'indicat, haec quia manifesta sunt, non libet plu- ribus persequi.   \pend
\phantomsection
\addcontentsline{toc}{subsection}{\textit{Namnomen Dei. }}
\subsection*{\textit{Namnomen Dei. }}\pstart Vides Paulum in hoc Cap. de occulta impietate et ficta sanclitate loqui, propter quam nemo male audit inter homines, sicut hactenus in Monachis ex- perti sumus, nemo non rectissime de ipsis sensit. Quid ergo dicemus ad Paulum, cum illam occultam hypo- crisin non iudicet mundus? Respondeo. Ista sancti - tas simulata, non potest se continere, quin aliquan- do prodat se suis fructibus: sicut hodie nemo non ha- bet in ore Monachos, nemo non ridet eorum san- ctitatem, reuelata est iam ignominia huius meretri- cis, ut et pueri in plateis ipsam rideant uulgaeribus Ficta sancti- tas.   \pend
\vspace{0.5cm}\noindent
\vspace{0.2cm}\rule{1cm}{0.2pt}\\ 
\hspace{0.2cm}\textit{mg}
\footnotesize Qui non eru- diuntur spiri- tu, nil nisi pec- cata cernunt.  
\normalsize| 
\hspace{0.2cm}\textit{mg}
\footnotesize 
\normalsize| 
\section*{COMMENT. POMER.  }\pstart quibusdam cantilenis: quod et accidit ludaeis, quem admodum Prophetae testantur. Est uero hic locus ex Ezechielis. 36. desumptus, tametsi quidam uelint ex Esaiae. 52. ipsum prolatum, ubi describuntur Gen- tium irrisiones, et uerba quaedam mimitica, quibus amarulente contra Iudaeos utebantur. Nam cum spe- ciem pietatis praetexebant, et pia docerent, tamen impia quaedam et scelerata Gentes in ipsis uidebant, quae mala esse iudicabant lege naturali. Sic autem ha- bet Ezech. Polluerunt nomem sanctum meum, cum di- ceretur de eis, populus Domini iste est. Ita et usura, luxus, auaritia, tyrannis, odia, inuidiae, etiam mun- do prodiderunt illam fictam sanctitatem Papisticorum Monachorum et Pontificum: hoc ergo est blasphema- ri nomen Dei in Gentibus. Caueamus et nos hodie hocipsum, qui maxime uolumus uideri Euangelici, ut non intus lateat sordidissima sentina omnium utiorum.   \pend
\phantomsection
\addcontentsline{toc}{subsection}{\textit{Nam circuncisio. }}
\subsection*{\textit{Nam circuncisio. }}\pstart Iam exponit quod dixit, quod externa sanctitas ni- hil sit, sine iustitia et sanctitate cordis: ubi illa est, ibi et externa sunt sancta: sine hac, omnia sunt hypo- crisis maledicta. Ergo hic eis supercilium detrahit Pau- lus, dicens: O Iudaei, nihil prodest uestra circuncisio externa, sine circuncisione cordis, in corde seruatur lex, non externis signis. Quemadmodum et baptis- mus inutilis est, nisi credas in Christumite regeneran  \pend
\vspace{0.5cm}\noindent
\vspace{0.2cm}\rule{1cm}{0.2pt}\\ 
\hspace{0.2cm}\textit{mg}
\footnotesize Polluere no- men Dei. 
\normalsize| 
\hspace{0.2cm}\textit{mg}
\footnotesize Circuncisio externa nihil prodest.  
\normalsize| 
\section*{23 IN EPIST. AD ROM.  }
\marginpar{[ p.]}\pstart tem et abluentem igni et Spiritusancto, nam etiam baptismus aquae, est baptismus Christi. Prodest cir- cuncisio, et illam magnifacio, sed absque re illam ha- bere quid prodest?  \pend\pstart Id est, si legem in corde non seruaris. CIR- CVNCISIO TVA. Scilicet, quae incute est.   \pend
\phantomsection
\addcontentsline{toc}{subsection}{\textit{Transgressor. }}
\subsection*{\textit{Transgressor. }}
\phantomsection
\addcontentsline{toc}{subsection}{\textit{In praeputium. }}
\subsection*{\textit{In praeputium. }}\pstart Id est, nihil est coram Deo, perinde est, ac sies- ses Gentilis.  \pend\pstart Hoc est, quod supra dixit, nam cum Gentes quae legem etc. id est, si Ethnicus dilexerit proximum, et fecerit ea quae mandat lex, is iustus coram Deo re- putabitur. Non tamen ex iis locis conuinces, Gen- tes natur ali lege, facere posse legem, et consequi iu- stificationem, quae ex solo Christo est. Sic in Euange- cit, quod possit seruare mandata etc.   \pend
\phantomsection
\addcontentsline{toc}{subsection}{\textit{Ergo. }}
\subsection*{\textit{Ergo. }}
\phantomsection
\addcontentsline{toc}{subsection}{\textit{Imputabitur . }}
\subsection*{\textit{Imputabitur . }}\pstart Id est, habebitur coram Deo pro uero Iudaeo,  \pend
\phantomsection
\addcontentsline{toc}{subsection}{\textit{Et iudicabit. }}
\subsection*{\textit{Et iudicabit. }}\pstart Periphrasis est Gentilis, eius qui non est ex carne Abrahae natus. SI LEGEM SERVA.  Ia das muss da bey seyn.   \pend
\vspace{0.5cm}\noindent
\vspace{0.2cm}\rule{1cm}{0.2pt}\\ 
\hspace{0.2cm}\textit{mg}
\footnotesize Iustificatio non- consequitur le- naturali. 
\normalsize| 
\section*{COMMENT. POMER.  }\pstart Id est, damnabit te. PER LITERAM- Qui in ipsa litera es transgressor, nam non habes in corde tuo quod litera exigit, lex igitur erit tibi in dam nationem, quae debebat esse in salutem. Sic omnem fidu- ciam aufert a sanctulis, ut cognoscant se habere cor im purum. Sic Christus de regina Austri loquitur.  \pend
\phantomsection
\addcontentsline{toc}{subsection}{\textit{Iudicabit. }}
\subsection*{\textit{Iudicabit. }}\pstart Coram Deo iustums est, non Iudaeus coram mundo.  Multi clamant, Domine Domine, sed in corde non agnoscunt Deum patrem. Iudaeus, confitentem significat, non labiis, sed corde confiteamur nomen Domini-  \pend
\phantomsection
\addcontentsline{toc}{subsection}{\textit{ludaeus. }}
\subsection*{\textit{ludaeus. }}
\phantomsection
\addcontentsline{toc}{subsection}{\textit{IN MANIFESTO. In cute coram hominibus.  }}
\subsection*{\textit{IN MANIFESTO. In cute coram hominibus.  }}\pstart Id est, corde per fidem: qui confitetur in corde Chri- stum, saluus erit. Rom. 1o.Itaque ludaei non intelligebant uoluntatem Dei in circuncisione, quam tamen Moses in Deuteronomio interpretatus est , dicebat: Cir- cuncidite praeputia cordium uestrorum. Et Hieremi- as, Incredulos, incircuncisos adpellauit.   \pend
\phantomsection
\addcontentsline{toc}{subsection}{\textit{In occulto. }}
\subsection*{\textit{In occulto. }}\pstart Quae est omne quod foris agi potest: Spums autem est motio Dei quae uiuit in corde, qui uere circuncidit cor: et hanc circuncisionem hoies non uident, non est ergo regnum Dei in obseruationibus externis et elementis mum- di. Idcirco falluntur, qui Christianum ab externis iudi-  \pend
\phantomsection
\addcontentsline{toc}{subsection}{\textit{Non litera. }}
\subsection*{\textit{Non litera. }}
\vspace{0.5cm}\noindent
\vspace{0.2cm}\rule{1cm}{0.2pt}\\ 
\hspace{0.2cm}\textit{mg}
\footnotesize Iudeus confi- tentem signi- ficat.  
\normalsize| 
\section*{EPIST. AD ROM.  }
\marginpar{[ p.24 ]}\pstart cant. Constat Deo iustitia uera, non hoibus, et Deus illam laudat et probat, non mundus, qui totus contra eam insanit, quemadmodum hodie uidemus, nam nihil est ipsi Belial cum Christo filio aeterno Dei, qui est filius lucis.   \pend
\endnumbering\beginnumbering\section{CAP .III.}\pstart aec parerga sunt, et digressio quaedam, in qua ob- L I iectioni impiorum respondet. Ex praedictis enim aliquis potuisset inferre, Si Gentes aequantur ludaeis, er- go nihil ipsi habent praerogatiui prae illis, cum Deus non fecerit taliter omni nationi, elegit eos in facerdotium re- gale, et peculiariter eis commisit Verbum, igitur aut tu impie doces, aut omnis Scriptura, quae praefert Iudaeos Gentibus, est mendax. His respondet Paulus.   \pend
\phantomsection
\addcontentsline{toc}{subsection}{\textit{\huge\textbf{Gentibus, est mendax . Hisrespondet Paulus . }\normalsize Quae utilitas? }}
\subsection*{\textit{\huge\textbf{Gentibus, est mendax . Hisrespondet Paulus . }\normalsize Quae utilitas? }}\pstart Scilicet prae Gentibus. MVLTVM. Respondet ad obiecta. Primu, iquit, eis promissiones diuiae commissae sunt: magna haec em gratia, quam saepe Prophetae decantant, quae olim ludaeis ostensa est, ut hodie nobis per Euan- gelium. Quia non potest frustra Verbum Dei emitti, Esa- test emitti.  55.ut maxime non omnes illud suscipiant, et quemadmodum no ctu in tenebris nemo uidet, nisi ad candelam accensam.  Ita omnes nihil uidemus sine Verbo Dei, quod uera est lucerna:et sicut lucerna non omnibus est utilis, utpote- somnolentis et caecis, sed solum euigilantibus: Sic Ver- bum Det excaecatis et cessatoribus nihil prodest, sed tan- tum iis quilampadibus accensis uocem sponsi expectant.   \pend
\vspace{0.5cm}\noindent
\vspace{0.2cm}\rule{1cm}{0.2pt}\\ 
\hspace{0.2cm}\textit{mg}
\footnotesize Verbum Dei frustra non po- test emitti.  
\normalsize| 
\section*{COMMENT. POMER.  }\pstart Non ferunt, ne hodie quidem Iudaei et hypocritae, ut sibi conferantur meretrices et publicant, ut maxime sexcenties dicant, Mea culpa, Mea maxima culpa.  Indignantur, rumpi uolunt quoties audiunt, Mere- trices et publicani, praecedent uos in regno Dei. Per- tinet autem proprie ad scopum huius Epistolae, quod Paulus Iudaeos Gentibus aequat, ut tandem inferat, om- nes homines esse peccatores. Cur uero hic praeferat Iudaeos Gentibus, similitudine rectius intelliges. Pri- mum, aliud est promittere, aliud praestare promissa.  Iudaeis et Gentibus sunt promissiones factae, ut in Ge- nesi, In semine tuo benedicentur omnes Gentes. At Iudaei, habebant id obsignatum et munitum literis Sa- cris. Gentes uero ignorabant ad se pertinere promis- siones, cum non haberent sacras literas in testimonium ut Iudaei, quemadmodum uidemus in Mose. Sicut pro- mitto duobus fratribus centum aureos nummos, uni- do literas meas in testimonium promissionis, alter ue- ro promissionem factam ignorat. Tempore designato do illi suam partem, do et alteri qui ignorabat pro- mnissum, nonne is melius habuit, cui dederam meas literas in testimonium, altero, quem fugiebat promis- sio? ut maxime tantum acceperit, quantum ille. Se- cundum promissionem hunc praefero, secundum rem promissam nihil habet praerogatiui prae altero, et fit, ut secundus gratissimus sit, alter ingratus, quia dicit, tu promisisti, ergo debebas mihe Sic Gentes, nihilta  \pend
\vspace{0.5cm}\noindent
\vspace{0.2cm}\rule{1cm}{0.2pt}\\ 
\hspace{0.2cm}\textit{mg}
\footnotesize Cur Paulus praeferat Iu- daeos Genti- bus.  
\normalsize| 
\section*{IN EPIST. AD ROM.  }
\marginpar{[ p.25 ]}\pstart le sperantes, promissiones acceperunt, ideo sunt gratis- simae et Deum laudant. Iudaei uero et iustitiarii putant se suis meritis accipere, ut parabola Euangelica de con- ductis operariis indicat: ergo tustitiarii nequaquam sunt meliores publicanis et meretricibus, nam si me- liores essent, id esset ex Verbo Dei, quod eis erat com- missum. Non tamen solum Iudaeis erat praestandum, sed et Gentibus, quibus aeque erat promissum, sed ignorauerunt, ne igitur uideretur Paulus contemne- re Verbum Dei, in quo praeferri possent Iudaei Gen- tibus, haec praeoccupauit.   \pend\pstart Scilicet, est in quo sunt praeferendi, habebant pro- missiones salutis: at ubi salus praestanda erat, et Gen- tibus haec contigit promissio.  \pend
\phantomsection
\addcontentsline{toc}{subsection}{\textit{Illud. }}
\subsection*{\textit{Illud. }}
\phantomsection
\addcontentsline{toc}{subsection}{\textit{Quid enim si quidam. }}
\subsection*{\textit{Quid enim si quidam. }}\pstart Alia est occupatio, ne quis diceret, quid profue- runt eis oracula Dei, cum maior pars Iudaeorum non crediderit illis, imo grauius propter haec damnati sunt et abiecti, nam patres in deserto perierunt, et tan tum Iosue et Caleb ingressi sunt terram: deinde qui bus data est terra, ii rursus eiecti sunt, et tandem ubi redirent ex captiuitate Babylonica, occiderunt Pro- phetas et Messiam. Vide orationem Stephani Act.7.  QVID. Scilicet dican-  \pend
\vspace{0.5cm}\noindent
\vspace{0.2cm}\rule{1cm}{0.2pt}\\ 
\hspace{0.2cm}\textit{mg}
\footnotesize Iusticiarii, no meliores pu- blicanis. 
\normalsize| 
\section*{COMMENT. POMER.  }\pstart Id est, fidem in promissiones Dei. Itaque hic obser ua, fides Dei ea est, qua Deo credimus, et eius Ver bo, non qua Deus credit. Nemo autem credit Deo, nisi loquatur, sicque uerbum materiam (ut ita dicam) fidei exhibet, in qua exercetur. Sic Christus loquitur de fide Mar.11.Habete fidem Deietc. ILLO- RVM. Scilicet, qui noluerunt credere.  \pend
\phantomsection
\addcontentsline{toc}{subsection}{\textit{Fidem Dei. }}
\subsection*{\textit{Fidem Dei. }}\pstart Tametsi Deus promiserit, et Iudaei non credide- rint, non sequitur Deum esse mendacem, ut caro in- fert. Deus non potest mentiri, promisit credentibus, non incredulis. Nam promissio ibi tantum est promis- sio, ubi est fides. Manet ergo Deus uerax, ut infert.   \pend
\phantomsection
\addcontentsline{toc}{subsection}{\textit{Absit. }}
\subsection*{\textit{Absit. }}\pstart Id est, maneat. Nos ex natura sumus mendaces et peccatum . Quicquid de pictate loquimur, menda- cium est et hypocrisis. Deus non potest fieri mem- dax, ettam si totus mundus periret, credentes acci- piunt, quod incredult non merentur.   \pend
\phantomsection
\addcontentsline{toc}{subsection}{\textit{Sit. }}
\subsection*{\textit{Sit. }}
\phantomsection
\addcontentsline{toc}{subsection}{\textit{Vt iustificeris. }}
\subsection*{\textit{Vt iustificeris. }}\pstart Id est, O Deus, ut tustus inueniaris.  \pend\pstart Scilicet, sentientes contra tuam ueritatem, scilicet, te non inueniri iustum.  \pend
\phantomsection
\addcontentsline{toc}{subsection}{\textit{Vincas. }}
\subsection*{\textit{Vincas. }}
\vspace{0.5cm}\noindent
\vspace{0.2cm}\rule{1cm}{0.2pt}\\ 
\hspace{0.2cm}\textit{mg}
\footnotesize Natura su- mus menda- ces.  
\normalsize| 
\section*{26 IN EPIST. AD ROM.  }
\marginpar{[ p.]}\pstart Scilicet ab eis. Admirabilis est haec sententia. Vi- de Annotationes nostras, et D. Martini in hunc Psal mum. Sic enim habet Psalmus, Tibi soli. Id est, in- oculis tuis sum peccator, ut maxime sim purus coram hominibus. Hoc ego tibi conqueror o Deus. Ad quid? ut tu tustus inueniaris, et uincas et ostendas te uera- cem, ubi ab aliis mendax iudicaris. Summa est, ut Deus uerax sit, necesse est me esse peccatorem. Tota Scriptura clamat nos esse peccatores, et per solum Christum iustificari, ergo necesse est, ut praecedat pec- catum, et declaretur iusticia Dei. Quomodo enim mihi Deus remitteret peccatum, cum non haberem peccatum, Esaiae. 4. Si abluerit Dominus sordes etc.  Iusticiarii cogitant, neminem occidi, sic de aliis pec- catis. Natura enim non nouit fontem et radicem pec- catorum, non nouit nisi fructus, id est, crassa peccata quae coram hominibus sunt peccata, ut Cap. 1. Pec- catores autem dicunt, tu promisisti te iustificaturum in iustos, ecce ego intustus sum, peccator sum, manife- sta in me gloriam tuam, et ueritatem tuam. Sic igi- tur iusticiarii tudicant Deum, et faciunt ipsum mem- dacem, cum se non putant opus habere Det iusti- cia. Verum Deus suo tempore, et ipsos uincet in- tentatione, scilicet, et horrore mortis inferni et peccatt, in qua pereunt, quia non aedificarunt su- pra petram, sed supra arenam, Matthaei septimo.      \pend
\phantomsection
\addcontentsline{toc}{subsection}{\textit{ludicaris. }}
\subsection*{\textit{ludicaris. }}\pstart  \pend
\vspace{0.5cm}\noindent
\vspace{0.2cm}\rule{1cm}{0.2pt}\\ 
\hspace{0.2cm}\textit{mg}
\footnotesize 
\normalsize| 
\section*{COMMENT. POMER.  }\pstart Tum confusione confundentur, neque erit eis lux, nam non possunt inuocare nomen Domini.  \pend\pstart Hoc omnino est extra propositum, sed mos est Pau- lo, propter citationem aliquam ex Scriptura, adduce- re quaedam extra propositum: tamen ea ut necessaria tractat, ne quis ex citata Scriptura aliquid minus in- telligat. Dixerat enim ex Psalm. Vt iustificeris, Quod. intelligere non potes, nisi legas priora, Tibi soli pec- caut . Alias enim flatim inferret caro,Ergo pecce- mus, nam in hoc faciemus Deo maximum honorem, quo grauius fuerit nostrum peccatum, tanto erit illu- strior Det gloria. Certe uerum est, Deum ex nostris peccatis glorificari. Sed quod infertur, Ergo pecca- bimus, hoc malum est. Simile exemplum Scripturae citatae per Paulum, habes Ephes . 4. ubi citarat ex- Psalm. Ascendens in altum, captiuam duxit captiui- tatem, dedit dona hominibus. Non ibi propter pri- mam aut fecundam particulam citatur Scriptura, ue- rum propter ultimam, quae est de donis, de quibus in eo loco Paulus agit. Sic et hoc loco, adducit Pauluis uersum, qui potuisset occasio esse impiis, ut dicerent, Ergo peccemus fortiter. Ideo praeoccupat et exponit.   \pend
\phantomsection
\addcontentsline{toc}{subsection}{\textit{Quod si iniusticia. }}
\subsection*{\textit{Quod si iniusticia. }}
\phantomsection
\addcontentsline{toc}{subsection}{\textit{Iniusticia. }}
\subsection*{\textit{Iniusticia. }}\pstart  \pend
\phantomsection
\addcontentsline{toc}{subsection}{\textit{}}
\subsection*{\textit{}}\pstart  \pend
\vspace{0.5cm}\noindent
\vspace{0.2cm}\rule{1cm}{0.2pt}\\ 
\hspace{0.2cm}\textit{mg}
\footnotesize Verus sensus Scripturae.  
\normalsize| 
\section*{17 EPIST. AD ROM.  }
\marginpar{[ p.]}\pstart Scilicet, contra nostram iniusticiam, qui tudicabit opera nostrae impietatis in extremo iudicio. Certe si- iudicabit, iniuste faciet, nam id ertt contra suum ho- norem. Dum peccauimus, non contra eum fecimus, sed pro etus honore. Iniustus equidem esset, si ob id nos puniret, ex quo illi oritur gloria.   \pend
\phantomsection
\addcontentsline{toc}{subsection}{\textit{Humano more. }}
\subsection*{\textit{Humano more. }}\pstart Temperat hanc blasphemiam. q. d. carnalis ho- minis et impii hoc est argumentum. ABSIT.  Quasi diceret: Blasphemia est.  \pend
\phantomsection
\addcontentsline{toc}{subsection}{\textit{Nam quomodo. }}
\subsection*{\textit{Nam quomodo. }}\pstart Qui uult arguere alios, non debet iis esse obno- xius, propter quae alios arguit. Si hoc decet hominem maxime Deum. Necesse est igitur eum irreprebensi- Irreprehens- bilem esse, qui alios uult damnare et iudicare. Sed bilis sit, qui ne quis dubitet hoc esse argumentum carnale, ampli- alios vuli ar- ficat ipsum, dicens.   \pend
\phantomsection
\addcontentsline{toc}{subsection}{\textit{Etenim si ueritas Dei, per meum mendacium. }}
\subsection*{\textit{Etenim si ueritas Dei, per meum mendacium. }}\pstart Id est, peccatum (quod ante dixerat iniustitiam)  \pend
\phantomsection
\addcontentsline{toc}{subsection}{\textit{Excelluit. }}
\subsection*{\textit{Excelluit. }}\pstart Id est, gloriosa apparuit, ad hoc, ut ipse inde glorificaretur.   \pend
\phantomsection
\addcontentsline{toc}{subsection}{\textit{Quid iudicor. }}
\subsection*{\textit{Quid iudicor. }}\pstart  \pend
\vspace{0.5cm}\noindent
\vspace{0.2cm}\rule{1cm}{0.2pt}\\ 
\hspace{0.2cm}\textit{mg}
\footnotesize  guere.  
\normalsize| 
\section*{COMMENT. POMER.  }\pstart Id est, a Deo, nam de iudicio Dei, non hominum, dicit, imo deberet me peccatorem exaltare, non dan- nare, cum hoc agam, ut gloria eius manifestior fiat.  An ueluti peccator iudicor, id est, ueluti admiserim aliquid contra Deum, cum pro eo fecerim . Hoc ar- gumentum carnis quisque apud se inueniet, ubi Spiri- tui sancto non permiserit se docendum. Praedestina- tio offendit carnem, nam non potest dicere cum Chri- sto, Fiat uoluntas tua. AC NON POTIVS.  Vuer das nit besser.   \pend
\phantomsection
\addcontentsline{toc}{subsection}{\textit{( Quemadmodum de nobis.) }}
\subsection*{\textit{( Quemadmodum de nobis.) }}\pstart Quando praedicatur Euangelium, non possunt car- nales non dicere, Ecce Apostoli et Euangelistae do- cent mala facere, nam prohibent bona opera. Id quod et nos audire cogimur hodie, dum tamen non opera, sed fiduciam in opera prohibeamus. Est autem haec sententia in parenthesi posita. DICERE.  Id est, docere.   \pend
\phantomsection
\addcontentsline{toc}{subsection}{\textit{Faciamus mala. }}
\subsection*{\textit{Faciamus mala. }}
\phantomsection
\addcontentsline{toc}{subsection}{\textit{Nam illud multo melius est. VT VENI- ANT BONA. Id est, oriatur inde gloria Dei- }}
\subsection*{\textit{Nam illud multo melius est. VT VENI- ANT BONA. Id est, oriatur inde gloria Dei- }}
\phantomsection
\addcontentsline{toc}{subsection}{\textit{Quorum damnatio. }}
\subsection*{\textit{Quorum damnatio. }}\pstart Haec est plena solutio argumenti carnalis, tametsi quidam putent Apostolum id contempsisse, et non soluisse-q. d-qui dicunt, si gloria Dei illustrior fit, peccabimus, illorum damnatio est iusta, quia nolunt resipi  \pend\pstart  \pend
\vspace{0.5cm}\noindent
\vspace{0.2cm}\rule{1cm}{0.2pt}\\ 
\hspace{0.2cm}\textit{mg}
\footnotesize Praedestina- tio offendu carnem.  
\normalsize| 
\hspace{0.2cm}\textit{mg}
\footnotesize Bona opera non sunt pro- hibenda.  
\normalsize| 
\section*{ IN EPIST. AD ROM.  }
\marginpar{[ p.]}\pstart scere, quasi non satis peccarint: et cum nihil aliud sint q̃ peccatum, putant se non satis peccasse, quemadmo- dum illi in Malachia. Tales autem non desyderant re- missionem peccatorum, quando peccatis gaudent, id- circo id consequuntur, ut pereant iusto iudicio Dei.   \pend\pstart Hic redit ad institutum, quia primo Cap. omnes Gentes peccatt arguebat, nec in hoc esse excusabiles, quod non habuerint legem scriptam, ostendebat. In secundo, et Iudaeos peccatores pronunciarat, ut - cunque simulerent iusticiam. Ideoque ; nec ipsos excusa- ri posse coram Deo. Nam quae ad literam tantum fiunt, ula nihil sunt coram Deo, sed quae sunt in spiritu. Ex iis ergo duabus propositionibus infertur tertia, Ergo omnes homines sunt peccatores: itaque inquit, Quid igitur praecellimus, scilicet, nos Iudaer (siquidem et Paulus suerat ludaeus). EOS. Scilicet Gentes.   \pend
\phantomsection
\addcontentsline{toc}{subsection}{\textit{Quid igitur. }}
\subsection*{\textit{Quid igitur. }}
\phantomsection
\addcontentsline{toc}{subsection}{\textit{Nullo modo. }}
\subsection*{\textit{Nullo modo. }}\pstart Quon ergo antea dixerat, Mulltum per omnem etc.  Verum est, quod ad promissiones factas attinet, melio- res sunt ludei: sed si respicias iusticiam, et Iudaei et Gen- tes destituuntur gloria Dei, id est, omnes sunt pctonres-  \pend\pstart Id est, in primo et secundo Capit-ubi non tan- tum caussas ostendit, uerumetiam dixit, ipsos excusationem non habere, sic enim seipsum interpretatur Paulus.   \pend
\phantomsection
\addcontentsline{toc}{subsection}{\textit{Ante. }}
\subsection*{\textit{Ante. }}
\section*{COMMENT. POMER.  }\pstart Hactenus ergo hanc propositionem tractauit Paulus, probabit et confirmabit eam iam quoque Scripturis cer- tissimis, ne qua elabi possimus: atque ex hoc disces, qui sit ordo praedicandi Euangelii. Primum, per le- gem sunt hominum peccata detegen ' a, ut sibiipsis in- cipiant displicere, et desperare in suis uiribus. Dein- de subiicienda est consolatio, Christum scilicet esse iustificatorem nostrum, remitti nobis peccata, quia ille placeat patri. Nemo utitur pharmacis, nisi mor- bum ante cognorit. Sic quomodo non intelliges Deum misericordem esse, nisi uideas te filium irae et damna- tionis esse? Euangelium reuelat cogitationes cordium, id est aperit, quanta sit in nobis iniquitas, ut tandem confugiamus ad Christum, adblandientem nobis, quem admodum mater suo filto blandiri solet. Sic Christus quando praedicabat, primum dicebat, Poenitentiam agite, agnoscite uos peccatores, et peruentet in uos regnum Dei. Quare prudenter Paulus praemisit dan- nationem nostri, priusquam iustitiam per Christum describat. Sed diceret Caro, Nunquid uerum est, om- nes homines esse peccatores? nonne multi sunt boni et deuoti (ut illi loquuntur) in Monasteriis patres? Haec praeuidens Paulus, ne putent esse somnium hu- manum, profert manifestas Scripturas contra eos, qui bus tamen uix conscientiae nostrae conuinci possunt, usq adeo sumus excaecati.   \pend
\phantomsection
\addcontentsline{toc}{subsection}{\textit{Omnes peccato. }}
\subsection*{\textit{Omnes peccato. }}\pstart  \pend
\vspace{0.5cm}\noindent
\vspace{0.2cm}\rule{1cm}{0.2pt}\\ 
\hspace{0.2cm}\textit{mg}
\footnotesize Ordo praedi- candi.  
\normalsize| 
\section*{IN EPIST. AD ROM.  }
\marginpar{[ p.19 ]}\pstart Scilicet in Prophetis. q. d. non est mea sententia.  Sciolis hoc imposuit, ut omnia congererent in Psalm.  13. cum tamen in eo non legantur apud Hebr. Sed in nouo Testamento id frequens est, ut citentur in uno lo- co ex diuersis scriptoribus sententiae, quod uides in principio Marci. Sicut scriptum est in Prophetis. In hoc loco quaedam ex Esaia, pleraque ex Psalmis de- prompta sunt, nonnulla ex Prouerbiis.   \pend
\phantomsection
\addcontentsline{toc}{subsection}{\textit{Sicut scriptum est. }}
\subsection*{\textit{Sicut scriptum est. }}
\phantomsection
\addcontentsline{toc}{subsection}{\textit{Non est iustus. }}
\subsection*{\textit{Non est iustus. }}\pstart Scilicet, in terra, quia Psalm. dixerat, Deus prospexit e coelo, num esset iustus quispiam: sed uidit omnes iniustos, ergo egent iusticia Dei. Quidam di- cunt exceptum Christum, per hoc quod dicitur, vsque  ad unum, sed haec non est sententia.   \pend
\phantomsection
\addcontentsline{toc}{subsection}{\textit{Intelligat. }}
\subsection*{\textit{Intelligat. }}\pstart Id est, cognoscat Deum, uel discernat bonum a malo, Esaiae. 5.Vae qui dicitis malum, bonum.  \pend
\phantomsection
\addcontentsline{toc}{subsection}{\textit{EXQVIRAT. Diligat Deum. }}
\subsection*{\textit{EXQVIRAT. Diligat Deum. }}
\phantomsection
\addcontentsline{toc}{subsection}{\textit{Omnes deflexerunt. }}
\subsection*{\textit{Omnes deflexerunt. }}\pstart Semper dicit, omnes, an non hoc intelligant nostri Iusticiarii. Esa. Omnes nos quasi oues errauimus, unusquisque declinauit in uiam suam, neglecta uia Dei.  \pend
\phantomsection
\addcontentsline{toc}{subsection}{\textit{Simul. }}
\subsection*{\textit{Simul. }}\pstart Id est, omnes inutiles sunt. Inutile est, quod non  \pend
\vspace{0.5cm}\noindent
\vspace{0.2cm}\rule{1cm}{0.2pt}\\ 
\hspace{0.2cm}\textit{mg}
\footnotesize Omnes ege- nus tusticia Dei- 
\normalsize| 
\section*{COMMENT. POMER.  }\pstart consequitur eum finem, ad quem creatum est, id est, uon apprehendunt salutem.   \pend
\phantomsection
\addcontentsline{toc}{subsection}{\textit{Exerceat. }}
\subsection*{\textit{Exerceat. }}\pstart Id est, faciat bonum, id est, non sunt candidi, Aliud in corde, aliud in ore habent. Aliud foris simae lant, aliud intus cogitant.   \pend
\phantomsection
\addcontentsline{toc}{subsection}{\textit{Sepulchrum. }}
\subsection*{\textit{Sepulchrum. }}\pstart Id est, nihil aliud proferunt, quam mendacia, ma- ledicta, foetorem illum subinde exhalant ex ore: nam cadauer est intus in corde, cuius foetor per os ascen- dit: ut in Euangelio dicitur, Ex abundantia cordis, os loquitur. Item, Quomodo potestis bona loqui, cum sitis mali, ut sequitur.   \pend
\phantomsection
\addcontentsline{toc}{subsection}{\textit{Ad dolum. }}
\subsection*{\textit{Ad dolum. }}\pstart Id est, ad decipiendum: non solum in rebus exter- nis id faciunt, sed et in iis quae salutis sunt: uo- lunt docere uiam salutis, et docent uiam perditio- nis. Dolum hunc optima specie tegunt, uocant enim opera bona, satisfactoria, interim nesciunt quid do- ceant, etiam conscientia eorum teste.   \pend
\phantomsection
\addcontentsline{toc}{subsection}{\textit{Venenum aspidum. }}
\subsection*{\textit{Venenum aspidum. }}\pstart Vide quomodo connectat Scripturam. Tale est il- lud uenenum, ut nulla medicina illud tollat.   \pend
\phantomsection
\addcontentsline{toc}{subsection}{\textit{Sub labiis. }}
\subsection*{\textit{Sub labiis. }}
\section*{IN EPIST. AD ROM.  30 }
\marginpar{[ p.]}\pstart Scilicet, eorum qui perdunt animas. Videntur optima docere, cum interim seducant: optima pro- raittunt, sed uenenum immittunt, non lac et mel.   \pend
\phantomsection
\addcontentsline{toc}{subsection}{\textit{Execratione. }}
\subsection*{\textit{Execratione. }}
\phantomsection
\addcontentsline{toc}{subsection}{\textit{Amarulentia. }}
\subsection*{\textit{Amarulentia. }}\pstart Id est, in rebus quae ad proximum attinent. Haec omnia est uidere in regno Papae, ubt regnat Excom- municatio, Aggrauatio, Reaggrauatio. Illa enim imperiosissime semper intonant et clamant Nul- la hic est Verbi Dei contio, praeterea male loquuntur de proximo.  \pend
\phantomsection
\addcontentsline{toc}{subsection}{\textit{Pedes }}
\subsection*{\textit{Pedes }}\pstart Sunt adfectus, id est, ad hoc-nituntur et cur- runt, ut effundant sanguinem, nulla prorsus in eis est commiseratio et misericordia, ad perdendum sunt efficaces, non sanandum. Videamus id in nostris san- ctis Papistis hodie, an non uelint Euangelii Praedi- catores, occidi, cremari, et perdi? hanc impieta- tem hodie Deus per Euangelium reuelat: uerum, quia traditi sunt in reprobum sensum, id non uident, uellent illi nos perditos uno momento. Exemplum habes de Isaac et ismaele in Galatis: non uocat hic Scri- ptura, uiros sanguinum, homicidus et latrones:  \pend\pstart  \pend
\vspace{0.5cm}\noindent
\vspace{0.2cm}\rule{1cm}{0.2pt}\\ 
\hspace{0.2cm}\textit{mg}
\footnotesize Pedes sunt af- fectus.  
\normalsize| 
\section*{TOMMENT. POMER.  }\pstart fed Hypocritas et huiusmodi sanctos Sathanae, qua - les omnes sumus a natura, nisi renascamur per Spi- ritum sanctum. Isaac non persequitur Ismaelem, Chri stiam non male praecantur suis aduersariis, sed pro- eis orant . Illi contra, non possunt non velle occidere, et quanto sunt sanctiores, tanto magis id cupiunt, et tantum abest, ut se in hoc peccare putent, ut etian- se obsequium Deo praestare arburentur . Cum hypo- crisi ergo est negocium Scripturae, non cum crassis peccatoribus, hos enim mundus iudicat.   \pend\pstart Confractio ad mentem refertur, sumptum ab iis rebus quas conterimus.   \pend
\phantomsection
\addcontentsline{toc}{subsection}{\textit{Contritio. }}
\subsection*{\textit{Contritio. }}
\phantomsection
\addcontentsline{toc}{subsection}{\textit{Calamitas in uiis. }}
\subsection*{\textit{Calamitas in uiis. }}\pstart Miseria est in omnibus iis quae ipsi tentant, a- gunt, et inueniunt, nihil aliud excogitant, quam uexationes conscientiarum, ut in regno Papae exper- ti sumus.   \pend\pstart Vides duplicem esse viam, impiorum, et pioru.  Confer haec cum praecedentibus, et intelliges quae sit illa contritio et calamitas. Via impiorum est, qua- se et alios misere perdunt. Via pacis, qua solum Deum agnoscimus, patrem, cuius misericordia pac antur nostrae conscientiae.   \pend
\phantomsection
\addcontentsline{toc}{subsection}{\textit{Viam pacis. }}
\subsection*{\textit{Viam pacis. }}
\vspace{0.5cm}\noindent
\vspace{0.2cm}\rule{1cm}{0.2pt}\\ 
\hspace{0.2cm}\textit{mg}
\footnotesize Christian pro aduersa riis orare de bent.  
\normalsize| 
\hspace{0.2cm}\textit{mg}
\footnotesize Duplex uia.  
\normalsize| 
\section*{IN EPIST. AD ROM.  31 }
\marginpar{[ p.]}\pstart Haec est radix omnium malorum. Non timent De- um, quid igitur facerent boni? Timent tantum, ne tol- lantur ab eis dignitates, census etc.ut olim Iudaet, Ne forte ueniant Romani et c. Talia omnes timemus, et non Deum, nisi per spiritum Dei eximantur. Haec omnia ex lege naturae nobis aperta esse deberent, ni- si per impietatem obtenebratum esset insciens cor- nostrum.   \pend
\phantomsection
\addcontentsline{toc}{subsection}{\textit{Non est timor . }}
\subsection*{\textit{Non est timor . }}
\phantomsection
\addcontentsline{toc}{subsection}{\textit{Scimus aute, quod quaecunque . }}
\subsection*{\textit{Scimus aute, quod quaecunque . }}\pstart Alia est occupatio contra ludaeos et sanctulos, namubi talia legunt (quemadmodum quidam hodie, qui maxime Euangelici dici uolunt) non putant ea de- se scripta esse, neque ad se pertinere, sed oscitanter transeunt omnia, quasi scriptura tantum de crassis pec- catis loquatur, manent ita in suis peccatis, iuxta cor- durum et poenitere nescium, quae interim illis non im- putantur qui resipiscunt. Itaque Iudaei qui se populum sanctum, regale sacerdotium, in Scripturis appella- ri sciebant, cum de excaecatis et reprobis legerent, (nempe, quod non sit qui requirat Deum etc.) omnia ad Gentes referebant: ita interpretabantur, Non est iustus, scilicet, inter Gentes, imo ne unus quidem. Quare Paulus hic eorum conscientiam tan git. Item obserua, legem hic dici totam sacram Scri- pturam Veteris Testamentt. Nam ex Esaia et Psalm.   \pend
\vspace{0.5cm}\noindent
\vspace{0.2cm}\rule{1cm}{0.2pt}\\ 
\hspace{0.2cm}\textit{mg}
\footnotesize Timor hypo- critarum.  
\normalsize| 
\hspace{0.2cm}\textit{mg}
\footnotesize Sinistre Iudaei interpretan- tur scripturas 
\normalsize| 
\section*{COMMENT. POMER.  }\pstart quaedam citarat, quae tamen non sunt lex Most- HIS. Scilicet, Iudaeis.   \pend
\phantomsection
\addcontentsline{toc}{subsection}{\textit{Qui in lege sunt, dicat. }}
\subsection*{\textit{Qui in lege sunt, dicat. }}\pstart Quasi diceret. Vos dicitis ea pertinere ad Gentes: atqui Scriptura Iudaeis, non Gentibus est edita, quare eam ludaei, non Gentes, legunt: siquidem, non fecit ta- liter omnt nationi. Et alibi Paulus ex Deuter. Non alligabis os boui trituranti: non propter boues lex scri- pta est, sed propter nos. Sic hic quoque , ut uides, ora- cula ludaeis prodita, leguntur, non Gentibus. Si igi- tur tales scripturas de Iudaeis scriptas conuincimus, cur de Gentibus interpretabuntur? quandoquidem Gen- tes non negant se esse peccatores, sentiunt se abiisse post deos alienos, et adorasse Idola, quare non opus habent lege scripta. At Iudaei illa opus habent, quia non possunt agnoscere suam hypocrisin.   \pend
\phantomsection
\addcontentsline{toc}{subsection}{\textit{Os obturetur. }}
\subsection*{\textit{Os obturetur. }}\pstart Gentium os obturatum est, non audent mutire.  Item Iudaeorum per Scripturam.  \pend
\phantomsection
\addcontentsline{toc}{subsection}{\textit{Totus mundus. }}
\subsection*{\textit{Totus mundus. }}\pstart Non solum Gentes, sed et iudaei, id est, ut om - nes coram Deo peccatores esse inueniantur.   \pend
\phantomsection
\addcontentsline{toc}{subsection}{\textit{Omnis caro. }}
\subsection*{\textit{Omnis caro. }}
\section*{IN EPIST. AD ROM.  58 }
\marginpar{[ p.]}\pstart Hic nemo excipitur, nullus homo. At quidam in- ter Gentes sancte et honeste uixerunt, ut Socrates, Camillus, Plato. Item inter Iudaeos, quidam se iu - stificabant operibus legis? Respondet Paulus. To- tus mundus, omnis caro est in peccato. Si itaque illi ex operibus diuinae legis, non sunt iustificati, quid tribuemus nostris traditionibus humanis?  \pend
\phantomsection
\addcontentsline{toc}{subsection}{\textit{In conspectu eius. }}
\subsection*{\textit{In conspectu eius. }}\pstart Id est, Dei. Tametsi fingant sibi aliquid iustitiae co- ram mundo, tamen coram Deo sunt peccatores. Ope- ra legis sunt, quae fiunt absque adfectu cordis uolunta- rio ubi iustificati sumus, et facti Christiani: bona opera quidem facimus, sed non tam magna et tam multa, ut per ea quidquam iustitiae consequamur, id quod etiam fieri non potest. Itaque nulla opera neque  ante, neque post fidem ad salutem prosunt, quia nun quam certi reddimur ex illis de nostra iusticia.   \pend
\phantomsection
\addcontentsline{toc}{subsection}{\textit{Per legem. }}
\subsection*{\textit{Per legem. }}\pstart Alia est occupatio. Ergo inutilis est lex? eur igitur Deus eam dedit? Simile est argumen- tum, Sol et Luna, item stellae, non iustificant, ergo sunt inutiles. Vide ineptias Iusticiariorum, qui quid opponant nesciunt, neque quod stulte op- ponunt Christo, intelligunt . Hic erit opus helle- boro ineptissimis hominibus: quia bona opera non  \pend
\vspace{0.5cm}\noindent
\vspace{0.2cm}\rule{1cm}{0.2pt}\\ 
\hspace{0.2cm}\textit{mg}
\footnotesize Nibil iu stitiae consequimur nostris operi- bus.  
\normalsize| 
\section*{COMMENT. POMER.  }\pstart iustificant, ideo inutilia erunt? Maximum sane erat beneficium legis, quod nobis errorem nostrum et cae citatem indicarit, non in ea requirimus iusticiam, quam dare nequit, sed peccati reuelationem, cor nostram erat obtenebratum et obscuratum. Contra, lex ignis quidam, quo illabente in cor nostrum, peccatum re- uelatum est, id erat in lege requirendum, et non aliud-  \pend
\phantomsection
\addcontentsline{toc}{subsection}{\textit{Nunc uero absque lege iustitia. }}
\subsection*{\textit{Nunc uero absque lege iustitia. }}\pstart Hactenus primam Verbi Dei partem, nempe le- gem tractavit Apostolus, nunc secundam partem ad greditur, scilicet, per solam fidem in Christum esse iustuciam . At hic est Epistolae slatus : et quis non sunt mo desyderio ista scire cuperet, cum omnes iusticiam quaeramus, eam quae Dei est, id est, quae per Deum in nos uenit, non qua Deus iustus est. MANI- FESTATA. Scilicet, per Euangelium.   \pend
\phantomsection
\addcontentsline{toc}{subsection}{\textit{Dum comprobatur. }}
\subsection*{\textit{Dum comprobatur. }}\pstart Id est, testimonium accipit. Haec necessario Pau- lus dixit, ut nobis iterum Scripturam commendaret contra novos Prophetas, qui hodie contemnunt Scri- pturam, et dicere audent. Quid mihi eum Veteri Te- stamento? spiritu opus est, qui me omnia doceat. Re cte quidem. Sed interim miseri non uident, ea quae docet Spiritus sanctus, esse confirmata, et iam olim praedicta in Mose et Prophetis: nam quod ho- die docemus, non est noua doctrina. An Christus  \pend
\vspace{0.5cm}\noindent
\vspace{0.2cm}\rule{1cm}{0.2pt}\\ 
\hspace{0.2cm}\textit{mg}
\footnotesize Vera insticia est per solam fidem in Chri- stum.  
\normalsize| 
\hspace{0.2cm}\textit{mg}
\footnotesize Doctrina spi- ritus sancti quae? 
\normalsize| 
\section*{EPIST. AD ROM.  33 }
\marginpar{[ p.]}\pstart quid aliud est, q̃ quod Adae, postea Abrabae et patri- bus promissam est? Ideo dixerat Paulus, Manifesta- ta, id est, non noua facta, uerum semper fuit, licet latuerit nos.   \pend
\phantomsection
\addcontentsline{toc}{subsection}{\textit{Iusticia Dei per fidem Christi. }}
\subsection*{\textit{Iusticia Dei per fidem Christi. }}\pstart Diligenter omnia exponit Paulus. Somniarant Iu- daei sibi quandam iusticiam, sed non respiciebant in mediatorem Christum, de quo erat dictum, In semi- ne tuo benedicentur omnes gentes. Item, Semen con- teret etc. Legebant illa quidem in Scrip. at non in- telligebant, credebant in Deum (ut putabant) sed non per Christum. Iohan. 14. dicit Christus, Creditis in Deum, et in me credite. Haec quoq sapientibus no- stris opponenda sunt, qui tantum ex Deo uolunt quae- rere iusticiam, Christum abiicientes, qui est is, per quem a Deo in nos uenit omnis iusticia.  Sola iusticia per Christum.   \pend
\phantomsection
\addcontentsline{toc}{subsection}{\textit{In omnes. }}
\subsection*{\textit{In omnes. }}\pstart Id est, quae pertinet ad omnes, super omnes. Co pia tantum est uerborum.   \pend
\phantomsection
\addcontentsline{toc}{subsection}{\textit{Non est enim distinctio. }}
\subsection*{\textit{Non est enim distinctio. }}\pstart Id est, siue iudaei sint, siue Gentes. Ecce Deus ma nifeste aequat illis crassos peccatores, qui in speciem sunt sanctissimi. q. d. Non est Iusticiarius melior co- ram Deo, quam latro, qui ad supplicium ducitur.  At quis hominum secundum naturam, sic se potest abii- cere, ut dicat se non meliorem tali, qui crasso delicto  \pend
\vspace{0.5cm}\noindent
\vspace{0.2cm}\rule{1cm}{0.2pt}\\ 
\hspace{0.2cm}\textit{mg}
\footnotesize 
\normalsize| 
\section*{COMMENT. POMER.  }\pstart infamisest. Itaque cum simus natura omnes filii irae, eripimur per solam gratiam Christi, in regnum lu- cis translati . Quicunque damnantur, propter increduli- tatem damnantur, quicunq saluantur, propter fiden- in solum Christum saluantur, nullo personae respectu.  \pend\pstart Id est, abalienati sunt a gloria Dei, carent illa, id est, non glorificant Deum, et Iudaei et Gentes sunt per naturam impii. Quid ergo habebunt Iusticiarii iusticiae, cum careant primo praecepto? quo docetur, quomodo Deus debeat glorificari. Quomodo colent Deum, quem non nouerunt esse suum Deum?  \pend
\phantomsection
\addcontentsline{toc}{subsection}{\textit{Destituuntur. }}
\subsection*{\textit{Destituuntur. }}\pstart Et Gentes et Iudaei gratis, id est, sine operibus le- gis, et quid illud sit, gratis, exponit.   \pend
\phantomsection
\addcontentsline{toc}{subsection}{\textit{Iustificantur. }}
\subsection*{\textit{Iustificantur. }}\pstart (Inquit) scilicet, Dei gratiam et fauorem, quia nobis bene uoluit propter suum Christum. Sic om- nes patres sancti, iustificatisunt a mundo condito, quia scilicet credebant in Christum uenturum (qui dile - ctus erat patri) quemadmodum hodie omnes iustifi- cantur, credentes in Christum iam exhibitum mundo.   \pend
\phantomsection
\addcontentsline{toc}{subsection}{\textit{Per illius. }}
\subsection*{\textit{Per illius. }}
\phantomsection
\addcontentsline{toc}{subsection}{\textit{Per redemptionem, }}
\subsection*{\textit{Per redemptionem, }}
\vspace{0.5cm}\noindent
\vspace{0.2cm}\rule{1cm}{0.2pt}\\ 
\hspace{0.2cm}\textit{mg}
\footnotesize salutis et d nationis cau- sae. 
\normalsize| 
\section*{IN EPIST. AD ROM.  }
\marginpar{[ p.34 ]}\pstart Fauor Dei primum est, quo saluamur. Alterum, quod Christus satisfecerit pro peccatis nostris, ut ui1 deas quam diligenter Apostolus exponat uocem, gra- tis, non omittens circunstantias, quibus depinga- gratiam, rhetorico more. Vnde est gratia? a Deo- Per quem? per Christum. Quibus contingit? creden- tibus, sine personarum respectu . Per quid? per san- guinem Iesu Christi. Quare? ut scilicet sic reuelare- tur iusticia Dei, et nostra iniusticia, nam nihil ac- ceptum est patri, nisi per illum. Hoc non uident lu- dei, sua opera tantum magnifaciunt, quod est se- ctari iusticiam legis, uerum ad iusticiam non perue- nire: hanc iusticiam humanam arguet Spiritus san- ctus suo aduentu, ut Christus dicit, ut sic reuele- tur Dei iusticia, et nostra omnia damnentur, ma- neatque solus Christus nostra iusticia. Sed quamuis haec diligenter inculcentur, scilicet, fauor Del, et Chri- sti satisfactio, remurmurat semper conscientia no- stra infirma. Tamen oportet esse satisfactionem pro peccatis, nec quiescit, nisi fuerit satisfactum Hinc finxerunt satisfactiones pro peccatis, quibus uidebantur conscientiae ad tempus pacificatae, sed fru- stra. Nam postea non desierunt peccata mordere con scientiam. Ergo, nunquam illis somniis potuerunt con- scientiae certificari. Certificantur autem sola et unica Christi satisfactione et redemptione. Beati qui baec intelligunt et credunt.   \pend
\vspace{0.5cm}\noindent
\vspace{0.2cm}\rule{1cm}{0.2pt}\\ 
\hspace{0.2cm}\textit{mg}
\footnotesize Satisfactionum humanarum origo origo. 
\normalsize| 
\section*{COMMENT. POMER. }
\phantomsection
\addcontentsline{toc}{subsection}{\textit{Proposuit Deus reconciliatorem. }}
\subsection*{\textit{Proposuit Deus reconciliatorem. }}\pstart Id est, mediatorem inter Deum et homines. Ecce quid opera faciunt? Solus Christus est reconciliator, ut maxime sexcenties sis Chartusianus, sine hoc Chri sto, non eris Deo acceptus. Rectius tamen legeris, propiciatorium, et erit eadem sententia, nisi quod per hoc in sanctissimam illam figuram respicere licet, nempe in propiciatorium arcae, super quod promise- rat se Deus sessurum, et exauditurum.   \pend\pstart Quem. Scilicet, Christum.  \pend\pstart Id est, dum credimus in illum. Ista est solutio eius, quod impii opponunt. Si Christus est mediator et re Sanguis Christi (inquit) in se non estinefficax, ve- rum pro amnium peccatis satisfecut, sed quod non om- rum pro omnium peccatis satisfecit, fed quod non om nes saluentur, in caussa est impietas: nemini enim est utilis et salutaris sanguis Christi, nisi credentibus: qui bibunt hunc sanguinem Christi, id est, credunt pro ipsorum peccatis fusum, illi salui sunt. Sed Tur- cae cum id negligunt, perduntur. Vbi non est fides, ibi necessario est, omnis maledictio, damnatio et pec catum .   \pend
\phantomsection
\addcontentsline{toc}{subsection}{\textit{Per fidem. }}
\subsection*{\textit{Per fidem. }}
\phantomsection
\addcontentsline{toc}{subsection}{\textit{Interueniente iplius sanguine. }}
\subsection*{\textit{Interueniente iplius sanguine. }}\pstart Hoc est, illud precium redemptionis, hic sanguis  \pend
\phantomsection
\addcontentsline{toc}{subsection}{\textit{}}
\subsection*{\textit{}}
\vspace{0.5cm}\noindent
\vspace{0.2cm}\rule{1cm}{0.2pt}\\ 
\hspace{0.2cm}\textit{mg}
\footnotesize Solus Chri- stus mediatos inter Deum e homines. 
\normalsize| 
\hspace{0.2cm}\textit{mg}
\footnotesize Cur Turca non saluentur. 
\normalsize| 
\section*{IN EPIST. AD ROM. 35 }
\marginpar{[ p.]}\pstart uerus, cuius figura erat sanguis uitulorum in Veteri Testamento, utuides in uerbis Christi, cum Sacra mentum hoc instituit. Sanguis ille Veteris Testamen- tt, non auferebat peccata, uerum quicunque crede- bant alium, scilicet, Christi sanguinem per illum signi ficari, illi non in sanguine uitulorum, sed Christi sal uabantur, qui effundendus expectabatur.  \pend
\phantomsection
\addcontentsline{toc}{subsection}{\textit{Ad ostensionem iusticiae suae. }}
\subsection*{\textit{Ad ostensionem iusticiae suae. }}\pstart Omnia sunt talia, quae non tantum non oscitan- ter legenda sunt, uerum summa diligentia. q.d.  Hactenus fuit quaedam iusticia mundo inculcata per Philosophos et Iusticiarios, sed est ficta, inefficax coram Deo, ideoque longe praecellentiore opus est, nempe Dei. Conati sunt passim Philosophi tot libris instituere rectam uitam, ut homines pii essent et iu - sti, uerum id non potuerunt praestare-Sic item Iudaei neglectis Dei promissionibus tantum in externa ope- ra oculos defigebant, quae Deus non ideo praecepe- rat, ut iustificarent, sed ut essent signa tantum inter nae fidei, et thesauri Verbi Dei in corde absconditi, Hoc est, quod Deus in Prophetis adeo damnet et exc cretur, etiam opera legis per ipsum institutae: exige bat fidem, non opera. Illi tamen (quia caeci sumus na tura in iis, quae Deus exigit a nobis) tantum opera spectantes, uoluntatem Dei negligebant: fuit ergo omnis illa iusticia, ficta iusticia. Ideoque mala, quemad Externa ope- ra, non iusti- ficant.  \pend
\vspace{0.5cm}\noindent
\vspace{0.2cm}\rule{1cm}{0.2pt}\\ 
\hspace{0.2cm}\textit{mg}
\footnotesize 
\normalsize| 
\section*{COMMENT. POMER.  }\pstart modum et iudicium erat malum, ob id, quia sapientes mundi, et tusti secundum carnem ludaei, non potuerunt di- cere, se errare, opus erat, ut ueniret Spiritus san- ctus, qui acutius cerneret et iudicaret, diceretque , uos Phulosophi, stulte scribitis de iusticia. Vos ludaei, caeci estis in iusticia operum. Omnis iusticia uestra, quantumuis bona in oculis uestris, est blasphemia Dei, et irrisio. Itaq manifestata est Dei iusticia, ut hu- mana omnis iusticia confunderetur, uidereturque ; ni- bil aliud esse, quam stercus et pannus menstruatae mulieris coram Deo, quantumuis oculo carnis ni- teat et adblandiatur.  \pend
\phantomsection
\addcontentsline{toc}{subsection}{\textit{Propter remissionem. }}
\subsection*{\textit{Propter remissionem. }}\pstart Id est, ut remitterentur praeterita peccata. Textus sic cohaereat. Quem proposuit Deus reconcillatorem propter remissionem etc.   \pend
\phantomsection
\addcontentsline{toc}{subsection}{\textit{Praeteritorum peccatorum. }}
\subsection*{\textit{Praeteritorum peccatorum. }}\pstart Hic ineptissime quidam de praeteritis peccatis lo- quitntur. Sunt enim, qui se in omnes formas uertunt calumniandi Scripturam, etiam contra conscientiam.  Alii dicunt, Christus pro omnibus peccatis ante illum factis, satissecit . Alii ea tantum quae erant ante bap- tismum in nobis, expiavit sua pabione . Caetera vero post baptismum facta, nobis per nostra bona opera reliquit abstergenda. O impias nugas, hoc est, foe- dare humano stercore Christi faciem, et dicere Scri-  \pend
\vspace{0.5cm}\noindent
\vspace{0.2cm}\rule{1cm}{0.2pt}\\ 
\hspace{0.2cm}\textit{mg}
\footnotesize Praeterita tan tum peccata remitti.  
\normalsize| 
\section*{36 IN EPIST. AD ROM.  }
\marginpar{[ p.]}\pstart pturam mendacem. Verum, non sunt alia peccata ni£ si praeterita, quae remittantur. Id quod uel ratio ui- dere potest . Cuius tu crassum accipe exemplum.  Si mihi centum aureos deberes, et non esses soluen- do, remitto igitur tibi debitum. Dic quaeso, an non ti- bi praeteritum debitum remittam? non profecto futu- rum, hoc inquam dimitto, quod mihi debes, non hoc quod nondum debes. Ita Deus quae alia peccata remitteret, quam praeterita, et quae habes? Hinc oramus, Dimitte nobis debita nostra, scilicet, quae iam nos premunt, nunq sumus sine peccatis, ideo semper oramus, ut dimittantur. Sed aduerte et hic stultissimam stulticiam Papisticam, quae primum nostris temporibus inualuit, olim non ita ineptiebant, ut de futuris peccatis indulgentias promitterent, id quod hodie faciunt Papistae, dant Bullas et diplomata qui bus a futuris absoluant peccatis.O homines impuden- tissimos, qui quid adfirment nesciunt. Deus solus re- mittit peccata, ideo ut non pecces, non ut porro pec- ces. Quae enim haec esset remissio? uolo liber esse a peccatis, et tamen manere in peccatis. Vt stulte fece- rit is, qui uehementer cupiat e carcere liberari, et ta- men manere in carcere, sic is quoq extremae insaniae accusandus, qui in hoc bullas acceperit. Deus quia fragilis es, tibi semper remittit peccatum, modo tibi displiceat peccatum.   \pend
\phantomsection
\addcontentsline{toc}{subsection}{\textit{Quae Deus tolerauit. }}
\subsection*{\textit{Quae Deus tolerauit. }}
\vspace{0.5cm}\noindent
\vspace{0.2cm}\rule{1cm}{0.2pt}\\ 
\hspace{0.2cm}\textit{mg}
\footnotesize Bullae pro fu- turis pecca- tis. 
\normalsize| 
\section*{COMMENT. POMER.  }\pstart Non statim uindicauit, quia erat nos olim purga- turus sanguine Christi . Inuitauit nos ad resipiscen- tiam Cap. 2. Sed quare?  \pend\pstart Id est, ut palam faceret mundo suam iusticiam.  Hic iterum posset oggamnire caro, Ergo peccabimus? fed habet suam solutionem supra, quorum damnatio tusta est: non solum autem tolerauit Deus, sed remi- sit quoque , ut palam fieret, quod ex sola eius misericor dia saluemur, quae, inquam, iam per Euangelium prae- dicatum, ostensa est, et olim descriptain Hiere. 31.  Peccatorum eorum amplius non recordabor. Haec iusticia non solum a lege et Prophetis habet testimo- nium, uerum et a cordibus nostris, in quibus Spiritus Dei illud testatur, etiam nolentibus nobis.   \pend
\phantomsection
\addcontentsline{toc}{subsection}{\textit{Ad ostendendam: }}
\subsection*{\textit{Ad ostendendam: }}\pstart Scilicet, tempore gratiae et reuelati Euangelii.  \pend
\phantomsection
\addcontentsline{toc}{subsection}{\textit{In praesenti. }}
\subsection*{\textit{In praesenti. }}\pstart Est enim ipse fons iusticiae uere, imo ipsissima iu- sticia, ipse solus iustus erat, cum omnes essemus dan- nati. Quare intellige, ut sit iustus, id est, declaretur esse iustus, secundum illud, Vt iustificeris, id est, iustus declareris hominibus.   \pend
\phantomsection
\addcontentsline{toc}{subsection}{\textit{Vt ipsesit iustus. }}
\subsection*{\textit{Vt ipsesit iustus. }}
\phantomsection
\addcontentsline{toc}{subsection}{\textit{Iustificans eum. }}
\subsection*{\textit{Iustificans eum. }}\pstart Ex hoc siquidem fonte omnes haurimus iusticiam,  \pend
\vspace{0.5cm}\noindent
\vspace{0.2cm}\rule{1cm}{0.2pt}\\ 
\hspace{0.2cm}\textit{mg}
\footnotesize Misericordia reuelatur per Euangelium.  
\normalsize| 
\section*{IN EPIST. AD ROM.  37 }
\marginpar{[ p.]}\pstart iuxta dictum Esaiae. 12. Haurietis aquam de fontibus saluatoris.   \pend\pstart Id est, qui credit in Christum, est Paulina phrasis.   \pend
\phantomsection
\addcontentsline{toc}{subsection}{\textit{Eum qui est. }}
\subsection*{\textit{Eum qui est. }}
\phantomsection
\addcontentsline{toc}{subsection}{\textit{Vbi igitur gloriatio.  }}
\subsection*{\textit{Vbi igitur gloriatio.  }}\pstart Hic omnino carnem confundit, aufert enim omnem gloriam, cum a Gentibus, tum a Iudaeis, non potest- tam gloriari Ethnicus de sua ciuili iusticia, neque Iu- daeus de operibus legis, nihil in te inuenies, quo te- praepares ad gratiam, ut illi somniarunt. Ideo respon- det Paulus  \pend
\phantomsection
\addcontentsline{toc}{subsection}{\textit{Exclusa est. }}
\subsection*{\textit{Exclusa est. }}\pstart Id est, nulla est, siue huc, siue illuc respicias.   \pend
\phantomsection
\addcontentsline{toc}{subsection}{\textit{Per quam legem. }}
\subsection*{\textit{Per quam legem. }}\pstart Id est, quo iure est exclusa. OPERVM.  Id est, quia tu opera quaedam sectaris, quae Deus in lege praecepit.  \pend
\phantomsection
\addcontentsline{toc}{subsection}{\textit{Per legem fidei. }}
\subsection*{\textit{Per legem fidei. }}\pstart Improprie sic dicit Apostolus, id est, per spiritum quo credimus in Iesum Christum. Sicut iusticiarii se- adstringunt ad sua opera, sic pii ad fidem in Chri- stum, sola enim fides efficit, ut possimus gloriari de Deo, non nostra opera. Haec facit me esse unum cum eo, omnia ex illo accipio, ergo ego nihil habeo de quo glorter, hoc hypocritae non uident neque possunt.  \pend
\section*{COMMENT. POMER.  }\pstart Nos dicimus, quia ante probauimus et Gentes et Iudaeos esse peccatores, idque etiam teste Scriptura. Item intulimus, neminem saluari, nisi per tusticiam Dei in Christum. Concludimus iam, nostra opera nihil esse coram Deo, sed ipsum tantum fidem cordium respice re. Cuius exemplum habemus latronem in cruce, qui cum socio suo etiam blasphemus erat, tamen paulo post cognouit peccatum suum, nihil fecit bonorum operum. Imo ne potuit quidem, cum manibus et pe- dibus suffixus penderet. Sed dixit tantum, Memento mei etc. quae sunt uerba fidei admirandae. Et sicut il- le est saluatus, sic nos saluari oportet. Hic non fuit opus, facere bona opera, per quae impetraret salutem, cum iam fide saluatus esset, quod Verba Christi, ipsi respondentis testantur.   \pend
\phantomsection
\addcontentsline{toc}{subsection}{\textit{Arbitramur. }}
\subsection*{\textit{Arbitramur. }}\pstart Non est Deus tantum Iudaeorum, qui uolunt glo- riari de lege. Illa ipsa opera legis sunt peccata, qua- propter, iustificabit et Gentes et Iudaeos gratis Deus, cum omnes natura sint filii irae.   \pend
\phantomsection
\addcontentsline{toc}{subsection}{\textit{An ludaeorum. }}
\subsection*{\textit{An ludaeorum. }}
\phantomsection
\addcontentsline{toc}{subsection}{\textit{Circuncisionem. }}
\subsection*{\textit{Circuncisionem. }}\pstart Id est, Iudaeos ex fide, et non ex operibus legis.  \pend
\phantomsection
\addcontentsline{toc}{subsection}{\textit{Legemigitur. }}
\subsection*{\textit{Legemigitur. }}\pstart Occupatio est. Tu Paule abiicis legem Dei? Re-  \pend
\vspace{0.5cm}\noindent
\vspace{0.2cm}\rule{1cm}{0.2pt}\\ 
\hspace{0.2cm}\textit{mg}
\footnotesize Fide fola sal. namur : 
\normalsize| 
\section*{IN EPIST. AD ROM.  }
\marginpar{[ p.38 ]}\pstart pondet. Non abiicio, imo stabilio eam, et doceo quomodo ea sit adimplenda, Non enim operibus, sed fide adimpletur. Lex dicit, Dilige Deum. Verum, hoc nemo potest secundum naturam, ergo lex inuti- liter docetur, si non praedicetur fides, unde discemus agnoscere nostram impossibilitatem. Per fidem Deus lar- gitur suum spiritum, qui diffusus in corda nostra, sup- peditat nobis utres, ut ultro faciamus legem, nemo enim satis diligit sine spiritu fidet, quare frustra lex praedicatur et inculcatur. Hoc si nostri uidissent, non tantam carnificinam bonorum operum, et sancte et religiose uiuendi in Monasteriis instituissent.   \pend
\endnumbering\beginnumbering\section{CAPA IIII.}\pstart \huge\textbf{H}\normalsize IC diligenter lector superiora obseruet, ne I I Pauli ordine negligenter animaduerso, et se quentia non capiat. Confirmattonis hic locus est, nam propositionem, scilicet, solam fidem in Christum iusti ficare, hic munit certis argumentis, quae ordine tudebimus.   \pend
\phantomsection
\addcontentsline{toc}{subsection}{\textit{Quid igitur? }}
\subsection*{\textit{Quid igitur? }}\pstart Iudaei et Iusticiarii audientes, ex fide tan- tum nos iustificari, obiiciunt. Num bona ope- ra nihil sunt, certe Abraham pater noster bona opera fecit, putas illa Deum non respexisse?  \pend
\vspace{0.5cm}\noindent
\vspace{0.2cm}\rule{1cm}{0.2pt}\\ 
\hspace{0.2cm}\textit{mg}
\footnotesize Lex fidem im- pletur.  
\normalsize| 
\section*{COMMENT. POMER.  }\pstart His Paulus respondet septem argumentis, quae et in plura discerpi possent, in quibus uidebis clarius pro- prietatem et circumstantias fidet et operum per collationem.  \pend\pstart Primum est argumentum, Abraham ex fide iu- stificatus est, non ex operibus, ergo et nos ex fi- de iustificamur, non ex operibus, nam est omnino par ratio, nemo aliter iustificatur . Hoc erat necesse Iudaeis probare, qui de lege digladiabantur.  \pend
\phantomsection
\addcontentsline{toc}{subsection}{\textit{Nam si Abraham. }}
\subsection*{\textit{Nam si Abraham. }}\pstart Quasi aliquid meruerit, ueluti mercenarius quis piam gloriatur de suo praemio. Itaque si ex operibus est tustificatus, certum est illam iusticiam tantum humanam fuisse, quia opera erant humana, gloriabitur igitur de ea iustitia tantum coram hominibus, non coram Deo.   \pend
\phantomsection
\addcontentsline{toc}{subsection}{\textit{Fsabet quod glorietur. }}
\subsection*{\textit{Fsabet quod glorietur. }}
\phantomsection
\addcontentsline{toc}{subsection}{\textit{Credidit Abraham. }}
\subsection*{\textit{Credidit Abraham. }}\pstart Fuit homo peccator ut nos, et iustificatus est, non quia non peccauerit, aut quod satisfecerit pro pecca- tis, sed quod crediderit promittenti Deo.   \pend\pstart Qui opus aliquod facit, mercedem a me accipit, Non dici potest, quod illi a me gratia contingat, sed quia operatus est, do mercedem debitam.   \pend
\phantomsection
\addcontentsline{toc}{subsection}{\textit{Qui operatur, }}
\subsection*{\textit{Qui operatur, }}
\phantomsection
\addcontentsline{toc}{subsection}{\textit{Qui iustificat. }}
\subsection*{\textit{Qui iustificat. }}\pstart Periphrasis est Dei, id est, propterea is est iustus,  \pend
\vspace{0.5cm}\noindent
\vspace{0.2cm}\rule{1cm}{0.2pt}\\ 
\hspace{0.2cm}\textit{mg}
\footnotesize Argument, iustificationi 
\normalsize| 
\section*{IN EPIST. AD ROM.  }
\marginpar{[ p.39 ]}\pstart quia credit iustificanti Deo, ut uideas solam gratiam et misericordiam, meritum omnino excluditur: nec est ullum, si tamen aliquod fingitur, non saluamur ex gratia. Idque est primum argumentum, nos fide iustificari.  \pend
\phantomsection
\addcontentsline{toc}{subsection}{\textit{Quemadmodum et Dauid. }}
\subsection*{\textit{Quemadmodum et Dauid. }}\pstart Secundum est argumentum. Iusticia est per im- putationem.   \pend
\phantomsection
\addcontentsline{toc}{subsection}{\textit{Beatificationem hominis. }}
\subsection*{\textit{Beatificationem hominis. }}\pstart Id est, qua homo fit iustus et beatus.   \pend\pstart Dicit hunc beatum, quem Deus iustificat, non qui facit bona opera. Acceptat nos pro iustis, tametsi secundum naturam simus iniusti: sicut ei dono debitum, qui non est soluendo, is iam liber est a debito, non quia soluerit, sed quia solutionem imputo, alioqui si non im- putassem, mansisset debitor. Sic non ex nobis, sed ex Deo tusti sumus. Audis ergo hic solam gratiam.   \pend
\phantomsection
\addcontentsline{toc}{subsection}{\textit{Cui Deus. }}
\subsection*{\textit{Cui Deus. }}
\phantomsection
\addcontentsline{toc}{subsection}{\textit{Obtecta sunt. }}
\subsection*{\textit{Obtecta sunt. }}\pstart Scilicet a Deo. Non dicit, qui non habent, Deus non imputat peccatum, sed iusticiam, ut est in Hiere, Et peccatorum eorum, non memor ero amplius.   \pend
\phantomsection
\addcontentsline{toc}{subsection}{\textit{Dicimus. }}
\subsection*{\textit{Dicimus. }}\pstart Scilicet ex scriptura, non ego. Tertium est ar- gumentum, Abraham ante circuncisionem est tustifi-  \pend
\section*{COMMENT. POMER.  }\pstart catus, ergo et nos ante omnia opera iustificamur.  Circuncisio primum est opus legis, sine quo reliqua commode fiert nequeunt, qualis et apud nos est bap- tismus. Praeterea est tantum signaculum iusticiae (non ipsa iusticia) quae est per fidem. Vide quae in Genesi de iis dicuntur, ubi non pauci anni intercesserunt in- ter hanc historiam fidet, et inter circuncisionem, Nam priusquam Isaac ex Sara, et Ismaelem ex Agar susciperet, de ipso dixerat Scriptura. Credidit Abra- ham, et reputatum etc. Ideoque ; cum circuncisio pri- mum sit opus legis, et tamen ante eam iustificatus sit Abraham, quid est, in quo haereamus? Ex hoc loco po- tes deprehendere errorem Iacobi Epistolae, in qua im- pium uides argumentum, praeterquam quod ridicule colligit. Citat Scripturam contra Scripturam, quod Spiritus sanctus ferre nequit, non potest igitur reli- quis adnumerari libris, qui tusticiam fidei praedi- cant. Epistola ad Hebraeos, longe dignius tractat baec exempla.   \pend
\phantomsection
\addcontentsline{toc}{subsection}{\textit{Non in circuncisione. }}
\subsection*{\textit{Non in circuncisione. }}\pstart Sententia est, ista beatificat io, qua Deus iustificat, non solum ad Iudaeos pertinet, sed ad Gentes, imo maxime ad Gentes, quia Abraham fuit iustificatus in praeputio, non in circuncisione. Itaque propter fidem in Dei promissionem, non propter opera est iustifica- tus. Haec constantissme semper opponenda sunt tusti  \pend
\vspace{0.5cm}\noindent
\vspace{0.2cm}\rule{1cm}{0.2pt}\\ 
\hspace{0.2cm}\textit{mg}
\footnotesize Empie conha S. Tacoh een 
\normalsize| 
\section*{IN EPIST. AD ROM.  40 }
\marginpar{[ p.]}\pstart ciariis Papisticis, ut confundantur cum suis traditioni bus, nam si non est iusticia ulla ex operibus diuinae le- gis, ut euidentissime Paulus probat, multo minus ex traditionibus humanis.   \pend
\phantomsection
\addcontentsline{toc}{subsection}{\textit{Et signum. }}
\subsection*{\textit{Et signum. }}\pstart Occupatio est, Ad quid ergo erat utilis circuncisio Abrahae, si ex ea non fuit tustificatus? Habes hic insig- nem locum, de signis externis, quales fuerunt fere olim omnes ceremoniae Iudaeorum, et hodie nostra Sacramenta, quae non iustificant per se, sicut nec uer- bum externe auditum per se, sed sunt certissima sig- na internae iusticiae-Sic Abraham accepit signum il- lud externum in praecisione pelliculae suae, in signum illius iusticiae, quam intus habebat in corde. Sgnifica- batur ipsi Abrahae tantum illo signo iusticia uera, sed per se non erat iusticia, Item per illud agnoscebatur populus Dei. Verum quia huiusmodi signa saepe fal- lere possunt, fit ut hypocritas plerunque faciant. Sed tunc signa dici possunt, non autem signa Dei. Quan quam nec signa quidem uere sunt, si nihil intus sit, quod uere significent, ut uideamus Sacramenta Dei idem efficere, quod uerbum externum, loquun- tur enim nobis, quasi Verbum Dei, testificantia bo- siam uoluntatem Dei erga nos, ut non dicam, quod sine Verbo Dei nihil essent: fide fasceptum  \pend\pstart  \pend
\vspace{0.5cm}\noindent
\vspace{0.2cm}\rule{1cm}{0.2pt}\\ 
\hspace{0.2cm}\textit{mg}
\footnotesize Stguorum usus. 
\normalsize| 
\section*{COMMENT. POMER.  }\pstart Verbum Dei iustificat, et fide suscepta sacramenta iu stificant. Nam fide iustificamur, fides autem est ex auditu, auditus autem per Verbum Dei. Sunt itaque  Sacramenta sigilla et signa et pactum, inter nos pec catores et beneuolentem Deum, ex manu enim eius ea accipimus . Secundaris autem per talia significa- mus aliis, ut nos pro fidelibus habeant . Signum et pactum inter me et Deus (si pro tali hoc habeam) non me fallit : Signum autem coram hominibus tantum saepe fallit, et hypocrisis est. Sic ubi peccata mor- dent nostras conscientias, sacramentum Eucharistiae accipimus, ut certificemus eas de remissione peccato rum, quae tantum est ex fide in Christum. Ex qua tan dem fide, fluunt bona opera, quibus etiam me ipsum de mea fide certifico. Nam alians ea non facerem, nsi ad talia Spiritussanctus me impelleret.  \pend
\phantomsection
\addcontentsline{toc}{subsection}{\textit{Quae fuerat in praeputio. }}
\subsection*{\textit{Quae fuerat in praeputio. }}\pstart Hic habes quoque locum, contra quem Noui Spiri- tus ne mutire quidem audent, pro paruulis, quos bap- tisamus hodie, Deus mandarat, ut octauo die mascu- lus circuncideretur, at circuncisio ex sese non erat fi- des aut iusticia, sed signum, quo significabatur fides in- tus esse. Siitaque fidem rationi humanae acceptamre- fers, quae tamen donum Dei est, dic mihi quomodo paruuli habuerint rationem, qui accipiebant circum cisionem signum internae iufticiae? hos profecto Deus  \pend
\vspace{0.5cm}\noindent
\vspace{0.2cm}\rule{1cm}{0.2pt}\\ 
\hspace{0.2cm}\textit{mg}
\footnotesize Sacramentu cur accipien- dum.  
\normalsize| 
\hspace{0.2cm}\textit{mg}
\footnotesize Loeus contra Nouos Spiri- tus, de bap- tismo paruu- lorum. 
\normalsize| 
\section*{IN EPIST. AD ROM.  43 }
\marginpar{[ p.]}\pstart rensebat in suum populum, cum diceret, Masculus cuius praeputium non fuerit circuncisum, eius anima peribit ex populo suo. Ita et nostri paruuli, ratio- nem non habent, tamen baptisantur, et signum ac- cipiunt fidei. Deus haec omnia facit sua gratia, noli haec metiri tua ratione. Nos offerimus Deo paruulos et animas paruulorum, non porcos, pro quibus paruu- lis, sanguis Christi effusus est. Noli igitur contra gra- tiam disputare, Christus uult ad se uenire paruulos, et dicit, Illorum est regnum coelorum. Garriant interim illi quicquid uelint, nos credemus, non inquiremums ratione humana.   \pend
\phantomsection
\addcontentsline{toc}{subsection}{\textit{Credentium. }}
\subsection*{\textit{Credentium. }}\pstart Id est, pater est omnium qui credunt diuinis pro missionibus.   \pend
\phantomsection
\addcontentsline{toc}{subsection}{\textit{lis qui. }}
\subsection*{\textit{lis qui. }}\pstart Exponit se, est pater credentium Iudaeorum et Gentium. Abraham enim non est pater Iudaeorum co- ram Deo, propter circuncisionem, neque Gentium pro- pter praeputium, sed utrorumque propter fidem.   \pend
\phantomsection
\addcontentsline{toc}{subsection}{\textit{Quae. }}
\subsection*{\textit{Quae. }}\pstart Scilicet, fides fuit in Abraham, antequam circun- cideretur. Est ergo fidei pater, non operum legis.   \pend
\phantomsection
\addcontentsline{toc}{subsection}{\textit{Non enim per legem. }}
\subsection*{\textit{Non enim per legem. }}\pstart Hoc argumento et pluribus, utitur in Galatis, Abraham iustificatus est ex promissione, et non ex  \pend
\vspace{0.5cm}\noindent
\vspace{0.2cm}\rule{1cm}{0.2pt}\\ 
\hspace{0.2cm}\textit{mg}
\footnotesize IIII. Argu- mentum.  
\normalsize| 
\section*{ }\pstart                              \pend
\vspace{0.5cm}\noindent
\vspace{0.2cm}\rule{1cm}{0.2pt}\\ 
\hspace{0.2cm}\textit{mg}
\footnotesize 
\normalsize| 
\section*{  }
\marginpar{[ p.]}
\phantomsection
\addcontentsline{toc}{subsection}{\textit{ }}
\subsection*{\textit{ }}\pstart    \pend
\phantomsection
\addcontentsline{toc}{subsection}{\textit{ }}
\subsection*{\textit{ }}\pstart    \pend
\phantomsection
\addcontentsline{toc}{subsection}{\textit{ }}
\subsection*{\textit{ }}\pstart              \pend
\phantomsection
\addcontentsline{toc}{subsection}{\textit{ }}
\subsection*{\textit{ }}\pstart   \pend
\phantomsection
\addcontentsline{toc}{subsection}{\textit{ }}
\subsection*{\textit{ }}
\phantomsection
\addcontentsline{toc}{subsection}{\textit{ }}
\subsection*{\textit{ }}
\phantomsection
\addcontentsline{toc}{subsection}{\textit{ }}
\subsection*{\textit{ }}\pstart   \pend
\vspace{0.5cm}\noindent
\vspace{0.2cm}\rule{1cm}{0.2pt}\\ 
\hspace{0.2cm}\textit{mg}
\footnotesize   
\normalsize| 
\section*{COMMENT. POMER.  }\pstart Deum esse mendacem, imo et stultum, si promitteret et dare uellet eam iusticiam, quam ante haberemus.  Est autem locus argumenti, ex absurdo desumptus.   \pend\pstart Iterum probat, fieri non posse, ut per legem iusti- ficemur, tantum abest, ut uel ullo modo nos Deo con ciliare possit. Imo ubi suum officium in nobis execu- ta fuerit ipsa, terret conscientias nostras, et excitat in nobis iram et indignationem contra Deum, uellemus tum in eo horrore et terrore, legem non esse, imo ip- sum Deum non esse Deum. Est hoc honorare Deum et seruare legem? facit ea, ut uideamus nostri dan- nationem, ubi hoc fit, fiert non potest, ut natura non irascatur et murmuret insolentissime contra Deum.  Sic fit, ut grautus fiat peccatum per legem, prouocat nos contra Deum, quod in septimo Cap. declaratur.  Maximus et superbissimus sit hypocrita necesse est, qui hoc in corde suo non uidet. Lex ergo magis aggra- uat conscientiam nostram, quam liberet uel iustificet.   \pend
\phantomsection
\addcontentsline{toc}{subsection}{\textit{Nam lex iram. }}
\subsection*{\textit{Nam lex iram. }}\pstart Non dicit, ibi nec est peccatum, quod diligenter obserua, ut maxime non transgrediaris legem externe, tamen peccato non cares. Ante legis reuelationem, nec peccatum, nec transgressionem uidemus. Qui per- ignorantiam peccant, tametsi non careant peccato, carent tamen transgressione, quam non uident. Putant  \pend
\phantomsection
\addcontentsline{toc}{subsection}{\textit{Siquidem. }}
\subsection*{\textit{Siquidem. }}
\vspace{0.5cm}\noindent
\vspace{0.2cm}\rule{1cm}{0.2pt}\\ 
\hspace{0.2cm}\textit{mg}
\footnotesize VI. Argu- mentum. 
\normalsize| 
\hspace{0.2cm}\textit{mg}
\footnotesize Legis opus.  
\normalsize| 
\section*{IN EPIST . AD ROM .  }
\marginpar{[ p.42 ]}\pstart enim se bonos et iustos esse uiros, coram Deo tamen sunt transgressores legis, id est, ipsam non implent ex corde, ut maxime coram mundo non accusentur impietatis. Ergo transgressio ad agnoscentes peccatum pertinet, non ad eos qui ignorant, sic Paulus dicit, Vbi lex uenit, mortuus sum.  \pend\pstart Irata Deo conscientia, non glorificat Deum, postq̃ se uidet abiici reuelato peccato. Nam in inferno quis confitebitur tibi etc? Ita augetur peccatum et fit ma nifesta transgressio. Si itaq iusticia esset ex operibus, nemo posset pacare conscientiam suam. At quia Deus gratis dat, non possumus esse incerti de iusticia no- stra, unde et conscientia tranquilla fit . Vchemens est hoc argumentum, nemo illud contemnat-  \pend
\phantomsection
\addcontentsline{toc}{subsection}{\textit{Idcirco ex fide. }}
\subsection*{\textit{Idcirco ex fide. }}
\phantomsection
\addcontentsline{toc}{subsection}{\textit{Vt firma sit promissio uniuer. so semini. }}
\subsection*{\textit{Vt firma sit promissio uniuer. so semini. }}
\phantomsection
\addcontentsline{toc}{subsection}{\textit{Vt secundum gratiam. Id est, gratis sit . Ad quid? }}
\subsection*{\textit{Vt secundum gratiam. Id est, gratis sit . Ad quid? }}\pstart Id est, nobis omnibus, ne quis nostrum dubitet de sua salute. Nam cum in manu Dei est, certi de ea sumus. Si autem in nostra esset, essemus miserrimae creaturae, siquidem de ea nunquam essemus certi- Quod contra illos est, qui dicunt nos non debere esse certos de salute, quo nihil aliud uolunt, quam non es se credendum, ac ita perierit plane omnis spes salu-  \pend
\vspace{0.5cm}\noindent
\vspace{0.2cm}\rule{1cm}{0.2pt}\\ 
\hspace{0.2cm}\textit{mg}
\footnotesize VII. Argu- nentum. 
\normalsize| 
\hspace{0.2cm}\textit{mg}
\footnotesize Certi sumus de salute .  
\normalsize| 
\section*{}\pstart  \pend\pstart  \pend\pstart  \pend\pstart  \pend\pstart  \pend
\phantomsection
\addcontentsline{toc}{subsection}{\textit{}}
\subsection*{\textit{}}
\phantomsection
\addcontentsline{toc}{subsection}{\textit{}}
\subsection*{\textit{}}
\phantomsection
\addcontentsline{toc}{subsection}{\textit{}}
\subsection*{\textit{}}
\phantomsection
\addcontentsline{toc}{subsection}{\textit{}}
\subsection*{\textit{}}
\phantomsection
\addcontentsline{toc}{subsection}{\textit{}}
\subsection*{\textit{}}
\vspace{0.5cm}\noindent
\vspace{0.2cm}\rule{1cm}{0.2pt}\\ 
\hspace{0.2cm}\textit{mg}
\footnotesize 
\normalsize| 
\section*{44 IN EPIST. AD ROM.  }
\marginpar{[ p.]}\pstart Coram Deo erat transferendum. Erasmus sic in- tellexit: quemadmodum Deus est communis pater om- tium, ita et Abrahamum communem omnium patrem constituit. Verum licet haec non mala, tamen alia est sententia hoc loco. Constitui te coram Deo patrem, id est, in regno inuisibili uel spirituali, quod est reg- num fidei, in quo tantum Deus est rex. Estque ; sim- plex quaedam Occupatiuncula, ne quis dicat: Quo modo est pater Gentium? Respondet: Non secundum carnem, sed per spiritum uel fidem est eorum pater.   \pend
\phantomsection
\addcontentsline{toc}{subsection}{\textit{QVI VITAE. Vel ad uitam.  }}
\subsection*{\textit{QVI VITAE. Vel ad uitam.  }}\pstart Scilicet ad se, ut adpareant. Nemo negat solum Deum suscitare eos, qui mortui sunt. Similiter uocauit omnes creaturas, cum diceret: Fiat lux etc. Vocauit inquam, cum uon essent, quasi fuissent, ut adfisterent et obedirent uocanti Deo. At Paulus hic longius respe- xit. Commendat enim nobis gratiam, quod Deus eos, ed est, Gentes qui perierant, fecit filios Abrahae, quod item uidemus in historia Lucae 15. de filio perdito nobis significari, quem pater iam mortuum putabat. Mortuae erant Gentes, id est, ad Idololatriam defecerant, ut taccan- de peccato Adae, Deus tamen resuscitauit eas, et re- uocauit ad uerum cultum, in quo facti sunt filii Abra- hae. Nemini igitur desperandum est, et si omnia desperata uideantur. Suscitat Deus mortuos, et quae nonsint, uocat tanam sint. Hoc facit Verbo suo. Gentium uo- catio.   \pend
\phantomsection
\addcontentsline{toc}{subsection}{\textit{Vocat. }}
\subsection*{\textit{Vocat. }}
\vspace{0.5cm}\noindent
\vspace{0.2cm}\rule{1cm}{0.2pt}\\ 
\hspace{0.2cm}\textit{mg}
\footnotesize 
\normalsize| 
\section*{COMMENT. POMER.  }\pstart filii sumus Dei per Verbum, domini omnium creatu- rarum per Verbum, cohaeredes Christi per Verbum, atque adeo nihil est, quod uirtute illius nobis non im- pertiatur. Lapides eramus, iam autem, ut Iohannes dixit, filii Det ex lapidibus. Haec tam insignis gratiae commendatio, omnes uires, omne Liberum Arb. no- strum tollit. Quid enim potest mortuus, qui nihil est? Quid emortua uulua? quid Isaac immolandus? quae omnia uides in historia Abrahae. Fides est, quae concepit et peperit Isaac, fides immolabat illum, ex quo erat posteritas promissa propaganda. Pudeat nos nostrae impietatis, qui tam aperta habemus ex- empla fidei.   \pend
\phantomsection
\addcontentsline{toc}{subsection}{\textit{Qui praeter spem. }}
\subsection*{\textit{Qui praeter spem. }}\pstart Haec est fidei natura, unde et finitio haberi po- test certissima, credere in spe, praeter spem, tam con- stanter nos credere oportet, contra omnem naturae sensum, ita, ut si Deus ignem frigidum dicat esse, credam sine controuersia, tametsi secundum naturam sit impossibile. Muliebria sustinere Sara non potuit, tamen praeter spem, Abraham credidit se ex ea suscep- turum Isaac.   \pend
\phantomsection
\addcontentsline{toc}{subsection}{\textit{Sic erit semen. }}
\subsection*{\textit{Sic erit semen. }}
\phantomsection
\addcontentsline{toc}{subsection}{\textit{Id est, innumerabile, ut stellae coeli.  Non infirmatus. Sed constanter credens.  }}
\subsection*{\textit{Id est, innumerabile, ut stellae coeli.  Non infirmatus. Sed constanter credens.  }}
\vspace{0.5cm}\noindent
\vspace{0.2cm}\rule{1cm}{0.2pt}\\ 
\hspace{0.2cm}\textit{mg}
\footnotesize Liber. Arb 
\normalsize| 
\section*{IN EPIST. AD ROM.  }
\marginpar{[ p.45 ]}\pstart Desierant iam illi menstrua fluere, id quod magis mirandumm est, quam Abrahae senectus.   \pend
\phantomsection
\addcontentsline{toc}{subsection}{\textit{Emortuam uuluam. }}
\subsection*{\textit{Emortuam uuluam. }}
\phantomsection
\addcontentsline{toc}{subsection}{\textit{Ob incredulitatem.  }}
\subsection*{\textit{Ob incredulitatem.  }}\pstart Non disceptabat secum in animo, forte ita est, for- te non est, sicut nostra conscientia, uerbo non firma- ta. Ah hoc est peccatum, hoc non est peccatum, non potest esse disceptatio, ubi radicem fixit promissio-  \pend
\phantomsection
\addcontentsline{toc}{subsection}{\textit{Non haesitabat.  }}
\subsection*{\textit{Non haesitabat.  }}
\phantomsection
\addcontentsline{toc}{subsection}{\textit{Tribuens gloriam. }}
\subsection*{\textit{Tribuens gloriam. }}\pstart Credere omni uerbo Dei, est dare gloriam Deo, ut maxime id sit contra omnem sensum naturae. Non enim est apud Deum impossibile omne Verbum, ut ait Angelus in Luca ad Mariam. Hanc ergo solam glo- riam a nobis exigit, et praeter eam nihil. Quare in nullis tribulationibus et angustiis est nobis desperan- dum, neque in tentationibus spiritualibus, neque in pec- catis. Esto, nobis omnia adpareant grauia et horren- da, demus gloriam Deo, quod uerax sit, quod promi- serit remissionen peccatorum. Noli, quaeso, O ho- mo, tua peccata maiora facere, diuinis promissioni- bus, alioqui eum mendacii argues, et desperabis ut Iudas, ut maxime sis scortator, homicida etc. Vi- de ne peius pecces desperando, quando tibi reuela-  \pend
\vspace{0.5cm}\noindent
\vspace{0.2cm}\rule{1cm}{0.2pt}\\ 
\hspace{0.2cm}\textit{mg}
\footnotesize Fideiuis et amplitudo.  
\normalsize| 
\section*{COMMENT. POMER.  }\pstart tur peccatum. Deus tibi peccatum non imputabit, mo- do ipsum ueracem esse credas. Quod tamen interim Sathan tibi non permittit, facit peccata esse multo maiora Dei promissionibus. Hic erit fortiter confisten- dum, et cum hoc hoste congrediendum: uide ut illius conteras robur. Quod tum fiet, si praeter omnem spem, praeter omnem sensum, credas Deo, qui te eripiet ex inferno suo tempore, non sinet te absorberi undis aquarum, tibi imminentium. Sic in omni necessitate externa, etiam penuria uentris, sentias oportet. Est omnipotens Deus per suum Verbum, quicquid promit- tit, hoc praestat, modo nos clamare ex corde posse- mus: Dimitte nobis debita etc. Sed quo minus hanc gloriam resignemus Deo, nos remoratur caro.   \pend
\phantomsection
\addcontentsline{toc}{subsection}{\textit{Certa persuasione. }}
\subsection*{\textit{Certa persuasione. }}\pstart Contra eos hoc dictum est, qui dicunt nos non cer- tos esse posse de nostra salute, satis esse bonam inten- tionem, quae tamen nulla est coram Deo. Petri inten- tio bona erat, cum diceret. Si oportuerit me mori te cum, non te negabo. Non equidem in ulio quod no- strum est, neque in bona nostra intentione consistit no- stra salus, sed in Det promissione. In hac sit nostra intentio et certitudo, hanc a nobis requirit gloriam, potens est eam praestare. Non indigebis, ut interces soribus utaris, Georgio, Christophoro, Barbara, aut diua Catharina. QVAPROPTER. Idest, quia dedit hunc honorem ucrutatis Deo.   \pend
\vspace{0.5cm}\noindent
\vspace{0.2cm}\rule{1cm}{0.2pt}\\ 
\hspace{0.2cm}\textit{mg}
\footnotesize Bona intentio. 
\normalsize| 
\section*{IN EPIST. AD ROM.  46 }
\marginpar{[ p.]}
\phantomsection
\addcontentsline{toc}{subsection}{\textit{Nonscriptum est. }}
\subsection*{\textit{Nonscriptum est. }}\pstart Hic inculcat diligenter nobis, quod omnia ( quod ad iustificationem tantum attinet) de Abraham dicta, ettam de nobis dicantur. Non dicit de operibus Abrabae, sed iustificatione, qua per Christum est iustificatus, ne quis diceret, ut infirmae solent conscientiae: Hoc de Abraham scribitur, ad me nihil attinet. De omni- bus dicitur: Beatus uir, cui non imputauit Dominus peccatum, non tantum de Abraham. Est igitur par- ratio iustificandi ex fide, omnium et Abrahae. Itaque  omnes ex fide, sine operibus iustificabimur.   \pend
\phantomsection
\addcontentsline{toc}{subsection}{\textit{Illi. }}
\subsection*{\textit{Illi. }}
\phantomsection
\addcontentsline{toc}{subsection}{\textit{Scilicet, ad iusticiam quod credidit. QVIBVS.  Scilicet nobis. IMPVTABITVR. Scilicet, ad iusticiam. IN EVM. Scilicet Deum patrem.  }}
\subsection*{\textit{Scilicet, ad iusticiam quod credidit. QVIBVS.  Scilicet nobis. IMPVTABITVR. Scilicet, ad iusticiam. IN EVM. Scilicet Deum patrem.  }}
\phantomsection
\addcontentsline{toc}{subsection}{\textit{Qui excitauit. }}
\subsection*{\textit{Qui excitauit. }}\pstart His concludit, Christum esse mediatorem, et iusti- ciam nostram per fidem.   \pend
\phantomsection
\addcontentsline{toc}{subsection}{\textit{Traditus. }}
\subsection*{\textit{Traditus. }}\pstart Nam non pepercit et Deus, tradidit illum, ut sa- tisfaceret pro peccatis nostris.   \pend
\phantomsection
\addcontentsline{toc}{subsection}{\textit{Excitatus. }}
\subsection*{\textit{Excitatus. }}\pstart Scilicet, ut nos maneremus in iusticia et nouita te uitae. Quidam hic differentiam suspicantur, quod non dicit, excitatus est propter peccata nostra, seu- remissionem peccatorumScriptura sic loqui solet: Estque   \pend\pstart  \pend
\section*{COMMENT. POMER.  }\pstart eadem sententia. Saepe tantum de morte loquitur, et tacet resurrectionem, nos tamen eam intelligimus, et rursus resurrectionem praedicat, ubi nos et mortem eius intelligimus. Alterum ex altero intelligitur. Qui enim liberaret nos a morte, si ipse non resurgeret? aut qui resurgeret, si non mortuus fuisset? Ideo mortuus est propter peccata nostra, ut excitaretur in nostri iu stificationem. Itaque hactenus uidimus comprobatum scripturis et testimoniis ex lege, Christum solum es se nostram iusticiam, nos uero natura damnatos, et filios irae-  \pend
\endnumbering\beginnumbering\section{CAP.V.}\pstart P  \pend\pstart Raefatur ab initio Capitis, in totum Caput, di- cens, Iustificati per fidem, pacem habemus- Visa est autem hactenus Epistola difficilis, quo et mul- ti terrebantur, uerum si ordo eius obseruatus fuerit, nihil inerit difficultatis. In superioribus docuit, tan- tum ex fide esse iusticiam, nunc consequitur, ex ea fi- de, inostras conscientias habere pacem erga Deum, et praeter id in hoc Capite nihil est. Hunc fructum fi- det si obseruaueris in hoc Cap. non uidebuntur tibi uaria hic tractari. Et ut rem totam compendio tracte- mus, uisum est quaedam praefari. Nonne omnes iusti- ciam quaerimus et per opera, et alia quaedam, in hoc, ut habeamus pacem erga Deum? Hac uia ommes ir-  \pend
\section*{EPIST. AD ROM.  49 }
\marginpar{[ p.]}\pstart cedimus, et olim quando uiuebamus in iusticiis no- stris, confitebamur, ieiunebamus, non ob aliud, ni- si ut auferentur peccata nostra, et conscientiae no- strae pacificarentur. Sed idem accidit nobis, quod isra elitis, ut Cap. nono uides, qui sectantes iusticiam, ad eam non peruenerunt. Putabamus et nos ex ope- ribus nobis futuram iusticiam, uerum decepti sumus: nihil est, ut hic fingas Liber. Arb. iusticiam operum, non consequeris pacem conscientiae, ad tempus falli potes, ad tempus item fingere potes sanctitatem et pacem, sed redeunte conscientia peccati, ubi est tua pax et sanctitas ficta? Redit plane eadem turbatio et inquietudo conscientiae, quam antea senseras, dixe- ras et somniaueras tibi pacem, sed tum uides non esse pacem, quia non est pax impiis (quales sunt omnes qui iusticiam tribuunt operibus, et non fidei in Chri- stum). Est itaque insignis commendatio, iusticiae fidei Fidel iusticia.  in hoc loco. q. d. Paulus. Experire et quaere, in uenies pacem, quicunque credis Christum satisfecisse pro peccatis tuis. Habes magistram ipsam experien tiam in corde tuo, quae te conuincit, tusticiam esse ex fide, etiam renitente ratione tua et carne, coge- ris fateri, et dicere, Ita se res habet. Quid hic plu- ribus testibus eges? noster spiritus, spiritui Dei dat te- stimonium, quod non sit alia iusticia, qua pacem ha- bet conscientia, quam fidei, ubi sentimus et credi- mus peccatorum remissionem per Christum. Iam con-  \pend
\vspace{0.5cm}\noindent
\vspace{0.2cm}\rule{1cm}{0.2pt}\\ 
\hspace{0.2cm}\textit{mg}
\footnotesize 
\normalsize| 
\section*{COMMENT. POMER.  }\pstart scientia exultat et in conspectum Dei, multa securi- tate carmen triumphale canit, hoc scilicet: Christus est mea iusticia, mea fortitudo, deleuit Chirographum peccati, uicit Sathanam principem inferorum, iam nihil est quod timeam, omnes hostes et aduersarii mei, funt uelut stipula, quae statim absumatur igni, et ue- lut puluis, quem difpergit uentus. Non potest quis- quam uerbis consequi, quanta haec sint gaudia, quan ta praedicatio hutus pacis in conscientia, quae uidit et sensit remissionem peccatorum. Experiatur aliquis, in omnibus aliis, ex quibus sperat caro salutem, num- possit habere pacem hanc et laeticiam conscientiae, ex illis: nunq profecto certus esse potuerit, de paterna erga se Dei uoluntate. Itaque uides hic describi nobi- lissimos fidei fructus, qui non corrumpi, nec interire possunt. Primo, quod certus es de gloria Dei, unde et pacem habes. Secundo, ut nequeas a spe excide- re in aduersis, quantumuis sint horrenda, neque mors neque peccatum te prosternet. Eris enim petra fortissima. Tertio, ut ne quid dubites de gratia (qua st indignus sis ) nam indignis, non dignis et san- ctis datur, quia Christus nos dilexit, cum essemus impii. Quarto, conscientia nostra est certa, ut di- cat, maiorem esse gratiam per Christum, quam pec- catum ex Adam in nos propagatum. Si igitur tam magna est condemnatio propter Adam, maiorem scus esse gratiam Dei per Christum. Atque haec est  \pend
\vspace{0.5cm}\noindent
\vspace{0.2cm}\rule{1cm}{0.2pt}\\ 
\hspace{0.2cm}\textit{mg}
\footnotesize Fidei fructus 
\normalsize| 
\section*{48 IN EPIST. AD ROM.  }
\marginpar{[ p.]}\pstart collatio, consolationis plena, Adae et Christi. Reli- quauide inuerbis Pauli.   \pend\pstart Id est, ubi consecuti sumus remissionem peccato- rum per fidem, conscientiae nostrae pacem habent: nam nullum peccatum iam nos turbat. Habemus le- gendum, non habeamus. In hoc loco lapsi sunt Columnae ccclesiae, ut putabantur, Origenes et caeteri, qui dicebant, quaedam moralia hic docere Paulum.  Honc admoneo, ne simus ingrati Deo, pro sua gratia, quam hodie nobis in suo Euangelio dedit. Nam uel pueri hodie id cernunt, imo in plateis clamant, quod tanti Doctores uiderant. Non loquitur Apostolus, de Ciuili quadam pace inter homines, sed de pace erga Deum, quam homines non uident, et qua solae con- scientiae adflictae fruuntur, ut non timeant Deum iu- dicem, sed sciant ipsum esse patrem. Hunc fructum fidei ubi habemus, statim sequuntur bona opera.   \pend
\phantomsection
\addcontentsline{toc}{subsection}{\textit{lustificati. }}
\subsection*{\textit{lustificati. }}
\phantomsection
\addcontentsline{toc}{subsection}{\textit{Per Christum. }}
\subsection*{\textit{Per Christum. }}\pstart Sine hoc mediatore Christo, nihil nobis donatur aDeo. CONTIGIT. Id est, nobis datum est.   \pend\pstart In gratiam hanc. Id est, hoc interpellatore in fauorem patris restituti sumus, ipse adduxit nos ad pa- trem, in cuuus conspectum antea, quia filiiirae, non eramus  \pend
\section*{COMMENT. POMER.  }\pstart ausi prodire, et hoc uerum esse scimus et credimus. Itaque non est iudex, sed pater.  \pend\pstart Scilicet, Dei gratia et fauore, scimus quod De- us nobis bene uelit.   \pend\pstart Id est, nunc regnamus, et sumus reges, quibus omnia subiecta sunt, cum antea iaceremus prostrati in terram, non audentes adspicere in coelum, imo de tenti eramus in carcere, ubi perpetua erat seruitus et calamitas.   \pend
\phantomsection
\addcontentsline{toc}{subsection}{\textit{In qua. }}
\subsection*{\textit{In qua. }}
\phantomsection
\addcontentsline{toc}{subsection}{\textit{Stamus. }}
\subsection*{\textit{Stamus. }}\pstart Verbum est certitudinis, nam nemo gloriatur, ni- si de eare, de qua certus est.   \pend
\phantomsection
\addcontentsline{toc}{subsection}{\textit{Gloriamur. }}
\subsection*{\textit{Gloriamur. }}\pstart Quomodo inquit gloriamur? sub spe? id est, quicquid iam habemus uitae aeternae, aut gloriae, id ha- bemus in spe. Certi sumus, quod Deus nos glorifica- bit. Sumus enim sancti et filii Dei. Huc refer quae in Octauo dicuntur de expectatione creaturarum. Item. I Iohan. 3. Nunc sumus filii Det, inquit, et nondum apparuit, quod futuri sumus etc. Iam contempti su- mus apud omnes, item apud nosmetipsos, quia uide- mus peccata nostra. Sed interim tamen salui sumus perspen-1-d. Res quidem, id est, uita eternaiam  \pend
\phantomsection
\addcontentsline{toc}{subsection}{\textit{Sub spe gloriae. }}
\subsection*{\textit{Sub spe gloriae. }}
\section*{IN EPIST. AD ROM.  49 }
\marginpar{[ p.]}\pstart adest, sed nondum uidetur, nondum est reuelata, Corpus hoc mortale, cum immortalitatem induerit, uidebitur, id quod hodie speramus, ut in Philip . 3.  Glorificabit hoc corpus humile.   \pend\pstart Id est, non solum gloriamur in spe gloriae futurae, et quod simus filii Dei reuelandi: sed etiam certi su- mus per illam pacem fidei, in cordibus nostris, quod nec quaecunque tentatio, non mors, nullum periculum- possit nos superare: in tantum exundat illa pax in cor- dibus nostris, ut aduersitates et tribulationes, alant ac fortiorem reddant eam. Sicut Christi crux et mors uita est et fortitudo.  Christi mors vita est.  \pend
\phantomsection
\addcontentsline{toc}{subsection}{\textit{Nec id solum. }}
\subsection*{\textit{Nec id solum. }}
\phantomsection
\addcontentsline{toc}{subsection}{\textit{De afflictionibus. }}
\subsection*{\textit{De afflictionibus. }}\pstart Gloriamur de cruce, quam Deus nobis proposuit (non de nostris bonis operibus) nam omnibus qui pie uolunt uiuere in Christo, non deerit crux: sicut quan do Sathan tentabit auferre ex corde promissiones Dei, quando sentiemus fidet fragilitatem, et peccatorum onus, iacebimus oppressi, et uidebimus, quod ex no- bis nihil sumus.   \pend
\phantomsection
\addcontentsline{toc}{subsection}{\textit{Scientes quod afflictio. }}
\subsection*{\textit{Scientes quod afflictio. }}\pstart Gradatio est rhetorica, commendatur nobis patien- tia, sed unde esset, nisi affligeremur a Deo. Crux quam Deus non imponit, non est crux: hinc macera- tiones et castigationes Monachorum, non sunt crux-  \pend
\vspace{0.5cm}\noindent
\vspace{0.2cm}\rule{1cm}{0.2pt}\\ 
\hspace{0.2cm}\textit{mg}
\footnotesize 
\normalsize| 
\section*{COMMENT. POMER.  }\pstart Christi, quia non impositae a Deo. Christus calicem quem bibit, a Deo patre habuit, ideo orauit: Aufer a me calicem, quod Monachi non possunt dicere, nam quis eos cogit, ut talia patiantur, est eorum potesta tis ferre, uel non ferre. Sicut nec crux fuit quando Chri stus quadraginta diebus teiunauit, nam erat in sua potestate, ideoque nec orauit clamando ad patrem, si- cut in horto.  \pend\pstart Experimentum, siquidem in patientia ipsa experi- mur, num credamus Deo: dum adfligimur exercetur patientia, in ea certi reddimur apud nos, quod Deus det uirtutem, ut possimus perdurare.   \pend
\phantomsection
\addcontentsline{toc}{subsection}{\textit{Probationem. }}
\subsection*{\textit{Probationem. }}\pstart Cur ultimo loco hic posuit spem, cum eam antea, caussam istorum omnium dixerit, quae uides in textu- Respondeo: Necesse est spem esse, quae fidem conser uet, ne tentationes opprimant. At quod hic dicit, pro- batio spem parit, uult Apostolus, tentationibus cer- tiorem et fortiorem eam fieri, id autem quisque apud se experitur, qui semel atq iterum uicit in arena ad uersarium, is maiore est animo et constantiore, ad plura pericula. Sic pius animus aliquoties tentatus et adflictus, perdurat in aduersis, ne Christum suum amittat, qui eius unica est spes-  \pend
\phantomsection
\addcontentsline{toc}{subsection}{\textit{Spem. }}
\subsection*{\textit{Spem. }}
\phantomsection
\addcontentsline{toc}{subsection}{\textit{Porro spes. }}
\subsection*{\textit{Porro spes. }}
\vspace{0.5cm}\noindent
\vspace{0.2cm}\rule{1cm}{0.2pt}\\ 
\hspace{0.2cm}\textit{mg}
\footnotesize Crux fidei ro- bur.  
\normalsize| 
\section*{IEPIST. AD ROM.  }
\marginpar{[ p.50 ]}\pstart Scilicet illa, de qua dixi, quam per perseueran- tiam adsecuti sumus, non eam dico, quam Hypocritae somniant.  \pend\pstart Scilicet, in conspectu Dei, et in tua conscientia, nam quicquid petieris, speraueris id adsequêre, ut cum Paulo dicere possis: Spe salui facti sumus.   \pend
\phantomsection
\addcontentsline{toc}{subsection}{\textit{Non pudefacit. }}
\subsection*{\textit{Non pudefacit. }}
\phantomsection
\addcontentsline{toc}{subsection}{\textit{Dilectio Dei. }}
\subsection*{\textit{Dilectio Dei. }}\pstart Scilicet, qua diligimus Deum: nam dicit de dile- ctione nobis data.   \pend\pstart Eam emeruit nobis Christus. Estque ; illa, ipse Spiri- tus sanctus. Quon autem illum accipimus? aut quando? eum Verbum praedicatur, id creditum, statim Spiritus san- ctus comitatur, qui hanc ipsam fidem in nobis ope- ratur, qua fit, ut suauissimo, non anxio corde, id prae stemus, quod lex exigit, nempe, ut diligamus Deum ex toto corde, complectamur proximum candido adfectu, Alians enim secundum naturam id non possumus. Nam omnia praecepta legis diuinae, sunt nobis impossibilia- Haec igitur est pax, quam non labefacit ulla aduer- sitas, non expugnatur caducarum rerum dulcedine. Quod si dicas te sentire adhuc peccatum, horro- rem mortis, ut maxime sentias, tamen habes quod statim obiicias. Quid? Remissionem omnium pec- catorum per Christum. Haec ubi firmiter tuo cordi-  \pend
\phantomsection
\addcontentsline{toc}{subsection}{\textit{Effusa est. }}
\subsection*{\textit{Effusa est. }}\pstart  \pend
\section*{COMMENT. POMER.  }\pstart insederit, quid faciet peccatum? nihil te moueat car nis cogitatio, quae tibi dicit, hanc gratiam sancti qui- dem adsecuti sunt, uerum non fuerunt peccatores, ut ego. Et esto, peccatores eos fiusse, tamen non tam magni, quantus ego sum. Respondeo: Tum primum te huiusmodi sentire, quando coeperis agnoscere te pec catorem, haec Sathanas excitat in corde tuo, qui- bus magnifaciat tibi peccatum, ut gratiam illam non agnoscas, nititur te impedire, quo minus dicas, Ma ior est gratia, quam peccatum meum. Ergo non in- dignos nos putemus ad hanc-  \pend
\phantomsection
\addcontentsline{toc}{subsection}{\textit{Quia Christus, cum essemus ad huc infirmi. }}
\subsection*{\textit{Quia Christus, cum essemus ad huc infirmi. }}\pstart Id est, quando eramus natura impii et indigni, ad hanc gratiam, redemit nos. Hic aestimare potes, quid sit Liberum Arbitrium, quae uox contraria est gratiae.   \pend
\phantomsection
\addcontentsline{toc}{subsection}{\textit{luxta temporis rationem. }}
\subsection*{\textit{luxta temporis rationem. }}\pstart Nos in tempore sumus, Deus extra tempus, er- go secundum praedestinationem sumus sancti et bea- ti, quia dilexit nos ante constitutionem mundi in Chri- sto Ephes. I. Verum dum sumus in carne, sumus in tempore, ideoque ; in ipsa impietate, non agnoscimus Deum et gratiam eius. Errantes oues sumus, et ex eo errore non eripimur, nisi ipse tractus ea dilectione qua ab aeterno nos dilexit, nos eripiat. Hic iterum ui-  \pend
\vspace{0.5cm}\noindent
\vspace{0.2cm}\rule{1cm}{0.2pt}\\ 
\hspace{0.2cm}\textit{mg}
\footnotesize Gratiae com mendatio. 
\normalsize| 
\hspace{0.2cm}\textit{mg}
\footnotesize Liberum Arb.  
\normalsize| 
\section*{    }
\marginpar{[ p.]}
\phantomsection
\addcontentsline{toc}{subsection}{\textit{ }}
\subsection*{\textit{ }}\pstart            \pend\pstart            \pend
\phantomsection
\addcontentsline{toc}{subsection}{\textit{ }}
\subsection*{\textit{ }}\pstart    \pend\pstart      \pend
\vspace{0.5cm}\noindent
\vspace{0.2cm}\rule{1cm}{0.2pt}\\ 
\hspace{0.2cm}\textit{mg}
\footnotesize 
\normalsize| 
\section*{COMMENT. POMER.  }
\phantomsection
\addcontentsline{toc}{subsection}{\textit{Christus pro nobis mortuus fuit. }}
\subsection*{\textit{Christus pro nobis mortuus fuit. }}\pstart Vt sic sua morte nos liberaret a peccatis, eramus impii et damnati, et nisi nobis mors Christi subue- nisset, periissemus.   \pend\pstart Est Argumentum a maiori, si suscepit et dilexit nos inimicos et malos, iam magis tuebitur, et dili- get amicos, et sanguine Christi redemptos. Benefi- cus fuit inimicis, quomodo non bene uellet amicis?  \pend
\phantomsection
\addcontentsline{toc}{subsection}{\textit{Multo igitur magis. }}
\subsection*{\textit{Multo igitur magis. }}
\phantomsection
\addcontentsline{toc}{subsection}{\textit{Nam cum inimici. }}
\subsection*{\textit{Nam cum inimici. }}\pstart Scilicet, Christi, quia iam uiuit, et nos iam ereg- no mortis, in regnum eius uitae traducti sumus.   \pend
\phantomsection
\addcontentsline{toc}{subsection}{\textit{Per uitam. }}
\subsection*{\textit{Per uitam. }}
\phantomsection
\addcontentsline{toc}{subsection}{\textit{Non solum autem hoc. }}
\subsection*{\textit{Non solum autem hoc. }}\pstart Non tantum gloriamur sub spe gloriae Dei, et de cruce, qua nos Deus probat, sed est tertia quoq glo- riatio, qua gloriamur, quod Deus noster est: non so- lum gloriamur, quod iste rex Deus det, et daturus sit multa bona, quibus nos ditet. Est hoc omnium ma- ximum, quod is in manu nostrasit, et noster diciue lit. Nos Dei sumus, et ipse cum omnibus suis uult esse noster, idque per Christum Iesum mediatorem. Sed haec exuperant omnem sensum, quae omnia eo perti- rent, ut pacem conscientiis nostris infirmis, confir-  \pend
\vspace{0.5cm}\noindent
\vspace{0.2cm}\rule{1cm}{0.2pt}\\ 
\hspace{0.2cm}\textit{mg}
\footnotesize Gloriatio no- stra.  
\normalsize| 
\section*{INTEPIST. AD ROM.  }
\marginpar{[ p.52 ]}\pstart ment et persuadeant, nam semper in dubium uoca- mur, per cogitationes de peccatis, quibus nihil est pe- stilentius, contra hanc pacem. Ideoque Paulus ait, ma- torem esse gratiam peccato, salutem non in eo consi- stere, quod nos sentimus et cogitamus, sed in fauo- re Dei per Christum, in quo solo obtines illam pacem.  Totus est Paulus in hoc, ut nobis ob oculos ponat uim gratiae et peccati, quam ubi agnouerimus, nihil est quod nos a Deo auocet.   \pend
\phantomsection
\addcontentsline{toc}{subsection}{\textit{Quemadmodum per unum hominem. }}
\subsection*{\textit{Quemadmodum per unum hominem. }}\pstart Obseruabis hic orationis constitutionem, nam non- statim eam absoluit, cum quaedam inspargit, ex qui- bus inclarescat uis gratiae et peccati, per collationem.   \pend
\phantomsection
\addcontentsline{toc}{subsection}{\textit{Per unum. }}
\subsection*{\textit{Per unum. }}\pstart Id est, Adam. Loquitur uero de peccato quod uo- cant, originali, quo omnes damnati sumus, huius effi- catia reliqua peccata crassa emanant, tanquam e fon- te. Rursus si hoc non esset in nobis, non haberemus reliqua omnia, et sicut omnia bona, quae in nobis sunt, debemus iusticiae Christi, sic omnia mala qui- bus sumus obnoxii, debemus peccato originali.   \pend
\phantomsection
\addcontentsline{toc}{subsection}{\textit{In mundum. }}
\subsection*{\textit{In mundum. }}
\phantomsection
\addcontentsline{toc}{subsection}{\textit{Id est, omnes homines.  }}
\subsection*{\textit{Id est, omnes homines.  }}
\vspace{0.5cm}\noindent
\vspace{0.2cm}\rule{1cm}{0.2pt}\\ 
\hspace{0.2cm}\textit{mg}
\footnotesize Originalis noxa.  
\normalsize| 
\section*{COMMENT. POMER.  }
\phantomsection
\addcontentsline{toc}{subsection}{\textit{Peccatum introiit. }}
\subsection*{\textit{Peccatum introiit. }}\pstart Id est, mors, iuxta hoc: In quacunque die comede- ris, id est, peccaueris, morte morieris.   \pend
\phantomsection
\addcontentsline{toc}{subsection}{\textit{In omnes homines. }}
\subsection*{\textit{In omnes homines. }}\pstart Id est, totam Adae posteritatem: propagamur enim omnes ab Adam, et eodem peccato inquinamur, quo ipse propagator.  \pend
\phantomsection
\addcontentsline{toc}{subsection}{\textit{Vsque ad legem. }}
\subsection*{\textit{Vsque ad legem. }}\pstart Scilicet, Mosi: Occupatio est, tamen peccatum non uidebatur, ergo non erat peccatum ullum: non consequitur. Sicut absurdum est et hoc: thesaurus non uidetur, ergo non est, nam forte alicubi in terra de- fossus est, ut uideri nequeat.   \pend
\phantomsection
\addcontentsline{toc}{subsection}{\textit{Peccatum non reputatur. }}
\subsection*{\textit{Peccatum non reputatur. }}\pstart Man acht seyn nicht. Per legem et prohibitio- nem peccatt, uidemus peccatum, hac sublata, tametsi peccatum sit et maneat, non uidetur peccatum. Con feratur hic locus cum superiore, ubi non est lex, ibi nec transgressio, eodem enim pertinet quod hic dici tur, peccatum non habebatur pro peccato ante legem, Sic hodie quoque eramus excaecati, ut etiam amissa cog- nitione naturalis legis, ignorauerimus peccatum.   \pend
\phantomsection
\addcontentsline{toc}{subsection}{\textit{Imo regnauit mors. }}
\subsection*{\textit{Imo regnauit mors. }}\pstart Idest, imperium tenuit, etiam ante legem latam.   \pend
\section*{EPIST. AD ROM.  }
\marginpar{[ p.53 ]}
\phantomsection
\addcontentsline{toc}{subsection}{\textit{In eos quoque qui non peccaue- rant ad similitudinem transgressio- nis Adam. }}
\subsection*{\textit{In eos quoque qui non peccaue- rant ad similitudinem transgressio- nis Adam. }}\pstart Id est, in eos quoque , qui non fecerant hoc peccatum, quod Adam fecit, in arbore uetita. Omnes qui sumus posteri Adae, facti sumus rei peccati Adae, erat is corruptus, et nos propagati sumus ab eo, et illius carne, nihil aliud sumus, quam ea caro peccati, quam Adam habuit: aggrauauit nos ille alieno, id est suo pec- cato, factumque est eius peccatum, nostrum peccatum, quia eadem sumus caro, cum ipsius carne, quae facta est nostra caro. Hic si quis potest sanctulorum, obstru- at os Paulo.  adae peccatii nostrum est-  \pend
\phantomsection
\addcontentsline{toc}{subsection}{\textit{Illius futuri. }}
\subsection*{\textit{Illius futuri. }}\pstart Id est, Christi, qui tum expectabatur, cuius ille habuit quandam figuram.   \pend
\phantomsection
\addcontentsline{toc}{subsection}{\textit{At nonut peccatu, ita et donum. }}
\subsection*{\textit{At nonut peccatu, ita et donum. }}\pstart Id est, Adam habet quidem similitudinem quan- dam Christi, sed longissimo interuallo, interposita dissimilitudine, ipsos disiungente: in hoc quiddam si- mile est, quod Adam sit propagator, et Christus, ue- rum dissimile est in re propagata, ille peccatum, hic iusticiam in nos propagauit, peccatum est nobis in dan- nationem, donum, id est, Spiritussanctus, Gratia, fauor Dei patris, in iusticiam: Sicut Adam damnauit  \pend
\vspace{0.5cm}\noindent
\vspace{0.2cm}\rule{1cm}{0.2pt}\\ 
\hspace{0.2cm}\textit{mg}
\footnotesize 
\normalsize| 
\section*{COMMENT. POMER.  }\pstart nos alieno, id est, suo peecato.Sic Christus nos libe- rauit aliena, id est, sua, non nostra iusticia.   \pend
\phantomsection
\addcontentsline{toc}{subsection}{\textit{Nam si unius delicto. }}
\subsection*{\textit{Nam si unius delicto. }}\pstart Exponit quod iam dixit: efficax sane fuituis pec- cati Adae, in omnium hominum damnationem, sed effi- cacior, iusticia Christi, in omnium hominum libera- tionem. Ideoque nostrae conscientiae respicere debent, non in peccatum, sed in gratiam, quae est multo for- tior et excellentior ipsa fortitudine peccati: quanto homine Deus est praestantior et melior, tanto uali- dior et nobilioris efficaciae gratia, peccato.   \pend
\phantomsection
\addcontentsline{toc}{subsection}{\textit{Gratia Dei. }}
\subsection*{\textit{Gratia Dei. }}\pstart Id est, fauor Dei quo nobis bene uult.   \pend
\phantomsection
\addcontentsline{toc}{subsection}{\textit{Et donum. }}
\subsection*{\textit{Et donum. }}
\phantomsection
\addcontentsline{toc}{subsection}{\textit{Quae fuit unius Christi. }}
\subsection*{\textit{Quae fuit unius Christi. }}\pstart Deus pater in hoc solo homine sibi complacuit, in bunc unice, sua benignitate et dilectione respiciebat, adeo, ut ipsum haeredem faceret totius mundi, colloca- ret ad dexteram suam, nihil dignitatis et potentiae in eum non transfunderet, et postquam ille sangu i- nem suum pro nobis effudit, quomodo non reconci- liaret nos patri? Itaque gratia Dei per Christum in nos est propagata: non potuisset Christus reddere no-  \pend
\vspace{0.5cm}\noindent
\vspace{0.2cm}\rule{1cm}{0.2pt}\\ 
\hspace{0.2cm}\textit{mg}
\footnotesize Gratiae uis. 
\normalsize| 
\section*{IN EPIST. AD ROM.  }
\marginpar{[ p.54 ]}\pstart bis propitium suum patrem, nisi ipse eum benignum et propitium habuisset .  \pend
\phantomsection
\addcontentsline{toc}{subsection}{\textit{Exuberauit. }}
\subsection*{\textit{Exuberauit. }}\pstart Abunde uenit peccatum, sed abundantius gratia.   \pend
\phantomsection
\addcontentsline{toc}{subsection}{\textit{Et non sicut per unum. }}
\subsection*{\textit{Et non sicut per unum. }}\pstart Scilicet, Adam uenerat mors.   \pend
\phantomsection
\addcontentsline{toc}{subsection}{\textit{lta et donum. }}
\subsection*{\textit{lta et donum. }}\pstart Scilicet, uenit: expontt dissimilitudinem.   \pend
\phantomsection
\addcontentsline{toc}{subsection}{\textit{Nam condemnatio. }}
\subsection*{\textit{Nam condemnatio. }}\pstart Scilicet, omnium uenit. EX VNO DE- LICTO. Scilicet, unius hominis Adae peccato: et loquitur hic de peccato originali.   \pend
\phantomsection
\addcontentsline{toc}{subsection}{\textit{Donuim autem.  }}
\subsection*{\textit{Donuim autem.  }}\pstart Scilicet, uenit, id est, ipse Spiritus Dei datus est per Christum.   \pend
\phantomsection
\addcontentsline{toc}{subsection}{\textit{Ex multis delictis. }}
\subsection*{\textit{Ex multis delictis. }}\pstart Id est, omnium hominum peccatis.  \pend
\phantomsection
\addcontentsline{toc}{subsection}{\textit{Ad iustificationem. }}
\subsection*{\textit{Ad iustificationem. }}\pstart Id est, ut iustificaremur et absolueremur ab illis.  Sathan, qui hodie omnem mouet lapidem in suis mi- nistris, audet tam impudenter nugari, contra Euangelicam lucem, ut uel pueri illius mendacium deprehendant , dicit enim, Christus passus est,  \pend
\section*{COMMENT. POMER.  }\pstart sed pro illis tantum, qui ante ipsum fuerunt: proinde illius passionem, nobis tantum utilem esse, ad origt- nalis peccati ablutiomem, caetera uero peccata per nosipsos expianda, sed id quam impium sit commen- tum, conscientiae nostrae testes sunt. Est autem phra- sis mira quod dicit, Ex multis delictis, quam et ha- bet Esa. 40. Suscepit duplicia pro omnibus pecca- tis, ergo non pro meritis. Deus omnes inuenit in pec- catis, et quod supra dixit, pro ratione temporis, huc quoque pertinet. Quod igitur fingemus nobis me- ritum, nisi damnationis? Inimici eramus, iram, non fauorem Dei merueramus.   \pend
\phantomsection
\addcontentsline{toc}{subsection}{\textit{Et enim si per unius peccatum. }}
\subsection*{\textit{Et enim si per unius peccatum. }}\pstart Semper eadem repetit, propter infirmas conscien tias, quae semper putant peccatum esse fortius et ma- ius gratia Dei. Respondet ergo per Anti hesin, mul- to esse uberiorem gratiam peccato, illud est ex Deo, hoc in nobis. At quid nos possumus, cum simus foe- num, et solus Deus aeternus, uere potens omnia.   \pend
\phantomsection
\addcontentsline{toc}{subsection}{\textit{Regnauit. }}
\subsection*{\textit{Regnauit. }}\pstart Emphaticum est, quod ait, regnauit. Nunquid aliquis potuit suis uiribus resurgere ad uitam? nun- quid ullus Philosophorum, suis uirtutibus moralibus, po- tuit huius regni iugum excutere? an et Iudaei uitam per legem consequi potuerunt?  \pend
\vspace{0.5cm}\noindent
\vspace{0.2cm}\rule{1cm}{0.2pt}\\ 
\hspace{0.2cm}\textit{mg}
\footnotesize Gratiae am- plitudo. 
\normalsize| 
\section*{IN EPIST. AD ROM.  }
\marginpar{[ p.55 ]}
\phantomsection
\addcontentsline{toc}{subsection}{\textit{Qui exuberantiam gratiae. }}
\subsection*{\textit{Qui exuberantiam gratiae. }}\pstart Gratia in dante est, donum ad accipientem venit, ut quando pater dat filio cibum, uel tunicam: fauor uel gratia Dei efficit, ut donum, id est, Spiritum san- tum accipiamus, per quem consequimur fillationem et haereditatem.   \pend
\phantomsection
\addcontentsline{toc}{subsection}{\textit{Doni iusticiae. }}
\subsection*{\textit{Doni iusticiae. }}\pstart Id est, per quod iustificamur, estque ipsa remissio peccatorum.   \pend
\phantomsection
\addcontentsline{toc}{subsection}{\textit{Per uitam. }}
\subsection*{\textit{Per uitam. }}\pstart In uita rectius legeris, id est, uiuentes regna- bunt, non manebunt in morte.   \pend
\phantomsection
\addcontentsline{toc}{subsection}{\textit{Autore uno. }}
\subsection*{\textit{Autore uno. }}\pstart Cui omnia debemus, is est mediator noster, et scala Iacob, qua in coelum ascendimus. Ideoque , cum uideas Paulum ubique id nobis inculcare, non est quod foliciti, huiusmodi alios locos quaeramus.   \pend
\phantomsection
\addcontentsline{toc}{subsection}{\textit{Itaque sicut. }}
\subsection*{\textit{Itaque sicut. }}\pstart Hie primum absoluitur oratio, et repetit supe- riora. MALVM. Id est, peccatum.   \pend
\phantomsection
\addcontentsline{toc}{subsection}{\textit{Ad condemnationem. }}
\subsection*{\textit{Ad condemnationem. }}\pstart Vt condemnarentur. Haec omnia propter infir - mas conscientias toties repetit, quae non satis roborari  \pend
\vspace{0.5cm}\noindent
\vspace{0.2cm}\rule{1cm}{0.2pt}\\ 
\hspace{0.2cm}\textit{mg}
\footnotesize Scala lacob.  
\normalsize| 
\section*{OMMENT. POMER.  }\pstart possunt per inobedientiam unius hominis Adae, om- nes facti sumus inobedientes: per obedientiam Chri sti, qui usq ad mortem patri factus est obediens, fa- cti sumus obedientes, recepti in filiationem, qui era mus antea exclusi.  \pend\pstart Id est, quotquot credunt et Christum suscipiunt .  His omnibus tantum uult confirmare pacem, de qua dixit, ut simus certi plus posse gratiam, quam pecca- tum, non est nobis negligendum, quod aliena tu- sticia saluamur. Itaque quoties nos sollicitat peccatum in Christum nos coniiciamus, credamus illius omnia nostra esse, et per ipsius opus nos saluari. Cum respi- cimus in nostram iusticiam, occurrunt tot millia pec- catorum, quibus magis ac magis uexantur conscien- tiae nostrae, si uero in solum Christum respexerimus, pacem habebimus et gloriabimur, Deum esse no- strum Deum, quod sicut peccatum fuit potens, gra- tiam multo potentiorem fore.   \pend
\phantomsection
\addcontentsline{toc}{subsection}{\textit{Multi. }}
\subsection*{\textit{Multi. }}\pstart Occupatio necessaria est ad praedicta, quasi dice- ret, nos diximus interim de peccato et tusticia, morte et uita. Ad quid ergo lex, quae interim data est et non frustra, ut testantur conscientiae nostrae? Respondet. Subiit, uenit et intrauit ad nos obiter, idest, peccato regnante, et stante ab eo tempore,  \pend
\phantomsection
\addcontentsline{toc}{subsection}{\textit{Caeterum lex. }}
\subsection*{\textit{Caeterum lex. }}
\vspace{0.5cm}\noindent
\vspace{0.2cm}\rule{1cm}{0.2pt}\\ 
\hspace{0.2cm}\textit{mg}
\footnotesize Pax cordis un- de petenda.  
\normalsize| 
\section*{IN EPIST. AD ROM.  }
\marginpar{[ p.50 ]}\pstart quo in nobis per Adam suas radices fixit. Venit ita- que lex, cum tam olim regnarat peccatum, sed non suerat reputatum ab hominibus, Reyner hats uor- sunde geacht. Sicut hactenus putauimus summam iu- sticiam Papisticos errores. Vbi autem reuelata est ira Dei de coelo, uidemus esse huiusmodi impietatem, summam Christi abnegationem.   \pend
\phantomsection
\addcontentsline{toc}{subsection}{\textit{Vt abundaret. }}
\subsection*{\textit{Vt abundaret. }}\pstart Quando lex uenit, peccatum inuentum est in mundo, nam ab Adam regnarat dormientibus homi- nibus, et non citrantibus illud, uerum lex (ut est eius natura) statim illud reuelauit conscientiis no- stris, ideoque grauius multo et molestius factum est hominibus, quam antea fuerat, quia non sentieba- tur, nunc cum sentitur et cernitur in conscientiis eo- rum, parit mirum quendam horrorem mortis et de- sperationis, uelint legem plane non esse, quae no- estris cupiditatibus obsistit, imo murmurant contra Deum, qui nos lege damnat. Hic, hic est penitisst- mus adfectus cordium nostrorum, quem pauci uident.  Inquit ergo, tantum abest, ut lex iustificet, ideo non est respiciendum in opera eius, sed in Christum.  Lex in hoc primum data est, ut reuelet peccatum. Se- cundo, ut augeat illud: sed quid dices? Ergo lex ma- lum est uscpumorespondet ese soncton,ut vides  \pend
\vspace{0.5cm}\noindent
\vspace{0.2cm}\rule{1cm}{0.2pt}\\ 
\hspace{0.2cm}\textit{mg}
\footnotesize Lex robin peccati.  
\normalsize| 
\section*{COMMENT. POMER.  }\pstart Sed ita agit Deus nobiscum cum lege, quemadmodum medicus cum egroto per medicinam, haec per se bona est et salutaris, attamen materiam delinquendi sepe praebet: ut si mandetur aegroto, ne uescatur carne suil- la, statim illi quidam adpetitus mouetur edendi suillam, siquidem natura nititur in uetitum, ita et lex multis est peccandi occasio, cum tamen Deus per eam nos urgeri uoluerit, ut uidentes nostrum peccatum et damnationem, adurgeremus ad quaerendam eius gra- tiam, nam quanto grauius et intolerabilius est deli- ctum, tanto citius quaerimus Christum, ut quo quis plus diligeret, hoc ei plus remitteretur, ut Christus in Luca dicit, in historia de muliere peccatrice.   \pend
\phantomsection
\addcontentsline{toc}{subsection}{\textit{Vbi uero exuberauit. }}
\subsection*{\textit{Vbi uero exuberauit. }}\pstart Haec quoque ad pacis confirmationem pertinent, etiamsi commisissem omnia peccata mundi solus, nul- la essent prorsus, si gratiam consequerer in Christo.  Deus sinit suos cadere, etiam grauissime quandoque , ut ardentius resurgant, et gratificentur Deo, sed carnales huiusmodi praedicatione Euangelica abutun- tur, dicentes: Ergo faciamus grauia peccata. Con- temptores sunt (quos Deus non recipit, et quibus Euangelium non praedicatur) non peccatores, id est, ii, qui agnoscant sui damnationem propter pec- catum. Hinc in Psalm. ait. Longe a peccatoribus sa- lus, quoniam non exquisierunt iustificationes tuas-  \pend
\vspace{0.5cm}\noindent
\vspace{0.2cm}\rule{1cm}{0.2pt}\\ 
\hspace{0.2cm}\textit{mg}
\footnotesize Lex cur data.  
\normalsize| 
\section*{IN EPIST. AD ROM.  }
\marginpar{[ p.57 ]}\pstart Non quia peccauerunt, ab eis longe est salus, sed quia iustificationes tuas non exquisierunt, nolebant per te- iustificari. Huc pertinet quod Christus ad Petrum di- cit, qui tamen negarat Christum, Petre amas me? etc.   \pend
\phantomsection
\addcontentsline{toc}{subsection}{\textit{In morte. }}
\subsection*{\textit{In morte. }}\pstart Id est, condemnatione hominum.   \pend
\phantomsection
\addcontentsline{toc}{subsection}{\textit{Per lesum, }}
\subsection*{\textit{Per lesum, }}\pstart Semper Christum interponit, ut sciamus per quem nobis talia beneficia contigerint.   \pend\pstart Actenus iusticiam ueram, per solum Christum me- diatorem peccatoribus, a Deo exhiberi: Item, etus iusticiae fructum esse, pacem conscientiarum do- cuit. Hic ea confutat argumenta, quae acumen ratio nis humanae, contemptricis gratiae et misericordiae Dei, solet obiicere. Nouit enim Paulus, qui et ipse caro fuit, quid soleat illa, in Det contemptum machinari, ubi gratiae praedicationem audierit, siquidem offendi- tur, cum audit: Vbi abundauit delictum, abunda- uit et gratia, statim infert: Si peccatum tam bonum est, ut sit occasio gratiae Dei in nobis, ergo pecce- mus? Item si sumus liberi, quiduis nobis impune licet.  Itaque summa huius Capitis, nihil aliud continet, quam dilutionem quandam huiusmodi ineptiarum, rationis  \pend
\section*{COMMENT. POMER.  }\pstart humanae. Simile in tertio Cap-uidistt, ad quod re- spondebat Paulus: Quorum damnatio iusta est. Hic fusius omnia agit, quo magis imbecilles conscientias erigat, nam ubi ratio obiicit : praestat ergo pecca- re, sacere quidvis . Respondet : Quando alii facti su- mus, necesse est, ut aliter vivamus, quando trans- lati sumus e regno tenebrarum, in lucis regnum, pu- deat nos eorum scelerum, et operum tenebrarum. quae lux iam ferre nequit. Absurdissima sunt haec, si- quidem inde sequeretur, esse scortandum, occiden- dum etc. cum tamen ratio interim haec pruden- tissime obiecisse sibi uidetur, putat Dei consilia stul- ta et inepta, sua uero sapientissima, et neutiquam reprehendenda. Sed ubi animaduerterit huiusmodi ratiocinationes non bene cedere, putat omnia praece- dentia de fide et gratia, esse ementita. Preinde hu- manis quoque utitur argumentis, praeter ea, quae e Sacris fontibus petita sunt, dicens: Minime est pec- candum, cum iam peccatorum remissio praedice- tur, non in hoc, ut porro peccemus, sed ut pecca- tum odiamus, et peccatum aboleatur : resuscitati sumus in uitam, non in mortem. In fine Capitis do- cet, quomodo id fieri possit, nempe adiutorio spiri- tus Dei, qui regnat in nobis: alias eni secundum naturam, peccatum quomodo odio haberemus? Odit Spiritus sanctus illud, is impellit et mouet no- stros adfectus in eiusdem odium: in hoc ille nobis da-  \pend
\vspace{0.5cm}\noindent
\vspace{0.2cm}\rule{1cm}{0.2pt}\\ 
\hspace{0.2cm}\textit{mg}
\footnotesize Ratio of- fenditur gra- tiae praedica- tione.  
\normalsize| 
\section*{}
\marginpar{[ p.]}\pstart  \pend
\phantomsection
\addcontentsline{toc}{subsection}{\textit{}}
\subsection*{\textit{}}\pstart  \pend
\vspace{0.5cm}\noindent
\vspace{0.2cm}\rule{1cm}{0.2pt}\\ 
\hspace{0.2cm}\textit{mg}
\footnotesize 
\normalsize| 
\section*{COMMENT. POMER.  }\pstart cato esset tribuenda salus, malti, imo omnes impii saluarentur.  \pend\pstart Id est, e regno Sathanae, translatt in Christi reg- num. Hic disce quam breuiter doceat Apostolus, in quo consistat Christianismus. Est is nihil aliud, quam mori peccato, et uiuere iusticiae: mortui autem su- mus illi, per mortem Christi, resuscitati per illius re- surrectionem.   \pend
\phantomsection
\addcontentsline{toc}{subsection}{\textit{Qui mortui sumus. }}
\subsection*{\textit{Qui mortui sumus. }}\pstart Scilicet, per Christum. Mortui sunt peccato, qui non habent de eo conscientiam, ut maxime illud sen- tiant, tamen et non consentiunt, sciunt se nomen de- disse iusticiae: Item ipsorum nomina scripta in coelis, item se pertinere ad iusticiae, non peccati regnum, ideo et ipsis non imputatur peccatum, ut aliis im- piis, qui peccato obsequuntur, grauantur quidem peccato, sed quamprimum uident in uultum benigni patris, euanescit.   \pend
\phantomsection
\addcontentsline{toc}{subsection}{\textit{Peccato. }}
\subsection*{\textit{Peccato. }}
\phantomsection
\addcontentsline{toc}{subsection}{\textit{Quomodo uiuemus. }}
\subsection*{\textit{Quomodo uiuemus. }}\pstart Absurdum esset dicere: Viuendum est peccato, cum mortui simus peccato, quod probat per symbo- lum baptismi nostri.   \pend
\phantomsection
\addcontentsline{toc}{subsection}{\textit{An ignoratis, quod qui cum. }}
\subsection*{\textit{An ignoratis, quod qui cum. }}\pstart Nunquid tam stupidi estis, ut ignoretis, nos unum iam esse cum Christo, item nos pertinere ad illius,  \pend
\vspace{0.5cm}\noindent
\vspace{0.2cm}\rule{1cm}{0.2pt}\\ 
\hspace{0.2cm}\textit{mg}
\footnotesize Christianis- mus quid. 
\normalsize| 
\section*{IN EPIST. AD ROM.  }
\marginpar{[ p.59 ]}\pstart non peccati eorpus, membra sumus Christi, non Sathanae.   \pend\pstart Id est, ita mortuus est Christus pro nobis, ut cum ipso simus quoq mortut omnes. Omnes eramus obno- xii morti propter peccatum, mori debebamus, sed Christus hoc debitum pro nobis soluit. Infirmiores eramus, quam ut mortem uinceremus, absorpsisset haec nos, ac pessundedisset: Verum Christus, fortior morte, absorpsit eam, non sibi, sed nobis: non era- mus soluendo, ille diluit debitum tam efficaciter, ut, tametsi nos non satisfecerimus, iam morti satisfactum sit. Mortut ergo sumus illius morte, sicut qui pro me soluit debitum, ut maxime ego non numerauerim pe- cuntas, tamen tam sum liber a debito, quia is pro me satisfecerit: ita iam perinde mortuus sum, atque Christus.   \pend
\phantomsection
\addcontentsline{toc}{subsection}{\textit{In mortem eius. }}
\subsection*{\textit{In mortem eius. }}
\phantomsection
\addcontentsline{toc}{subsection}{\textit{Sepulti igitur. }}
\subsection*{\textit{Sepulti igitur. }}\pstart Id est, sicut ipse sepultus est, Vuyr seyn insey- nen todt getrochen. Quare?  \pend
\phantomsection
\addcontentsline{toc}{subsection}{\textit{Vt quemadmodum excitatus Chri- stus ex mortuis per gloriam. }}
\subsection*{\textit{Vt quemadmodum excitatus Chri- stus ex mortuis per gloriam. }}\pstart Id est, quemadmodum Deus pater ex ignominio- sa morte, traxit suum Christum, et in gloriam collo- cauit, iuxta illud Iohan. 17. Pater glorifica filium tuum, deditque illi nouam uitam, quam et nobis cum  \pend
\vspace{0.5cm}\noindent
\vspace{0.2cm}\rule{1cm}{0.2pt}\\ 
\hspace{0.2cm}\textit{mg}
\footnotesize Christi satisfa- ctio nostra.  
\normalsize| 
\section*{COMMENT. POMER.  }\pstart eo communem fecit, ita et nos e morte extracti, et in nouam vitam constituti, perseveremus in ea, non repetamus mortem . Semel ille mortuus est, et iam semper vivit : Mors illi ultra non dominabitur .  Ita et nos, qui in ipso reuiximus, non cupiamus ite- rum mori: iam regnat ille in regno uitae, non mor- tis, idque hodie in gloriam eius praedicatur per totum orbem. Summo dedecore ipsum adfecerimus, ubi con- tempta noua uita, ad mortem deficere uoluerimus, quod ii faciunt, qui inferunt: ergo peccabimus, ma- lunt illi mortem, quam uitam, uide quam absurda haec sint, peccatum remittitur, non ut pecces, sed quia Deus odio habet peccatum. Itaque ; peccatum ad Christi regnum, et nouam uitam non pertinet. Be- nignitate alicuius principis e carcere soluor, stultus fuerim, qui iam benignitatem putem non libertatem, sed captiuitatem fore, non ideo solutus sum, ut ma- neam in carcere, sedut sim liber.   \pend
\phantomsection
\addcontentsline{toc}{subsection}{\textit{Nam si insiticii. }}
\subsection*{\textit{Nam si insiticii. }}\pstart Id est, si inserti sumus illi, et unum factum cum eo, si cum eo uere sumus mortui, non allegorice, sicut quidam causantur, certe et resurreximus cum eo. Volunt hic quidam Magistri carnis, similitudi- nem mortis, nescio quibus tenebris obfuscare, sed ne- queunt. Aliter enim nos docet Spiritus sanctus Dei,  \pend
\vspace{0.5cm}\noindent
\vspace{0.2cm}\rule{1cm}{0.2pt}\\ 
\hspace{0.2cm}\textit{mg}
\footnotesize Exemplo Chri sti moriamur et nos: 
\normalsize| 
\section*{60 IN EPIST. AD ROM.  }
\marginpar{[ p.]}\pstart qui regnat in cordibus credentium, quare non fa- cienda est hic allegoria mortis. Similiter, et uere nos mortui sumus, quando ipse mortuus est pro nobis: illius mors, nostra est mors, in conspectu Det, ut maxime ratio humana id non capiat. Hoc uerum est fundamentum, cui innititur nostra salus, hic uerus usus mortis Christi, quem illi, qui somniant allego.  ricam mortem, non uident, praescribunt nobis tan- tum imitationem Christi, sed interim haec negligunt, quae inprimis scire retulerat. Itaque que certus sim necesse, me mortuum esse in morte Christi, me resurrexisse re- surrectione Christi, regnare in regno Christi, sede - re ad dexteram patris in Christo: In summa, omnia quae Christi sunt, mea esse, Ephesios primo. Non est ergo similitudo haec : Christus mortuus est, ergo et nos libenter moriamur. Sed illa: Christus mor- tuus est, ergo et nos tain sumus mortui. Christus re- surrexit: Ergo et nos resurreximus, hoc qui credit, facile et hilari corde imitabitur Christum, bonis ope- ribus, imo et audebit pro illo mori, et mortificare membra, quae sunt super terram.   \pend
\phantomsection
\addcontentsline{toc}{subsection}{\textit{Illud scientes , quod uetus. }}
\subsection*{\textit{Illud scientes , quod uetus. }}\pstart Duplicem describit hominem, Veterem et nonum, vetus est, qu ex carne nascitur, ut Christus dicit.  \pend
\vspace{0.5cm}\noindent
\vspace{0.2cm}\rule{1cm}{0.2pt}\\ 
\hspace{0.2cm}\textit{mg}
\footnotesize Mortis Chri- sti usus.  
\normalsize| 
\section*{COMMENT. POMER.  }\pstart Quod natum est ex carne, caro est, id est, omnia quae pertinent ad hominem et uires eius: totum quod na- scitur, anima et corpus, caro est et peccatum, neque  intelliges haec de ratione et sensualitate. Hinc Chri- stus dicit: Oportet uos renasci, adeo nihil est in ue- stra illa prima natiuitate quod bonum sit, omnia sunt damnata, non potestis uelle et cogitare quae Deisunt, hinc est, quod lex Dei uobis displiceat, et contra eam murmuretis, uelitis eam non esse, neutiquam potest uestra uoluntas, cum Dei uoluntate conspirare, er- go quicquid uultis, quicquid cogitatis, aut concupi- scitis id contra Deum. Nouus uero homo est, in quo nihil est naturae, nihil uirium humanarum, sed spiri- tus Dei, qui ipsum mouet et in pellit, qui quiequid agit, scit certissime se agere Dei bona uoluntate, scit omnia illi placere, nihil esse in se quod ille odiat et damnet.   \pend
\phantomsection
\addcontentsline{toc}{subsection}{\textit{Crucifixus est. }}
\subsection*{\textit{Crucifixus est. }}\pstart Sicut antea mortuus, ita nunc crucifixus dicitur, ni- si quod hoc intersit, quod qui crucifixus est, nondum plane mortuus sit, uiuere potest adhuc diebus aliquot, interim tamen, dum uiuit, pro mortuo habetur, neque  spes ulla est uitae illius, cum plane sit adiudicatus mor- ti. Sic nos, utcunque sentiamus peccatum, crucifixi sumus, et morti iam destinati, imo sumus in ipsa mor- te, iam coept in nobis mort, et moritur uere peccati et caro nostra.   \pend
\section*{IN EPIST. AD ROM.  }
\marginpar{[ p.61 ]}\pstart Id est, ut non maneret in peccato, quod contra il- lud est quod dicunt: Ergo peccabimus. His ostendit carnem non significare corpus tantum, uerum quid- quid uirium est in homine, nempe, anima, intellectus etc. Siquidem non potest concupiscere corpus sine anima. Hoc totum, ait, per Christi mortem est abolen- dum, reformandumque iusticia Dei. Spiritus baptisat nos toto tempore uitae nostrae, donec totum illud cor- pus moriatur: utcunque adhuc sentiamus illud pecca- tum, tamen crucifixum est, non regnat, sed subigi- tur spiritui, ergo corpus peceati, nihil aliud est, quam quod peccat.  \pend
\phantomsection
\addcontentsline{toc}{subsection}{\textit{Vt aboleretur corpus. }}
\subsection*{\textit{Vt aboleretur corpus. }}
\phantomsection
\addcontentsline{toc}{subsection}{\textit{Etenim qui mortuus est, iustifi catus est. }}
\subsection*{\textit{Etenim qui mortuus est, iustifi catus est. }}\pstart Non potest accusari peccati, uel quod peccet, is qui mortuus est. Et ut simili declarem, si quispiam accusaret me adulterii heri perpetrati, conuincerem autem illum testibus, eo tempore me fuisse mortuum, nonne accusator mendax repudiaretur per iudicem? Quamprimum ergo totum corpus peccati aboletur et moritur, mortuum est omne peccatum, et uis eius.   \pend
\phantomsection
\addcontentsline{toc}{subsection}{\textit{Viuemus. }}
\subsection*{\textit{Viuemus. }}\pstart Scilicet, ipsare, et uere cum illo in resurrectio- ne eius.   \pend
\vspace{0.5cm}\noindent
\vspace{0.2cm}\rule{1cm}{0.2pt}\\ 
\hspace{0.2cm}\textit{mg}
\footnotesize Mori pec- cato.  
\normalsize| 
\section*{COMMENT. POMER.  }
\phantomsection
\addcontentsline{toc}{subsection}{\textit{Scientes quod. }}
\subsection*{\textit{Scientes quod. }}\pstart Mors semel Christo dominata est, ut Petrus in Actis dicit. Impossibile erat eum teneri in morte, qui morti obnoxius non erat. Diligentissime haec semper inculcat eadem Apostolus, ut opprimantur uirulen- tae linguae canes, quae dicunt, ergo peccabimus: quo modo peccare uellemus, translati in regnum iusticiae?  \pend
\phantomsection
\addcontentsline{toc}{subsection}{\textit{Nam quod mortuus fuit, pec- cato etc. }}
\subsection*{\textit{Nam quod mortuus fuit, pec- cato etc. }}\pstart Peccatum ipse portauit, non quia fecerat pecca- tum, sed ut peccatum in nobis mortificaretur, uiue- remusque ; noua uita Deo.   \pend
\phantomsection
\addcontentsline{toc}{subsection}{\textit{Ita et uos existimate. }}
\subsection*{\textit{Ita et uos existimate. }}\pstart Credite uos sic mortuos in Christo, ut non pecee- tis, ubi autem solicitarit uos caro ad peccatum, cogi- tate et credite, uobis nihil esse cum peccato, et ex eius regno uos esse ereptos, uiuitis Deo, non pecca- ti dominio.  \pend
\phantomsection
\addcontentsline{toc}{subsection}{\textit{Per Christum. }}
\subsection*{\textit{Per Christum. }}\pstart Haec adiicit, ut intelligas mediatorem ipsum esse, per quem haec omnia nobis contigerint tam opulen- te gratiae.   \pend
\phantomsection
\addcontentsline{toc}{subsection}{\textit{Ne regnet igitur. }}
\subsection*{\textit{Ne regnet igitur. }}
\section*{IN EPIST. AD ROM.  62 }
\marginpar{[ p.]}\pstart Est Occupatio: quid faciemus? inquient pii, carnis imbecillitati adhuc obnoxii, non sumus Angeli, habe- mus tale corpus, cui procliuum est, ut tantum pec- cet . Respondet Paulus: Vos quidem portatis morta Ie corpus, sentitis peccatum, et carnis iniurias, id non potestis effugere, dum in mundo estis. Neque enim dico, ut non sit in uobis peccatum, sed ne reg- net in uobis, non obediatis peccato: Spiritus sanctus repugnet carni, et contundat illius robor, sed haec ipse clare exponit. ILLI . Scilicet, corpori.   \pend
\phantomsection
\addcontentsline{toc}{subsection}{\textit{Per cupiditates. }}
\subsection*{\textit{Per cupiditates. }}\pstart Concupiscentia in nobis est, ergo totum pecca- tum: Verum non imputatur, quia uellemus illud ab esse. Sic igitur nobis cum Deo conuenit, qui queque  uult abolitum esse peccatum. Sed hic constantissima fi- de opus est, non opinione carnis, qua nostri somni- ant fidem.   \pend
\phantomsection
\addcontentsline{toc}{subsection}{\textit{Neque accommodetis. }}
\subsection*{\textit{Neque accommodetis. }}\pstart Id est, non abutimini membris uestris, quae iam sunt instrumenta sanctificationis, nam sunt membra corporis Christi. ICor. 6. non porro peccandum est, quum erepti sumus e regno peccati.   \pend
\phantomsection
\addcontentsline{toc}{subsection}{\textit{Sed accommodetis uosmetipsos Deo. }}
\subsection*{\textit{Sed accommodetis uosmetipsos Deo. }}
\section*{COMMENT. PIOMER.  }\pstart Vt infra Duodecimo docet: Offeramus nos Deo et corpore et spiritu, id est, nihil sit in uobis, quod non seruiat iusticiae.   \pend\pstart Id est, iam resuscitati estis, et uiuitis: facite er- go opera uiua non mortua.   \pend
\phantomsection
\addcontentsline{toc}{subsection}{\textit{Velut ex mortuis uiuen. }}
\subsection*{\textit{Velut ex mortuis uiuen. }}\pstart Id est, ad Dei honorem, quod ubique Paulus admo- net. Est enim caro propensissima, audita praedicatio ne Euangelii, in illius abusum, quo blasphe met Deum.  Hactenus itaque ratio et efficacia baptismi explica- ta est, qua monemur, ne porro peccemus: et ut bre uius repetamus, Signum illud externum duo habet, im mersionem et exiractionem, quae sunt duae partes scripturae, siue duo officia praedicationis, id est, lex et Eudngelium. Lege immergimur, non solum cor- pore, sed et anima, id est, toti, ut significetur to- tum nostrum hominem, quantus quantus est, opor- tere mori: neque tamen solum id significari, uerum ipsa re fieri, reuera enim per baptismum eum Christo morimur. Extrahimur autem ex hoc baptismate, id est, resurgimus cum Christo, crescente in nobis spiri- tu: fit igitur, ut simul et moriamur et uiuamus: moriamur per legem, uiuamus autem per Euangelium: nam utrunque in ratione baptismi deprehenditur. Non est igitur quod causetur caro, nobis liberum esse, at  \pend
\phantomsection
\addcontentsline{toc}{subsection}{\textit{Deo. }}
\subsection*{\textit{Deo. }}
\vspace{0.5cm}\noindent
\vspace{0.2cm}\rule{1cm}{0.2pt}\\ 
\hspace{0.2cm}\textit{mg}
\footnotesize Baptismi sig- num quid.  
\normalsize| 
\section*{IN EPIST. AD ROM.  }
\marginpar{[ p.63 ]}\pstart peccemus. Sed quia haec nihil quod Dei est intelli- git, non potest aliud, quam Dei gratiam et miseri- cordiam blasphemare, id quod hodie in multis cum dolore est cernere.   \pend
\phantomsection
\addcontentsline{toc}{subsection}{\textit{Peccatum uobis non dominabitur. }}
\subsection*{\textit{Peccatum uobis non dominabitur. }}\pstart Et hic locus probat non porro peccandum: est ue- ro insignis, quia legis abrogationem tractat, in quam caro crasso iudicio trruit, et nihil non audet. Est au tem argumentum a facili. q. d. Antea quando ad- huc eratis sine fide, iusticia Dei, ducebamint quasi uin- cti ad peccandum, nam tum uobis dominabatur pec- catum, nunc uero quando estis sub gratia, habetis De- um patrem, non iudicem, non potestis peccare: non potuistis olim non peccare, imo necessario ad id fere bamini, nunc quia restituti estis, per Dei gratiam, non peccatis. Hic quae de Libero Arb. dicunt illi, sunt respuenda, damnamus ea et execramur, Dei uero praedicamus potentiam, quae in nobis quiduis potest efficere. Eramus autem usqueadeo obnoxii legi olim, ut simul omnes morti uorandos obiiceret, quemadmo dum is, qui hominem occiderit, ad mortem, per legem de homicidis latam urgetur, sic nos quia peccauera- mus omnes, morti adiudicabamur per legem, iuxta illud: In quacunque die comederis etc. Haec erat ty- rannis, hoc intolerabilissimum dominium legis, quo uexabamur, nunc tandem Dei fauore sumus liberatt, Contra Libe- rum Arb.   \pend
\vspace{0.5cm}\noindent
\vspace{0.2cm}\rule{1cm}{0.2pt}\\ 
\hspace{0.2cm}\textit{mg}
\footnotesize 
\normalsize| 
\section*{COMMENT. POMER.  }\pstart habemusque iam dominium in legem, non ex nostris uiri- bus, sed Deo, fitque ut non peccemus: nam quomodo peccaremus,  cum nec lex nec peccatum, nobis domi- netur. Quare diligenter spectandum est, quid sit es- se sub lege, quid sub gratia, nam si neglectum id fue rit, facile usu et fructu huius loci, de abrogatione le- gis frustraberis. Sub lege esse, est damnari in consci- entia, illius sententia, quia peccauimus. Sub gratia es se, est conscientiam certissimam esse de Deo, quod sit pater, illiusque benignitate omne remissum pecca- tum. Vbi igitur erit peccatum, cum remissum sit? quis damnabit, cum lex non possit damnare? Sed et hoc monendum est, multos hodie abuti libertate Chri- stiana, qui dicunt: Cum abrogata sit lex, nihil est praecipiendum Christianis: Item, ietunium Papisti- cum nihil est, ergo non ieiunabimus: Item, oratio- nes institutae per Papam, nihil prosunt, ergo non est orandum: sic illi iudicant, et abutuntur Christiana libertate. Illis non secus accidit, ac Iudaeis et Gen- tibus quibusdam olim, qui uidebantur credere Euan- gelio, cum tamen non crederent: abrogationem et libertatem nondum sunt experti in cordibus suis, ideoque oculos tantum in externa defigentes, et cras- sa errant. Verum qui sentiunt in conscientia sua le- gis uim sublatam, ii uere norunt, quae sit libertas.  Id autem fit per Euangelii praedicationem, qua re- missio peccatorum annunciatur, quemadmodum  \pend
\vspace{0.5cm}\noindent
\vspace{0.2cm}\rule{1cm}{0.2pt}\\ 
\hspace{0.2cm}\textit{mg}
\footnotesize Sub lege, au gratia esse quid.  
\normalsize| 
\hspace{0.2cm}\textit{mg}
\footnotesize Libertate quo modo abutan- tur quidam.  
\normalsize| 
\section*{EPIST. AD ROM.  64 }
\marginpar{[ p.]}\pstart Gentium conscientiae olim erant alligatae quibusdam culti- bus externis obseruationi dierum. Praedicato autem Euangelio, deprehendebant errorem, illaque ; esse pror- sus ridicula: idem in ludaeis quibusdam uisum est, qui quoque suis constitutionibus uexabantur, reuelato ue- ro Euangelio, intellexerunt esse fabulas. Sic nos quoq.  iim uulgato Euangelio per orbem, scire oportebat, conscientias nostras liberatas per illud, a torturis et uexa ationibus Papisticis, non per id datam occasionem stul tae carni, ut qui duis praesumat, quam insaniam Paulus et Apostoli ubique damnant. Norunt autem tantum ueri Christiant, legis abrogationem, hanc admirantur, et adorant in spiritu, cum sit ipsa Dei misericordia. Mul- tis nunq̃ potuit persuaderi totam legem abrogatam, quare tantum Ceremonialem legem sublatam praedicabant : Proin de dicebant, si moralis lex est abrogata, ergo licet scor- tari, furari etc. quae sunt contra Deum et proximum.  Vsque adeo haeserunt in ratione humana, ideoque ; uani fa- cti sunt, non uiderunt neq intellexerunt Det diuitias, no- bis per Christum exhibitas. Quid potuit hoc stultius di- ci? si eni Ceremonialis lex abrogata esset, ergo tantum Iudaeis, non aliis hoibus abrogata esset: cum tamen dicant Apostoli, neminem potuisse sustinere onus legis, ideoque  ridiculum est haec audire. Sed hoc omnt dementta inepti us, ꝙ cum nos per Christum asserant liberatos a lege Cere moniaerum, interim tamen multo grauioribus nos onera- rint, ut eorum Decreta indicant. Est hoc praedicare  \pend
\section*{COMMENT. ROMER.  }\pstart Euangelium, aut Dei gratiam, et intelligere liberta- tem Christianam? Non igitur hic, et similes loci, de abrogatione Ceremonialis legis tantum intelliguntur, uerum de untuersa lege, qua sublata est quoque ma- ceries, tota mandatorum lex, Chirographum, quod erat contra nos, sublatum, Ephe. 2. Col. 2. Haec pro secto de morali lege decem Praeceptorum dicta sunt, per quae praecepta, mea conscientia, uacua gloria Dei, ob ligatur, damnatur, quod sciam et iam agnoscam me non dilexisse, neque posse diligere Deum et proximum ex corde: ubiuero agnouero Christum mihi datum, qui pro me legi satisfecit, aboletur tota lex, et qui fueram seruus et captiuus legis, spiritum filiorum, qui est libertatis, consequor. Itaque libertate illa legis, Primo remissionem omnis peccati obtinemus, quo quis non uidet, omne ius legis abrogari? nam nihil amplius in nobis illa inuenit quod condemnet. Quid plura requiris? nihil hic reliquum quod debeas legi, soluta sunt omnia per remissionem. Secundo, dona- tur nobis Spiritus sanctus, per quem libere facimus quidquid exigit lex, non tamen lege cogente, et id quod facimus bonum est, quia sine lege facimus, spi- ritu intus nos ad id trahente. Si tamen quid contrarium facere nos contigerit, dolemus eodem spiritu, et non imputabitur, quia in remissionem peccatorum consti- tuti sumus. Hoc igitur est uere facere legem, non ty- rannide compulsos, sed spiritus mra suauitate et fa-  \pend
\vspace{0.5cm}\noindent
\vspace{0.2cm}\rule{1cm}{0.2pt}\\ 
\hspace{0.2cm}\textit{mg}
\footnotesize Lex quatenus abrogata.  
\normalsize| 
\section*{IN EPIST. AD ROM.  65 }
\marginpar{[ p.]}\pstart eilitate, ad illud perducti. Tertio, consequimur, ut peccatum quod reliquum est in nobis, et quod soli- citat nos, non imputetur. Siquidem Spiritus sanctus nobis donatus est, sumusque filii, et Deus pater est.  Hic nihil habet turis lex, non potest nos: damnare, quia in nobis abrogata est, Deus est supra legem, qua- re potest efficere, ut non imputetur ipsum peccatum: et hoc est quod Paulus dicit: Vbi Spiritus domini, ibt libertas, quam qui habent, non ea abutuntur: at qui eam sibi fingunt, abutuntur. Et in hanc sententiam bic Paulus loquitur: Non estis sub lege, id est, non habet lex ius uos damnandi, tametsi uerum sit, ubi in uos respieitis, peccatum habetis, ubi uero in gra- tiam, non habetis, nam remissa sunt omnia peccata, et ea tantum facitis, quae sunt uoluntatis diuinae.  lectionem. Dilectionem Dei non possunt non habere, qui Christiant sunt, ideoque de illa non praecipitur, sed tantum ea quae est in proximum. Exiis omnibus stulticiam eorum uides, qui inferunt: ergo peccabi- mus: quae esset illa liberatto a lege, si adhuc captiui per eam teneremur? A morte, a lege, a peccato su- mus liberati, quare minime est peccandum: stultus is fuerit, qui dixerit, e carcere sum liberatus, ergo in grediar carcerem.   \pend
\phantomsection
\addcontentsline{toc}{subsection}{\textit{An nescitis. }}
\subsection*{\textit{An nescitis. }}\pstart Hie respondet carnalibus, qui inferebant: ergo  \pend
\section*{COMMENT. POMER.  }\pstart peccandum. Est autem talis responsio, quam ratio no- uit: respondebit autem fortius in sequenti Cap. ꝗ.  d. Si uultis peccare, ergo non estis liberati a pecca- to. Estque generalis sententia: Si ego alicui principi- me addico, seruus illius sum, sic qui se peccato accon modat, seruus peccati efficitur, nam illi obedit.  ECCATI. Scilicet, serut. AD MORTEM. Id est, condemnationem.  OBEDIEN. Scilicet, serui.   \pend\pstart Id est, iustificationem, ut sitis iusti. Sententia est olim pcton seruiebatis, iam alium dominum habetis Chri- stum, illi seruite, nam in illius regnum estis traducti.   \pend
\phantomsection
\addcontentsline{toc}{subsection}{\textit{lusticiam. }}
\subsection*{\textit{lusticiam. }}\pstart Scilicet, sit Deo, quod liberati sitis a peccatis, uon quod fueritis peccatores.   \pend
\phantomsection
\addcontentsline{toc}{subsection}{\textit{Gratia. }}
\subsection*{\textit{Gratia. }}\pstart Id est, suscepistis hanc doctrinam per fidem.   \pend
\phantomsection
\addcontentsline{toc}{subsection}{\textit{Obedistis. }}
\subsection*{\textit{Obedistis. }}\pstart Id est, non carnaliter, aut hypocritice. DO- CTRINAE. Scilicet, Euangelicae-  \pend
\phantomsection
\addcontentsline{toc}{subsection}{\textit{Ex animo. }}
\subsection*{\textit{Ex animo. }}\pstart Scilicet, a Deo, quasi erepti ex alio, et in dliud traducti. Hic uides, quid relinquatur Libero Arb- cum traducamur per illud: non potes hanc doctrinam siscpere, cum sit ex sola fide.   \pend
\phantomsection
\addcontentsline{toc}{subsection}{\textit{Traducti. }}
\subsection*{\textit{Traducti. }}
\section*{IN EPIST. AD ROM.  66 }
\marginpar{[ p.]}
\phantomsection
\addcontentsline{toc}{subsection}{\textit{Humanum, }}
\subsection*{\textit{Humanum, }}\pstart Haec ad sequentia refero, non ad superiora, ut qui dam, quasi non satis sit Christianum, dici regnum iusticiae seruitutem: Quae quidem sententia bona est, sed malo argumenta esse rationis, quae pia uidentur, at si quis expendat, non satis consistent. Siquidem dicit Paulus: Loquor quemadmodum soletis uos lo- qui, de iis rebus, attamen intimius illa sunt scrutan- da- Non satis est dixisse: Si potuistis in regno mortis facere malum, potestis et iam in regno iusticiae face- re bonum et iustum. Paulus enim hoc interim con- temnit, quod hodie nostri Theologi se in Scrip. inue nisse gloriantur. Nam et ex hoc loco Liberum Arb.  statuunt miseri, non uidentes, quod Paulus dicat se lo- qui humano more, id est, uti rationibus et argumen- tis rationis. Nam subiungit.  \pend
\phantomsection
\addcontentsline{toc}{subsection}{\textit{Propter infirmitatem uestram. }}
\subsection*{\textit{Propter infirmitatem uestram. }}\pstart Id est, forte aliter non potestis capere, utar argu- mento rationi uestrae obuio. q.d-olim fecistis malum, ta benefacite, humana sunt haec et insufficientia ad fa- ciendam legem, ideoque ; hoc diligenter spectandum est inte- rim, Paulum hic praesupponere: Nos esse iam sub gratia, non sub lege, sub qua non aliud q̃ bonum facere possumus, alians frustra haec diceret, quia id non est in nostra pote- state, uerum in Spiritus, qui in nobis est, omnia faciens quae bona sunt. SERVA. Id est, serurentia.   \pend
\phantomsection
\addcontentsline{toc}{subsection}{\textit{Immundi. }}
\subsection*{\textit{Immundi. }}
\phantomsection
\addcontentsline{toc}{subsection}{\textit{}}
\subsection*{\textit{}}
\section*{COMMENT. POMER.  }\pstart Hac intelligit, scortationem, adulteria, stupra, qui bus in proprium corpus peccatur: iniquitatis uocabu- lo, reliqua uitia omnia intelligit, quae sunt extra cor- pus nostrum . 1. Corinth. 12.   \pend
\phantomsection
\addcontentsline{toc}{subsection}{\textit{Serua iusticiae. }}
\subsection*{\textit{Serua iusticiae. }}\pstart Est Antithesis-q. d. Nolite iam abuti membris ue- stris. Nam in hoc estis erepti a peccatis, ne posthac uiuatis in eis, membra uestra seruiant iusticiae ad san- ctificationem uestri, quae indies in uobis crescit, ut de fide Paulus dixit Rom. I.   \pend
\phantomsection
\addcontentsline{toc}{subsection}{\textit{Cum enim serui essetis. }}
\subsection*{\textit{Cum enim serui essetis. }}\pstart Hic duplicem habes seruitutem et libertatem, bo- nam alteram, alteram malam. Libertas iusticiae, haec pessima est, et improprie dicta, Vuen eyner uont Gott frey ist, und des Teufels eygen, neque medium est, aut peccatum, aut iusticia. Simile est hoc Psal- mi-87. Factus sum sicut homo sine adiutorio, in- ter mortuos liber, id est, derelictus, destitutus omni auxilio, tametsi interpretes alio detorserint. Sic hic Paulus: Liberi eratis iustitiae, id est, derelicti a iusti- cia, non pertinebatis ad regnum illius-Sumus hodie multi liberi iustieia, id est, nolumus habere dominum tusticiam, uerum Sathanam, cui miserrimam seruimus feruitutem. Haec omnia in naturae nostrae condemna tionem dicuntur, que nihil alud facere potest, quam peccare.   \pend
\section*{IN EPIST. AD ROM.  }
\marginpar{[ p.67 ]}\pstart Hic urget conscientias ipsorum, sed tamen ratio- nibus humanis. q. d. quid profuit, quam utilitatem inde habuistis? nunc erubescitis, quoties illa recur- runt uobis in memoriam (quemadmodum et nos ho- die, nonne pudet fuisse in tanta caecitate et errore?) Quid igitur infertis: Ergo peccabimus, cum erube- scatis? Esset hoc regnum iusticiae. ILLORVM.  Scilicet, quae faciebatis antea. MORS.  Id est, aeterna damnatio. MANVMISSI.  Id est, liberi facti. A PECCATO. Id est, seruitute peccati. IN SANCTIFICA- TIONEM. Scilicet, uestri:iterum Antithesi ex ponit, Scilicet, VITAM. Hanc opponit morti-  \pend
\phantomsection
\addcontentsline{toc}{subsection}{\textit{Quem igitur fructum? }}
\subsection*{\textit{Quem igitur fructum? }}
\phantomsection
\addcontentsline{toc}{subsection}{\textit{Etenim autoramenta. }}
\subsection*{\textit{Etenim autoramenta. }}\pstart Id est, praemium, stipendium, quo aliquis obno- xius fit seruituti: non potest aliud praemium confequi peccatum, quam mortem. In quacunque die comederis.   \pend
\phantomsection
\addcontentsline{toc}{subsection}{\textit{Donum. }}
\subsection*{\textit{Donum. }}\pstart Non dicit, bona opera sunt caussa aeternae uitae, sed Det donum, quia remissa sunt nobis pec- cata per Christum mediatorem no- strum. Haec omnia docent, non porro peccandum.   \pend
\section*{COMMENT. POMER.  }
\endnumbering\beginnumbering\section{CAP.VII.}\pstart \huge\textbf{}\normalsize Oc caput, cum multis fuerit difficile intelle- ctu, paucique post Apostolorum seculum sen- tentiam illius attigerint, satis indicat, solum Christum habere clauem scientiae, qua nobis omnia aperit et reuelat. Reuelata enim ratione fidei, simul omnia quoque mysteria regni coelorum aperuit. Esto, sint aliqua in Scrip. quae non satis intelligam, ex clarto ribus locis petam. Sunt uero aperta et clara uerba Capitis, si quis uideat, eademque res multis uerbis re- petitur, tametsi docti quidam nostris temporibus suis tur pars Capitis similitudinem de contugio habet, qua, quae dicta sunt in superiori Cap. declarantur. Estque ea fere sententia. Quemadmodum mulier, mortuo ma- rito, libera est a iure illius, ita, ut et liceat ipsi, al- tert nubere: Sic nos per Christi mortem, legi quoque  sumus mortui, et eius iuri, id est, liberi sumus a le- ge, insitique Christo, hoc est, nouo cuidam marito, nempe Christo nupsimus, et in eius regnum, quod iusticiae est, translati sumus. Decet igitur, ut iam, quae nouo marito, non mortuo, placcant (quid enim illi iam uita functo placeret?) faciamus, utpote iusti- ciae fructus. Quid igitur stulte causabitur caro, li- beri sumus a lege, ergo peccandum? Sed haec textus clarius dicet.  \pend
\vspace{0.5cm}\noindent
\vspace{0.2cm}\rule{1cm}{0.2pt}\\ 
\hspace{0.2cm}\textit{mg}
\footnotesize Mors Christi a lege liberat 
\normalsize| 
\section*{IN EPIST. AD ROM.  }
\marginpar{[ p.68 ]}
\phantomsection
\addcontentsline{toc}{subsection}{\textit{Scientibus. }}
\subsection*{\textit{Scientibus. }}\pstart Id est, illis qui tenent et intelligunt legem, qua- les quoque tum erant Iudaei per gratiam Det Christiani facti  \pend
\phantomsection
\addcontentsline{toc}{subsection}{\textit{Quoad. }}
\subsection*{\textit{Quoad. }}\pstart Non est haec sententia, sed quoad homo is uixe- rit. Lex quid sit, uel ex superioribus intelligere po- tes. Est sententia quaedam generalis, in uiuentes tan- tum, non in mortuos lata, sicuti furibus uiuentibus praecipitur: Non furaberis, uita functis, non potest praecipi. Hic habes exemplum de lege coniugii, qua dicitur: Sub uiri potestate eris, qua mulier premi- tur, quamdiu in uiuis est maritus.   \pend
\phantomsection
\addcontentsline{toc}{subsection}{\textit{Adultera iudicatur. }}
\subsection*{\textit{Adultera iudicatur. }}
\phantomsection
\addcontentsline{toc}{subsection}{\textit{Scilicet, per legem. ALTERI. Scilicet, adultero.  }}
\subsection*{\textit{Scilicet, per legem. ALTERI. Scilicet, adultero.  }}
\phantomsection
\addcontentsline{toc}{subsection}{\textit{Itaque fratres. }}
\subsection*{\textit{Itaque fratres. }}\pstart Hic similitudinem accommodat: uideti haec simili- tudo mirabilis et non satis adposite adducta, cum ta- men sit adpositissima. Sic enim habet, cum sumus mor- tui, per mortem Christi, etiam lex et peccatum no- bis mortua sunt: Sicut, quando maritus mortuus est, mortua est et mulier iuri illius, non haec amplius est legi obnoxia, cum et ipsa lex, qua dicebatur, sub uirt potestate eris, sit cum marito mortua, neq iam potest illi progignere liberos, nupta alteri, nouae cuipiam legi est mancipata, prioris oblita, facit ergo quaecunque uiuus uult, non mortuus ille,  \pend
\section*{COMMENT. PO MAER.  }\pstart et dans liberis procreandis operam, iam nec priori, sed isti nouo parit. Pari ratione Christiana uita habet, peccatum, lex, nobis mortua sunt, et nos illis per Chri stum. Liberi igitur sumus a lege, tanquam priore ma rito, at quia marito carere non poteramus, elegi- mus nobis multo beatiorem priore, Christum Iesum, huius legi nos subiecimus, cum illo unum facti su- mus, et inspiramur uno spiritu. Par igitur est, ut illi quoque sobolem proferamus, nempe iusticiam et fru- ctus spiritus, non respicientes in demortuum mari- tum, legem. Quid ergo potuit efficacius dici ad prae sentem locum, quo monet, Christianis non peccandum.  \pend
\phantomsection
\addcontentsline{toc}{subsection}{\textit{Mlortificati estis legi. }}
\subsection*{\textit{Mlortificati estis legi. }}\pstart Id est, nullum ius iam habet in uos damnandi.   \pend\pstart Illi enim insiti estis, et ablati a corpore pecca- ti, iam e carcere et tenebris, in lucem producti estis: habebatis uirum mortis, et tenebrarum principem, iam autem autorem lucis, Christum. NIMI- RVM EI. Est periphrasis Christi.   \pend
\phantomsection
\addcontentsline{toc}{subsection}{\textit{Per corpus. }}
\subsection*{\textit{Per corpus. }}
\phantomsection
\addcontentsline{toc}{subsection}{\textit{Fructificemus Deo. }}
\subsection*{\textit{Fructificemus Deo. }}\pstart Id est, Deo pariamus iam sobolem, cui iuncti su- mus, obliuiscamur prioris mariti, id est, legis, quae nos damnabat.   \pend
\phantomsection
\addcontentsline{toc}{subsection}{\textit{Cum enim essemus. }}
\subsection*{\textit{Cum enim essemus. }}\pstart Exponit diligentius, quod dixerat, nos legi mor-  \pend\pstart  \pend
\vspace{0.5cm}\noindent
\vspace{0.2cm}\rule{1cm}{0.2pt}\\ 
\hspace{0.2cm}\textit{mg}
\footnotesize Liberi sumu a lege. 
\normalsize| 
\section*{IN EPIST. AD ROM.  71 }
\marginpar{[ p.]}\pstart tuos. Duplicem uitam describit, alteram carnis, id est, ueteris hominis: alteram sptritus, id est, nout ho- minis. Haec ad peccatum pertinet, illa ad iusticiam: in carne esse, est spiritu nondum renatum esse, Ioham . 3.   \pend\pstart Id est, irritantur per legem. Primum, ipsi adfe- ctus per legem, quam mali sint, indicatur. Secundo per illam, incipiunt insantre, adeo, ut irritentur ma- gis peioresque fiant, tantum abest, ut iustificare pos- sit lex.   \pend
\phantomsection
\addcontentsline{toc}{subsection}{\textit{Qui sunt. }}
\subsection*{\textit{Qui sunt. }}
\phantomsection
\addcontentsline{toc}{subsection}{\textit{Affectus peccatorum. }}
\subsection*{\textit{Affectus peccatorum. }}\pstart .Hic uides quam stulte Hypocrisis in externa ope- ra respiciat, quae dum simulat, putat se satisfacere le- gi Sed Paulus hic adfectus peccatorum urget, osten ditque concupiscentiam ex lege, haec non est in manu mea, sed in corde, uiget in tota nostra natura, per Adam corrupta. Haec libentius dico propter eos, qui multa de legis abrogatione scripserunt, cum tamen nibil horum, quae hic scribit Paulus, intellexerint, Ideoque frustra legeris eorum commentarios, et maxime Origenis in hanc Epistolam: Non solum ex- ternum adulterium prohibetur, sed internum, quod intus solicitat: ita et de aliis peccatis dicendum.  Operib . non- satis sit legi.  \pend
\phantomsection
\addcontentsline{toc}{subsection}{\textit{Vigebant. }}
\subsection*{\textit{Vigebant. }}\pstart Id est, fortes erant, ut maxime opus externum non perficiamus, tamen ipsi adfectus praesto sunt, et sen- per peccant.   \pend
\vspace{0.5cm}\noindent
\vspace{0.2cm}\rule{1cm}{0.2pt}\\ 
\hspace{0.2cm}\textit{mg}
\footnotesize 
\normalsize| 
\section*{COMMENT. POMER.  }
\phantomsection
\addcontentsline{toc}{subsection}{\textit{Ad fructificandum morti. }}
\subsection*{\textit{Ad fructificandum morti. }}\pstart Id est, priori marito. Hic nos per uim legis, pec- catis foecundabat, ideoque ; non alios fructus, non aliam prolem, q̃ peccatum, ex nobis sustulit. NVNC VERO LIBERI SVMVS. Per mortem Chri- sti, in qua et nos mortui sumus a lege. EI. Legi.   \pend
\phantomsection
\addcontentsline{toc}{subsection}{\textit{In qua detinebamur. }}
\subsection*{\textit{In qua detinebamur. }}\pstart Eramus uincti, non erat ulla libertas, quae iam per Christum data est nobis-Semper enim nos incre- pabat, die noctuque ; hanc uocem in aures nostras in- tonans: Dilige Deum, dilige proximum ex corde, quod cum cerneremus, non posse nos facere, dice- bat : Ergo damnatus es, tra Dei super te. Quid face- remus, tam insuaut et inexorabili marito obnoxii? desperandum erat, nunc uero uix dici potest, quanta lae- titia cor nostrum exundet, cum hunc mortuum uidea- mus, nouoque ́, id est, Christo desponsos. O foelicem hanc maritorum commutationem. SERVIAMVS.   \pend\pstart Scilicet, Deo. PER NOVITATEM SPIRITVS. Id est, non Hypocritice.   \pend
\phantomsection
\addcontentsline{toc}{subsection}{\textit{Non per uetustatem literae. }}
\subsection*{\textit{Non per uetustatem literae. }}\pstart Vt Hypocritae solent. Hic non solum aufert ex- terna peccata, sed etiam hypocrisin, et omne specio- sum. De litera et sprritu, clarus est locus, in. 2. Co- rinth. 3. et supra Cap. 2. Curcuncisio nihil est etc-  \pend
\section*{IN EPIST. AD ROM.  }
\marginpar{[ p.70 ]}\pstart Est uero litera, quicquid a nobis dici, praedicari, fieri potest operibus externis. Sic circuncisionem uocat li- teram, quod nisi de externis operibus intelligas, ab- surdum erit. Quae enim litera est in praepucio, aut quid ibi legitur? Spiritus est, corSpiritu sancto inno- uatum Iohan. 4. In spiritu, id est, non in externis adoratio fit, et in ueritate, id est, non hypocrisi- Praeterea et uetustas ipsa, est litera, id est, quicquid uiribus nostris facimus: uetustatem uocamus, et li- teram, namut si Christus non uenisset, potuissem aut radi, aut non, ieiunare, ire Romam. Cur igitur uenisset Christus, si haec tustificassent me, meis, ut il- li dicunt, uiribus facta? Sed quia haec erant opera tan- tum mortis, non uitae, peccati, non iusticiae: Ideoque ́ necesse erat, aliam iusticiam, id est, Christum exhi- beri, qui non per uetustatem literae, id est, tantum hypocritice impleret, sed sprritu, id est, uere. Nam cum lex diceret: Non furaberis, nos quidem legebamus et audiebamus, sed non implebamus, deerat eni cor purum, quod habere, non erat nostrae potestatis, simulabamus igitur, nos non furari opere externo, cum tamen in corde, cogeremur semper furari: uerum Christus, non tantum au- diuit et legit: Non furaberis, sed ipso corde impleuit, id est, adfectui furti non fuit obnoxius, quemadmodum nos: uicit illum, per hoc quod diligebat proximum, quod nos natura non poteramus, necdum possumus. Sic opponen da sunt uetustas literae et nouitas spiritus. Quid igitur inferemus : Ergo peccabimus?  \pend
\vspace{0.5cm}\noindent
\vspace{0.2cm}\rule{1cm}{0.2pt}\\ 
\hspace{0.2cm}\textit{mg}
\footnotesize Litera et spi- ritus.  
\normalsize| 
\section*{COMMENT. POMER.  }\pstart Ab hoc loco, qui supertoribus cohaeret, usque ad finem Capitis, haec est sententia Pauli. Caro dicit, li- beri sumus a lege, ergo licet quiduis facere, darem inquit Apostolus, hoc tibi, si totus uetus Adam mor- tuus esset in nobis. Ex hoc loco commodissime intelli ges carnis uetustatem et gratiam, nos quidem, qui spiritum Dei habemus, ut uelimus esse absque peccato, et dolemus uehementer in nobis esse peccatum, sen- tirique malos adfectus, liberi esse uelimus, sed fieri ne quit, quia nondum sumus perfecte mortui, quantum ergo in nobis peccati et ueteris Adae, tantum habet adhuc in nos imperii lex. Ideoque non perfecte a lege sumus absoluti, quia non sumus perfecte mortui: at haec nonne contradicunt supertoribus? nequaquam.  Per gratiam enim Dei, id est, quia Deus mihi fauet, et quod sciam me illius esse filtum, sum saluatus, non per hoc, quod non sentiam peccatum. Accidit saepe, ut filius in domo patris peccet, non tamen eiicitur, ut seruus. Sunt enim ii gradus, ut dicam, gratiae Det, quae in nos peruenit. Primo, suscepti sumus, eum impii essemus, remissaque nobis omnia peccata- Secundo, spiritum accepimus, quo bene et recte ex- corde uelle coepimus. Tertio, quod non imputetur re- ltquum, quod adhuc in nobis est peccatum, quia spi- ritu contranitimur peceato, et illud odimus. Atq baec summa est gratiae. Liceret igitur nobis quiduis fa  \pend
\phantomsection
\addcontentsline{toc}{subsection}{\textit{Quid ergo dicemus? }}
\subsection*{\textit{Quid ergo dicemus? }}\pstart  \pend
\vspace{0.5cm}\noindent
\vspace{0.2cm}\rule{1cm}{0.2pt}\\ 
\hspace{0.2cm}\textit{mg}
\footnotesize Legis tyran- nis sublata.  
\normalsize| 
\section*{IN. EPST. AD ROM.  }
\marginpar{[ p.69 ]}\pstart cere, si uetus Adam totus esset mortuus in cruce, at quia mortuus non est, custodiendus est diligenter, ne de- cruce descendat, siquidem nihil magis optat. Est nobis assidua cum illo pugna per totam uitam: non solum ad Christum, sed ad nos quoque dicunt Pharisei: sifi- lius Deies, descende de cruce. Nos dicemus, quia fi- lii Dei sumus, manebimus in cruce, quamdiu caro no- stra uixerit. Hic certandum est, hoc opus, hic labor: uult libera esse caro, non suffigi in crucem, hocra- tio, hoc doctrina Sathanae clamat, sed non est au- scultandum illis, in finem perdurandum est, donec to tus moriatur Adam, tum licebit, ut de cruce tollatur: non tamen tunc descendet, quia mortuùs est, sed quia donata est nobis uita aeterna: non per nos, sed per Deum, quare non quae nobis, sed Deo placent, face re debemus-Videbis igitur ex hoc loco, totam illam pugnam carnis et spiritus, quae durat per totam ui- tam. Habesque eam et in Gal. 5. multa enim uolu- mus secundum carnem, quae non sunt facienda, uo- lumus et quaedam secundum spiritum, quae sunt fa- cienda, ex quibus et superior, item, praesens locus intelligi potest, nam inferebant prumum, quiduis fa- ciendum est, quia sumus sub gratia, Recte: quaecun que ex spiritu et gratia Dei uolumus, sunt facien- da, quae uero inde non fluunt, minime sunt facienda.  Similiter et hic ex opposito gratiae inferunt, quia a lege sumus liberi, ergo quiduis faciendum. Respond  \pend
\vspace{0.5cm}\noindent
\vspace{0.2cm}\rule{1cm}{0.2pt}\\ 
\hspace{0.2cm}\textit{mg}
\footnotesize Spiritus pug- na cum carne.  
\normalsize| 
\section*{COMMENT. POMER.  }\pstart det Paulus: Si totus Adam uester mortuus esset, id est, si nihil eorum uelletis, quae caro uult, affirmarem uos liberos a lege. Sed id conscientia uestra teste, nondum adsecuti estis. Est autem uix clarior locus in Scripturis, qui gratiam nobis commendet, et aperiat, quam is est quo discimus, nos iustos esse Det gratia, quantum- uis etiam interim in nobis, aestuet peccatum, quod non uideremus, nisi lex id nobis ostenderet. Quod au- diens caro, iterum dicet: Ergo lex malum quiddam est, nam uiget in affectibus, et irritat eos. Respon- det Apostolus: Lex Dei est, non potest esse mala, sed ipsum peccatum, in nobis est peccatum, quod non cerni- mus sine lege, et quemadmodum aegrotus, cut medi- cus praecipit, ne suilla uescatur, ex ipso praecepto, su- mit occasionem appetendi suillam, quod antea non fe- cerat. Vbi uides concupiscentiam esse illud malum, quod magis per praeceptum in se bonum excitatum est : simili modo nos per legem peccandi occasionem arri- pimus, et quod antea peccatum esse non putasse- mus, uidemus esse peccatum per legem, quae bona est: Nitimur in uetitum etc. Quis enim sciuisset con- cupiscere alterius rem esse peccatum? sanctos nos ha- ctenus esse putauimus, quia faciebamus externa bo- na opera, non uidimus concupiscentiam esse pecea- tum: nunc cum reuelatur Verbum, quod iudicat om- nem carnem, et uires nostras, relicta fiducta ope- rum desperamus: nam uidemus, quod ooncupiscen-  \pend
\vspace{0.5cm}\noindent
\vspace{0.2cm}\rule{1cm}{0.2pt}\\ 
\hspace{0.2cm}\textit{mg}
\footnotesize Gratiae com- mendatio.  
\normalsize| 
\section*{IN EPIST. AD ROM.  }
\marginpar{[ p.72 ]}\pstart tia sit uerum peccatum-Clamamus: Lex Dei, Ver- bum Dei, malum est, quia facit nos peccatores. Ah- stultos nos et miseros, nonne morbum tantum dan- nat, et malum esse dicit aegrotus, et gaudet de me- dici tudicio, quod agnoscit morbum? Pudeat nos no- strae stultitiae, qua Deum accusamus, aperit nobis ex sua misericordia peccatum, in bonum, per legem (quae nihil aliud est, quam conscientia reuelata) quod si nobis non bene uellet, non aperiret, relin- queret nos nostro sensui, sicque fieret, ut nunquam agnosceremus ipsum patrem, auxilioque eius nobis opus esse, uerum aliud non possumus natura, quam indignari Deo et legi eius, cum exigit, ut sim pu- ro corde, candidoque ; adfectu in proximum, nam id non in se inuenit: Itaque blasphemat bonitatem Dei, qui impossibilia exigat. Sic lex uitio nostro mala est, in se sanctissima. Sed haec interim Iusticia- rii et sancti illi homines, non uident, quibus uela - men Mosi, per Spiritum Christi, a facie non est amo- tum. Caue haec Philosophice interpreteris, laba- risque in crassiores Cymeriis tenebras . Est ergo Oc- cupatio, cum dicitur: Quid dicemus, lex pecca- tum est? Cui Apostolus respondet, Absit: imo non est peccatum, sed cognitio tantum peccati. Sed in- terim et hoc obserua, duo esse hominum gene- ra, que male de lege tudicant: Alterum a dextris.  \pend
\vspace{0.5cm}\noindent
\vspace{0.2cm}\rule{1cm}{0.2pt}\\ 
\hspace{0.2cm}\textit{mg}
\footnotesize Lex sanda. rationi mala uidetur.  
\normalsize| 
\section*{COMMENT. POMER.  }\pstart quod est Iusticiariorum, qui legi iusticiam tribuunt, atque ubl eis dicitur, lex auget peccatum, tantum ab- est, ut iustificet, inferunt: Ergo est inutilis, si non iusti- ficat. Alterum est impiorum, a sinistris, qui ubi au- diunt nos liberatos a lege, dicunt: Ergo peccabimus, Praeterea, ut maxime ostendat peccatum, tamen dan- nat nos, ergo est mala et peccatum quibus dicitur, non probari peccatum, quod ostenditur per legem, sed cognitionem peccati: quemadmodum in aegrota- res habet, morbus est malus, at cognitio morbi bo- na, uerum uide haec in textu.   \pend
\phantomsection
\addcontentsline{toc}{subsection}{\textit{Sed peccatum non cognovi. }}
\subsection*{\textit{Sed peccatum non cognovi. }}\pstart Id est, nunquam ipsum quod uere peccatum est, et radix crassorum peccatorum, in externis operi- bus sensissem, nist id me lex docuisset.   \pend
\phantomsection
\addcontentsline{toc}{subsection}{\textit{Nam et concupiscentiam. }}
\subsection*{\textit{Nam et concupiscentiam. }}\pstart Lex enim non tantum dicit: Non furaberis, id est, manu ne sustuleris aliena, et non adulteraberis, id est, ipso opere carnis ne inquineris. Verum etiam quod omnium maximum est: Non concupiscas rem et uxorem proximi, id est, cor tuum illa non spectet, aut appetat, secundum naturae crassitiem: quis hoc dixisset esse peccatum, cum omnes respiciamus tantum in externa opera? Quid igitur sunt uires et ingeniam rostra, si caput damnationis nostrae, id est, concu- piscentiam in corde, non cernimus et deprehendi-  \pend
\vspace{0.5cm}\noindent
\vspace{0.2cm}\rule{1cm}{0.2pt}\\ 
\hspace{0.2cm}\textit{mg}
\footnotesize Lex ad adfe- ctus pertinet.  
\normalsize| 
\section*{73 IN EPIST. AD ROM.  }
\marginpar{[ p.]}\pstart mus? Ex his tenebris et ignorantia, quis emerget sine Spiritus Dei tudicio, qui solus nobis ostendit, concu- piscentiam esse peccatum. Eant nunc Theologt Ari- stotelici, et iactent suum Aristotelem, lumen na- turae, hic uident omnem suam sapientiam, et acumem pudefieri, tametsi multos libros, de Septem peccatis mortalibus, Decem Praeceptis, et Quinque Sensibus: Item, de Summo Bono, et Beatitudine conscripserint, in iis uere impletum est, quod Dominus in Propheta dixit: Perdam sapientiam sapientum. Nonne perdi- ta est et confusa? quae maior esse potuit Sathanae im- postura, quam peccatum cordis, id est, concupiscen tiam, ex omnium hominum oculis et cognitione tol- li? quapropter laudandus Deus in secula, qui nobis reuelauit hodie, per suum Euangelium (quod est in- terpres legis) concupiscentiam esse peccatum. Sed utinam, id multi nostrum, non tantum lingua et ore confiterentur, uerum et in corde sentirent illam con- cupiscentiam reuelatam: nam uel ex hoc discerent, non abuti libertate Euangelii, sed humiliati et con- fracti spiritu, sese submitterent: id quod non faciunt, qui Scripturam pro sua libidine torquent, ut non et ipsi posteriores in Euangelii ministerio esse uideantur, ad quod eos illa caecitas, et summa impietas cordis urget et impellit. Siquidem Christiantsmi ratio non Qιλαυτία, aut ambitione fastuque ́, sed magis sui con- temptu et ταπενώσι constat: sed uiderint ipsi,  \pend
\vspace{0.5cm}\noindent
\vspace{0.2cm}\rule{1cm}{0.2pt}\\ 
\hspace{0.2cm}\textit{mg}
\footnotesize In schoeerme rios.  
\normalsize| 
\section*{COMMENT. POMER.  }\pstart nam Dominus nos et illos iudicaturus est. Non est res tam aperta et facilis, sentire concupiscentiam esse pec- catum: multi desperare uolebant, ubi in externa pec- cata incidissent, usq adeo diuexabantur eorum con- Icientiae: quis horror, quis luctus, quis cruciatus con- cientiam uexabit, ubi fontem malorum et peccatorum omnium, ipsum cor, in quo residet tanq̃ regina, con- cupiscentia, esse perspexeris? grauia sane sunt pctam, fructus externi concupiscentiae, sed multo grauius, est ipsa concupiscentia horum mater. Quae spes ergo nobis poterit ulla esse salutis, si non primum condisca- mus matrem, id est, concupiscentiam damnare? fru- stra te adflixeris flagris propter filios, qui te solicitant, id est, externa peccata, si non primum matrem ip- sam caedas, faciasque ́, ne tales liberos pariat. Ridicu- lum est, si quispiam arborum suarum, quas forte in horto habet plurimas, fructus, cum mali sint, et pu- tres natura, damnet, si non magis ipsas arbores euel- lat, excindatque ́: Hoc est, quod hic Paulus tantopere sudet in descriptione, ueri peccati, quod nec natura, nec hypocrisis uidet. Voluit ille, nam ipsum uereba- tur futurum, ne, ubi tantum in fructus peccati defi- xissemus oculos, in hypocrisin prouolueremur: idque ́ hodie, nostro periculo experti sumus. Vide quot sectas Iusticiariorum et Hypocritarum, haec caecitas protu- lerit, omnes tantum in externa opera respiciebant, nemo in cordis malitiam. Quisque conabatur effuge-  \pend
\vspace{0.5cm}\noindent
\vspace{0.2cm}\rule{1cm}{0.2pt}\\ 
\hspace{0.2cm}\textit{mg}
\footnotesize Concupiscen- tiae peccatum.  
\normalsize| 
\section*{IN EPIST. AD ROM.  74 }
\marginpar{[ p.]}\pstart re et uitare crassa peceata, hoc ubi speciosissime fie- bat, statim se sanctum et dignum, qui coronaretur a Deo, putabat: at si quis iis neglectis, in cor spectas- set, illudque neutiquam mutabile suis uiribus uidisset, non equidem, in has laruas iusticiae carnis lapsus es- set, uerum desperasset, opemque ; tantum benigni pa- tris implorasset. Haec obiter propter infirmitatem no- stram monut, quae non ueritatem ipsam, sed umbras tantum metitur, sed ad textum redeamus.   \pend
\phantomsection
\addcontentsline{toc}{subsection}{\textit{Sed occasione accepta, peccatum. }}
\subsection*{\textit{Sed occasione accepta, peccatum. }}
\phantomsection
\addcontentsline{toc}{subsection}{\textit{Scilicet, quod erat in me, et latebat in corde meo, me ignorante- GENVIT. Id est, excitauit.  PER PRAECEPTVM. Scilicet, illud: Non concupisces.  }}
\subsection*{\textit{Scilicet, quod erat in me, et latebat in corde meo, me ignorante- GENVIT. Id est, excitauit.  PER PRAECEPTVM. Scilicet, illud: Non concupisces.  }}
\phantomsection
\addcontentsline{toc}{subsection}{\textit{Omnem concupiscentiam. }}
\subsection*{\textit{Omnem concupiscentiam. }}\pstart Id est, Tum coepi habere conscientiam de peccato concupiscentiae, nam antea nesciebam esse peccatum, putabani tantum per legem opus externum prohi- beri, at cum uidi et adfectus damnari, coepi eam odio persequi. Hic ridendi ueniunt magni illi scri- pturae interpretes, qui dixerunt tantum Cere- moniarum legem per Christum abrogatam.  Quod esset hoc Christi beneficium? inutile sane: per Christum omnis lex est abrogata, et nisi ho- lie quoque in nobis abrogetur per fidem, detinebimur  \pend\pstart  \pend
\section*{COMMENT. POMER.  }\pstart in illius damnatione perpetuo. Lex dicit: Non concu- pisces, hanc legem, nisi in te mortuam uideas et ab rogatam per Christum (quod fit, ubi integro et can- dido es animo in proximum, et in quem si fieri pos- set te totum uelis profundere omni propensione et of- ficio) semper odio habebis Deum, dicesque ; ipsum ma- lam tulisse legem.   \pend
\phantomsection
\addcontentsline{toc}{subsection}{\textit{Siquidem absque lege. }}
\subsection*{\textit{Siquidem absque lege. }}\pstart Id est, cum mihi lex nondum esset reuelata in cor- de, quae non est lex literae: lex enim tum est uere lex, quando mihi est lex, id est, cum reuelat mihi iram Dei in corde propter peccatum.   \pend
\phantomsection
\addcontentsline{toc}{subsection}{\textit{Peccatum erat mortuum. }}
\subsection*{\textit{Peccatum erat mortuum. }}\pstart Non agnoscebatur per me, nec eram eius mihi con- scius, nempe peccati concupiscentiae.   \pend
\phantomsection
\addcontentsline{toc}{subsection}{\textit{Ego autem uiuebam line lege quondam. }}
\subsection*{\textit{Ego autem uiuebam line lege quondam. }}\pstart Id est, absque reuelatione peccati, uide in Philip- et Gal. Viuebam inquam in carnis iusticiis, eramque ́ optimus Iudaeus, et in speciemuir sanctus. Attamen id fiebat sine lege, non uidebam illos malos cordis ad fectus in me. Tales nos hactenus fuimus Hypocritae, putabamus legem tantum externa opera bona exige- re, non cor integrum et purum: quisq in seipsum dirigat aciem, et intelliget se olim non habuisse ullam conscientiam de concupiscentia.   \pend
\section*{75 IN EPIST. AD ROM.  }
\marginpar{[ p.]}\pstart Veniente mandato. Id est, ubi reuelata est mihi lex in corde.   \pend
\phantomsection
\addcontentsline{toc}{subsection}{\textit{Reuixit peccatum. }}
\subsection*{\textit{Reuixit peccatum. }}
\phantomsection
\addcontentsline{toc}{subsection}{\textit{Ego uero mortuus sum. }}
\subsection*{\textit{Ego uero mortuus sum. }}\pstart Id est, agnoui per legem peccatum et damnatio nem: antea sancte uiuere uidebar, cum sine lege ui- uerem, nunc reuelata lege, non uiuo, sed mortuus sum, id est, uideo me Verbi Det esse contemptorem. Nam cum lex dicit: Dilige ex corde, non possum- Consequitur igitur me esse mortuum, id est, nihil ha- bere iusticiae, et plane damnatum esse.   \pend
\phantomsection
\addcontentsline{toc}{subsection}{\textit{Et repertum est mandatum, quod institutum erat ad uitam , mihi ce- dere ad mortem. }}
\subsection*{\textit{Et repertum est mandatum, quod institutum erat ad uitam , mihi ce- dere ad mortem. }}\pstart Id est, lex, quae praecipiebat, ut cauerem malum, ut Adamo dictum est: Quacunque die comederis, morte nos reos fa- morieris, malum prohibet, et exigit bonum, sed ui- cit.  re hoc faciendi non praestat, et quo magis id urget in nobis, hoc magis damnatio nostri ob oculos nobis ponitur: hinc recte sceptrum exactoris uocat eam Pro- pheta: Exigit bonum a nobis, atqui nihil boni in no- bis est, prohibet malum, et nihil aliud sumus, quam malum. Conuincimur ipso Chirographo peccati, id est, testimonio conscientiae nostrae, nos non facere bo  \pend\pstart  \pend
\vspace{0.5cm}\noindent
\vspace{0.2cm}\rule{1cm}{0.2pt}\\ 
\hspace{0.2cm}\textit{mg}
\footnotesize Lex omnes 
\normalsize| 
\section*{COMMENT. POMER.  }\pstart num. Quo palam fit, ut lex, ex qua uitam speraram (nam Iusticiarii eam quaerunt in illa) me occidet et non uiuificet: mandando nunquam efficere potest, ut quod deliquimus, opere aliquo bono restituatur.  Opus est ergo ad uitam non legis operibus, quae face- re non possum, sed condonatione peccati et gratia Dei. Quemadmodum homicida quispiam, quomodo resuscitabit occisum suis manibus? Ab eo peccato, bic coram mundo liberari nequit, nisi ipsi homicidium condonetur, ita maliun nostrum quod facimus, nunq per legis opera expiabimus, imo per eam magis dan- nabimur.   \pend
\phantomsection
\addcontentsline{toc}{subsection}{\textit{Nam peccatum. }}
\subsection*{\textit{Nam peccatum. }}\pstart Experientia doctus sum, meo uitio mandatum mi- hi cessisse ad mortem.   \pend
\phantomsection
\addcontentsline{toc}{subsection}{\textit{Per praeceptum. Id est, legem in corde reuelatam.  }}
\subsection*{\textit{Per praeceptum. Id est, legem in corde reuelatam.  }}
\phantomsection
\addcontentsline{toc}{subsection}{\textit{Decepit me. }}
\subsection*{\textit{Decepit me. }}\pstart Non lex me decepit, sed meum peccatum, meum malum et uitium, tametsi ipsum praeceptum, occa- sionem decipiendi me, peccato dederit, in hoc ta- men non lex, sed ego sum accusandus. Vt aeger occa- sionem peccandi, id est, edendi quae non conducunt, ex mandato medici praescripto accipit, ita nostrum peccatum, per legem irritatum insolescit.   \pend
\vspace{0.5cm}\noindent
\vspace{0.2cm}\rule{1cm}{0.2pt}\\ 
\hspace{0.2cm}\textit{mg}
\footnotesize Lex occasio peccati.  
\normalsize| 
\section*{IN EPIST. AD ROM.  }
\marginpar{[ p.70 ]}\pstart Scilicet, peccatum, occidit me lex, quia ipsa est uirtus peccati, habetque ius occidendi peccatores.   \pend
\phantomsection
\addcontentsline{toc}{subsection}{\textit{Et per illud. }}
\subsection*{\textit{Et per illud. }}\pstart Quia Dei est, est enim peccati agnitio, uetatque ́ malum, ideoque mala esse nequit.   \pend
\phantomsection
\addcontentsline{toc}{subsection}{\textit{Itaque lex sancta. }}
\subsection*{\textit{Itaque lex sancta. }}
\phantomsection
\addcontentsline{toc}{subsection}{\textit{Ergo quod bonum. }}
\subsection*{\textit{Ergo quod bonum. }}\pstart Alia est Occupatio, sic caro infert: Lex cum sit bo- na, mihi facta est mors et damnatio, quae sane non est nisi pessimares, ergo non potest esse, aut dici bona, Respondet Apostolus  \pend\pstart Quasi diceret, impietatis haec sunt argumenta, et Dei blasphemia.   \pend
\phantomsection
\addcontentsline{toc}{subsection}{\textit{Absit. }}
\subsection*{\textit{Absit. }}
\phantomsection
\addcontentsline{toc}{subsection}{\textit{Imo peccatum. }}
\subsection*{\textit{Imo peccatum. }}\pstart Est mihi mortis caussa, per legem, quae bona est, in hoc, dum ostendit mihi ipsum peccatum, quod latebat in corde, me nesciente: uoluit lex prohibere et ostendere mihi concupiscentiam, dumque id egis- set, coepi insanire, cognouique per eam me damnari, tametsi id non agebat primum, at magis ut peccatum cognoscerem: meo uitio ergo factum quod lex ma- la uideri poterat. Quemadmodum culter, qui fa- ctus est in bonum usum, scindendi panis, bonus est, at si eo abutar, me hoc occidendo, malus di- ciposset, cum tamen malus non sit. Sic de legis  \pend
\vspace{0.5cm}\noindent
\vspace{0.2cm}\rule{1cm}{0.2pt}\\ 
\hspace{0.2cm}\textit{mg}
\footnotesize Peccatum mor- tis caussa.  
\normalsize| 
\section*{COMMENT. POMER.  }\pstart officio iudicandum erit, non est haec mala, at si abu- tar, mala uideri potest. Est enim finis et officium eius, reuelare peccatum, non damnare quenquam.  At si te damnet, non fit, nisi tuo uitio et caecitate, qui non vides sub ipso ossicio legis, id est, revelatione peccati, iucundibimum quiddam latere,  nempe ur- geri ad gratiam. Qui igitur sic lege utuntur, sancte utuntur, ut uides in Dauide, Ezechia, item Petro.  Qui autem tantum in damnationem figunt oculos, non obseruantes gratiam, sub reuelatione abscondi- tam, ii pereant necesse est, nam lege abutuntur, ut in Iuda, Saule, Cayno cernimus: non enim in hoc suis reuelat peccatum Dominus, ut excidant ab eius misericordia, sed ut discant misericordiam agnoscere.  Itaq ait recte Apostolus: Ipsum peccatum est mihi operatum mortem per legem, quae hoc faciebat, ut appareret mihi, et notum fieret peccatum Ante hanc reuelationem, non putabam esse tantam uim concupi scentiae, nunc uideo et sentio uehementer potentem, ac prorsus insanientem in carne mea. Hic succumbo, et incido in Dei blasphemias, uelimque legem non es- se: non aliud uideo, quam damnationem, et Dei iram in me, quo magis enim peccatum consydero, eo fit fortius et potentius. Hoc necesse est, omnes homines sentiant, dum reuelatur eis peccatum, et nisi monen te Dei Spiritu, obseruarint, sub illa reuelatione sub- esse gratiam, ficri non potest, ut saluifiant, nam de  \pend
\vspace{0.5cm}\noindent
\vspace{0.2cm}\rule{1cm}{0.2pt}\\ 
\hspace{0.2cm}\textit{mg}
\footnotesize Legis officium.  
\normalsize| 
\marginpar{[ p.K5 ]}
\marginpar{[ p.]}\pstart glutiteos infernus, et uincit magnitudo peccati.   \pend
\phantomsection
\addcontentsline{toc}{subsection}{\textit{Scimus quod lex spiritualis est. }}
\subsection*{\textit{Scimus quod lex spiritualis est. }}\pstart Caussam reddit superiorum, cum dixerat, non conuenire inter legem et homines. Sed interim non praetereunda est ineptissima quorundam huius loci in terpretatio. Nam sic exponebant: Lex est spiritua- lis, id est, libri Sacri debent intelligi spirituali- ter, id est, Allegorice: similiter, cum Mosen expo- nerent: Creauit Deus coelum, id est, bonos, et ter- ram, id est, malos. Quae sunt haec portenta? nonne κυβιςwν κματαιολόγαν, quibus non exponitur, sed destruitur uniuersa Scriptura. Haud aliud conantur nostri quidam noui spiritus hodie, quibus quoque uolupe est, inuertere quiduis in Scrip. Erat simplex Gran- matica, et uis uocabulorum obseruanda, at quia il- li tantum spiritus, non Grammatici esse uolunt, ideo contemnunt et nihilifaciunt huiusmodi minutias. Por ro, si quis hunc locum hoc pacto interpretabitur, quid consequetur fructus ex toto hoc Capite? Imo, totius Epistolae lectione frustrabitur, quare mirum non est hanc semper obscuram fuisse, quippe, cum ipsi qui de- bebant lucem addere, interpretatione aperta, offunde rint tenebras. Verum haec est sententia: lex est spiri- tualis, id est, requirit spiritualia, non carnalia, aut externa, non impletur, quando opere ipso abstines ab alienare, aut uxore proximi, sed requirit spiritum,  \pend
\section*{COMMENT. POMER.  }\pstart id est, ne quid tale in corde tuo sentias, quo cupias rem alienam, atque puro et integro sis adfectu in omnes: id ubi natura, quia carnalis es, audierit, terretur, et necesse est, ut emoriatur. Proinde, si spiritus im perio cogatur obsequi, crucifigitur, inuitaque ; quicquid est fpiritus legis fert. Haec ergo est illa pugna carnis et spiritus, quam nemo non sensit sanctorum, pro- ponitque seipsum exemplum: non autem audimus, qui dicunt: non loquitur de se Paulus, sed in persona car- nalis hominis. Quid aliud hii faciunt, quam quod sa- xa et truncos ex sanctis facere uolunt, et quasi nun- quam senserint peccatum? Sed apertus est Textus, clare uidebis de carnali homine non posse intelligi- Ego uiuebam sine lege: Item, in me repertum est et c- Quare non nisi de Paulo intelligitur hic locus, et in eo de omnibus, in cuius exemplo, nos aestimare pos- sumus reliquos sanctos, cum sumus omnes eiusdem massae: non carent illi peccato, dum in hac carne degunt.   \pend
\phantomsection
\addcontentsline{toc}{subsection}{\textit{At ego carnalis sum. }}
\subsection*{\textit{At ego carnalis sum. }}\pstart Id est, non sum Dei amans, sed mihi, quaero mea, cum lex iubeat, ut non nostra, sed aliorum et Dei quaeramus, quare meo uitio, lex mihi mala est .   \pend
\phantomsection
\addcontentsline{toc}{subsection}{\textit{Venditus sub peccatum. }}
\subsection*{\textit{Venditus sub peccatum. }}\pstart Durissimum dominium sustineo, factus manci- pium : alludit autem ad seruos, qui olim uendebantur, et nunc alicubi, in perpetuam seruitutem. Vides igi-  \pend
\vspace{0.5cm}\noindent
\vspace{0.2cm}\rule{1cm}{0.2pt}\\ 
\hspace{0.2cm}\textit{mg}
\footnotesize Pugna carnis et spiritus.  
\normalsize| 
\section*{IN EPIST. AD ROM.  78 }
\marginpar{[ p.]}\pstart tur, quam multa adhuc sanctis desint. Sanctus erat Paulus, at quid in seipso desyderet, uides uel ipsius te- stimonio.  \pend
\phantomsection
\addcontentsline{toc}{subsection}{\textit{Quod enim ago, non probo. }}
\subsection*{\textit{Quod enim ago, non probo. }}\pstart Declarat hanc seruitutem suam, id est, quicquid ago, hoc mihi displicet, nam uideo nihil fieri sine re- liquiis peccati: nihil mihi conscius sum, tamen in hoc non tustificor.   \pend
\phantomsection
\addcontentsline{toc}{subsection}{\textit{Non enim quod uolo. }}
\subsection*{\textit{Non enim quod uolo. }}\pstart Spiritu. HOC FACIO. Scilicet, carne: nam caro pugnat aduersus spiritum, Gal. 3. Nam quicquid est natum, ex muliere et Adam, peccatum est.  \pend\pstart Id est, nolo secundum spiritum. HOC AGO.  Scilicet, per concupiscentiam, quae est in me, id est, uelim nolim, cogor esse peccator, quod sanctus sum, ex gratia habeo, non meo merito.   \pend
\phantomsection
\addcontentsline{toc}{subsection}{\textit{Quod odi. }}
\subsection*{\textit{Quod odi. }}
\phantomsection
\addcontentsline{toc}{subsection}{\textit{Si uero quod non uolo. }}
\subsection*{\textit{Si uero quod non uolo. }}\pstart Secundum spiritum . HOC FACIO.  Scilicet, carne . CONSENTIO LEGI.  Id est, tum ipsa experientia cordis nostri docemur legem esse bonam, nam spiritualis homo conscientiam pcti- sentit, et uellet abesse, et in hoc quod uelit abesse, consentit legi, id est, dicit legem esse bonam, siquidem cogi- tat, legem debitorem facere: at non facio, tamen mi- hi non desperandum est, sed dicendum quotidie: Dimitte nobis debita nostra.   \pend
\section*{COMMENT. POMER.  }
\phantomsection
\addcontentsline{toc}{subsection}{\textit{Nunc autem non iam ego per- }}
\subsection*{\textit{Nunc autem non iam ego per- }}
\phantomsection
\addcontentsline{toc}{subsection}{\textit{petro illud. }}
\subsection*{\textit{petro illud. }}\pstart Id est, quia ego non consentio peccato, uelim il- lud abesse: mordet enim peccatum illud naturae me, uelut pediculus quidam, quem inuitus sustineo.   \pend
\phantomsection
\addcontentsline{toc}{subsection}{\textit{Sed inhabitans in me peccatum. }}
\subsection*{\textit{Sed inhabitans in me peccatum. }}\pstart Id est, morbus ille naturae, seu quod peccatum uo- cant originale.   \pend
\phantomsection
\addcontentsline{toc}{subsection}{\textit{Noui. }}
\subsection*{\textit{Noui. }}\pstart Id est, ego, qui spiritum Dei habeo, et cuius iam beneficio mihi lex reuelauit peccatum.   \pend
\phantomsection
\addcontentsline{toc}{subsection}{\textit{Quod non habitet in me, hoc est, in carne. }}
\subsection*{\textit{Quod non habitet in me, hoc est, in carne. }}\pstart Video me, quantus quantus sum, esse nihil aliud, quam peccatum. Ecce hic totum hominem damnat, ubi iam inuenies Libus Arb. et praeparationem ad gra- tiam? nunquid ex turbido fonte, puram aquam hau- ries, num ex mala arbore, fructum bonum petes?  \pend
\phantomsection
\addcontentsline{toc}{subsection}{\textit{Nam uelle adest mihi. }}
\subsection*{\textit{Nam uelle adest mihi. }}\pstart Id est, per gratiam quidem id habeo, ut quae bona sunt, uelim facere.   \pend
\phantomsection
\addcontentsline{toc}{subsection}{\textit{At ut faciam bonum non reperio. }}
\subsection*{\textit{At ut faciam bonum non reperio. }}\pstart Id est, in meis uiribus non est : itaque non loqui- tur hic de manu, sed ipso corde, quod non uult ex na- tura bonum faccre, ergo bonum in hac uita est, ut-  \pend
\section*{IN EPIST. AD ROM.   }
\marginpar{[ p.]}\pstart uerc pura conscientia et corde, omni momento, ut in nullo deficias. Sed quis est ille, et laudabimus eum? Non enim quod uolo, scilicet, spiritu, facio bonum, scilicet, in carne.   \pend
\phantomsection
\addcontentsline{toc}{subsection}{\textit{Sed quod non uolo, malum. }}
\subsection*{\textit{Sed quod non uolo, malum. }}\pstart Scilicet, spiritu, hoc ago, scilicet, carne, perpe- tuo reluctor, et crucifigo carnem, neque erit finis col- luctationis huius, priusquam tota emoriatur.   \pend
\phantomsection
\addcontentsline{toc}{subsection}{\textit{Reperio igitur. }}
\subsection*{\textit{Reperio igitur. }}\pstart Duo inquit in me reperio, carnem tolentem ma- lum et spiritum repugnantem carni. Vana sunt igi- tur commenta eorum, qui somniant hic Libus Arb.  tribui homini, quo uelit bonum, cum Paulus non de ullis uiribus, in nobis intelligi uelit, sed ipso Spiritu Dei, qui uult bonum, qui operatur in nobis, et uel- le et perficere, pro sua in nos bona uoluntate. Prae- terea insaniunt et illi, qui hic sensualitatem fingunt, quae tantum uult malum. Nonne est hoc expuere in fa- ciem Iesu, et alapis caedere? Totum hominem, quan- tus est secundum naturam, dicit nihil boni posse, uelle, siue dicas Libus Arb. siue rationem. Ingenium et quicquid est in homine, omnia sunt maledictioni obnoxia, quaecunque cogitas, uis, aut contendis ex na- tura tua, aut sapientia carnis: Ait enim Paulus: Re- perio, id est, perspicio tandem per legem Dei, quae mihi bene uult (nam tantum in hoc mihi reuelat pec- Libus Arb.  Natura male dicta est.   \pend
\vspace{0.5cm}\noindent
\vspace{0.2cm}\rule{1cm}{0.2pt}\\ 
\hspace{0.2cm}\textit{mg}
\footnotesize 
\normalsize| 
\hspace{0.2cm}\textit{mg}
\footnotesize 
\normalsize| 
\section*{COMMENT. POMER.  }\pstart catum, ut discam diffidere meis uiribus, et Dei mi- sericordiam inuocare) quod dum bonum facere uolo (quod non ex me, sed spiritu Dei fit) mihi adiunctum sit, id est, innatum et insitum in carnem meam to- tam, malum, id est, peccatum, hoc est, illud tantum me uelle uideo secundum meas uires, quod est contra Deum: et, ut in summa dicam, partim sum spiritua lis, partim carnalis, neutiquam conuenit inter uolun tatem spiritus et uoluntatem carnis.  \pend\pstart Haec non sunt carnalis hominis uerba, sed spiritua- lis, quomodo enim delectaretur in lege Dei, caro? cui maxime displicet, intelligit ergo per internum ho- minem, nouam creaturam, de qua Christus: Nisi quis renatus fuerit ex spiritu sancto.   \pend
\phantomsection
\addcontentsline{toc}{subsection}{\textit{Delectat me lex. }}
\subsection*{\textit{Delectat me lex. }}
\phantomsection
\addcontentsline{toc}{subsection}{\textit{Sed uideo aliam legem. }}
\subsection*{\textit{Sed uideo aliam legem. }}\pstart Id est, est alia vis in mea carne, quae repugnat legi mentis, id est, spirituisancto in corde meo, quae vis me captiuum tenet sub lege peccati, id est, video pec- catum in me mondum esse totum mortuum, vivit, et carnis indulgentia illi adblanditur. Habet caro sua desyderia, quae hic Paulus legem uocat: habet et spi- ritus, inter quae duo, nunq conuenit. Spiritus dicit: Non concupiscas. Caro dicit: Ego concupiscam. Haec cum sic se in nobis habeant, nonne ingemiscemus cum Paulo dicentes?  \pend
\vspace{0.5cm}\noindent
\vspace{0.2cm}\rule{1cm}{0.2pt}\\ 
\hspace{0.2cm}\textit{mg}
\footnotesize Peccatum non plane emori- tur in nobis.  
\normalsize| 
\section*{IN EPIST. AD ROM.  80 }
\marginpar{[ p.]}
\marginpar{[ p.]}\pstart Est exclamatio. EX HOC CORPO.  re, id est, hac carne, ab hoc peccato, nam uideo me- morti per hoc obnoxium fieri: per peccatum enim mors introiit in mundum. Vide quam doleat Paulus, et homo spiritualis est sub lege carnis, uelit ab hac liberart, sic suspirant ad Deum omnes sancti.   \pend
\phantomsection
\addcontentsline{toc}{subsection}{\textit{Miser ego homo. }}
\subsection*{\textit{Miser ego homo. }}
\phantomsection
\addcontentsline{toc}{subsection}{\textit{Gratias ago Deo. }}
\subsection*{\textit{Gratias ago Deo. }}\pstart Antea cum dixerat: Miser ego homo, respexerat Paulus in suum peccatum, hic respicit in gratiam Dei, q. d. utcunque saeuiat damnatio, mors et peccatum, tamen scio me in gratiam Dei, susceptum per Chri- stum: quod ad me attinet, peccatum non possum su- perare, cogor desperare: uerum ubi uideo Deum es- se patrem Christi crucifixi, et ita esse patrem illius, ut per hoc quoque meus sit pater, quomodo non gestiam, quomodo non gaudeam et exultem? siquidem iam ui- deo omne illud peccatum esse abolitum et non im- putari. Lubens ergo, quia Deus uult, sustinebo in me peccatum, contranitar, Dei uirtute, non permittam illi ut dominetur. Haec omnia eo spectant, ut intellt- gas, quod dixerat, peccatum in nobis non- dum plane esse mortuum, sed perpetuam pugnam nobis esse cum reli- quiis, quemadmodum lu- daeis olim cum Cha- nanaeis.   \pend
\vspace{0.5cm}\noindent
\vspace{0.2cm}\rule{1cm}{0.2pt}\\ 
\hspace{0.2cm}\textit{mg}
\footnotesize Gratiae uis maior pecca- to.  
\normalsize| 
\section*{COMMENT. POMER.  }\pstart IC longius persequitur, quod antea dixerat- Gratias ago Deo per Christum. q. d. Pec- catum in me est, in quod dum uideo, concutior horro- re, sentioque ; damnationem: at cum in Dei gratiam ocu- los dirigo, uideo nullam damnationem Estque ; hoc pro- prium sanctorum, qui non negant se habere pecca- tum, et peccatores esse, quod magna laetitia spiritus dicant: Nulla est in nobis condemnatio, peccatum non potest nos damnare, ut maxime in nobis sit. Siqui - dem est spiritus Christi in nobis, cut displicet pecca- tum, non posset autem displicere, nisi id ex Deo ha- beremus. Natura ita sumus adfecti, ut oblectatio quaedam sit nobis in peccato, secundum spiritum ue- ro, est summus cruciatus et displicentia. Stulti ergo funt, qui sanctis peccatum adimunt: est uere in illis peccatum, sed sine condemnatione: sunt filii Dei, non serui, ut olim sub Mose, quod in resurrectione pa- lam fiet. Non sunt ergo praedicatores gratiae, sed con temptores, quicunque docent, sanctos absque pecca- to fuisse, contra fidem sacrarum historiarum. Quo- modo enim non resiliremus a Dei misericordia, ui- dentes illos quasi non peccatores, saluos esse factos? Vna sumus cum his massa et caro peccati, quod ut glorificetur Dominus, nos confiteri aduersus Satha- nam oportet.   \pend
\endnumbering\beginnumbering\section{CAP.VIII.}
\vspace{0.5cm}\noindent
\vspace{0.2cm}\rule{1cm}{0.2pt}\\ 
\hspace{0.2cm}\textit{mg}
\footnotesize Sancti quoq peccatores sunt. 
\normalsize| 
\section*{IN EPIST. AD ROM.  81 }
\marginpar{[ p.]}
\phantomsection
\addcontentsline{toc}{subsection}{\textit{Nulla igitur nunc est condemnatio. }}
\subsection*{\textit{Nulla igitur nunc est condemnatio. }}\pstart Quamuis peccatum est reliquum, quod non con- demnat nos, nam ut hostem odimus illud .   \pend
\phantomsection
\addcontentsline{toc}{subsection}{\textit{His qui insiti sunt Christo. }}
\subsection*{\textit{His qui insiti sunt Christo. }}\pstart Id est, qui sunt facti uiuum corpus per fidem, nam ipse caput est, nos membra, cui insiti sumus, ut uxor marito, ubi ex duobus corporibus una fit caro, om- nia item fiunt communia, quae mariti sunt, ad uxo- rem quoque pertinent, et quae sunt uxoris, et in ma- ritum transferuntur. Itaque illam nostri cum Christo unionem, uix rectius capies, quam ex contugii ratio- ne: cum Christo per fidem coniungor, iam iusticia, fortitudo, gloria, uictoria, et quaecumque sunt in Chri- sto, mea sunt, et quae mea, sunt illius, peccatum meum in se transfert, absorbet et consumit omnia mea ma- la, suo bono, nempe meam iniusticiam, sua iusticia tollit, mortem meam, sua uita superat: In summa, qui Spiritum Christi habent, illi sunt eius, quemadmodum inter contuges, unus spiritus et fides, facit unam car nem: Sic qui Domino adhaerent, unum sunt cum eo corpus, sed hoc declarat sequentibus.   \pend
\phantomsection
\addcontentsline{toc}{subsection}{\textit{Qui non iuxta carnem ambulant, }}
\subsection*{\textit{Qui non iuxta carnem ambulant, }}\pstart Non dico, qui non habent peccatum, aut illud non sentiunt, sed qui id habent et uere sentiunt, at non obsequuntur, imo pugnant contra illud, odiunt maxime, iuxta illud quod supra dixit: Non quod uo-  \pend
\vspace{0.5cm}\noindent
\vspace{0.2cm}\rule{1cm}{0.2pt}\\ 
\hspace{0.2cm}\textit{mg}
\footnotesize Vnum sumus cum Christo.  
\normalsize| 
\section*{COMMENT. POMER.  }\pstart lo, hoc ago. Est uero iuxta carnem ambulare, nihil aliud, quam Deum non timere, nullam habere salu- tis, aut Det rationem: quales sunt ettam, qui Euan- gelici hodie uidentur, omnia contemnentes et con- fundentes, cum interim nihil aliud quaerant, quam sua, non quae Dei sunt, aut proximi, ut spiritualis ho mo, qui et Dei et proximi summam habet curam, igitur addit  \pend\pstart Id est, qui legi Dei consentiunt, et dicunt: Lex oris tui, super millia auri et argenti, tametsi caro repugnet illis: sed sciunt hanc uitam esse militiam, in qua nihil aliud expectandum est a carne, q̃ perpetua pugna: insultatio et insidiae Sathanae non quiescunt-  \pend
\phantomsection
\addcontentsline{toc}{subsection}{\textit{luxta spiritum : }}
\subsection*{\textit{luxta spiritum : }}
\phantomsection
\addcontentsline{toc}{subsection}{\textit{Nam lex, spiritus uitae. }}
\subsection*{\textit{Nam lex, spiritus uitae. }}\pstart Duplicem hic legem tangit: Alteram, spiritus ui- tae, id est, spiritus uiuificantis: Alteram literae, id est, mortis et damnationis-Vtilissimus hic locus est imbe- cillibus et adflictis conscientiis, quae eum uident, se- urgert per legem ad peccatum, per peccatum ad mor- tem, ingemiscunt, et fere succumbunt per desperatio nem: tanta est tyrannis illius legis, quae excitat pec- catum, et post peccatum excitatum, reuelat aeternam damnationem, et nisi illucesceret eis Christus (qui spi- ritum patris sui nobis impetrauit) deglutirentur ab in ferno et Sathana. Ihs igitur spiritus, est noua lex,  \pend
\vspace{0.5cm}\noindent
\vspace{0.2cm}\rule{1cm}{0.2pt}\\ 
\hspace{0.2cm}\textit{mg}
\footnotesize Legis terror.  
\normalsize| 
\hspace{0.2cm}\textit{mg}
\footnotesize Spiritus con- folatio.  
\normalsize| 
\section*{ST. AD ROM.  }
\marginpar{[ p.82 ]}\pstart non exigens, non crucians conscientias, sed securas ip sas reddens, dicitque ́. Quid est, quod affligamini? ut ma£ xime peccatum in nobis sit, tamen illud non obest, est remissum, non potest uos occidere. Sic fit, ut ta- les conscientiae ex lege mortis, et ipsa morte, per le gem spiritus uiuificantis, in uitam reuocentur. Ideoque ́ conuenientissimo uocabulo, spiritum uitae, eo quod ui uificet cor, appellat. Itaque cum ait: Quare non est condemnatio in me? Respondet, quia lex spiritus ui uificantis me, et illuminantis conscientiam meam, libe rauit me a iure peccati et mortis, id est, a sententia et lege damnationis, cut debebam esse obnoxius pro pter peccatum.  \pend
\phantomsection
\addcontentsline{toc}{subsection}{\textit{Etenim quod lex. }}
\subsection*{\textit{Etenim quod lex. }}\pstart Hic nequaquam ferendum est, ut transferas, ea par- te, uerum sic. Etenim quod lex praestare non pote- rat, quando (aut si mauis, quia) imbecillis erat etc- Hic locus uel solus satis esset, ad conuincendos Lib- Arb. adsertores, quod lex etiam diuina, imbecillis erat propter carnem et naturam nostram, nihil illa potuit praestare ad salutem . QVOD. Id est, quain iustificationem nostri.   \pend
\phantomsection
\addcontentsline{toc}{subsection}{\textit{Quia imbecillis erat. }}
\subsection*{\textit{Quia imbecillis erat. }}\pstart Hebraice dictum est, id est, non implebatur: exl gebat illa quidem, at uires non suppeditabat, ut im- pleretur . In tertio Cap. contrariae sententiae est Tropus  \pend
\vspace{0.5cm}\noindent
\vspace{0.2cm}\rule{1cm}{0.2pt}\\ 
\hspace{0.2cm}\textit{mg}
\footnotesize Libus Arb.  
\normalsize| 
\section*{COMMENT. POMER.  }\pstart Legem stabilimus, id est, docemus legem impleri per fidem. Ex eo loco rectius et hunc intelliges: tam praesta tur, quod lex exigit, olim uero non praestabatur, cum adhuc essemus sub dominio peccati, et affectuum no- strorum. HOC, Id est, hanc iustificationem praestitit nobis Deus per Christum, quem misitin hunc mundum. Hic uide, quae possit in nobis esse damna- tio, cum satisfactum sit pro peccatis nostris, per Chri- stum, datus est Dei spiritus, quem et ipse emeruit no- bis, huius uirtute implentur omnia, quae lex exigit: peccatum dimittitur, manetque ; in nobis uera iusticia, per solam gratiam non nostrum meritum. Porro et ex hoc loco, quae sit libertas Christiana et legis ab- rogatio, clarum est: non sic abrogata est, ut quid - uis facias. At tum uere est abrogata, quando nihil in te inuenerit, quod accuset aut damnet . Praeterea magna quaedam est misericordissima in nos dilectio, quod proprium filium miserit, nam si quid aliud cha- rius et nobilius habuisset, misisset utique . Est et in hoc aestimanda peccati uis et potentia, quae si alio quo- dam modo, aut per ullam aliam creaturam, aut per nos tolli potuisset, cur ipsum filium misisset? Is igitur solus est nostra iusticia et certissima anchora, ad quam confugiamus.   \pend\pstart His circunstantiis Christi exinanitio, proinde et di- gnitas quam pro nobis habuit, declaratur:nostra, inquit,  \pend
\phantomsection
\addcontentsline{toc}{subsection}{\textit{Sub specie carnis. }}
\subsection*{\textit{Sub specie carnis. }}
\vspace{0.5cm}\noindent
\vspace{0.2cm}\rule{1cm}{0.2pt}\\ 
\hspace{0.2cm}\textit{mg}
\footnotesize Libertatis Christianae usus.  
\normalsize| 
\section*{IN EPIST. AD ROM.  83 }
\marginpar{[ p.]}\pstart carnem Christus habuit ( quae semper est obnoxia pec- cato) attamen cum esset in Christo, absque peccato sen- per fuit. Erat itaque caro illius, nostra caro, habens etiam speciem peccati, tametsi in ea illud non fuerit, tamen propter nos, nostro peccato, et totius mundi, hic unus homo Christus Iesus, grauissime coram Deo obnoxius factus est, cum non eramus soluendo-  \pend
\phantomsection
\addcontentsline{toc}{subsection}{\textit{Ar de peccato. }}
\subsection*{\textit{Ar de peccato. }}
\phantomsection
\addcontentsline{toc}{subsection}{\textit{Scilicet, quod ipse factus est pro nobis . 2. Co- rinth. 5.  }}
\subsection*{\textit{Scilicet, quod ipse factus est pro nobis . 2. Co- rinth. 5.  }}
\phantomsection
\addcontentsline{toc}{subsection}{\textit{Non factus est peccator. }}
\subsection*{\textit{Non factus est peccator. }}\pstart Nam nunquam peccauerat, sed ipsum factus est peccatum, et maledictio, ut nos a peecato libera- remur, ut in Esa. 53.uides: Deus in eo posuit ini- quitates omnium, Et, si posuerit peccatum animam suam, id est, si suam uitam, se totum, quantus quantus est, fecerit peccatum, uel ut sit nihil aliud, quam pec- catum, uidebit semen longeuum etc Ideoque ; ait, de peccato, aut per peccatum, quod ipse factus erat, condemnauit peecatum, scilicet nostrum, idque per carnem, scilicet, suam. In summa, factus est pecca- tum, ut nos innocentia fieremus: factus est mors, ut nos uita essemus: factus est maledictum, ut nos esse- mus benedicty, iuxta promissionem: In semine tuo be- nedicentur omnes Gentes-Sed haec omnia Textus ha- bet. Quare haec factus est Christus?  \pend\pstart  \pend
\vspace{0.5cm}\noindent
\vspace{0.2cm}\rule{1cm}{0.2pt}\\ 
\hspace{0.2cm}\textit{mg}
\footnotesize Christus fa- ctus peccatum. 
\normalsize| 
\section*{COMMENT. POMER.  }
\phantomsection
\addcontentsline{toc}{subsection}{\textit{Vt iustificatio legis impleretur in nobis. }}
\subsection*{\textit{Vt iustificatio legis impleretur in nobis. }}\pstart Hanc iustificationem quidem lex exigebat, sed non praestabat, nam ex operibus legis, non est iusticia: Ergo illa impletio, non est nostrorum uirium, sed per Christum facta.   \pend
\phantomsection
\addcontentsline{toc}{subsection}{\textit{Qui non secundum carnem. }}
\subsection*{\textit{Qui non secundum carnem. }}\pstart Ne quid sapientia, aut temeritas carnis ; hic au- deat, in nobis, inquit, impletur iustificatio legis, qui su- mus assidui in excubiis spiritus, contra insidias carnis, aut qui nihil permittimus carni, uerum omnia per il- lum spiritum, nobis per Christum donatum, fiunt- Nemo ergo carens hoc spiritu, implet legem: praecipi- tur, ut diligas proximum, non potes ex tuis uiribus, potes autem ubi hunc habueris.   \pend
\phantomsection
\addcontentsline{toc}{subsection}{\textit{Nam qui carnales sunt. }}
\subsection*{\textit{Nam qui carnales sunt. }}\pstart Id est, qui non habent spiritum Dei, ea quae car- nis sunt curant. Vides Paulum, non de corpore no- stro, aut carne tantum, quemadmodum quidam som- niant, ubi carnis uocabulum audierint, uerum de co- toto, quod nos sumus, id est, de omnibus nostris ui- ribus, quae sunt in corpore et anima nostra: non in appositum ergo fuerit si legas, procurant, sapiunt, id est, ipso adfectu sentiunt. QVAN CARNIS SVNT. Id est, mundi, nempe carnalia, in qui- bus sese totus homo oblectat.   \pend
\vspace{0.5cm}\noindent
\vspace{0.2cm}\rule{1cm}{0.2pt}\\ 
\hspace{0.2cm}\textit{mg}
\footnotesize Caro quid in Scrip. signifi- cet.  
\normalsize| 
\section*{IN EPIST. AD ROM.  }
\marginpar{[ p.84 ]}
\marginpar{[ p.]}\pstart Id est, spiritum Christi habent, sapiunt ea quae Christi sunt et illius spiritus. Damnantur ergo hic op- timae quaeq intentiones, ut uocant, et omnia nostra quantumuis speciosa studia, ut uerum sit, quod infra dicit Paulus: Quicquid non est ex fide, peccatum est.   \pend
\phantomsection
\addcontentsline{toc}{subsection}{\textit{At qui spirituales sunt. }}
\subsection*{\textit{At qui spirituales sunt. }}
\phantomsection
\addcontentsline{toc}{subsection}{\textit{Nam adfectus carnis, mors est. }}
\subsection*{\textit{Nam adfectus carnis, mors est. }}\pstart Id est, illa prudentia quae est in nobis, ratio huma- na, omnes uires hominis, sunt mors: hoc est, nihil ualent in iis, quae Dei sunt, non potes harum praesi- dio saluari. At quid igitur, utilis ratio humana aut Qρόνκμα? Ad gubernandam Remp. ciuiles mores formandos, colendum agrum, ad aedificandum, quae omnia ita fiant, ne in illis confidas, aut conquiescas- Non enim licet Christiano homini in creaturis pone- re spem suam. Si uero quis uoluerit iis uti, uelut in- strumentis uitae aeternae et salutis, hic opera mortis fa- cit, quae coram Deo non uiuunt, nam hac parte, nihil aliud sunt q̃ mors, humani conatus.   \pend
\phantomsection
\addcontentsline{toc}{subsection}{\textit{Adfectus uero spiritus, uita. }}
\subsection*{\textit{Adfectus uero spiritus, uita. }}\pstart Id est, prudentia et sapientia, quae est ex solo Deo per gratiam, quae et facit nos sapere spiri- tualiter, id est, ea quae Dei tantum sunt, non carnis nostrae, illa inquam, nihil aliud est quam uita et pax, nam per eam, qui eramus mortui per peccatum, reuiuiscimus, habemusque pacem et tranquillitatem in conscientia  \pend
\vspace{0.5cm}\noindent
\vspace{0.2cm}\rule{1cm}{0.2pt}\\ 
\hspace{0.2cm}\textit{mg}
\footnotesize Vires huma- nae nihil sunt.  
\normalsize| 
\section*{COMMENT. POMER.  }\pstart nostra, iam non mordet nos et affligit stimulus mor tis, ipsum peccatum .  \pend
\phantomsection
\addcontentsline{toc}{subsection}{\textit{Propterea quod adfectus carnis. }}
\subsection*{\textit{Propterea quod adfectus carnis. }}\pstart Ecce quicquid sapit caro, quando sine spiritu Dei et Verbo sapis, carnaliter sapis, et ut maxime Ver- bum habeas, audias, abuteris co, si spiritu uacas: ergo quicquid homo intelligit per se et rationem hu- manam, id totum est contra Deum: neque hic fingere potes ullum medium, dicit enim crassissime: Immicitia est aduersus Deum, et seipsum clarius interpretatur,  \pend
\phantomsection
\addcontentsline{toc}{subsection}{\textit{Nam legi Dei non subditur, neque  enim potest. }}
\subsection*{\textit{Nam legi Dei non subditur, neque  enim potest. }}\pstart Egregium hoc est Libus Arb Encomium: nouit Paulus futurum, ut totis buccis quidam disputaturt es- sent de Libus Arb. defensurique eius partes: non potest, inquit, facere et uelle ea quae Dei sunt, caro et eius sapientia, uelit nolit, cogitur esse inimica aduersus Deum, nam dicit Christus: Qui uult uenire post me, abneget semetipsum: Non dicit carnem, ne et ibi tor queas Christi Verbum, quo significat, quicquid est omnis home, id totum mecum pugnat, inde oramus: Fiat uoluntas tua, non nostra. Itaque stultitia fuit, quod hactenus Monachi, ducti sua sapientia et bona in- tentione, finxerunt secundum suas cogitationes Deum, cum tamen non uelut fingi, aut forman a nobis, ue-  \pend
\vspace{0.5cm}\noindent
\vspace{0.2cm}\rule{1cm}{0.2pt}\\ 
\hspace{0.2cm}\textit{mg}
\footnotesize Carnis sapien- tia.  
\normalsize| 
\hspace{0.2cm}\textit{mg}
\footnotesize Lib . Arb.  
\normalsize| 
\section*{COMMENT. POMER.  }\pstart non uisa carnali oculo, aut audita aure humana. Subli- mia sunt, et soli spiritui cognita: quapropter Tex- tum uideamus.   \pend
\phantomsection
\addcontentsline{toc}{subsection}{\textit{Debitores sumus, non carni. }}
\subsection*{\textit{Debitores sumus, non carni. }}\pstart id est, non carnem iam sequamur, aut eius pruden- tiam, quae mors est, sed spiritum, quo agatur, buic omnia debemus, carni nihil: noli interim negligere Paulinam Phrasin, qua hic utitur.   \pend
\phantomsection
\addcontentsline{toc}{subsection}{\textit{Moriemini. }}
\subsection*{\textit{Moriemini. }}\pstart Id est, si secuti fueritis carnis sapientiam, obnoxii eritis damnationi. Alludit uero ad legem Adae da- tam, ubi dicebatur: In quacunque die comederitis, mo- riemini. Hic interim obserua, quid Paulus (nam pro pter crassos haec moneo) secundum carnem uiuere dicat, ne uel qui de Ciuilibus moribus praecepta scri- pserunt, tibi imponant. Illi enim putabant, eos non carnaliter uiuere, qui non scortarentur helluarentur, et decoquerent omnia. Est hoc sapientiae humanae, non Verbt, quod adfectus et carnis prudentiam, carnalem uitam uocat, non tantum scortationem, et ebrietatem, id quod uides in Eua, quae non fa- cto, sed ipsa carnis prudentia decepta est. Nam fic cogitabat, audito diuino mandato: Haec arbor cunt sit tam pulchra et grata, qui fieri potest, ut mali et noxiieius sint fructus, nequaquam uenenum habent, ut maxime Deus edi uetuerit, creatura est Dei, ut ego,  \pend
\vspace{0.5cm}\noindent
\vspace{0.2cm}\rule{1cm}{0.2pt}\\ 
\hspace{0.2cm}\textit{mg}
\footnotesize Carnaliter u uere quid. 
\normalsize| 
\section*{EPIST. AD ROM.  88 }
\marginpar{[ p.]}\pstart quomodo non potest esse bona? profecto gustabo, ni- hil oberit gustus. In his peccauit Eua, ut maxime non dum pomum attigerat. Siquidem Deum et spiritum eius, mendacii accusarat. Vbi hic opus, cui adscri- bas peccatum? ergo peccatum carnis, est uiuere, cogi- tare, agereque secundum prudentiam carnis, ne qua elabantur Monachi et Sanctuli qui causantur se nihil unq̃ peccatorum opere fecisse, ergo esse inculpatos, et dignos qui sedem suam ad thronum altissimi col- locent.   \pend
\phantomsection
\addcontentsline{toc}{subsection}{\textit{Quod si spiritu. }}
\subsection*{\textit{Quod si spiritu. }}\pstart Scilicet, Dei uel Christi . FACTA COR.  PORIS. Id est, carnis. Semper autem id ob ocu- los habeto, Paulum non externa, sed interioremcon cupiscentiam intelligere. Ea enim est summum et pri- mum opus corporis uel carnis, per eam quiduis agit.  Pii, ubi hoc opere solicitantur, concupiscentiae resi- stunt, nomen Dei inuocant, precantur, ut ab ea liberentur.  Concupiscen- tia  \pend
\phantomsection
\addcontentsline{toc}{subsection}{\textit{Viuetis. }}
\subsection*{\textit{Viuetis. }}\pstart Scilicet aeterna uita, quemadmodum, in moriemint, aeternam mortem intellexit.   \pend
\phantomsection
\addcontentsline{toc}{subsection}{\textit{Etenim quicunque spiritu Dei du- cuntur. }}
\subsection*{\textit{Etenim quicunque spiritu Dei du- cuntur. }}
\vspace{0.5cm}\noindent
\vspace{0.2cm}\rule{1cm}{0.2pt}\\ 
\hspace{0.2cm}\textit{mg}
\footnotesize 
\normalsize| 
\section*{COMMENT. POMER.  }\pstart trahimur desyderio, quemadmodum spiritu carnis trahimur et ferimur in carnalia, ut ille dixit: Tra- hit sua quenque uoluptas. Verum ne quis dicat: St tra- himur, ergo non est libertas? Recte-Haec tractio mi- ro quodam fit gaudio et cordis dilatatione, in Deo: tametsi interim quoq non per omnia Adae ueteri pla- ceat, ideoque ; magis ad hunc quod trahatur referes, quan- ad Dei spiritum, aut nouum hominem, nam is sua sponte in omne bonum fertur, neque etiam nisi desy- derantem trahit Deus per illum spiritum, qui dum traheris, tibi desyderium in Deum simul ingessit, sed hoc caret caro, habet autem et suum desyderium, quo fertur, sicut suum spiritum, quod ex diametro pugnat cum sancto hoc desyderio spiritus.   \pend
\phantomsection
\addcontentsline{toc}{subsection}{\textit{Spiritu Dei. }}
\subsection*{\textit{Spiritu Dei. }}\pstart Scilicet, non suo. Hic iterum uides, quid sit car - naliter uiuere, nempe suo spiritu agi, et ambulare secundum adfectus qui sunt in carne. Ergo optima quae que opera in speciem, si non e Dei spiritu profluant, impia sunt. Nisi ille te doceat, et moneat quid faci- as, nihil tua ratione efficies.   \pend
\phantomsection
\addcontentsline{toc}{subsection}{\textit{Hi sunt filii Dei. }}
\subsection*{\textit{Hi sunt filii Dei. }}\pstart Propter spiritum inhabitantem, quem eundem ha- bent cum Deo. Quemadmodum filius carnalis unam habet carnem cum patre, et uxor et maritus sunt una caro: Sic, Cuius spiritus nos et Deum unum fa-  \pend
\section*{IN EPIST. AD ROM.  89 }
\marginpar{[ p.]}\pstart cit, id est, ut nos simus filii eius, ipse uero pater noster.  \pend
\phantomsection
\addcontentsline{toc}{subsection}{\textit{Non enim accepistis spiritum seruitutis. }}
\subsection*{\textit{Non enim accepistis spiritum seruitutis. }}\pstart Vide, alius spiritus est in seruis, alius in filiis, Eyn- ander mut und synn: filii sciunt se liberos, non timent eiectionem sui e domo patris, sunt securi de haeredita- te, certo seiunt se habere patrem, et omnia quae pa- tris sunt, ad se quoque pertinere: Sciunt praeterea se non meruisse, ullo in patrem officio, quod haec simul cum eo possideant, uerum, ex hoc tantum quod filii sint. Sed huius haereditatis et libertatis ius, nemo carna- lis exprimet, nisi qui ipse se filium Dei esse norit.  Contra, serui nihil habent libertatis, nesciunt qua ho- ra e domo eiiciantur. Non sunt enim filii, nihilque quod in domo est, ad eos pertinet, faciunt omnia coacti et territi, premunturque tyrannide molestissima legis, quae semper dicit, fac fac, dilige dilige, ac nequeunt, cum hoc ipsum filii non audiant, nam per se faciunt omnia et diligunt Deum: quomodo enim non dilige- rent patrem? Proinde, hic aduerte quam rotundus sit Paulus in argumentis breuissimis, seque mutuo com- plectentibus. Primum: Non est peccandum, quia non estis debitores carnis. Secundo, Si secuti carnem fue- ritis, mortemini. Tertio, Si spiritu ambulaueritis, ui uetis. Quarto, Hoc spiritu estis filii Dei, non serui, Filii sumus non serui.   \pend
\vspace{0.5cm}\noindent
\vspace{0.2cm}\rule{1cm}{0.2pt}\\ 
\hspace{0.2cm}\textit{mg}
\footnotesize 
\normalsize| 
\section*{COMMENT. POMER.  }\pstart trahimur desyderio, quemadmodum spiritu carnis trahimur et ferimur in carnalia, ut ille dixit: Tra- hit sua quenque uoluptas. Verum ne quis dicat: Si tra- himur, ergo non est libertas? Recte-Haec tractio mi- ro quodam fit gaudio et cordis dilatatione, in Deo: tametsi interim quoque non per omnia Adae ueteri pla- ceat, ideoque ; magis ad hunc quod trahatur referes, quam ad Dei spiritum, aut nouum hominem, nam is sua sponte in omne bonum fertur, neque etiam nisi desy- derantem trahit Deus per illum spiritum, qui dum traheris, tibi desyderium in Deum simul ingessit, sed hoc caret caro, habet autem et suum desyderium, quo fertur, sicut suum spiritum, quod ex diametro pugnat cum sancto hoc desyderio spiritus.   \pend
\phantomsection
\addcontentsline{toc}{subsection}{\textit{Spiritu Dei. }}
\subsection*{\textit{Spiritu Dei. }}\pstart Scilicet, non suo. Hic iterum uides, quid sit car- naliter uiuere, nempe suo spiritu agi, et ambulare secundum adfectus qui sunt in carne. Ergo optima quae que opera in speciem, si non e Dei spiritu profluant, impia sunt . Nisi ille te doceat, et moneat quid faci- as, nihil tua ratione efficies.   \pend
\phantomsection
\addcontentsline{toc}{subsection}{\textit{Hi sunt filii Dei. }}
\subsection*{\textit{Hi sunt filii Dei. }}\pstart Propter spiritum inhabitantem, quem eundem ha- bent cum Deo. Quemadmodum filius carnalis unam habet carnem cum patre, et uxor et maritus sunt 'una caro: Sic, Cuius spiritus nos et Deum unum fa-  \pend
\section*{EPIST. AD ROM.  }
\marginpar{[ p. ]}\pstart cit, id est, ut nos simus filii eius, ipse uero pater uoster.  \pend
\phantomsection
\addcontentsline{toc}{subsection}{\textit{Non enim accepistis spiritum seruitutis. }}
\subsection*{\textit{Non enim accepistis spiritum seruitutis. }}\pstart Vide, alius spiritus est in seruis, alius in filiis, Eyn- ander mut und synn: filii sciunt se liberos, non timent eiectionemn sui e domo patris, sunt securi de haeredita- te, certo seiunt se habere patrem, et omnia quae pa- tris sunt, ad se quoque pertinere: Sciunt praeterea se non meruisse, ullo in patrem officio, quod haec simul cum eo possideant, uerum, ex hoc tantum quod filii sint.  Sed huius haereditatis et libertatis ius, nemo carna- lis exprimet, nisi qui ipse se filium Dei esse norit.  Contra, serui nihil habent libertatis, nesciunt qua ho- ra e domo eiiciantur. Non sunt enim filii, nihilque ; quod in domo est, ad eos pertinet, faciunt omnia coacti et territi, premunturque tyrannide molestissima legis, quae semper dicit, fac fac, dilige dilige, ac nequeunt, cum hoc ipsum filii non audiant, nam per se faciunt omnia et diligunt Deum: quomodo enim non dilige- rent patrem? Proinde, hic aduerte quam rotundus sit Paulus in argumentis breuissimis, seque mutuo com- plectentibus. Primum: Non est peccandum, quia non estis debitores carnis. Secundo, Si secuti carnem fue- ritis, moriemini. Tertio, Si spiritu ambulaueritis, ui- uetis. Quarto, Hoc spiritu estis filii Dei, non serui,  \pend
\vspace{0.5cm}\noindent
\vspace{0.2cm}\rule{1cm}{0.2pt}\\ 
\hspace{0.2cm}\textit{mg}
\footnotesize Eilii sumus non servi. 
\normalsize| 
\section*{COMMENT. POMER.  }\pstart At haec in textu clarius uideamus, olim, inquit, dum serui eratis, et sub lege, impossibile erat uos non ti- mere Deum, quem sciebatis, non patrem, sed iudicem es- se, ita enim docebat uos mala conscientia, quae rea erat: nolebatis illum accedere cum fiducia, sed fugie batis magno horrore et trepidatione, sentiebatis onus peccati et iram legis, et dum ab ea uos absolutum iri, per externa opera et hypocrisin putabatis, maiorem timorem, imo desperationem incidistis, id experti et nos sumus. Nonne et nos ad bona opera agebamur propter peccata, statutis Pontificum et carnificum Romanensium: multa fecimus, ieiunando, orando, canendo, templa Diuorum aedificando, nunc Diuum Iacobum, nunc Veronicam, aut aliud portentum inuisendo, sed quid promouimus, non potuimus a peccatis absolui: Erat enim tantum spiritus seruitu- tis, nunquam pacificatam sensimus conscientiam, quae tamem, ubi spiritus est Dei, et filiorum, tranquillissi ma redditur. An non uerum est, quod de ludaeis in Nono dicit Apostolus: Sectando legem iusticiae, ad legem iusticiae non peruenimiis. Eundem timorem et in Adam uidemus faisse, qui post peccatum latuit, fugiens Dei conspectum, nam ante non erat filius, si quidem nunc habebat spiritum seruitutis. De hoc ti- more loquitur Iohannes in Epistola: Perfecta Chari- tas foris eiicit timorem. Non fert charitas, ut Deum timeas, quandoquidem scis te filium esse, accedis hi  \pend
\vspace{0.5cm}\noindent
\vspace{0.2cm}\rule{1cm}{0.2pt}\\ 
\hspace{0.2cm}\textit{mg}
\footnotesize Conscientiae carnificina le gis opera 
\normalsize| 
\section*{IN EPIST. AD ROM.  90 }
\marginpar{[ p.]}\pstart Iari animo, agnoscendo patrem. Non loquitur autem hie de timore eo, quem tantopere commendat Scriptura, qui nihil aliud est, quam cultus et Dei reuerentia, de quo scribitur: Timor Domini expellit peccatum, Item: Initium sapientiae, timor Domini. Hic timor ac- cedit Deum, alius ille seruilis fugit conspectum Dei, quod in nobis fiebat, ante praedicatum Christi Euan gelium. Nunc uero quia filiorum timorem habemus, accedimus magna cum conscientiarum laetitia. Op- ponit itaque hoc loco filiis seruos.  \pend
\phantomsection
\addcontentsline{toc}{subsection}{\textit{Sed accepistis spiritum adoptionis. }}
\subsection*{\textit{Sed accepistis spiritum adoptionis. }}\pstart Id est filiationis, si dicere liceret, per quem sumus adoptati in filios, et quemadmodum Christus natu- ra est filius, ita nos adoptione: neque minus tamen ha- bemus iuris aut haereditatis, quam ipse, omnes ex aequo complectitur nos una cum ipso pater, nihil ille ha bet quod nos non habeamus, nihil ille uult, quod nos detrectemus. PER QVEM. Scilicet, spiri- tum nobis datum,  \pend
\phantomsection
\addcontentsline{toc}{subsection}{\textit{Clamamus. }}
\subsection*{\textit{Clamamus. }}\pstart Scilicet, ad Deum patrem, fidelt corde, qui cla- mor fit etiam clauso ore: nam et Moses legitur clauso ore ad Deum clamasse.  \pend
\phantomsection
\addcontentsline{toc}{subsection}{\textit{Abba pater. }}
\subsection*{\textit{Abba pater. }}\pstart Syrum est Abba, quo et Hebraei utuntur, πα τίρ Graecum, Puto uero Paulum iis usum noibus, quae idem significant  \pend
\section*{COMMENT. POMER.  }\pstart propter Iudaeos et Gentes, quorum appellatione om- nes homines in scripturis intelliguntur. Haec est uox filiorum Dei, non hypocritarum, sentiunt enim per spiritum, ita uere in corde rem habere, non in uerbis tantum esse, quando enim dicunt: Pater, patrem ip- sum habent, non uoce tantum, aut corde fingunt. Hae sunt ipsae diuitiae, thesaurique ; diuinitatis et maie- Haereditas fi- statis Dei, hic filii et haeredes sumus, neque alicunde liorum Dei. nobilius quiddam est expectandum, ideoque ; iubet ora- pr nn2dre: Pater noster etc.q. d. Stin corde uere dixeritis me patrem uestrum, ero pater, non absens, sed prae sentissimus. Possidebo uos, et uos me: habebitis et in me omnia, nempe, sanctificationem nominis mei, regni mei potentiam, uoluntatis mei agnitionem, omnium necessariarum copiam, peccatorum remissio- nem, uictoriam Sathanae, et omnium tentationum: In summa, libertatem qua sentietis, uos neutiq̃ posse- fieri obnoxios ulli malo. O immensam Dei bonitatem, quae usqueadeo, suis diuitiis nos obruit, quas utinam agnosceremus. Vult esse pater, modo uelimus esse fi- lii: trahit, solicitat nos in sui amorem, suo spiritu: inuitat et externis beneficiis, sed quia corda nostra indurata sunt, neque audimus tam dulcemuocem, neque  uidemus ulla eius beneficia.   \pend
\phantomsection
\addcontentsline{toc}{subsection}{\textit{Idem spiritus testatur. }}
\subsection*{\textit{Idem spiritus testatur. }}\pstart Id est, ipse spiritus sanctus, certum reddit spiri-  \pend
\vspace{0.5cm}\noindent
\vspace{0.2cm}\rule{1cm}{0.2pt}\\ 
\hspace{0.2cm}\textit{mg}
\footnotesize 
\normalsize| 
\section*{ IN EPIST. AD ROM.  }
\marginpar{[ p.]}\pstart tum nostrum. Duplicem hic audis spiritum, nostrum et Dei, cum tamen idem sit spiritus, ubi illuminati su- mus, nam ante sumus uere absq spiritu, et carnales dicimur, inde noster dicitur, quia nobis datus: antea mens nostra et intellectus non capiebat Deum, iam uero per Spiritum sanctum, transiit in spiritum a sua car- nalitate, et factus est cum Deo unus spiritus, tra- hens paulatim post se totum hominem, de quo dicit: Qui spiritu Dei aguntur etc. Quid agit ille? nihil aliud, quam quod firmissime nobis persuadeat et inculcet Deum esse patrem nostrum, nosque ; illius fili- os, afficit nos dulcibus erga illum adfectibus, ne quid aut dubitemus, aut desyderemus, non sunt ergo ig- nari Christiani de uoluntate Dei patris erga se, scunt uere se esse filios: Sunt obsignati hoc charactere aeter no et inuiolabilt, quem intuentes, securi sunt de il- lius dilectione, id quod dicunt in Psalm. Signatum est lumen uultus tui Domine etc. Tametsi Papistae alium quendam characterem bestiae, de qua in Apoc. inue- nerunt. Non suntitaque haec simplicia uerba, cum di- cit: Spiritus, certos nos facit, nos esse Det filios, e coelo datus est ille in corda nostra, quae obsignauit: sumus consortes Christi: ergo inuncti hoc spiritu quo ille, quamuis ille plus inunctus fuerit, cum esset ca- put nostrum. Nemo est ergo Christianus, nisi credat in Deum per Christum, si credit, iam spiritum ha- bet per quem credit, qui regnans in corde, dat ilud  \pend
\vspace{0.5cm}\noindent
\vspace{0.2cm}\rule{1cm}{0.2pt}\\ 
\hspace{0.2cm}\textit{mg}
\footnotesize Duplex spi- ritus.  
\normalsize| 
\section*{COMMENT. POMER.  }\pstart testimonium, nam corde creditur ad iusticiam: Solius est igitur Spiritus sancti credere, quo non regnante in cordibus nostris, non credimus. Hic ridendi sunt, quicunque fidem sibi, nescio quam, fingunt, cum ta- mem certi, de Det erga se uoluntate bona, esse nequeant, quales sunt omnes, qui bonas sibi fingunt intentio- nes, et nescio quam pietatem stertentem. Tales igno- rant, imo semper dubii sunt de peccato remisso, nequam remissum sibi unquam sentire possunt. Sunt et non- nulli, qui hoc spiritu contempto, expectant, nescio quas e coelo reuelationes et motiones Spiritus sancti, uerum decipiuntur, Sathanas eius fascini auctor est, quo sic diuexantur, et illi, ut caro in huiusmudi est Spiritus quid? prona, diunt a Deo factum . Praeterea hic propter quosdam, de spiritu, anima, et corpore, quaedam in dicanda sunt, nam saepe illis utitur-Spiritus, est ex- cellentior hominis pars, id est, intellectus, illumina tus Spiritu sancto, et cum eo unum factus, de quo ani- males homines nesciunt, neque intelligunt quid fit, ut dictum . Hic ineomprehensibiles, inuisibiles, aeternasque res complectitur, estque plane domus fidei et Dei, quia Verbo Dei, et Spiritu sancto illustra- tur, alioqui, obscura domus impietatis est, sicareat Verbo. Anima uero est, ea uis omnis, qua uiuifica- tur totum corpus, fiuntque per eam omnes operationes in corpore, ettam dum dormis, haec operatur, et tunquam quiescit. Comprehendit autem illa non res  \pend
\vspace{0.5cm}\noindent
\vspace{0.2cm}\rule{1cm}{0.2pt}\\ 
\hspace{0.2cm}\textit{mg}
\footnotesize Anima quid? 
\normalsize| 
\hspace{0.2cm}\textit{mg}
\footnotesize 
\normalsize| 
\section*{IN EPIST. AD ROM.  92 }
\marginpar{[ p.]}\pstart incomprehensibiles, aut aeternas, sed eas tantum, quas ratio cognoscere potest et intueri: Idcirco et huius domus, ratio, lux esse dicitur, qua amota, nihil ani- mna ipsa intelligit, aut concupit. Corpus autem, to- tum illud externum est, constans pluribus membris, ha- betque suos motus et operationes, ex primis duabus partibus, anima et spiritu, nam omnia, ut per animam cognoscuntur, et spiritu creduntur, in illo fiunt. Et ut id clarius fiat, simili quodam ex Scripturis, taber- naculum Mosi proponemus. Atrii nomine, corpus significari potest, nam ut atrium illud, sub Diuo erat omnium oculis patens, sic corpus externum est et cer- nitur, tangitur: hic nihil est secretum, aut arcanum, neque tale quiddam hic geritur, sed omnia in propatulo erant. Verum omnia intus geruntur, aut in Sancto et ante fores eius, aut Sancto sanctorum: itaque Sanctum, animam significabit, in eo, et ant? fores eius, omnia uiribus humanis fiebant, caedebantur uictimae, expi- ationesque celebrabantur. Non secus in aniares habet, quaecumq per eam fiunt, sunt humanae uires, ex ratione et sapientia carnis profluunt, hic nihil adhuc est uere san- ctum, et uere diuinum, tametsi quaedam multam spetiem habe- ant sanctitatis. Quare, in Sanctum sanctorum est progre- diendum, ubt nihil humanarum est uirium, nihil naturalis- lucis, uerum tenebrae et caligo, quo significabatur hanc esse domum Dei, qui habitat in abscondito, quique ; non cernitur presidio humanae ullius lucis, aut sapientiae  \pend
\vspace{0.5cm}\noindent
\vspace{0.2cm}\rule{1cm}{0.2pt}\\ 
\hspace{0.2cm}\textit{mg}
\footnotesize Corpus quid? 
\normalsize| 
\section*{COMMENT. POMER.  }\pstart Huius ratio quodammodo, et in spiritu nostro depre- benditur, qui nihil externum admittit, uerum quic- quid continet per Verbum, hoc est solius Dei: estque ́ in eo lux illa inaccessibilis, quam neque oculus, neque  auris, aut cor humanum capiat. Sed diligenter est ui- dendum, ne hunc spiritum, aut Sanctum sanctorum nostrum, occupent Assirii, id est, humanae traditio- nes, quibus obnoxius iam non est spiritus Dei, sed Sa- thanae. Ergo orandum est, ut hic non ratio, aut quid dam carnale dominetur, sed solus Deus, cuius spiri- tus certum reddit nostrum spiritum, quod sumus filii Dei: Si filii, ergo et haeredes, id est, possessores oim quae sunt patris, nempe coeli, terrae, omniumque crea- turarum. Praeterea non tantum haeredes Det, sed et Christi cohaeredes, id est, quo iure et potestate ille omnia possidet, ita et nos. Quae igitur creatura plus bonoratur, quam Christianus, an maior possit alicui contingere gloria a Deo? Tamen fit persepe, ut tem- tatio et crux eam obscuret. Beatus ergo ille, qui in- finem perseueraturus est. Sed haec hactenus.   \pend
\phantomsection
\addcontentsline{toc}{subsection}{\textit{Siquidem simul. }}
\subsection*{\textit{Siquidem simul. }}\pstart Suauissimae hae sunt consolationes, ne aut in cru- ce succumbamus, aut peccato obsequamur: nos, in- quit, certi sumus, et in hoc, quod simus cohaeredes et fratres Christi, quia affligimur tota die, quemad modum ipse adflictus est, solicitamur per carnem ad  \pend
\vspace{0.5cm}\noindent
\vspace{0.2cm}\rule{1cm}{0.2pt}\\ 
\hspace{0.2cm}\textit{mg}
\footnotesize Crux saluta- ris. 
\normalsize| 
\section*{IN EPIST. AD ROM.  }
\marginpar{[ p.93 ]}\pstart peccandum, deiicimur ad inferos, omnes iniurias mum- di perferimus, non tamen aliam ob caussam, nisiut per omnia, illius imagini crucis, adsimilati, in eandem quoque gloriam cum illo constituamur. Est itaque opti- mum hoc ab exemplo argumentum, et utile: Si enim naturalis filius passus est, cur nos non pateremar, qui sumus adoptiui? is suo spiritu filius est Dei, tamen cruci suffixus est, pudeat nos temeritatis et insipien- tiae nostrae, qui crucem detrectamus, quos ille facere uult Dei filios, non nostro, sed suo spiritu.   \pend
\phantomsection
\addcontentsline{toc}{subsection}{\textit{Nam reputo. }}
\subsection*{\textit{Nam reputo. }}\pstart Argumentum est a praemio: tanta, inquit, gloria nos expectat, ut nulla cum ea res conferri queat, imo et crux quantumuis sit intolerabilis et calamitosa, uix umbra quaedam est in collatione huius gloriae, quae reuelabitur in nobis in resurrectione. Si igitur pares non sunt adflictiones, taceamus merita operum, mi- hil meremur, aut operibus nostris, aut cruce, omnia per gratiam nobis exhibentur. Nam si adflictiones pro- pter Christum susceptae, non sunt condignae, quando condigna forent nostra opera?  \pend
\phantomsection
\addcontentsline{toc}{subsection}{\textit{Etenim solicita. }}
\subsection*{\textit{Etenim solicita. }}\pstart Multis argumentis hactenus Paulus eam sententia comprobauit, qua dixit, nihil esse in nobis damnatio nis, ut maxime in nobis sit reliquum peccatum, et peccati sensus, carnisque adfectus, qui primum in re-  \pend
\vspace{0.5cm}\noindent
\vspace{0.2cm}\rule{1cm}{0.2pt}\\ 
\hspace{0.2cm}\textit{mg}
\footnotesize Merita no- stra nulla sunt 
\normalsize| 
\section*{COMMENT. POMER.  }\pstart surrectione abolebuntur, nunc tandem confirmat ani- mos nostros, aduersus crucem, et insultus Sathanae.  q' d- non poteritis consistere in ista tribulatione et adflictione, sine patientia: Ideoque exemplum patien- tiae sumite uobis a totius uniuersi (quod ipse crea- turam uocat) anxia expectatione, in qua, multa cum patientia sustinet, quod impii eo abutantur, et Sathanas pro sua libidine, ipsum cum mem- bris suis, id est, creaturis impiis in os et rictum pro trudat . Nolite, quaeso, desperare, non enim uos soli patimini externa, habetis consortes omnes crea- turas in cruce. Praeterea, quod hic Paulus creaturam uocat, Philosophi, uniuersum, et Moses coelum et terram dixere, illa, cum iam omnia sint corruptioni obnoxia, spem habent gloriae futurae, non propter se, sed propter nos: erit enim aliquando coelum no- uum, et terra noua, ut Esaias dixit: Item, Petrus in Epistola indicauit: Iam sunt in labore perpetuo, nun- quam quiescunt, seruiunt nobis, ferunt etiam male- dictionem propter nos, sed cum expectatione, sciunt enim fore, ut aliquando liberemur, quod tum fiet, quando nos liberabimur perfecte in resurrectione ul- tima a nostro peecato. Magnum hoc est Pauli myste- rium, quod tamen nondum comprehendimus, ue- rum sicuti resurrectio et gloria nostra, fide tan- tum comprehenditur, ita et illud. Docemur praete- rea, hoc arcano, nos uere dominos esse et tan-  \pend
\vspace{0.5cm}\noindent
\vspace{0.2cm}\rule{1cm}{0.2pt}\\ 
\hspace{0.2cm}\textit{mg}
\footnotesize Consolatio ne succumbamus in malis. 
\normalsize| 
\section*{IN EPIST. AD ROM.  }
\marginpar{[ p.94 ]}\pstart quam patrisfamilias in hoc mundo, reliquasque omnes creaturas, quasi familiam nostram, dolere nostram in foelicitatem, interim tamen cum respiciant in solam Dei uoluntatem, perferunt omnia, et uanitatem et corruptionem. Ideo addidit  \pend
\phantomsection
\addcontentsline{toc}{subsection}{\textit{Non uolens. }}
\subsection*{\textit{Non uolens. }}\pstart Quod ad creaturas attinebat, non uellent subii- ci iis malis, sed cum Dei sit uoluntas, illi se permit- tunt. Similtratione, et Sancti, uelint abesse pec - catum omne: At cum inteligant et sentiant hanc esse Deiuoluntatem, ut in carne sint hoc tempore, perferunt, tametsi inuiti, peccatum, clamant: Ne inducas nos in tentationem, sed libera nos a malo-  \pend
\phantomsection
\addcontentsline{toc}{subsection}{\textit{Propter eum: }}
\subsection*{\textit{Propter eum: }}
\phantomsection
\addcontentsline{toc}{subsection}{\textit{Scilicet, Deum.  }}
\subsection*{\textit{Scilicet, Deum.  }}
\phantomsection
\addcontentsline{toc}{subsection}{\textit{Qui subiecit illam sub spe. }}
\subsection*{\textit{Qui subiecit illam sub spe. }}\pstart Scilicet, liberationis, habet certissimam spem, quod olim debeat liberari: nihil enim est sub so- le, quod non mauelit esse liberum, beatumque sua conditione, quam rapi, discerpique , et in seruitutem redigt, uel in totum aboleri: nonne quaelibet arbor bona pro naturae suae bonitate libenter fructum fer- ret, fit autem ut nonnunquam aduersa tempestate, icta  \pend
\vspace{0.5cm}\noindent
\vspace{0.2cm}\rule{1cm}{0.2pt}\\ 
\hspace{0.2cm}\textit{mg}
\footnotesize Speseruamur in tentatione 
\normalsize| 
\section*{COMMENT. POMER.  }\pstart aut uermibus arrosa, nihil fructuum adferat. Cupit Sol et Luna, ut tempus uernum redeat, quo omnia reflorescant, secundum hoc Prophetae: praeceptum po suit et non praeteribit, accidit tamen, ut saepe quid ad uersi uideamus in eo tempore. Et mulier uehemen- ter cupit se liberam a partu, at qui fieri nequit, haec enim est Dei uoluntas, ut isto periculo laboret: in omnibus ergo creaturis quiddam crucis cernis et ama- ritudinis, qua premuntur et adfliguntur, omnia sunt uicissitudinibus, et defectibus obnoxia, ad seruien- dum interim homini: sed quemadmodum illae creatu- rae spem habent liberationis suae, ita et nos non suc- cumbamus in cruce, et aduersitatibus, nam adie- cit Paulus  \pend
\phantomsection
\addcontentsline{toc}{subsection}{\textit{Quoniam et ipsa creatura libe- rabitur a seruitute corruptionis, in libertatem gloriae filiorum Dei. }}
\subsection*{\textit{Quoniam et ipsa creatura libe- rabitur a seruitute corruptionis, in libertatem gloriae filiorum Dei. }}\pstart Id est, quando filii Dei glorifie abuntur, liberabuntur et creaturae a seruitute, neque eni tum earum ministerio indi- gebimus, uidebimus omnia noua facta: tunc fiet ser- mo qui scriptus est: Vbi est mors uictoria tua? Item il- lud Solomonis: Impii de terra perdentur, iusti autem habitabunt in ea, in seculum seculi, quia coelum et terra iustorum tantum sunt, cum quibus et glorifica- buntur, Psalm. 8.1. Timoth.4. non impiorum- Possumus huius mysterii exemplum quoddam, uxcunque   \pend
\vspace{0.5cm}\noindent
\vspace{0.2cm}\rule{1cm}{0.2pt}\\ 
\hspace{0.2cm}\textit{mg}
\footnotesize Glorificatio iustorum.  
\normalsize| 
\section*{99 IN EPIST. AD ROM.  }
\marginpar{[ p.]}\pstart iam in nobis praesentiscere. In tentatione, in qua omnes creaturae nobis contrariae uidentur, etiam sol ipse qui clarissimus est, nobis tenebrascit, secundum ea, quae in Prophetis legimus: Omnia quae iucunda fue- runt et dulcia, in amarorem nobis uersa sentimus, imo et ad strepitum decidentis folii ex arbore trepi- damus. At ubi spem aliquam in Deum in ea tentatio- ne rursus conceperimus, nihil est quod non stet a no- bis, nihil terret nos, arrident omnes creaturae, et no biscum gaudere uidentur, siquidem Deus pater arri- det. Ecce hic gustus est, primitiarum gloriae huius, hic experimur illud quod dicit Paulus: Diligentibus Deum, omnia cooperantur in bonum, id est, omnes creaturae et quicquid est rerum, etiam aduersarum, salutem nostram promouent, id quod in resurrectione uidebimus, cum omnis impietas extirpata fuerit  \pend
\phantomsection
\addcontentsline{toc}{subsection}{\textit{Scimus eni, quod omnis creatura, }}
\subsection*{\textit{Scimus eni, quod omnis creatura, }}\pstart Nemo haec uerba, nisi spiritualis sit, dicere po- test, caro enim ignorat, quod creaturae etiam pessi- mae, item peccatum et mors, piis cooperentur in bonum-  \pend
\phantomsection
\addcontentsline{toc}{subsection}{\textit{Congemiscit. }}
\subsection*{\textit{Congemiscit. }}
\phantomsection
\addcontentsline{toc}{subsection}{\textit{Gemit nobiscum, et cupit liberari a seruitute.  }}
\subsection*{\textit{Gemit nobiscum, et cupit liberari a seruitute.  }}
\phantomsection
\addcontentsline{toc}{subsection}{\textit{Dimulque nobiscum parturit. }}
\subsection*{\textit{Dimulque nobiscum parturit. }}\pstart Id est, in angustia est ut mulier parturiens, que uihil aliud ob oculos habet, quam mortem, ut inte- rim dolorem taccam, lohan . 16.   \pend\pstart  \pend\pstart 3 Scilicet, nostrum: nunquam cessat gemere crea- tura, donec liberetur a seruitute, quemadmodum mu lier a partu: debebant profecto conscientiae nostrae me ditatione huius mysterii exerceri, nam quaedam utique  comprehenderent, et uel ipsa experientia condiscerent.   \pend
\phantomsection
\addcontentsline{toc}{subsection}{\textit{\huge\textbf{COMMENT. POMER. }\normalsize Vlque ad hoc tempus. }}
\subsection*{\textit{\huge\textbf{COMMENT. POMER. }\normalsize Vlque ad hoc tempus. }}\pstart Argumentum est aliud, quo nos consolatur, supe- riori tamen cohaerens: creatura inquit ingemiscit, et expectat liberationem, sed non sola: est Spiritus san- ctus in cordibus nostris, qui ide. facit, hic ingemiscit, uelit quoque nos liberos a peccato, et carnis uersutia, sed fieri nequit: Sustinet, quemadmodum creaturae di- cens: Fiat uoluntas tua. Sunt uere gemitus parturien- tis, quando dicit Paulus: Quis me liberabit a corpo- re mortis huius? ah infoelix ego. Christus in horto dixit: Transfer a me calicem, quae omnia ex infirmi- tate carnis dicta uides. At spiritus, constanter hoc unum dixit: Fiat uoluntas tua, sed hoc nemo, nisi tem- tatione et spiritu edoctus, satis intelligit.  \pend
\phantomsection
\addcontentsline{toc}{subsection}{\textit{Non solum autem illa. }}
\subsection*{\textit{Non solum autem illa. }}
\phantomsection
\addcontentsline{toc}{subsection}{\textit{Qui primitias spiritus habemus. }}
\subsection*{\textit{Qui primitias spiritus habemus. }}\pstart His seipsum exponit, ne quis id de carnalibus spi- ritibus nostris dictum putet, qui non ingemiscunt, sed toti insaniunt in carne et spumant. Non sunt in- firmi, sed robusti et audientes omnia: Item aduentum domini non expectant, sed deprecantur: in omnibus  \pend
\vspace{0.5cm}\noindent
\vspace{0.2cm}\rule{1cm}{0.2pt}\\ 
\hspace{0.2cm}\textit{mg}
\footnotesize Creatiraege- munt, et ex pectant libe- rationem.  
\normalsize| 
\section*{IN EPIST. AD ROM.  96 }
\marginpar{[ p.]}\pstart enim tantum suam gloriam, et utilitatem quaerunt, hon Dei. Alludit autem hic Apostolus ad legis obla- tiones, ubi non totus fructus offerebatur, sed primi- tiae fructuum. Christus totum habet spiritum, non ad mensuram, nos autem primitias tantum et gustum quen dam spiritus, ad mensuram, quibus contra peccatum pugnamus, donec totum aboleatur: baptismus durat interim per omnem uitam, hinc addit  \pend
\phantomsection
\addcontentsline{toc}{subsection}{\textit{Et nos ipsi in nobis ipsis gemimus. }}
\subsection*{\textit{Et nos ipsi in nobis ipsis gemimus. }}\pstart Id est, intus in corde nostro, suntque ; hi gemitus Spiritus sancti, nam alians nequaquam, cum habea- mus ferreum cor, ullus gemitus propter peccatum edi posset. Sed quare gemimus?  \pend
\phantomsection
\addcontentsline{toc}{subsection}{\textit{Adoptionem filiorum expectantes. }}
\subsection*{\textit{Adoptionem filiorum expectantes. }}\pstart Non intellige quod dicat, nos nondum esse adop- tatos, uerum quod adoptio nostra nondum sit reuela- ta: sumus uere iam filii Dei per fidem, sed non palam factum, primum autem illud in resurrectione manife- stum fiet. Quid autem speramus futurum in illa reue- latione adoptionis? nihil aliud, quam  \pend
\phantomsection
\addcontentsline{toc}{subsection}{\textit{Redemptionem corporis nostri. }}
\subsection*{\textit{Redemptionem corporis nostri. }}\pstart Id est, tum redimetur hoc nostrum corpus a serui- tute et corruptione et peton, quod impossibile est fie- ri in hac uita. Quare, ut maxie doleamus et ingemisca- mus ob haec oia, malimusque este liberi, audienda em uox  \pend
\vspace{0.5cm}\noindent
\vspace{0.2cm}\rule{1cm}{0.2pt}\\ 
\hspace{0.2cm}\textit{mg}
\footnotesize In resurrectio ne palam fiet nos filios esse 
\normalsize| 
\section*{COMMENT. POMER.  }\pstart benigni et pii patris, qua quemuis in cruce et perse cutione stantem, solatur, dicens: Sufficit tibi gratia mea, ut maxime sis peccator, sustine, sufficit, haec est mea uoluntas. q. d. quod in te non habes, in me ha bes, perfer saltem ad tempus hanc crucem: sed haec ipse exponit.  \pend
\phantomsection
\addcontentsline{toc}{subsection}{\textit{Siquidem spe seruati sumus. }}
\subsection*{\textit{Siquidem spe seruati sumus. }}\pstart Saluati sumus, sed sub spe: sumus uere filii a Deo suscepti, at illud non uidetur, omnia sunt abscondita sub formis crucis, nihil aliud uidetur, quam ignomi- nia et opprobrium, spectaculum sumus mundo, et ipsius πρί ημα: Ideoque ́, ut maxime iam omnia pos sideamus, ut haeredes Dei, et Domini creaturarum, tamen reuelationem huius dominii et haereditatis no- strae expectamus: Ipse promisit, non fallet nos, pri- us coelum et terra transitura scimus, quam Verbum etus: quod dum credimus, habemus omnia quae no- bis in illo promittuntur: ergo expectandum est cum pa- tientia. Non in crucem et eius amaritudinem diriga- mus oculos, sed in solum Verbum, in quo omnes diui- tias et plenitudinem inueniemus, id quod sola fide fit, quae in iis, quae non palam uidentur, uersatur. Haec his quoq uerbis Iohannes expressit: Scimus quod si- mus filii Dei, sed nondum apparuit, id est, nisi ex fi- de essemus certi, nos esse filios, dominosque ; omnium cre aturarum, non possemus expectare reuelationem eius gloriae, que futura est.  Porro  \pend
\vspace{0.5cm}\noindent
\vspace{0.2cm}\rule{1cm}{0.2pt}\\ 
\hspace{0.2cm}\textit{mg}
\footnotesize Gloria nostra nondum reue- lata est.  
\normalsize| 
\section*{97 IN EPIST. AD ROM.  }
\marginpar{[ p.]}\pstart Id est, res quae speratur . SI VIDEATVR.  Scilicet, coram . NONEST SPES. Id est, spes tantum est de futuro bono, non praesen- ti. q.d. sumus adhuc in regno fidei, id quod sumus, nempe, filii Dei, nondum est reuelatum . QVOD ENIM QVIS CERNIT. Scilicet, ob ocu los positum . SI VERO QVOD NON VIDEMVS. Scilicet, adesse.   \pend
\phantomsection
\addcontentsline{toc}{subsection}{\textit{Porro spes. }}
\subsection*{\textit{Porro spes. }}
\phantomsection
\addcontentsline{toc}{subsection}{\textit{Speramus, id per patientiam expectamus. }}
\subsection*{\textit{Speramus, id per patientiam expectamus. }}\pstart Scilicet, crucis: interim opus est patientia. Simi- lia aduertis in carnalibus exemplis, cum speras te pul- chram et bonam uxorem ducturum, sustines omnia, Eurris sursum ac deorsum, imo et in feruitutem te coniicis propter eam, ut lacob fecit: Spes est quae te conseruat et patienter omnia ferre iubet. Praeterea fi quid magni tibi promissum fuerit, id expectas sum- mo desyderio, sed priusquam id ueniat, multa inte- rim sunt sustinenda, quae intercidunt. Quemadmo- dum pater puerum Paedagogo tradit imbuendum, is cogitur aliquandiu ferre tyrannidem et plagas, donec intelligat se et honore et doctrina auctum, quae res tum est admodum iucunda, licet priora peracerba uidebantur, quae et patientia uicit: ita Christiano non msi per patientiam in cruce est permanendum, do- Spes.   \pend
\vspace{0.5cm}\noindent
\vspace{0.2cm}\rule{1cm}{0.2pt}\\ 
\hspace{0.2cm}\textit{mg}
\footnotesize 
\normalsize| 
\section*{COMMENT. POMER.  }\pstart nec ueniat, qui uenturus est, quemadmodum uides in Abacuc Propheta.   \pend
\phantomsection
\addcontentsline{toc}{subsection}{\textit{Consimiliter autem et spiritus auxiliatur infirmitatibus nostris. }}
\subsection*{\textit{Consimiliter autem et spiritus auxiliatur infirmitatibus nostris. }}\pstart Consolatio haec est in primis memoranda, siquidem ducit argumentum Paulus a facili. q. d. Quamuis patiamur et Sathanae et mundi iniurias, idque ; omnium creaturarum consortto, tamen ne ulla desperationis co- gitatio nos uincat, qua ob oculos nobis ponitur, Deum non esse nostrum Deum, non esse patrem nostrum: Item, oblitum nostri, scitote, et ipsum Dei spiritum uobis auxiliari, non putate uos derelictos, etiam in perturbatissima cruce, siquidem non potest non esse uobiscum, etiam in morte, Dei spiritus: hic uos con- seruat, hic aduersus peccatum uos confirmat, nam cum sitis filii Dei, qui fieri potest, ut non omnia agat in uobis et sustineat? Est hic arrabo uobis, uos pertinere ad regnum Christi, hic solus certas reddit conscientias uestras, uos esse in Dei manu, omniaque  fieri Dei uoluntate. Multi uero, qui hoc carent spiri- tu, in tentatione deficiunt: Iudas et Cayn absorpti sunt ab inferno, nam hoc auxiliatore carebant. Vi- des hoc horrendum Dei iudicium, in Parabola, de se mine expositum, nam quantula pars est eorum, qui per terram bonam significantur? Tres sunt partes co- rum, in quibus uirtus Dei non operatur, quantumuis  \pend
\vspace{0.5cm}\noindent
\vspace{0.2cm}\rule{1cm}{0.2pt}\\ 
\hspace{0.2cm}\textit{mg}
\footnotesize Spiritus san- cti auxilium.  
\normalsize| 
\section*{IN EPIST. AD ROM.  98 }
\marginpar{[ p.]}\pstart etiam splendescant et niteant, mundi quadam po- tentia et sapientia. Primi conculcantur pedibus eo- rum, qui per uiam transeunt, deuorantur item a uo- lucribus coeli, quae aliud mihil sunt, quam impii do- ctores, siue Daemones, qui transeunt et uolant per mundum, omnia inquinantes, quae sancta sunt, etiam si quid reliquum est deuorantes, non enim permit- tunt, ut fructum in iis faciat hoc semen benedictum. Alteri, typo saxorum notantur, in quibas Verbum non potest altius radices agere, nam ubi crux im- minet, recedunt, negantque ; sanctum Det Euangelium. Postremi spinarum nomine indicantur, qui item mun- di huius solicitudine adflicti, nequeunt Deum inuoca- re, et illius bonttate fidere. Soli ergo pii fructum adferunt, sed in pattentia, quam nemo habet, nisi spiritus Dei adsit, quandoquidem natura, ubi- se flagellari uidet iusto Dei iudicio, statim oggan- nit, Deum sibi aduersari, trasct, ac tyrannum es- se, ideoque Spiritus sanctus in nobis, intelligens pa- ternam Dei uoluntatem, reprimit naturam, di- cens: Necesse est, ut patiaris, nam haec est bo- na Dei uoluntas. Videtur autem in cruce saepenu mero, hic sprritus se abscondere, et non adesse, uerum prodit tamen suo tempore, uult enim Deus, ut ores, et inuoces illius auxilium de sancto. Oeculte hic spiri- tus dicit: Fiat uoluntas tua, ut maxime non sentias.  Sit nobis exemplo, Christus in horto, qui dixit: Au-  \pend
\vspace{0.5cm}\noindent
\vspace{0.2cm}\rule{1cm}{0.2pt}\\ 
\hspace{0.2cm}\textit{mg}
\footnotesize Parabola se- minantis. 
\normalsize| 
\section*{COMMENT. POMER.  }\pstart fer a me calicem: Spiritus autem, Fiat uoluntas tua - Idem Susannae docet historia, cui cum mors intentare tur, nisi consentiret senibus : caro in summa fuit an- gustia, at spiritus dixit : Melius est, ut incidam et c.  et nemo est hodic, qui non time at mortem, quoties au- dit tumultuari Principes et Episcopos aduersus Euan gelium, hic cogitat caro, actum est de nobis, iam iam occidemur, ue nobis, quid futurum est nobiscum? Sed contra spiritus occulte dicit: Ah iniquum et im- pium facinus, negabis Dei Euangelium? Timebis bul- las carnis et Sathanae? quid facient nolente Deo? etiamsi omnes creaturae insultent, tamen nihil possunt efficere, consiste, consiste, sustine, haec Det est uo- luntas, satius est mori pro illius Verbo, quam uiuere, et eius misericordiam negare. Hanc spiritus uocem, audit Deus, non eam, quae carnis est, quae semper de- trectat crucem Tolerat autem Deus in nobis has in- firmitates, qui non uersamur secundum carnemt, sed qui primitias has spiritus habemus, uidet quod non sit nostrarum utrium perdurare in cruce, non tamen abii cit, propter eum spiritum clamantem in nobis, et agnoscentem bonam eius uoluntatem. Quem, quae- so, haec non consolarentur, ettam in profundissima inferni uoragine? sed exponit clarius quae dixit.   \pend
\phantomsection
\addcontentsline{toc}{subsection}{\textit{Siquidem hoc ipsum quod etc. }}
\subsection*{\textit{Siquidem hoc ipsum quod etc. }}
\vspace{0.5cm}\noindent
\vspace{0.2cm}\rule{1cm}{0.2pt}\\ 
\hspace{0.2cm}\textit{mg}
\footnotesize Carnis et sp ritus pugna in aduersis. 
\normalsize| 
\section*{99 IN EPIST. AD ROM.  }
\marginpar{[ p.]}\pstart Dignissimius hic locus est de uera oratione, quam et Deus exaudit, haec spiritus est, non carnis nostrae, fiquidem fit carne ipsam ignorante, tantum coram Deo- Quid igitur trepidabimus in periculis, aut morte? cumn spiritus uere inuocet nomen Dei, nunquam est audienda caro, quae omnia crucis declinat, siquidem non intelligit, quid sub ea lateat, nempe optima in nos Dei patris uoluntas, qui suos filios ita uult castigatos: cur desperabimus, cum non uideamus id nobiscum ge- ri, quod nos uelit secundum Dei uoluntatem in despe- rationem coniicere? Alia quidem est facies crucis, alta paternae uoluntatis ratio. Hanc caro ubi cernit et sentit, putat se perituram, illam uero cum fides, spi- ritus Det in nobis, cognorit, mirum in modum nos confortat, et contemnere facit quicquid est portarum inferi. Exhorruit et Dauid mortem, idque ; ex carnis in firmitate, sed quoties uides in Psalmis, in summa tem- tatione dici: Exaudiuit Dominus deprecationem me- am etc. et similia. In ipsa cruce orat spiritus, dum interim caro nil aliud uult, q̃ desperare, non enim uidet crucem nobis prodesse ad salutem, quod tamen spiritus probe nouit. Hinc uocat eam orationem, uel intercessionem spiritus pro nobis, gemitus inenarra- biles, id est, talem affectum occultissimum in nobis, quouere dicimus, Abba pater, et qui solus probat Dei uoluntatem, ut maxime in cruce et horrore sit sepultus: quo fit, ut recte gemitus dixerit, nam tam  \pend
\vspace{0.5cm}\noindent
\vspace{0.2cm}\rule{1cm}{0.2pt}\\ 
\hspace{0.2cm}\textit{mg}
\footnotesize Oratio uera.  
\normalsize| 
\section*{IN EPIST. AD ROM.  }
\marginpar{[ p.102 ]}\pstart dulcis in Deum adfectus erumpit, et se exerit, ex ipsa amaritudine, ita, ut uideas gaudium, in tristi- tia sepultum, tristitiam in gaudio, mortem in ui- ta, uitam in morte, desperationem spet permixtam, spem desperationi. Sed cur pluribus haec persequor, sciens ea in ipsa experientia, non ex chartis condisci-  \pend\pstart Deus qui nouit profunda cordium, audiens oran- tem spiritum, se statim illi accommodat, siquidem ui- det eum intercedere, non iuxta carnis uoluntatem, uidet suam et illius uoluntatem esse eandem. Ergo ne minem Deus exaudit, nisi oret secundum ipsius spiri- tum, qui in oratione errare non potest.   \pend
\phantomsection
\addcontentsline{toc}{subsection}{\textit{At ille qui scrutatur. }}
\subsection*{\textit{At ille qui scrutatur. }}
\phantomsection
\addcontentsline{toc}{subsection}{\textit{Intercedit pro sanctis. }}
\subsection*{\textit{Intercedit pro sanctis. }}\pstart Non impiis, in quibus spiritus Dei non est, hinc etiam non orant -  \pend
\phantomsection
\addcontentsline{toc}{subsection}{\textit{Scimus autem, quod his. }}
\subsection*{\textit{Scimus autem, quod his. }}\pstart Argumentum aliud ab utili, dicit enim: tan- tum abest, ut propter aduersitates, quas sustinemus, deficiamus a Christo, ut etiam adiuuent nos: nihil ex- cipitur, omnia mala sunt nobis commodo, nempe, gla- dius, mors, Sathanas, impii qui infestant, adeoque ; ip- sum peccatum. Mira itaque haec est commendatio gratiae, quod ea quae nos uelint damnare, cogantur nobis seruire ad salutem: et hoc ideo fieri, ꝙ simus electi, in quo di-  \pend
\vspace{0.5cm}\noindent
\vspace{0.2cm}\rule{1cm}{0.2pt}\\ 
\hspace{0.2cm}\textit{mg}
\footnotesize Gratie com- mendatio. 
\normalsize| 
\marginpar{[ p.N 4 ]}
\marginpar{[ p.]}\pstart scamus salutem nostram in electione consistere, nihil meriti in nobis inueniri, ut sola praedicetur gratia et Det mi- sericordia-Stergo, quae pessima sunt, cooperantur uobis in salutem, quid dubitabimus de reliquis crea- turis? Et Paulo datus erat stimulus Sathanae, qui mi- hil aliud erat, quam gratia, qua custodiretur, ne in- cideret in laqueos Sathanae, nempe, superbiam, odia, fraudes etc. Non deberemus indignius ferre, si quae durius in nos agit Deus, nam haec omnia ta- lia sunt, quibus nos seruet, et fidem nostram probet:  \pend
\phantomsection
\addcontentsline{toc}{subsection}{\textit{Qui diligunt. }}
\subsection*{\textit{Qui diligunt. }}\pstart Nam qui odiunt Deum, omnia eis cooperantur in malum et damnationem.  \pend
\phantomsection
\addcontentsline{toc}{subsection}{\textit{Nimirum his, qui iuxta propo situm uocati sunt. }}
\subsection*{\textit{Nimirum his, qui iuxta propo situm uocati sunt. }}\pstart His uerbis excluduntur Iusticiarii, qui se suis ope- ribus tustificant, non Det electione. Saluamur itaq ex proposito Dei, non nostro, quicunq saluamur.   \pend
\phantomsection
\addcontentsline{toc}{subsection}{\textit{Quoniam quos praesciuerat. }}
\subsection*{\textit{Quoniam quos praesciuerat. }}\pstart Locus hic de praedestinatione utilissimus est, iis qui adfliguntur in hac carne sanctis. Contra, impiis et carnali sapientiae inutilissimus, quemadmodum et hodie uulgo uidemus fieri inter pocula, quoties mentio fit rerum sacrarum, insigniter uideri uolunt quidam  \pend
\section*{COMMENT. POMER.  }\pstart disputare tantum de praedestinatione, cum tamen nunq didicerint, quando cogitanda et meditanda sit prae - destinatio, quare prorsus est talibus et tam curiosis hominibus inutilis. Est autem hoc caput praedestina- tionis, ut sciamus nos mortificandos, quemadmodum filius Dei mortificatus est: glorificandos item, quem admodum filius glorificatus est, hoc qui non iudet et credit, praedestinationis solerti cogitatu, seipsum se ducit et perdit, id quod ex ipso textu hic clarum est, nam dicit: Quos praesciuerat, scilicet ab aeterno, uel suo consilio elegerat (ut in Ephes- est: Dilecti sumus in Christo, ante mundi constitutionem) hos et praefi- niuit, id est, ordinauit ad ea omnia, quae erant futu- ra in Christo filio suo, ut et cruce et glorificatione similes Christo fierent, et filii Dei essent, sed per pri- mogenitum Christum. Hoc qui credunt, gaudent, nec quid sentiunt aduersi in sua conscientia. Nam quid est dulcius et iucundius, quam scire, nos diuina ele- ctione saluandos, non nostro merito, siquidem dixit: Vos non me elegistis, sed ego elegi uos. In hoc ubi re- spicimus, excidimus ab omni fiducia humana, sci- mus enim nos esse in manu Dei, ex qua nemo nos po- test rapere, ergo perire non possumus. Periissemus uero iam dudum, si in nostra manu salus collocata di- ceretur. Caueamus igitur nobis, ne in solicitudinem praedestinationis incidamus cum carnalibus, qui per totam uitam turbantur praedestinatione, quae, quia Dri  \pend
\vspace{0.5cm}\noindent
\vspace{0.2cm}\rule{1cm}{0.2pt}\\ 
\hspace{0.2cm}\textit{mg}
\footnotesize Praedestinatic nis utilitas. 
\normalsize| 
\section*{IN EPIST. AD ROM.  }
\marginpar{[ p.101 ]}\pstart est et spiritualis, a carne intelligi non potest, imo so- li spiritu, placet: Tonitru quoddam est in animo carna- lis hominis praedestinatio, quo uastantur et proster- nuntur omnia, id est, quo tollitur omnis Dei timor el- reuerentia, nam dicunt: Si non est in manu mea salus, quid- uis faciam: Si Deo est uisum, saluabor, etiansi nihil boni faciam, sin minus, peribo, ettansi sanctissima uita fulgeam. His argumentis caro nixa perit, nam sunt inextrica- bilia, et solui nequeunt, siquidem ad iudicia Dei per- tinent, quae nobis sunt incognita. Gratias ergo aga- mus Deo, qui dedit nobis in Euangelio agnitionem sui Christi, quo fidimus, et cui nos totos committimus, cer- to scientes, petam nostra esse illius sanguine sublata.  \pend
\phantomsection
\addcontentsline{toc}{subsection}{\textit{Vt ipse sit primogenitus. }}
\subsection*{\textit{Vt ipse sit primogenitus. }}\pstart Ipse natura est filius Dei, prior et ante secula dilectus patri, nos uero filii adoptione, facti quoq dilecti per illum patri: est tamen primogenitus, ut maxime sit fra- ter noster, non essemus facti filii Det, nisi ipse fuisset filius natura, et in quo sibi complacuisset pater. Ergo quocunque respicias, Christus est primogenitus, siue in diuinttatem, nam natura Deus est, nos electione Dii sumus, siue in resurrectionem mortuorum, ipse quoq est primogenitus ex mortuis, quando glorificato eo dictum est: Ego hodie genuite: Sede a dextris meis.  Si in creationem respicias, est primogenitus ante om- nem creaturam, ut est in Colos. Vt sit ipse primatum  \pend
\vspace{0.5cm}\noindent
\vspace{0.2cm}\rule{1cm}{0.2pt}\\ 
\hspace{0.2cm}\textit{mg}
\footnotesize Caro offendt tur praedesti- nationis sen- tentia.  
\normalsize| 
\section*{COMMENT. POMER.  }\pstart tenens in omnibus etc. Multi sunt Dii, sed per ip- sum Christum, multi sunt sancti, multi sacerdotes, sed per illum. Ergo si glorificari uolumus, sequamur Christum, in quo ad gloriam electi sumus, ita tamen ut per uiam crucis et patientiae ingrediamur.  \pend
\phantomsection
\addcontentsline{toc}{subsection}{\textit{Porro, quos praedefinierat. }}
\subsection*{\textit{Porro, quos praedefinierat. }}\pstart Quos statuerat saluandos, ut uideas, non esse uo- lentis aut currentis, sed miserentis, quae res tamen uehementer offendit carnem, quare hic sinito Deum plura scire, quam te. Gratiae est haec commendatio, non meriti, quod gratissimum est omnibus Christia- nis. Eramus enim inimici Dei pro tempore, nunc ue- ro facti per gratiam filii Dei, hinc dicit  \pend
\phantomsection
\addcontentsline{toc}{subsection}{\textit{Eosdem et uocauit. }}
\subsection*{\textit{Eosdem et uocauit. }}\pstart Scilicet, in tempore, quamuis ab aeterno electi- simus a Deo ad uitam, tamen permisit nos errare ad tempus in Dei blasphemia: at quando uoluit, uoca- uit nos. Itaque hunc locum non de uocatione aeterna, ed temporis intelliges.  \pend\pstart Scilicet, uocatione temporis, per suum Verbum, et reuelatum Euangelium,  \pend
\phantomsection
\addcontentsline{toc}{subsection}{\textit{Et quos uocauit. }}
\subsection*{\textit{Et quos uocauit. }}
\phantomsection
\addcontentsline{toc}{subsection}{\textit{Los et iustificauit. }}
\subsection*{\textit{Los et iustificauit. }}\pstart Quo iterum omne meritum tollit humanum: quot quot igitur uocat Deus, eos omnes inuenit in pecca- to. Etiam Iohannes Baptista, testatur se esse peccato-  \pend
\vspace{0.5cm}\noindent
\vspace{0.2cm}\rule{1cm}{0.2pt}\\ 
\hspace{0.2cm}\textit{mg}
\footnotesize Electio uel praedestina- tio, est ab eterno. 
\normalsize| 
\hspace{0.2cm}\textit{mg}
\footnotesize Vocatio, a tem pore. 
\normalsize| 
\section*{COMMENT. POMER.  }\pstart rem, cum dicit: Ego deberem a te baptizari.   \pend
\phantomsection
\addcontentsline{toc}{subsection}{\textit{Quos iustificauit, hos et glorifi. }}
\subsection*{\textit{Quos iustificauit, hos et glorifi. }}\pstart Hanc gloriam per spem iam habemus, accepturi eam olim re ipsa. Quare diligens lector, haec omnia ex admiratione quadam, per Paulum dicta intelligat, ne putemus nos posse excidere, qui Euangelio Dei non abutimur: quotquot autem abutuntur, non pos- funt esse certi de gratia. Dicit ergo Apostolus, quid est, quod crucem aegreferatis, cum sitis praedestinati, uocati, iustificati et glorificatt? haec est uia, qua ad Deum conscenditur, nolite timere Sathanam et im- pios, quantumuis ettam insaniant, nihil enim praeua lebunt, siquidem certissimi estis, de bona Dei erga uos uoluntate. Sed quia hic de uocatione fit mentio, obii- ciunt quidam ex Euangelio: Multi sunt uocatt, pau- ci uero electi. Sciant hos locos non posse inteligt de una uocatione, Christus enim in Euangelio loquitur de uocatione, uel inuitatione sanctorum, qui nos uo- cant, ut serui Dei ad Euangelium, id quod uides satis clarum in nupttarum ratione, nam quidam, tan- etsi uocati, non ueniunt, quidam ueniunt, sed sine ueste nuptiali, quales sunt omnes, qui abutuntur Euangelio, in libertate carnis. At Paulus de ea uo- catione loquitur hic, qua Deus nos trahit ad se per suum Verbum, id quod ministri Verbi non pos- sunt: Traheris hac et raperis, neq alio te con- uertere potes, fit enim Dei spiritu, qui fortis est,  \pend
\vspace{0.5cm}\noindent
\vspace{0.2cm}\rule{1cm}{0.2pt}\\ 
\hspace{0.2cm}\textit{mg}
\footnotesize Matth . 3.  
\normalsize| 
\hspace{0.2cm}\textit{mg}
\footnotesize Vocatio du- plex. 
\normalsize| 
\section*{COMMENT. POMER.  }\pstart et cui nemo potest resistere. Illa autem uocatio, uirtu- tis nihil habet, nisi credatur Verbo, per ministros praedicato. Sic ergo conciliabis hos locos, nam ibi serui dicuntur inuitare uel uocare, hic autem Deus: praeterea hic dicitur de iis uocatis, qui ab aeterno ad boc praedestinati sunt, hi omnes sunt electi et glo- rificabuntur.   \pend
\phantomsection
\addcontentsline{toc}{subsection}{\textit{Quid igitur dicemus? }}
\subsection*{\textit{Quid igitur dicemus? }}\pstart Quid est, quod soliciti simus et impatientes? ecce sumus in manu Dei, unde excidere nequimus. non possum ego effari, inquit, haec mira Det uerba, qui sua benignitate et fauore nos ad se allicit, uti pa- ter filios, tamen tanta est nostra impietas, et fidei imbecillitas, ut de hac eius in nos uoluntate dubite- mus-Confitentur omnes creaturae, ipsum esse omni potentem, solus homo semper fluctuat, et timet, ne- coelum ruat, cum tamen quotidie Verbi praedicatio- ne, per ministros certificetur. Hac igitur ratione ui- des, quod Christianos faciat dominos omnium crea- turarum, nam si Deus proets est, quis potest esse contra eos?  \pend
\phantomsection
\addcontentsline{toc}{subsection}{\textit{Qui proprio filio non pepercit. }}
\subsection*{\textit{Qui proprio filio non pepercit. }}\pstart Similis sententia est in Iohan.3.Sic Deus dilexit mnundum etc-ut omnino certi simus, hanc gratiam ex Deo profectam, ex dilectione, qua ipse prior nos dilexit, non qua nos illum, quia eramus inimici na-  \pend
\section*{IN EPIST. AD ROM.  }
\marginpar{[ p.103 ]}\pstart tura. Quid posset fortius contra merita et Liberum Arbitrium dici? Scribant nunc magnos libros Liberi Arbitrii adsertores, uides uocationem, electionem, iustificationem esse Dei, non nostram, cum nihil aliud sit in nobis, quam peccatum et damnatio: tam cer- tos nos facit Apostolus, de ista Dei gratia, ut nulla creatura nos a Deo possit separari. Si nulla, ergo nec Sathan, qui posset ergo sordidum peccatum? Domi- nus pugnat pro nobis, nos sabbatum agimus. Talium locorum plena est scriptura, adeo, ut fere ubique pug- net contra merita nostra, et Liberum Arbitrium. His qui non uult moueri, lapis sit, necesse est, maxime cum audit proprium filium Det, pro se datum, non adop- tiuum, aut Angelum, aut hominem aliquem: quae er- go maior erga nos possit esse Dei dilectio? TRA- DIDIT. Scilicet, in mortem.   \pend
\phantomsection
\addcontentsline{toc}{subsection}{\textit{Qui fieri potest? }}
\subsection*{\textit{Qui fieri potest? }}\pstart Certissimum argumentum. Si dedit in mortem fi- lium suum unigenitum, per quem omnia fecit, et qui haeres est omnium creaturarum, qui fieret, ut etiam non daret haereditatem? Contemptius sane quiddam haereditas est, haerede: Abraham contempsit hanc, cum haerede careret, deit enim: Quid mihi cum di- uitiis, quas mihi dedisti, haeres nullus est? Ergo con- stanter credendum, quod Deus dederit pro nobis, et no- bis, Esauae y suum filium: ergo cum illo, et in ilo  \pend
\vspace{0.5cm}\noindent
\vspace{0.2cm}\rule{1cm}{0.2pt}\\ 
\hspace{0.2cm}\textit{mg}
\footnotesize Contra Lib- Arb.  
\normalsize| 
\section*{COMMENT. POMER.  }\pstart omnia accepimus, Itaque Christiani domini sunt omnium: uerum quemadmodum ridet caro sententiam praede- stinationis, ita et ridet hos dominos, nam uidet eos despectos, et formis crucis obnoxios, dicit eni: Sunt hi reges? O egregios mundi principes, interim nos talia non curantes certi sumus, cum filio data nobis omnia, cuius pedibus omnia sunt subiecta, Psalm. S- Non igitur redigamur sub potestatem ullius hominis 2. Corin.7. aut creaturae, cum sint omnia nostra- Magna haec sunt, quae tamen Paulus, ut sentit non satis potest explicare.  \pend\pstart Quis quaeso nos accusabit? nam si quis uoluerit ac- cusare, quem habebit iudicem? Deus illum non au- diet, quomodo enim ferret accusationem in dilectos filios suos? Si filii sumus, ille pater est, non ergo tu dex. Vix maior possit imbecillibus conscientiis adhi- beri consolatio.   \pend
\phantomsection
\addcontentsline{toc}{subsection}{\textit{Quis intentabit. }}
\subsection*{\textit{Quis intentabit. }}
\phantomsection
\addcontentsline{toc}{subsection}{\textit{Deus est qui iustificat. }}
\subsection*{\textit{Deus est qui iustificat. }}\pstart Qui fauet, ille remisit omne peccatum, ergo non est quod accusetur. O nos stultos, qui in adflictione conscientiae, respicimus in merita, in peccata nostra: Inefficatia sunt omnia, non potest te accusare Satha- nas, quod minus bonorum operum, quam oportebat fece- ris, neque etiam peccatum te tugulabit, curi Deus per Christum sustulit illud. Subiungit autem et aliam interroga  \pend
\vspace{0.5cm}\noindent
\vspace{0.2cm}\rule{1cm}{0.2pt}\\ 
\hspace{0.2cm}\textit{mg}
\footnotesize Carnales ri- dent gloriam plorum.  
\normalsize| 
\section*{IN EPIST. AD ROM.  104 }
\marginpar{[ p.]}\pstart tionem, quo magis exponit, optimum et Apostolicum suum adfectum erga omnes Christianos, dicens-  \pend
\phantomsection
\addcontentsline{toc}{subsection}{\textit{Quis ille, qui condemnet? }}
\subsection*{\textit{Quis ille, qui condemnet? }}
\phantomsection
\addcontentsline{toc}{subsection}{\textit{Nemo . Ecce hic }}
\subsection*{\textit{Nemo . Ecce hic }}
\phantomsection
\addcontentsline{toc}{subsection}{\textit{Christus est, qui mortuus est, }}
\subsection*{\textit{Christus est, qui mortuus est, }}\pstart Nullum hic est peccatum in sanctis, quo conden- nentur, peccatum eorum omne Christi morte est sub- latum. Praeterea, ille ipse Christus, qui mortuus est, resuscitatus est, in iustificationem nostri, nam mori Chri- stum, nihil profuisset nobis, nisi et ex morte in uitam transisset, haecque est gloria et uictoria nostra. Insuper ille Christus sedet ad dexteram patris, id est, rex est et dominus omnium, uti pater, et intercedit pro no- bis, id est, summus est Pontifex, placans patrem quem offenderamus peccatis nostris. Quomodo ergo quis- piam condemnaret nos, quum Deus ipse non sit iu- dex, sed pater noster? Christus uero interpellator no- ster apud patrem. St hi non accusant et condemnant nos, nemo poterit nos accusare et condemnare.  Gloria nostra  \pend
\phantomsection
\addcontentsline{toc}{subsection}{\textit{Quis nos separabit a dilectione Dei? }}
\subsection*{\textit{Quis nos separabit a dilectione Dei? }}\pstart id est, qua nos ipse diligit, uel qua nos diligi- mus Deum, siquidem dat nobis gratiam, ut ipsum diligamus: malo tamen intelligi de Dei dilectione in nos, ut infra: Per eum qui dilexit nos, ut eadem sit sententia, quae hic est: Quis nos separabit a dile- ctione Dei? cum ista, nemo poterit nos rapere e  \pend
\vspace{0.5cm}\noindent
\vspace{0.2cm}\rule{1cm}{0.2pt}\\ 
\hspace{0.2cm}\textit{mg}
\footnotesize 
\normalsize| 
\section*{COMMENT. POMER.  }\pstart manu patris, nam nostra dilectio saepenumero fluctu at, uti Dauidis et Petri Apostoli.  \pend\pstart Scilicet, separabit, aut aliqua crux: Minime. NVM FAMES? Nequaquam  \pend
\phantomsection
\addcontentsline{toc}{subsection}{\textit{Num adflictio. }}
\subsection*{\textit{Num adflictio. }}\pstart Nullo pacto. Siquidem haec omnia sunt talia, quae ad uitam Christianorum pertineant, ideoque ex Psalmo uersiculum adscripsit, quo sanctorum animus et con- stantia in omnibus his declaratur. Dicunt igitur, quo modo non mortem, famem etc-toleraremus, siqui- dem in hoc nati sumus, ut per totam uitam patiamur, et ueluti oues in lanienam pertrahamur.  \pend
\phantomsection
\addcontentsline{toc}{subsection}{\textit{Num gladius? }}
\subsection*{\textit{Num gladius? }}\pstart Scilicet, O Deus, propter tuum nomen, quia tan topere nos diligis et conseruas, quod nos agnoscen- tes profitemur coram toto mundo, tuam uirtutem et misericordiam: hoc quamprimum facimus TRA- DIMVR MORTI. Idest, mundus nos per- sequitur. TOTA DIE. Id est, iugiter.   \pend
\phantomsection
\addcontentsline{toc}{subsection}{\textit{Propter te, }}
\subsection*{\textit{Propter te, }}
\phantomsection
\addcontentsline{toc}{subsection}{\textit{Habiti sumus uelut oues desti - natae mactationi. }}
\subsection*{\textit{Habiti sumus uelut oues desti - natae mactationi. }}\pstart Quae, scilicet, ab hominibus in id tantum seruan tur, ut mactentur. Aestimamur et nos pari ratione a mundo, indigni hac uita: Sunt autem omnia nobis  \pend
\section*{IN EPIST. AD ROM.  105 }
\marginpar{[ p.]}\pstart dulcia per Christum, quem eadem passum scimus, qui non alia uia, quam per crucem in gloriam in- gressus est, ideo subtexuit,  \pend
\phantomsection
\addcontentsline{toc}{subsection}{\textit{Verum in omnibus his superamus }}
\subsection*{\textit{Verum in omnibus his superamus }}\pstart Id est, per tota pericula, non deficimus a Deo meritis ne nostris? non, sed per eum, scilicet, Deum, qui nos dilexit, is enim est pro nobis, ideoque ; nihil no- bis potest obesse.   \pend
\phantomsection
\addcontentsline{toc}{subsection}{\textit{Nam mihi persuasum habeo. Id est, certus sum, siue uiuam, siue moriar- }}
\subsection*{\textit{Nam mihi persuasum habeo. Id est, certus sum, siue uiuam, siue moriar- }}
\phantomsection
\addcontentsline{toc}{subsection}{\textit{Quod neque mors, neque angeli. }}
\subsection*{\textit{Quod neque mors, neque angeli. }}\pstart Scilicet, mali-Dicit uero haec Paulus hyperboll côs, simul quoque bonos angelos intelligens, tametsi neminem a Deo abstrahant, quemadmodum et in Galatis loquitur: Etiamsi angelus de coelo etc. quod tamen fieri non potest, ut angelus bonus praedicet con- tra Dei Verbum. Significat uero in iis omnes creatu- ras, bonas et malas, quae omnes cooperantur no- bis in bonum, tantum abest, ut separent a Deo, tam etsi uideantur nonnunquam id uelle-  \pend
\phantomsection
\addcontentsline{toc}{subsection}{\textit{Principatus. }}
\subsection*{\textit{Principatus. }}\pstart Id est, omne quod magnum est in mundo, qualla sunt, sapientia et potentia humana.   \pend
\phantomsection
\addcontentsline{toc}{subsection}{\textit{Potestates. }}
\subsection*{\textit{Potestates. }}\pstart Scilicet, nequitiae huius aëris, aut mundi, Ephes . 6.   \pend
\section*{COMMENT. POMER.  }
\phantomsection
\addcontentsline{toc}{subsection}{\textit{Neque instantia. }}
\subsection*{\textit{Neque instantia. }}\pstart Id est, illa quae iam adsunt, etiamsi pessima sint.   \pend
\phantomsection
\addcontentsline{toc}{subsection}{\textit{Neque futura. }}
\subsection*{\textit{Neque futura. }}\pstart Id est, quicquid malorum comminatur nobis mundus.   \pend
\phantomsection
\addcontentsline{toc}{subsection}{\textit{Neque altitudo, neque profunditas. }}
\subsection*{\textit{Neque altitudo, neque profunditas. }}\pstart Id est, neq coelum, neque infernus, neque terra, neque gloriosum, neque contemptum in mundo.   \pend
\phantomsection
\addcontentsline{toc}{subsection}{\textit{A dilectione. }}
\subsection*{\textit{A dilectione. }}\pstart Scilicet, quam adsecuti sumus per Christum Ie- sum Dominum nostrum. Vides quae hactenus scrip- ta sunt esse ignita uerba, quae debebant permouere corda nostra, ne resiliremus in periculis, cum sit per seuerandum in illis, secundum Domini uoluntaten- Et nisi id fecerimus, non agnoscet nos suos. Siquidem persecutio seu crux, est tessera Christianorum, qua illius filii agnoscimur.   \pend
\phantomsection
\addcontentsline{toc}{subsection}{\textit{}}
\subsection*{\textit{}}\pstart BSOLVTA prima parte, in qua docet, A Y quomodo peccatum nostrum sit agnoscendum, et unde remedium, nempe, remissio peccatorum, sit petendum, tractanda est et altera pars, quanta fieri potest diligentia, discemus enim ex ea, ad quos pertineat gratia, uel peccatorum remissio: tametsi  \pend
\vspace{0.5cm}\noindent
\vspace{0.2cm}\rule{1cm}{0.2pt}\\ 
\hspace{0.2cm}\textit{mg}
\footnotesize Crux Christi anorum tesse ra.  
\normalsize| 
\section*{IN EPIST. AD ROM.   }
\marginpar{[ p.]}\pstart in superioribus quoque ; illa nobis descripta sit, uerum quia imbecilles et infirmi sumus in his intelligendis, fu- siore et frequentiore declaratione opus est. Si igi- tur non ex nostra iusticia, aut urribus, consequimur salutem, sed ex proposito Dei, consequitur, ut illi tantum saluentur, quos Deus uult: sed quamdiu ui- uimus, murmurat caro contra hanc sententiam prae destinationis. Spiritus uero, qui intelligit sensum Do- mini, non reluctatur, sed se submittit: quae res, cum ita nobiscum habeat, quomodo possumus secundum carnem, non impii et blasphemi esse in Deum? nam ubi uidemus, non ualere nostra opera, statim inferi mus, ergo damnabimur, nam si in Dei uoluntate est, qui sciemus, an nostri curam habeat? tyrannus est, iudex est, immitis, potens, nulia spes nobis penes il- lum: sic Iudas periit, et Cayn. Ex hac impietate, nemo nos potest liberare, nisi sola Dei gratia, neq unquam̃ orabimus: Fiat uoluntas tua, dum sic sentimus: Deo ut nos commendemus, et in nihilum transeamus, necesse est, qui si etiam uelit nos damnatos, praesto si- mus: at nemo consentiet, nisi sprritu Dei sit plenus.  Omnes dicimus haec uerba: Fiat uoluntas tua, sed per hypocrisin: qui uere dicit, non aliud quam hoc uult quod Deus, qui si hunc uelit perdere, non remurmu rat: si etiam saluare, non inant gloria, ut solet caro, attollitur: submittit se, amore ductus, uoluntati eius, uon autem desperatione, id quod in impiis cerni-  \pend
\vspace{0.5cm}\noindent
\vspace{0.2cm}\rule{1cm}{0.2pt}\\ 
\hspace{0.2cm}\textit{mg}
\footnotesize Caro praede- stinatione, quon offendatur.  
\normalsize| 
\section*{COMMENT. POMER.  }\pstart mus. Si amor carnalis multa dura et mala susti- net propter puellam, haecque non curat, nec sentit esse dura et mala, quomodo dilectio in Deum non susti- neret et perferret in nobis omnia? Et huius exemplum, se ipsum hic Paulus nobis proponit, cum dicit se uelle Ana- thema esse a Christo pro fratribus. Ecce quam pius adfectus, quem in Mose quoque uides, cum dicit ad Do- minum: Remitte hanc culpam populo, aut dele me de libro uiuentium. Haec dicta sunt per hos uiros co- ram Deo, non ex hypocrisi: hoc denique uere est ab- negare seipsum, sed horrendum est carni. Hunc ad- fectum qui senserit, nempe, quod potest se Dei uolun- tati committere, etiam si uelit damnare, damnari non potest, nam per omnia gloriam Dei, et salutem pro- ximi quaerit. Nostra fragilitas ad hunc adfectum nno peruenit, non audemus talia cogitare, nedum cupe- re. Sed interim haec nostra est consolatio, quod ui- demus multos in speciem bonos, damnatos esse, ut Sau- lem : mnultos abiectos, iustificatos, ut Dauidem ho- micidam et adulterum, cum tamen illi posset Saul praeferri humano iudicio. Ergo hic uides nobis pro- poni hoc quod non intelligimus, non quod curiosi- tas nostra possit indagare: Dei iudicia sunt et consi- lia, non nostra: illi, in omnibus quaecunque geruntur nobiscum, adscribamus necesse est, non nobis, nam nostra tantum est confusio, et illius gloria. Quando igitur turbamur ob praedestinationem, gaudere debe  \pend
\vspace{0.5cm}\noindent
\vspace{0.2cm}\rule{1cm}{0.2pt}\\ 
\hspace{0.2cm}\textit{mg}
\footnotesize Dilectio ue ra Dei. 
\normalsize| 
\section*{107 IN EPIST. AD ROM.  }
\marginpar{[ p.]}\pstart inus, nam ea turbatio, portio quaedam est salutis.  Estque plane illud, quod Christus in Matthaeo dixit: Esurire et sitire iustitiam. Qui uero sunt absq spiri- tu Dei, sententiam huiusmodi praedestinationis irri- dent. Vacant autem omnes Dei spiritu, quotquot De- um non timent, neq turbantur, ex agnito peceato, hi non sitiunt iusticiam, ideoque ; neque saturabuntur.  Si seruauerimus hunc ordinem, quem hic Paulus ser- uat, in Christianismi forma, utilis erit nobis praede- stinatio: At si alia uia ingrediemur, nempe, initium sumentes a praedestinatione, non a peccati agnitione in nobis, erit nobis damnosa omnis praedestinationis cogitatio, incidemus enim in laqueum, et excaeca- bimur. Sed uide, quaeso, Pauli ordinem, a peccato Gentium et Iudaeorum, Epistolam auspicabatur, mox iusticiam et gratiam, qua tolleretur peccatum, de- pingebat, hac aucta, pugnam nasci perpetuam in nobis contra peccatum, docebat, in qua mortificare tur, et destrueretur uetus homo (haec enim est crux) His omnibus primum diligenter inculcatis, coepit in posteriore loco, de praedestinatione agere. Quem ordinem non est, quod tantum requiras ex ipso tex- tu, imo ipsa conscientia te edocebit, ut in Octauo di- xit de spiritu. Ergo stultissime errant, qui dicunt: Si scirem me praedestinatum ad uitam, facerem bona ope- ra, melius instituerem uitam etc. Pii contra, qui sci- ue salutem in manu Deipositam, eum laudant et  \pend
\vspace{0.5cm}\noindent
\vspace{0.2cm}\rule{1cm}{0.2pt}\\ 
\hspace{0.2cm}\textit{mg}
\footnotesize praedestinatio consoletur nos 
\normalsize| 
\section*{COMMENT. POMER.  }\pstart praedicant, certi sunt se non e manu illius rapi posse- Sciunt se in ipsouiuere et mori: Liberii Arbitrii ad- sertores, qui eo non perueniunt, haec omnia conten- nunt, et stulticiam quandam vocant . Et quemadmo dum non potes satis eloqui, cur tibi ignotus quispiam donet mille aureos, ita neque intelligis, cur te Deus (siuoles in te totum respicere) saluet. Scimus ergo et non scimus, nos esse Dei filios, nam in Psalmo di- xit: et scitote, quoniam segregauit Dominus san- ctum etc. Deinde non scimus, quia uidemus in tota Scriptura per spiritum nos admoneri, de Dei in nos bona uoluntate, nihil enim frustra is agit: idpsum quanta diligentia hic quoque Paulus id fecerit, per eun- dem spiritum, iudicabit mens Christo obnoxia. Nihil autem, quam hanc et similes sententias, in tota illa disputatione de praedestinatione, explicauit: Qui uult me sequi, abneget semetipsum. Quare ignari sunt Scripturarum et omnis pietatis, quicunque hac Pau- lina disputatione offenduntur: multa quiritantur de Libero Arbitrio, et suis bonis operibus, sed non sunt operarii et cultores uineae Domini, uerum cants et porci, qui eam tantum depopulantur. Sit autem haec summa trium sequentium Capitum, Noni, Decimi et Vn decimi: Gratia, de qua hactenus diximus, datur iis qui non sunt meriti, siue ii sint ex Iudaeis, siue Genti- bus: qui autem somniant sibi merita, et sua bona ope- ra, illi abiiciuntur, qui enim quaerunt iusticiam, non in  \pend
\vspace{0.5cm}\noindent
\vspace{0.2cm}\rule{1cm}{0.2pt}\\ 
\hspace{0.2cm}\textit{mg}
\footnotesize Pii certiss.. praedestinati- one consola- tionem habent 
\normalsize| 
\section*{IN EPIST. AD ROM.  }
\marginpar{[ p.108 ]}\pstart ueniunt, et qui non quaerunt, inueniunt. Fuit autem maxime necessarium, ut illa doceret Paulus, ne qua dubitaremus de gratia.   \pend\pstart Occupatione Pathetica utitur, quo magis perspi- cuum fiat, quam uehementer doleat perire Iudaeos, ex quibus ipse quoque natus erat. Hunc adfectum ha- bent omnes Christiant, cum audiunt quosdam peritu- ros: at ubi uident id Dei negotium esse, committunt omnia ipsius iudicio, et dicunt: Fiat uoluntas tua- Videt uero libenter in nobis Dominus talem adfectum, quia testatur cor nostrum, ex aequo in omnes homi- nes, siue amicos, siue inimicos, ex summa dilectio - ne propensum. Non potest enim quenq odisse cha- ritas, quae docet, increpat, rogat, inuocat, omnia pro inimicis facit, dum sperat illorum salutem, et temporalem et aeternam. Sed fides, quae in Dei tan- tum uoluntatem dirigit oculos, dicit: Domine, tu es qui perdis et redimis, pereat, qui periturus est, glo rificetur, qui glorificandus est, tua enim est potentia, non nostra-Itaque dicit Paulus: Veritatem dico, id est, ue- rum est testante idipsum in me Christo. Vides hic Paulum iurare, at non in sua caussa, non decet eni Christianum, ut uel iuret, uel litiget, aut uindictam quaerat propter se Iurarunt Apli, sed propter frens, propter Dei glo- riam amplificandam: poenas sumit Magistratus de malis,  \pend
\phantomsection
\addcontentsline{toc}{subsection}{\textit{Veritatem dico. }}
\subsection*{\textit{Veritatem dico. }}
\vspace{0.5cm}\noindent
\vspace{0.2cm}\rule{1cm}{0.2pt}\\ 
\hspace{0.2cm}\textit{mg}
\footnotesize Charitatis sol- licitudo pro aliorum salute.  
\normalsize| 
\marginpar{[ p.non ]}
\phantomsection
\addcontentsline{toc}{subsection}{\textit{\huge\textbf{COMMENT. POMER. }\normalsize Attestante conscientia. }}
\subsection*{\textit{\huge\textbf{COMMENT. POMER. }\normalsize Attestante conscientia. }}\pstart Id est, non loquor contra bonam meam conscien- tiam. Sed quomodo esse potest bona conscientia? per Spiritum sanctum. q. d. haec mea conscientia me fal- lere non potest, nam testimonium dat illi Spiritus san- ctus quod non erret. Quare hic discas, non satis es- se nostram conscientiam, in operibus aut institutis qui buscunque piis, nisi ipsa per Verbum Dei rectificata sit: unde fit, quod blasphemiae fuerint nominis Dei, omnes nostrae Missae, omnes item Monachorum mace- rationes. Ergo impium est habere fidem conscientiae suae, nisiuideat eam esse Verbo Det conformem, et per omnia respondentem. Addit uero  \pend
\phantomsection
\addcontentsline{toc}{subsection}{\textit{Quod dolor magnus. }}
\subsection*{\textit{Quod dolor magnus. }}\pstart Quomodo non doleam, aut uehementer crucier, qui uideam perituros Iudaeos? nam omnia secundum Dei gratiam fiunt, non secundum opera legis, aut me rita humana, quibus tamen illi innituntur.   \pend
\phantomsection
\addcontentsline{toc}{subsection}{\textit{Optarim enim Anathema esse. }}
\subsection*{\textit{Optarim enim Anathema esse. }}\pstart Id est, excommunicatus, separatus, maledictus, et cui nihil sit commertii cum Christo, quod sane ni- hil aliud est, q̃ sese dedere aeternae damnationi. Is adfectus fuit in Mose, cum ait Exodi. 32. Dele me de libro, quem scripsisti: cui Dominus respondit: Qui peccauerit mihi, delebo eum de libro meo. q. d. ista charitas, qua te exponis pro fratribus, significat te  \pend
\vspace{0.5cm}\noindent
\vspace{0.2cm}\rule{1cm}{0.2pt}\\ 
\hspace{0.2cm}\textit{mg}
\footnotesize uisse. 
\normalsize| 
\section*{IN EPIST. AD ROM.  109 }
\marginpar{[ p.]}\pstart non esse perdendum, usqueadeo est impossibile, ut hi perdantur, qui hoc adfectu insignes sunt: at pau cos inuenias, in quibus hunc depraehendas.   \pend
\phantomsection
\addcontentsline{toc}{subsection}{\textit{Qui sunt Israelitae. }}
\subsection*{\textit{Qui sunt Israelitae. }}\pstart Nam ex Ifraël sunt oriundi.   \pend
\phantomsection
\addcontentsline{toc}{subsection}{\textit{Quorum est adoptio. }}
\subsection*{\textit{Quorum est adoptio. }}\pstart Id est, ad eos pertinet haereditas filiorum Dei, ta- men non saluantur, quia respiciunt in se et dona, non in donantem Deum. Si ergo tales abiiciuntur, quorum ea adoptio, quid futurum est nostris sanctis, qui tantum in sua stercorea iusticia gloriantur? Hor- rendum sane est contra Iudaeos, quod dixerunt Apo- stoli in Actis: Ecce conuertimur ad Gentes-  \pend
\phantomsection
\addcontentsline{toc}{subsection}{\textit{Et gloria. }}
\subsection*{\textit{Et gloria. }}\pstart Scilicet, a patre, quod scilicet essent filii Dei, quae sane ineffabilis est gloria.   \pend
\phantomsection
\addcontentsline{toc}{subsection}{\textit{Et testamenta. }}
\subsection*{\textit{Et testamenta. }}\pstart Scilicet Dei. ET LEGIS CONSTITV TIO. Id est, ipsa tota lex, quemadmodum Psalm.  dicit: Non fecit taliter omni nationt etc.   \pend
\phantomsection
\addcontentsline{toc}{subsection}{\textit{Et cultus, }}
\subsection*{\textit{Et cultus, }}\pstart Dei, nam reliquae Gentes in hoc errabant, cum ni- hil haberent certi a Deo praescriptum.   \pend
\phantomsection
\addcontentsline{toc}{subsection}{\textit{Et promissiones. }}
\subsection*{\textit{Et promissiones. }}\pstart Scilicet, et praesentis et aeternae uitae.   \pend
\section*{COMMENT. POMER.  }\pstart Patriarchae, Prophetae, quae est maxima eorum gloria.   \pend\pstart Hic clare uides dici, Christum Deum esse, quis id per totum Iohannis Euangelium tractetur, imo to- ta Scriptura passim nihil aliud agit, quam ut id nobis clarum faciat: unde curiositatis esse putauerim, unum atque alterum locum tantum in Scrip-in hoc obserua- ri. Sed redibunt Arriana tempora, et homines absque  fide, qui negabunt, et Christum et Spiritum, scili- cet esse Deum, quod quid aliud erit, quam negatio remissionis peccatorum?  \pend
\phantomsection
\addcontentsline{toc}{subsection}{\textit{Et ii, ex quibus Christus est. }}
\subsection*{\textit{Et ii, ex quibus Christus est. }}
\phantomsection
\addcontentsline{toc}{subsection}{\textit{Quorum sunt patres. }}
\subsection*{\textit{Quorum sunt patres. }}
\phantomsection
\addcontentsline{toc}{subsection}{\textit{Non autem hoc loquor. }}
\subsection*{\textit{Non autem hoc loquor. }}\pstart Vt maxime illi pereant, tamen ueritas et Dei pro- missio manet firma. Non hic agitur secundum praero gatiuas, quas habuerunt Iudaei secundum carnem, sed simpliciter secundum Dei uoluntatem, qui iusto suo iudi- cio, carnalem Israëlem abiicit, spiritualem autem su- scipit. Verax est Deus in suis promissis, quantumuis etiam multi interim pereant, pertinent enim ad cre- dentes, non ad impios, Dei promissiones.   \pend
\phantomsection
\addcontentsline{toc}{subsection}{\textit{Non enim omnes, qui sunt exlsrael. }}
\subsection*{\textit{Non enim omnes, qui sunt exlsrael. }}\pstart Est expositio sententiae praecedentis: nihil Deo cum Ifraële secundum carnem, nihil illi cum semine Abrahae secundum carnem, et Ismaël erat filius Abrahae primo genttus, non tamen a Deo pro filio Abrahae habetur:  \pend\pstart  \pend
\vspace{0.5cm}\noindent
\vspace{0.2cm}\rule{1cm}{0.2pt}\\ 
\hspace{0.2cm}\textit{mg}
\footnotesize Christus De- us- 
\normalsize| 
\hspace{0.2cm}\textit{mg}
\footnotesize niii Abrahae 
\normalsize| 
\section*{n0 IN EPIST. AD ROM.  }
\marginpar{[ p.]}\pstart uult ille tantum eos, qui per promissionem sunt Abra- hae filii, qualis erat Isaac. Ideoque et Christus Iudaeos in Iohanne reprobat, dicens: Si filii essetis Abrahae, utique faceretis opera Abrahae, id est, crederetis in Deum, quemadmodum ille credidit. Fides filios Det fa- cit, non caro: ergo omnes sunt impiorum filii, quot- quot iusticiis externis nituntur. Sie conciliabis lohan nem et Christum, dum uariis nominibus Iudaeos appel- lant: iam genimina uiperarum, nationem prauam et adulteram, nunc ex Diabolo patre natos. Ex Ismaëlis familia, sunt omnes Monachi et operarii carnis, qui su- perbiunt coram hoc mundo, et quibus primus, pro- pter sanctitatem, qua nitent externe, defertur honos.  Ex Isaac autem familia, qui ignorant opera bona, qui bus confidant, sed tantum Christum crucifixum nouerunt.  \pend
\phantomsection
\addcontentsline{toc}{subsection}{\textit{Sed qui sunt filii promiisionis. }}
\subsection*{\textit{Sed qui sunt filii promiisionis. }}\pstart Promissio Det, est fides nostra, id est, quicumque fide Deo promittenti credunt, hi sunt uerum semen Abrahae-  \pend
\phantomsection
\addcontentsline{toc}{subsection}{\textit{Promissionis enim sermo hic est. }}
\subsection*{\textit{Promissionis enim sermo hic est. }}
\phantomsection
\addcontentsline{toc}{subsection}{\textit{In tempore hoc. }}
\subsection*{\textit{In tempore hoc. }}\pstart Vt in Genesi legis. q. d. Dominus: Vbi ad te redie- ro, anno reuoluto, habebis ex Sara filium. Mirum est, si in infoecundam et uetulam mulierem, aut in senem Abrahamum respiciamus, sed hic foecundum est Ver- bum Dei, ut maxime infoecundi sint, Sara et Abra- ham, ergo fidei filius est isaac, non carnis :Deus carnem  \pend
\section*{COMMENT. POMER.  }\pstart non moratur, quicquid etiam feceris tuis uiribus. Itaque  satis clarum est, ex primo exemplo, de duobus filiis Abrahae, solam fidem et Dei promissionem facere fi- lios Dei.  \pend
\phantomsection
\addcontentsline{toc}{subsection}{\textit{Non solum autem hoc. }}
\subsection*{\textit{Non solum autem hoc. }}\pstart Alterum est exemplum, quo sententiam suam con- munit Apostolus, nempe, sola Dei electione nos fie- ri Dei filios, non nostris meritis aut uiribus. Potuis- set enim sapientia carnis sic contra obiicere: Non mi rum est Ismaëlem abiectum, quia fuerat natus ex an- cilla, ipsum uero isaac susceptum in haereditatem, ꝗ Sara domina eum enixa sit: hanc obiectionem Paulus diluens, alterum hoc opposuit exemplum de lacob et Esau, quo sane grauissime perstringit in uniuersum totam humant ingenii, ac rationis acrimoniam, ut palam uideas, hic nihil aliud in caussa esse, quod ma- ior abiiciatur, et minor seruetur, quam solam Dei, quae latet omnem sapientiam carnis, uoluntatem et electionem, tametsi etiam haereditas ad maiorem natu, secundum legem (Deuter. 13.) pertinebat. Hic non dixeris, Deum fecisse contra legem, nam legem de- dit hominibus, sub qua omnes concluderentur: ipse ue- ro sub ea concludi non uoluit, liberum est illi quiduis agere. Ergo non alia huius factiratio hic dari potest, quam praedestinatio.  \pend
\phantomsection
\addcontentsline{toc}{subsection}{\textit{Rebecca ex uno conceperat. }}
\subsection*{\textit{Rebecca ex uno conceperat. }}\pstart  \pend\pstart  \pend
\vspace{0.5cm}\noindent
\vspace{0.2cm}\rule{1cm}{0.2pt}\\ 
\hspace{0.2cm}\textit{mg}
\footnotesize Electione filii facti sumus, non merito.  
\normalsize| 
\section*{IN EPIST. AD ROM.  }
\marginpar{[ p.m ]}\pstart Scilicet, congressu. q. d.Paulus: uerum est, di uersas matres fuisse, diuersos congressus carnis, in primo exemplo, at in hoc, una est mater, unus con  \pend
\phantomsection
\addcontentsline{toc}{subsection}{\textit{Nondum enim natis pueris, cum nequeboni etc. }}
\subsection*{\textit{Nondum enim natis pueris, cum nequeboni etc. }}\pstart Ecce ubi hic est meritum, cum nondum essent na- ti, omnia aguntur occulto et admirabili Dei iudicio.  \pend
\phantomsection
\addcontentsline{toc}{subsection}{\textit{Vt secundum electionem, pro positum maneret Dei. }}
\subsection*{\textit{Vt secundum electionem, pro positum maneret Dei. }}\pstart Id est, illud tantum firmum esset, quod Deus se- cum statuerat: ut cum quaereretur, quare sic factum est, responderetur, Deus sic uoluit.   \pend
\phantomsection
\addcontentsline{toc}{subsection}{\textit{Non ex operibus. }}
\subsection*{\textit{Non ex operibus. }}\pstart Scilicet, quae nulla adhuc fecerant, quia nondum nati.   \pend
\phantomsection
\addcontentsline{toc}{subsection}{\textit{Ded ex uocante. }}
\subsection*{\textit{Ded ex uocante. }}\pstart Id est, uoluntate Dei . DICTVM EST IL- LI. Scilicet Rebeccae, non propter opera, quae non fecerant, sed propter Dei uoluntatem.   \pend
\phantomsection
\addcontentsline{toc}{subsection}{\textit{Naior seruiet minori. }}
\subsection*{\textit{Naior seruiet minori. }}\pstart Sic iudicauit Deus de illis, antequamnascerentur, etiam contra omne ius primogeniturae. Isaac uolens benedicere primogenitum Esau, impeditur, namnon  \pend
\vspace{0.5cm}\noindent
\vspace{0.2cm}\rule{1cm}{0.2pt}\\ 
\hspace{0.2cm}\textit{mg}
\footnotesize Merita no- stra nulla.  
\normalsize| 
\section*{COMMENT. POMER.  }\pstart hunc, sed ipsum Iacob primogenitum Deus habere uo lutt. SICVT SCRIPTVM EST. In Malachia.  \pend
\phantomsection
\addcontentsline{toc}{subsection}{\textit{lacob dilexi, Esau odio habui. }}
\subsection*{\textit{lacob dilexi, Esau odio habui. }}\pstart Quare hunc dilexerit, illum odierit, nescimus, sic illi placuit.   \pend\pstart Hic non potest non obmurmurare ratio contra Deum, dicit enim haec audiens: Ergo Deus est iniustus, qui ista facit, nam hos abiicit, illos saluat, sine omni uel merito uel demerito: si omnes saluaret, id esset immem sae eius bonitatis et misericordiae indicium, et hanc pu- tassem eius esse naturam propriam: quemadmodum lapis non posset sursum agi natura, ita nec Deum cum bonus sit, quenquam posse abiicere. Respondet Apo- stolus . ABSIT. Id est, non est iniustus, qui haec faciat, uerum manet iustus.   \pend
\phantomsection
\addcontentsline{toc}{subsection}{\textit{Quid igitur dicemus ? }}
\subsection*{\textit{Quid igitur dicemus ? }}
\phantomsection
\addcontentsline{toc}{subsection}{\textit{Nam Mosi dicit, miserebor. }}
\subsection*{\textit{Nam Mosi dicit, miserebor. }}\pstart Deus nulli hominum quidquam debet, ideoque ne - mo illum accusabit: quemadmodum si Petro darem cen- tum aureos, Iohannt nihil, et neutri deberem, hic dices, quid in caussa fuerit, quod Petro dederim? libe- ralitas: quid in caussa, quod Iohannem praeterierim? mea uoluntas, quod dare noluerim (nihil hic dico de charitate, quae non quaerit delectum). Quid ergo hic dicet ratio? garriet multa, at uihil huius negotii  \pend
\vspace{0.5cm}\noindent
\vspace{0.2cm}\rule{1cm}{0.2pt}\\ 
\hspace{0.2cm}\textit{mg}
\footnotesize Rationis ar- gumenta, co ista facit, tra praedesti- nationis sen- tentiam.  
\normalsize| 
\section*{IN EPIST. AD ROM.  }
\marginpar{[ p.112 ]}\pstart intelliget: uerum certo sciamus, nos omnes in dan- natione esse natos, nec debere nobis Deum, ut nos ex ea liberet, iuste ergo omnes damnamur, si autem qui saluantur, id misericordiae et beneficio Dei tri- buamus necesse est. Nullum hic audis meritum, nihi in nobis esse, quod Deum placatum nobis reddat, id quod satis est manifestum e posteriore exemplo. Ille itaque Deo placet, cui Deus placet, ut inquit Augusti nus, quique uidens tam diuersa eius iudicia, dicit: Fiat uoluntas tua. Hoc item in parabola Euangelica Mat- thei 2o. exprimi uidemus, ubi paterfamilians dicit: An non mihi licet facere quod uolo, nunqnid me mo- uebit, quod oculus tuus nequam sit, qui contra bene ficentiam meam toruam dirigit aciem? Blasphemus es, legem mihi uis praescribere, ast ego non quod tu, sed quod ego uolo, faciam.  \pend
\phantomsection
\addcontentsline{toc}{subsection}{\textit{Itaque non est uolentis. }}
\subsection*{\textit{Itaque non est uolentis. }}\pstart Scilicet, hominis salus, uel illud quo indiget.  \pend
\phantomsection
\addcontentsline{toc}{subsection}{\textit{Neque currentis. }}
\subsection*{\textit{Neque currentis. }}\pstart Id est, annitentis et conantis suis omnibus uiri- bus- Non Liberum Arbitrium, non nostra sancti- tas hic quidquam est, sed sola Dei misericordia-  \pend
\phantomsection
\addcontentsline{toc}{subsection}{\textit{Dicit enim Scriptura, in hocip- lum excitauite. }}
\subsection*{\textit{Dicit enim Scriptura, in hocip- lum excitauite. }}
\phantomsection
\addcontentsline{toc}{subsection}{\textit{Id est, magnumfeci.  }}
\subsection*{\textit{Id est, magnumfeci.  }}\pstart Vt ostendam in te potentiam meam.  \pend
\vspace{0.5cm}\noindent
\vspace{0.2cm}\rule{1cm}{0.2pt}\\ 
\hspace{0.2cm}\textit{mg}
\footnotesize ᗞamnati su- mus, et sola misericordia fervamur.  
\normalsize| 
\section*{CO.MMENT. POMER.  }\pstart Hic Deus abusus est (si dicere liceret nostro iudi- io) Pharaone, quem creauerat, et in amplissimui. regnum constituerat, ut notum faceret in eo suam iram et potentiam, non Pharaont, uerum toti mundo: Perdidit illum propter electos, unde discerent om- nes homines, in Dei solius manu omnia esse, non, in nostra, quo tam horrendo exemplo, debebant hodie moueri principes et tyranni nostri, qui haud aliter iu sto Dei iudicio, sunt indurati contra populum Dei, quem nolunt dimittere ex Aegypto-  \pend
\phantomsection
\addcontentsline{toc}{subsection}{\textit{Itaque cui uult miseretur. }}
\subsection*{\textit{Itaque cui uult miseretur. }}\pstart Conclusio est, omnia sunt in illius uoluntate, non nostra- DICES ERGO MIHI. Scilicet, o homo.  \pend
\phantomsection
\addcontentsline{toc}{subsection}{\textit{Quid adhuc conqueritur. }}
\subsection*{\textit{Quid adhuc conqueritur. }}\pstart Scilicet, Deus. Secundum est argumentum carnis, quod per occupationem infert: Si Deus est autor bo- norum et malorum, facitque nos peccare, si eitis uo- luntas est ut peccemus, quod iam in Pharaonis exem- plo deprehendimus, cur ergo nos accusat et dan- nat, cum ipse potius sit accusandus, qui haec factat? Haec est stultitia sapientiae nostrae, aliud eni inferre non potest, uultque alligare Deum suae rationi, ut faciat non quod uult, sed quod illa intelligit, suo carnali iu- dicio, et dum maxime uult sapere, maxime Deum blasphemat. Est ergo inimica Deo, Rom. 8. mouetque ́  \pend
\vspace{0.5cm}\noindent
\vspace{0.2cm}\rule{1cm}{0.2pt}\\ 
\hspace{0.2cm}\textit{mg}
\footnotesize Pharaonis horrendum exem plum.  
\normalsize| 
\marginpar{[ p.113 ]}\pstart Perpetuo bellum contra illum, ut quondam Gigantes contra iouem, ut fabulae aiunt. Est igitur omnino in solubile hoc argumentum, ut Paulus in fine Cap. n.  indicat: neque aliud habemus, quod dicamus, quam o profunditatem etc. Creat uero ob duas caussas, im- pios Deus . Primum, propter se, ut ostendat in illis suam potentiam. Secundo, propter electos, ut ui- deant quantum gratiae ab illo acceperint, et ut per illos exerceantur, et saluentur: nam plurimum ex- empla illorum nos exercent, ne scilicet ignari iu- diciorum Dei impingamus in illa. Ergo non frustra- nascuntur impii, Deus illis utitur, et Sathanae regno in suam gloriam manifestandam. Multus est in his Hiob, quem uiderl potes, ubi cum amicis suis uarie de his rebus disputat, quare Deus neutiquam est accusandus qui illos crearit, qui etus Verbum abiiciunt, nam pari- ratione posset obstrepi Deo, cur uenenatas bestias cre- auerit. Sed audi impietatis uocem, quae non uult reprimi-  \pend
\phantomsection
\addcontentsline{toc}{subsection}{\textit{Nam uoluntati eius, quis resistit? }}
\subsection*{\textit{Nam uoluntati eius, quis resistit? }}\pstart Cogor facere quod uult: si igitur non saluor ex bens factis, cur damnabor ex malefactis, nam non ego quod inalum perpetro, facio, sed ille-Sed respondet Paulus  \pend
\phantomsection
\addcontentsline{toc}{subsection}{\textit{Atqui ohomo. }}
\subsection*{\textit{Atqui ohomo. }}\pstart O miser, tu figmentum es, et creatura, uis illi praescribere legem quid faciat: ecce disceptas cum co, quo quid stultius dici potest? Nonne ridiculum est, fi figmentum se figulo oppondt, et dicat: Quare me  \pend
\vspace{0.5cm}\noindent
\vspace{0.2cm}\rule{1cm}{0.2pt}\\ 
\hspace{0.2cm}\textit{mg}
\footnotesize impiicur cre- entur.  
\normalsize| 
\section*{COMMENT. POMER.  }\pstart fic formasti? Sed quid ille respondet? Nunquid lutum in manu mea est, an non possum eo uti ad uasa for- manda honoris, id est, quae etiam principibus pos- fint digne apponi, et ad uasa ignominiae, id est, in quae recipiantur sordes, et abiiciantur.  \pend
\phantomsection
\addcontentsline{toc}{subsection}{\textit{Ideoque Deus uolens ostendere iram, tulit. }}
\subsection*{\textit{Ideoque Deus uolens ostendere iram, tulit. }}\pstart Id est, protulit, ne intelligas sustinuit, quod uer - bum satis horrendum est, profert enim impios, us figulus uasa sua.   \pend\pstart Non est ut ex hoc loco probes, Deum non liber- ter damnare impios, imo ideo profert, ut in illis ex- erteat jauut tuuiciun, sea poses proruaie iuos, fere pa- fieler torn blajphenaaet corelpra, ue Pharuonu, et patitur Deus ab illis, in electis patientibus etc. , Vt notas faceret diuitias gloriae. Id est, abundantiam suae beneficentiae-  \pend
\phantomsection
\addcontentsline{toc}{subsection}{\textit{Multa animi lenitate. }}
\subsection*{\textit{Multa animi lenitate. }}
\phantomsection
\addcontentsline{toc}{subsection}{\textit{Erga uasa misericordiae. }}
\subsection*{\textit{Erga uasa misericordiae. }}\pstart Id est, ipsis electis, ut intelligerent, experienturque  ipsum solum esse potentem, et omnia agere potenter.   \pend
\phantomsection
\addcontentsline{toc}{subsection}{\textit{Quos et uocauit , nimirumnos. }}
\subsection*{\textit{Quos et uocauit , nimirumnos. }}\pstart Iam quod coeperat, pergit clarius exponere, nen- pe, gratiam non contingere, nisi electis, uel praede stinatis: Iudaeos uero abiict, Gentes suscipi: et in  \pend
\vspace{0.5cm}\noindent
\vspace{0.2cm}\rule{1cm}{0.2pt}\\ 
\hspace{0.2cm}\textit{mg}
\footnotesize Electis tantum contingit gra- tia.  
\normalsize| 
\section*{IN EPIST. AD ROM.  }
\marginpar{[ p.114 ]}\pstart summa, Iudaeorum et Gentium salutem constare ele- ctione, et Dei uocatione. Ait ergo, nos scilicet Chri- stianos, uel Apostolos,  \pend
\phantomsection
\addcontentsline{toc}{subsection}{\textit{Non solum ex ludaeis. }}
\subsection*{\textit{Non solum ex ludaeis. }}\pstart Qui tamen soli uidebantur pertinere ad gratiam.  \pend
\phantomsection
\addcontentsline{toc}{subsection}{\textit{Verum etiam ex Gentibus. }}
\subsection*{\textit{Verum etiam ex Gentibus. }}\pstart Scilicet Idolatris, in quibus nemo quicquam me- riti uidere potuit, erant enim Dei hostes.   \pend
\phantomsection
\addcontentsline{toc}{subsection}{\textit{Quemadmodum et Oseae di- cit, Vocabo populum. }}
\subsection*{\textit{Quemadmodum et Oseae di- cit, Vocabo populum. }}\pstart Hic populum Det uidemus fieri, Dei uocatione, non meritis.  \pend
\phantomsection
\addcontentsline{toc}{subsection}{\textit{Qui meus non erat. }}
\subsection*{\textit{Qui meus non erat. }}\pstart Gentes enim non dicebantur populus Dei se- cundum rationem temporis, non censebantur tum in populo Dei: coram Deo uero, ab aeterno fuerunt Dei populus, quia dilexit eas ante mundi constitutionem.   \pend
\phantomsection
\addcontentsline{toc}{subsection}{\textit{Et eam quae dilecta non erat. }}
\subsection*{\textit{Et eam quae dilecta non erat. }}\pstart Id est, quam non diligebam, Die meyn bul nit mar, Sic enim Scriptura, quasi de muliere loquitur.  \pend
\phantomsection
\addcontentsline{toc}{subsection}{\textit{Et erit ubi dictum fuerat eis. }}
\subsection*{\textit{Et erit ubi dictum fuerat eis. }}
\phantomsection
\addcontentsline{toc}{subsection}{\textit{\huge\textbf{COMMENT. POMER. }\normalsize Esaias autem clamat. }}
\subsection*{\textit{\huge\textbf{COMMENT. POMER. }\normalsize Esaias autem clamat. }}\pstart Locum hunc ex Propheta de Iusticiariis abiectis adducit, quo testatius faciat, quod dixit, q. d. Innu- meri erunt operarii, et qui in suas uires respicient, non in Deum, sed ipsi perdentur, et tantum reli- quiae, id est, pauci salui erunt, quod factum est in Apo- stolis, et Iudaeis credentibus-  \pend
\phantomsection
\addcontentsline{toc}{subsection}{\textit{Sermonem enim perficiens: }}
\subsection*{\textit{Sermonem enim perficiens: }}\pstart Lege potius perficit et abbreuiat, pro partici- piis. Et pro, cum iusticia, in iusticia. Paulus hunc locum citauit ex Septuaginta Interpretibus, non ex Hebraeo- Siquidem Romanis gnaris Graecae linguae, in manu erant Graeca, non Hebraica biblia. Sunt in Hebraeo aliauerba, at non alia sententia: in Paulo est obscu- rior. Vulgo sic interpretantur, quod Deus antea sub Veteri Testamento, dederit multa uerba, ceremonias, leges et praecepta, quae omnia uoluit in Nouo Testa mento colligere in breue quoddam uerbum, nempe, Dilige proximum, et crede Verbo Dei. Nunc ergo abbreuiatum est, et consummatum, perfectumque ́, id est, quod solum perficit, et in nobis consummat Dei tusticiam. Nihil tusticiae consequebantur ii, qui uiue- bant secundum Legem Mosi, et externa etus ope- ra, ut maxime perpetuo multa praeciperet, et urge- ret. Videtur tamen alia posse colligi ex uerbis Pro- phetae sententia: quod hic uerbum habet, in Pro-  \pend
\vspace{0.5cm}\noindent
\vspace{0.2cm}\rule{1cm}{0.2pt}\\ 
\hspace{0.2cm}\textit{mg}
\footnotesize Verbum ab- breuiatum 
\normalsize| 
\section*{IN EPIST . AD ROM.  }
\marginpar{[ p.115 ]}\pstart pbeta consummaiionem dici scias. q.d. Consummatio, idest, finitio, uel finis aderit in hoc populo, in iusti- cia, scilicet, fidet, nam hoc solum perficit, et con- summat omnia: lex infirma est, non potest praestare utres faciendi, quae exigit. Paulus uero, non de Ver- bo breui, citat hunc locum, sed de reliquiis saluandis tantum, ut si uelis ex Esaia interpretart, haec sit sen- tentia. Sermonem perficit et abbreuiat in iusticia, id est, Dominus finiet, siue perficiet, et breuiabit numerum sui populi ludaici, ut paucis uerbis nume- rari possit: et hoc in iusticia, id est, ille numerus pau- cus reliquiarum erit iustus, siue fidelis, aliis infidelitate pereuntibus etc. Quemadmodum hoc etiam anno- tauimus in Esaiam: et plane hanc sententiam uult Paulus-q. d. Videntur tam multi esse, qui sint po- pulus Dei, uerum futurum est, ut paucissimi fiant, et qui statim numerart possint . Es voirt mit dem zeln- itzund behend zu gehn.  \pend
\phantomsection
\addcontentsline{toc}{subsection}{\textit{Et quemadmodum prius dixit. }}
\subsection*{\textit{Et quemadmodum prius dixit. }}\pstart Scilicet, Esaias in primo Capite.  \pend
\phantomsection
\addcontentsline{toc}{subsection}{\textit{Nisi Dominus Sabaoth reliquis- set nobis semen. }}
\subsection*{\textit{Nisi Dominus Sabaoth reliquis- set nobis semen. }}\pstart Id est, reliquias, nam semen pro filiis accipitur, qui sunt illae reliquiae. At scias cum emphasi dictum, reliquisset. q. d. Ipse uoluit ut aliqui salui fieremus, dians oës, ut Sodoma, periissemus, ubi nemo remansit,  \pend
\section*{COMMENT. POMER.  }\pstart Dominus ergo nos reliquos esse uoluit, non nostro merito, aut carnis sapientia-  \pend
\phantomsection
\addcontentsline{toc}{subsection}{\textit{Quid igitur dicemus ? }}
\subsection*{\textit{Quid igitur dicemus ? }}\pstart Id est, quae est nostra sententia ex his uerbis, quae tractamus, et quae summa? et respondet: haec est summa eorum quae tractamus  \pend
\phantomsection
\addcontentsline{toc}{subsection}{\textit{Quod gentes quae non secta- bantur iusticiam, apprehenderunt iusticiam. }}
\subsection*{\textit{Quod gentes quae non secta- bantur iusticiam, apprehenderunt iusticiam. }}\pstart Id est, illi, qui omnino erant impii et Idololatrae, inimici Dei, qui nihil boni fecerant, neque facere po- tuerant, facti sunt populus Det, et sunt iustificati per solam gratiam, eo quod Deus eis suum Verbum mi- serit, quo credentes sanatt sunt, quod exponit  \pend
\phantomsection
\addcontentsline{toc}{subsection}{\textit{Iusticiam aut eam quae est exfide, }}
\subsection*{\textit{Iusticiam aut eam quae est exfide, }}\pstart Id est, quae tantum coram Deo ualet, quae non est ex hominibus, sed ex mera Dei misericordia, qua tu- stificat quos uult, non eos, qui uidentur meruisse hanc tusticiam.   \pend
\phantomsection
\addcontentsline{toc}{subsection}{\textit{Contra Israel quae sectabatur. }}
\subsection*{\textit{Contra Israel quae sectabatur. }}\pstart Et hoc quoque ad summam praesentis disputationis pertinet, nempe, quod Iudaei qui uidebantur serua- re legem, facere omnia quae Dominus mandarat, nihil effecerint suis operibus coram Deo, imo sint abie- eti, et dicantur iam non esse populus Dei.   \pend
\section*{IN EPIST. AD ROM.  }
\marginpar{[ p.116 ]}
\phantomsection
\addcontentsline{toc}{subsection}{\textit{Propter quid ? }}
\subsection*{\textit{Propter quid ? }}\pstart Scilicet, id factum est, aut cur non peruenerunt ad legem iusticiae?  \pend
\phantomsection
\addcontentsline{toc}{subsection}{\textit{Quia non ex fide. }}
\subsection*{\textit{Quia non ex fide. }}
\phantomsection
\addcontentsline{toc}{subsection}{\textit{Scilicet, sectabantur iusticiam.  }}
\subsection*{\textit{Scilicet, sectabantur iusticiam.  }}
\phantomsection
\addcontentsline{toc}{subsection}{\textit{Sed tanquam ex operibus. }}
\subsection*{\textit{Sed tanquam ex operibus. }}\pstart Id est, putabant iusticiam, operibus comparari, quod tamen est contra totam scripturam: in hoc itaq falsi sunt . IMPEGERVNT ENIM.  Scilicet, ludaei.   \pend
\phantomsection
\addcontentsline{toc}{subsection}{\textit{In lapidem offendiculi. }}
\subsection*{\textit{In lapidem offendiculi. }}\pstart Scilicet, qui semper manet, Christum: quibusdam enim est salus, uita, et resurrectio: multis autem rui- na et offendiculum. Debebant uero Iudaei permoueri per hoc, quod Propheta dixerit duabus domibus Is- raël, et non Gentibus, id quod ipsi libenter tulissent.  Vide quae Simeon dicat in Luca de Christo, et quae hic sequantur.   \pend
\phantomsection
\addcontentsline{toc}{subsection}{\textit{Ecce ponoin Sion. }}
\subsection*{\textit{Ecce ponoin Sion. }}\pstart Id est, in Ecclesia mea, in populo meo, non ie Gentibus, quae adhuc populus non sunt, id quod horrendissimum est, nempe, in Christum offendere, qui proximi uidentur.   \pend
\phantomsection
\addcontentsline{toc}{subsection}{\textit{Et petram offensionis. }}
\subsection*{\textit{Et petram offensionis. }}\pstart Idest, in quem impingant, et ad quem offendantiae.  \pend
\vspace{0.5cm}\noindent
\vspace{0.2cm}\rule{1cm}{0.2pt}\\ 
\hspace{0.2cm}\textit{mg}
\footnotesize Christus aliis salus aliis of- fendiculum.  
\normalsize| 
\section*{COMMENT. POMER.  }
\phantomsection
\addcontentsline{toc}{subsection}{\textit{Et omnis qui creditin eo, non pudefiet, }}
\subsection*{\textit{Et omnis qui creditin eo, non pudefiet, }}\pstart Mirus est hic lapis: necesse est hunc esse Deum, quemadmodum et saepe in Scriptura petra uocatur Deus. Qui ergo crediderit in hunc lapidem, ille non pudefiet, non confundetur, consequetur omnua, quae sunt promissa semini ipsius Abrahae. Rursus uero: qui alio spectarint, illi confundentur, et peribunt.   \pend
\endnumbering\beginnumbering\section{CAP.X.}
\phantomsection
\addcontentsline{toc}{subsection}{\textit{Fratres, propensa quidem uo- luntas cordis mei. }}
\subsection*{\textit{Fratres, propensa quidem uo- luntas cordis mei. }}\pstart Id est, nihil altud magis desydero, ac obsecro in  \pend\pstart \huge\textbf{}\normalsize  \pend
\vspace{0.5cm}\noindent
\vspace{0.2cm}\rule{1cm}{0.2pt}\\ 
\hspace{0.2cm}\textit{mg}
\footnotesize Lapis offen- fionis. 
\normalsize| 
\section*{IN EPIST. AD ROM.  }
\marginpar{[ p.117 ]}\pstart torde meo, id est, coram Deo, quam ut saluentur Iu- daei: et nos hodie multa dicimus contra Clericos, quos uocant, damnantes eorum iusticiae opera, inte- rim tamen quaerimus eorum salutem, et ita erga eos adfecti sumus in corde, ut Paulus erga Iudaeos-  \pend
\phantomsection
\addcontentsline{toc}{subsection}{\textit{Testimonium enim illis perhi- beo, ꝙ studium Dei habent. }}
\subsection*{\textit{Testimonium enim illis perhi- beo, ꝙ studium Dei habent. }}\pstart Mirum est, carni tantum zelum non satis esse, Iudaei tribuentes iusticiam operibus suis, putabant se agere Dei caussam, pugnabant pro eis optima, ut eis uidebatur, conscientia, tamen contra Deum, et non pro Deo agebant. Sic hodie nostri, putant se de- fendere Dei gloriam, cum eos occidunt, qui operi- bus bonis nihil tribuunt, sed omnia gratiae Der: mordicus tenent suos patres, et traditiones, arbi- trantes nihil in his esse constitutum, quod non ex bo- na conscientia, et intentione profectum sit, sed ne- sciunt se misere falli. Non profecto sufficit bona con- scientia, si hoc tam ardens studium in Iudaeis abiici- citur. Impiissima est, et impia, et contra Deum sunt, quaecunque sine Dei Verbo ipsa statuit. Hoc zelo et bona intentione, Christus et Paulus occisi sunt: item alii innumeri, qui Dei tantum iusticiam praedi- cabant. Deus igitur custodiat corda nostra, ne et ho- die, atia speciositate huius occulti mali, seducamur: Videnaus enum qui noui spiritus surgant, iurantes se-  \pend
\vspace{0.5cm}\noindent
\vspace{0.2cm}\rule{1cm}{0.2pt}\\ 
\hspace{0.2cm}\textit{mg}
\footnotesize Conscienti- pona non oia fiunt quae ui- 
\normalsize| 
\section*{COMMENT. POMER.  }\pstart omnia docere bona conscientia, sed ipsi uiderint: Ideo et Paulus hic subiicit, studium habent, sed non secundum scientiam, id est, certissimam intelligentiam, quae ex solo Dei Verbo et spiritu est, non recte intelli gunt, neque norunt, quid uelit Dominus in Scripturis.  \pend
\phantomsection
\addcontentsline{toc}{subsection}{\textit{Nam ignorantes Dei iusticiam. }}
\subsection*{\textit{Nam ignorantes Dei iusticiam. }}\pstart Id est, fidet, qua solus Deus nos iustificat, unde et altas humanas uocamus.   \pend
\phantomsection
\addcontentsline{toc}{subsection}{\textit{Et propriam iusticiam quaerentes: }}
\subsection*{\textit{Et propriam iusticiam quaerentes: }}\pstart Scilicet, operum, uirium humanarum. Potes er- go hunc locum observare diligentius, ut scias auae sit differentia iusticiae Dei, et operum, quae coram Deo tantum ualeat, et quae tantum oculos hominum fal- lat Fiert autem non potest, quemadmodum hic dicit Apostolus, ut simul tuam statuas iusticiam, et Dei: superbus erat Cayn, superbi ludaet: nolebant sub- diti esse iusticiae Dei, id est, nolebant obedire Deo, nempe negabant solam Dei misericordiam, satis esse ad salutem: quod et hodie nostri faciunt, tam durae sunt ceruicis, et suaue iugum Christi, ut maxime com- modum sit, ferre nolunt. Sed Deus suo tempore, et species contemnet et abiiciet, addit ergo  \pend
\phantomsection
\addcontentsline{toc}{subsection}{\textit{NAM PERFECTIO. Id est, finis }}
\subsection*{\textit{NAM PERFECTIO. Id est, finis }}
\phantomsection
\addcontentsline{toc}{subsection}{\textit{Legis Christus. }}
\subsection*{\textit{Legis Christus. }}\pstart Id est, Nihil tibi legis deest, nihil est quod non im- pleueris, si Christum habeas.  \pend
\phantomsection
\addcontentsline{toc}{subsection}{\textit{}}
\subsection*{\textit{}}
\vspace{0.5cm}\noindent
\vspace{0.2cm}\rule{1cm}{0.2pt}\\ 
\hspace{0.2cm}\textit{mg}
\footnotesize Iusticiae Dei et humanae afferentia.  
\normalsize| 
\section*{ne IN EPIST. AD ROM.  }
\marginpar{[ p.]}
\phantomsection
\addcontentsline{toc}{subsection}{\textit{Ad iustificationem omni credenti. }}
\subsection*{\textit{Ad iustificationem omni credenti. }}\pstart Vt scilicet, Omnis qui credit in eum iustificetur: postea explicabit quid uelit, cum dicit: Omni cre- denti, ne Iusticiarii putent id quoque de se dictum: nam saepe ac multum iactant se credere in Deum, cum ta- men ne sciant quidem Deum.  \pend
\phantomsection
\addcontentsline{toc}{subsection}{\textit{Moses enim scribit de iusticia. In Leuitico, et Deutero. }}
\subsection*{\textit{Moses enim scribit de iusticia. In Leuitico, et Deutero. }}
\phantomsection
\addcontentsline{toc}{subsection}{\textit{Quae est ex lege, ꝙ qui fecerit ea homo, uiuet per illa. }}
\subsection*{\textit{Quae est ex lege, ꝙ qui fecerit ea homo, uiuet per illa. }}\pstart Qui illa, quae in Lege sunt expressa fecerit, non punietur, non occidetur, aut lapidabitur: uel, quem admodum iusticia illa est carnalis, sic tantum uitam obtinebit carnalem, nam Lex ubique habet suans poenas annexas, quas nemo euitare potuit, nisi qui faceret in hominumconspeclu, ea quae in ea praecipiebantur: habrbat et promissiones suas, si seruarent, ꝗ bene facturus esset eis benedictione, uberis agri, cella- rii. Emphasis est ergo, non in uerbo uiuet, sed in, fecerit.   \pend
\phantomsection
\addcontentsline{toc}{subsection}{\textit{Caeterum ea quae exfide. }}
\subsection*{\textit{Caeterum ea quae exfide. }}\pstart Ecce hic aliam iusticiam describit, quae est per fi- dem: estque iusticia Euangelii, qua tantum coram Deo, et non coram hominibus, per fidem iustificatus es, ut supra dictum, ex qua iusticia produ iste fructus, quam  \pend
\section*{COMMENT. POMER.  }\pstart sentis te spontaneo, puro, recto, et iusto corde di- ligere proximum: haec non prouenit ex Lege, sed ex Dei Spiritu, qui mouet, et agitat cor tuum, ut be- ne sentiat de proximo. Hic solus facit factores legis, non tantum auditores, quemadmodum Moses: neq unquam potest adimpleri Lex, nisi ille praesto sit, et omnia in te faciat, et operetur. Quare impossibile est, ut iusticia Legis et speciosissima hominum ope- ra, apud Deum ualeant: ideoq haec iusticia nulla pla- ne est, nam non potest fieri. Quomodo enim fieret, cum id in te desyderetur, per quod fieri debebat: hine Hieremias;1. dixit: Scribam non in saxeas tabulas, sed in corda: huc pertinet Petri dictum, Acto. 15.  Quod neque nos, neque patres etc. Deus tuste exigit a nobis hanc iusticiam, sed non praestamus, imo neq possumus, nisi ipse suo digito Legem ipsam in corda nostra scribat. Facile itaque obseruabis, in quo iusticia Legis et Euangelii confistat : haec tantum praeciptt, et nihil implet: illa uero nihil praecipit, et tamen omnia implet, quia Christum habet, qui est Legis impletio, quique ; spiritum suum donat, per quem remissionem peccatorum consequimur : is ani- mat et hilares nos facit, ad quoduis bonum facien- dum, et si quid peccatur, facit is ne imputetur pec- catum, nam perpetuo luctatur cum Carne et Sang uine-Sic itaque uides Christum quoque dicere nobis; hoc fac, et uues, iustus eris: uerum ita dicitur, ut  \pend
\vspace{0.5cm}\noindent
\vspace{0.2cm}\rule{1cm}{0.2pt}\\ 
\hspace{0.2cm}\textit{mg}
\footnotesize Legis et Euan gelii iusticia.  
\normalsize| 
\section*{19 IN EPIST. AD ROM.  }
\marginpar{[ p.]}\pstart non terrearis hac uoce, quemadmodum Mosis uoce territi sunt Iudaei. Non hic respicitur in uires huma- nas, ubi hanc uocem in corde tuo audis, sed in so- lum Deum, quem sinis in te operari, dum tu inte- rim Sabbathum agis. Sed audiamus Paulum-  \pend
\phantomsection
\addcontentsline{toc}{subsection}{\textit{Ea sic loquitur, ne dixeris in cor- de tuo. }}
\subsection*{\textit{Ea sic loquitur, ne dixeris in cor- de tuo. }}\pstart Id est, ne sic senseris. Estque hoc contra Iusticia- rios dictum, qui ore non negant Christum mortuum pro peccatis, attamen in corde sentiunt et dicunt, operibus parandam iusticiam. Vides iterum uerbum legis in ore, et saxeis tabulis tantum esse, Verbum autem Euangelii, in corde quo solo purificantur, et creantur in nobis munda corda. Non uult hic Moses hypocri- sin, quae tantum in externis haeret, sed ipsum cor, Qui ergo dixerit in corde suo  \pend
\phantomsection
\addcontentsline{toc}{subsection}{\textit{Quis ascendet in coelum? }}
\subsection*{\textit{Quis ascendet in coelum? }}\pstart Hic corde negauerit Christum ascendisse, et sede- re ad dextram patris. Includit autem, quis, interro gatiue negationem, id est, nullus ascendet.   \pend
\phantomsection
\addcontentsline{toc}{subsection}{\textit{Aut quis delcendet in abyssum, hoc est Christum etc. }}
\subsection*{\textit{Aut quis delcendet in abyssum, hoc est Christum etc. }}
\phantomsection
\addcontentsline{toc}{subsection}{\textit{}}
\subsection*{\textit{}}\pstart  \pend
\vspace{0.5cm}\noindent
\vspace{0.2cm}\rule{1cm}{0.2pt}\\ 
\hspace{0.2cm}\textit{mg}
\footnotesize iusticiarii.  
\normalsize| 
\section*{COMMENT. POMER.  }\pstart Id est, dicere in corde, nullus est qui descendit in in- fernum, in mortem: quod quid aliud est, quam negare Chri- stum fuisse in morte propter nos? Sic igitur erit subla- tus Christus, ad quem confugiemus tandem? quo freti, in coelum ascendamus, et mortem effugiamus? Non aliud proponemus nobis, quam nostra opera et iusticias externas, nam aliter fieri nequit, quam pri- mum negaueris Christum uictorem mortis, et in glorta patris per mortem restitutum. Vides ergo, quo con silio haec citarit Paulus, ut intelligamus Mosen non tantum de Legis aut Carnis iusticia locutum, sed de alia quoque , id est, ea quae est per fidem, ut obstrue- retur os Iusticiariorum, qui solum sua opera praedi- cant in corde, et Christum ore, qua non potest esse maior blasphemia in Deum. Sunt itaque diligenter uer- ba Pauli obseruanda, ne qua tibi uera sententia elaba- tur' Emphasis est, et tota uis huius loci in dictione, in corde-q'd. Nolo ut addubites de Dei in te mise- ricordia, et beneficentia, ipse est qui te iustificat et faluat: non respice in tuum Libus Arb. aut tua ope- ra, necesse est ut constantissimus sis in hac Petra, ne qua fiat, ut in arenam demissus, demergaris, et tempestatibus totus expugneris. Addit uero Paulus  \pend
\phantomsection
\addcontentsline{toc}{subsection}{\textit{Sed quid dicit? }}
\subsection*{\textit{Sed quid dicit? }}\pstart Scilicet, illa iusticiafidei . PROPE TE EST VERBVM. Scilicet Dei. IN ORE TVO.   \pend
\vspace{0.5cm}\noindent
\vspace{0.2cm}\rule{1cm}{0.2pt}\\ 
\hspace{0.2cm}\textit{mg}
\footnotesize Moses non tan tum de legis iusticia locu- tus.  
\normalsize| 
\section*{IN EPIST. AD ROM.  }
\marginpar{[ p.120 ]}\pstart Scilicet, per confessionem, et in corde per fidem q'd. Quid est quod uelis respicere in tua opera, qui bus non satisfacis legi? huc intende oculos cordis tui, in Verbum, ex quo solo est uita-  \pend
\phantomsection
\addcontentsline{toc}{subsection}{\textit{Hoc est uerbum fidei quod praedi, }}
\subsection*{\textit{Hoc est uerbum fidei quod praedi, }}\pstart Scilicet nos: non potest ergo esse uerbum legis, aut operum, sed fidei, de quo Hieremias dixit: In cor- de scribam eam etc-  \pend
\phantomsection
\addcontentsline{toc}{subsection}{\textit{Nempe si confessus fueris ore tuo.  }}
\subsection*{\textit{Nempe si confessus fueris ore tuo.  }}\pstart Corammundo, et Sathana. DOMINVM IESVM. Scilicet illum esse Dominum totius mun di, item inferorum, peccati, mortis: Iesum esse non nege- uerunt, at Christum esse non crediderunt . Huc pertinet locus ex17. lohan: Vt cognoscant quem misisti Iesum Christum, scilicet esse, id est, eum qui unctus est rex et sacerdos.   \pend
\phantomsection
\addcontentsline{toc}{subsection}{\textit{Corde enim creditur ad iusticiam. *1-4-5ebis Id est, per hoc quod corde credis, iustificaris.  }}
\subsection*{\textit{Corde enim creditur ad iusticiam. *1-4-5ebis Id est, per hoc quod corde credis, iustificaris.  }}
\phantomsection
\addcontentsline{toc}{subsection}{\textit{Ore autem confessio fit ad salutent. }}
\subsection*{\textit{Ore autem confessio fit ad salutent. }}\pstart Vt salutem consequaris, necesse est ut tuam fidem coram toto mundo, et Sathana palam facias. Qui- dam haec duo esse putant, ideoque esse diuidenda, quod quam faciant inepte, uel puer iudicare potest: unum est quod dicitur, non duo. Fidei tribuitur tota iustificatio, sed si eam negabo coram homini- bus, non est uera: ut ergo fides ucra sit, exigis  \pend
\phantomsection
\addcontentsline{toc}{subsection}{\textit{}}
\subsection*{\textit{}}
\vspace{0.5cm}\noindent
\vspace{0.2cm}\rule{1cm}{0.2pt}\\ 
\hspace{0.2cm}\textit{mg}
\footnotesize Fides ore con- fitenda.  
\normalsize| 
\section*{COMMENT. POMER.  }\pstart Deus confessionem eius coram hominibus, ita enim et Christus dixit': Qui me negauerit coram homini- bus, negabo et illum coram patre meo-  \pend
\phantomsection
\addcontentsline{toc}{subsection}{\textit{Dicit scriptura. Omnis qui fidit illi, non pudefiet. }}
\subsection*{\textit{Dicit scriptura. Omnis qui fidit illi, non pudefiet. }}\pstart Id est, siue sit ludaeus, siue Graecus: si crediderit, quodcunque uoluerit, adsequetur. Pudefiunt autem Iu- sticiarii, postulantes multa a Deo, uerum nihil acci- piunt. Non exaudtuntur qui osculantur manus suas.  \pend
\phantomsection
\addcontentsline{toc}{subsection}{\textit{Non enim est distinctio. }}
\subsection*{\textit{Non enim est distinctio. }}\pstart Id est, haec tusticia fidet ex aequo ad omnes homi- nes pertinet, non tantum ad Iudaeos- Vide Cap-;  \pend\pstart Nam idem Dominus omnium. diues in omnes.  \pend\pstart Id est, abundans misericordia et liberalitate su- per omnes. INVOCANTES SE. Idest, credentes: nihil retulerit, siue sint coram hominibus sancti, siue non, hinc addit.   \pend
\phantomsection
\addcontentsline{toc}{subsection}{\textit{Quisquis inuocauerit nomem Do- mini saluus erit. }}
\subsection*{\textit{Quisquis inuocauerit nomem Do- mini saluus erit. }}\pstart Hic nulla est exceptio, uel ludaei, uel Graeci: qui cunq in corde habet Dominum, ille quoque in corde eum inuocabit. Explicat autem id clarius gradatione pulchra, dicens:  \pend
\phantomsection
\addcontentsline{toc}{subsection}{\textit{Quon inuocabunt in quem non credunt. }}
\subsection*{\textit{Quon inuocabunt in quem non credunt. }}
\section*{IN EPIST. AD ROM.  }
\marginpar{[ p.121 ]}\pstart Quasi diceret, impossibile est, ut ad Deum con- fugias, tanquam patrem, nisi credas ipsum esse tuum pa- trem. Ex hoc concludimus, omnes Monachos regni Papistici, nunquam orasse, neque inuocasse, siquidem non credebant in Deum, qui erat inuocandus: credebant ex suis operibus iusticiam se consecuturos, quemadmo dum Iusticiarii Iudaet, quod quid aliud erat, quam abnega- re Deum? Multo rectius senserunt Iudaei de Deo in Euangelio, qui dixerunt: Num homo potest remitte re peccatum? Non audiunt nostri, quod dicitur Rom- 24. Quicquid non est ex fide, peccatum est: nam alioqui abiecissent Cappas et Blattas, comedissentque ́ feriis sextis carnes, aut saltem docuissent, in talibus non esse tusticiam. Cum interrogantur, num haec De- us mandarit? obmutescunt, non audent confiteri suam stultam impietatem: Praeterea, tametsi uideantur cla- mare in tribulatione, tamen non exaudiuntur, quia non fidunt Deo: quemadmodum Esau constitutus in angustia, liberart uoluit, cum Iacob benedictionem praeripuisset, at si quae amiserat, ipsi restituta fuis- sent, nullum curam de Deo habuisset.   \pend
\phantomsection
\addcontentsline{toc}{subsection}{\textit{Quomodo autem credent ei, de U'Th5Aa quo non audierunt ? }}
\subsection*{\textit{Quomodo autem credent ei, de U'Th5Aa quo non audierunt ? }}\pstart Noluerunt audire Iudaei, cum praedicaretur eis gra- tia Dei per Christum, quomodo igitur credent? Con- tra nouos hoc est Prophetas, qui nolunt ministros  \pend
\phantomsection
\addcontentsline{toc}{subsection}{\textit{}}
\subsection*{\textit{}}
\vspace{0.5cm}\noindent
\vspace{0.2cm}\rule{1cm}{0.2pt}\\ 
\hspace{0.2cm}\textit{mg}
\footnotesize Papistae non in vocant Deum 
\normalsize| 
\hspace{0.2cm}\textit{mg}
\footnotesize Prophetae no vitii.  
\normalsize| 
\section*{COMMENT. POMER.  }
\phantomsection
\addcontentsline{toc}{subsection}{\textit{* l.5 getas.  }}
\subsection*{\textit{* l.5 getas.  }}\pstart Verbi audire, sed tantum suos spiritus e coelis: utitur De- us in hoc ministerio, humana uoce, tamen eius ma- nus in hoc non potest dici abbreuiata: mandat ille, ut audiamus, dicunt isti: alia uia est incedendum, ni- hil est quod hunc uel illum audiam: sic Sathanas suo- rum animos dementat, ut non uideant, quod iis me- diis sit utendum, nam persuadet eis, si usus fueris his, iam Verbi maiestas, est per te contempta, alligasti enim te externis. O praeclarum commentum, sic agit ille impostor, ut et res, quae ex Deo nobis propo- nuntur, et ipsum Deum amittamus. Quare diligen- ter obseruandus est Sathanas, ne in hanc calamitatem et nos coniiciat. Subiumgit uero Apostolus plura, qui bus autoritatem faciat externo ministerio Verbi, dicens  \pend
\phantomsection
\addcontentsline{toc}{subsection}{\textit{Quomodo autem praedicabunt nisi missi fuerint ? }}
\subsection*{\textit{Quomodo autem praedicabunt nisi missi fuerint ? }}\pstart Estque contra eos dictum, qui non uocati accedunt ad Verbi ministerium: currunt non missi, hinc ni- hil boni efficiunt, agunt id uentris caussa, non prop- ter Det gloriam. Nos interim audiamus Dei Verbum, siquando is uoluerit, ut Verbum praedicemus, inue- ntet nos etiam in pistrino. Summe cauendum est, ne ambiamus tanta munera nostri caussa, quod si faci- mus, non docemus ex fide: si non ex fide docemus, quomodo ipsam fidem alios docebimus? non possu- mus Deo uentris curam committere, multo magis in  \pend
\vspace{0.5cm}\noindent
\vspace{0.2cm}\rule{1cm}{0.2pt}\\ 
\hspace{0.2cm}\textit{mg}
\footnotesize Vocatio expe ctanda- 
\normalsize| 
\section*{12 IN EPIST. AD ROM.  }
\marginpar{[ p.]}\pstart re tam magna erimus impii. Sed diceres: Charitas urget, ut praedicem, Respondeo: uide ne ficta sit, si- tiera est, non potest esse malum, quod eius impulsu facis, sed etiam atque etiam uide, quod charitas non facit contra Dei ordinationem: quomodo praedica- bunt, nisi mittantur? Et si tali occasione uenias in ec- clesiam aliquam, quod tuum est facito, si te audie- rint, potesque ́, si opus fuerit, uictum et amictum ab eis repetere, quod si nolint audire, neque uictum neq amictum acceperis ab eis, sed excutieris puluerem pedum, alio ibis: itaque tantum missi a Deo, uere praedicant-  \pend
\phantomsection
\addcontentsline{toc}{subsection}{\textit{Sicut scriptum est, Quam ama- biles pedes. }}
\subsection*{\textit{Sicut scriptum est, Quam ama- biles pedes. }}\pstart Scilicet, sunt eorum, qui pacem, scilicet, conscientia rum annunciant. Vidit tum Esaias in spiritu Euange- lii praedicatores, currentes per totum orbem, cum tamen eo tempore, tantum in iudaea aliquot ueri es- sent praedicatores, gauisus est ergo quod uiderit futu- rum, quod toti mundo esset patefacienda Dei gratia- Super montes, dicit, ut intelligas ipsum praenouisse, Euangelium praedicandum extra urbes, et in montibus, et in desertis, id est, ubique , etiam in Gentibus. Annunciat autem pacem is tantum, cut mandauit Dominus, ut amabiliter et cum silentio ueniat, non horrendo stre- pitu pedum legis. Ecce hic uide, quis audebit se di- cere nuncium Domint, cui Verbum a Deo non  \pend
\vspace{0.5cm}\noindent
\vspace{0.2cm}\rule{1cm}{0.2pt}\\ 
\hspace{0.2cm}\textit{mg}
\footnotesize Vocati tantum pacem praedi cant .  
\normalsize| 
\section*{COMMENT. POMER.  }\pstart sit datum et commissum? Ad idem pertinet quod ad- dit: Annunciantium bona, non damnationem ex le- gis sententia, sed consolationem Euangelii, quae est optima-  \pend
\phantomsection
\addcontentsline{toc}{subsection}{\textit{At non omnes obedierunt Euan gelio. }}
\subsection*{\textit{At non omnes obedierunt Euan gelio. }}\pstart Occupatio est: Sicut non omnes sunt reiecti , ita non omnes crediderunt, quamuis etiam per totam Iu- daeam tum erat Euangelium praedicatum, neque hodie omnes credunt qui audiunt.  \pend
\phantomsection
\addcontentsline{toc}{subsection}{\textit{Elaias enim dicit, Domine quis credit auditui nostro? }}
\subsection*{\textit{Elaias enim dicit, Domine quis credit auditui nostro? }}\pstart Quasi diceret, quis est, uel quotusquisque est, qui credat his uerbis, quae ex nobis audiunt, nam ea gra- tia quam praedicamus, tota est obscurata formis cru- cis, non aliud hic cernitur, quam quod maxime odiat natura. Habet autem illud Caput Prophetae Christi passionem, ut uideas non aliud uelle Esaiam signifi- care, quam multos esse, qui nolint audire Euange- lium, quia sit de Christo despecto et infami in boc mundo.   \pend
\phantomsection
\addcontentsline{toc}{subsection}{\textit{Ergo fides ex auditu est. }}
\subsection*{\textit{Ergo fides ex auditu est. }}\pstart Id est, ex praedicatione quae auditur, est autem Hebraismus. Ex hoc loco clarum est, auditum et prae- dicatores, requiri ad Euangelium, contra nouos hu- ius temporis Prophetas, qui aliud non possunt quan contemnere, quaecunque talia sunt.   \pend
\vspace{0.5cm}\noindent
\vspace{0.2cm}\rule{1cm}{0.2pt}\\ 
\hspace{0.2cm}\textit{mg}
\footnotesize Caussa, cur non credatur Euangelio.  
\normalsize| 
\section*{IN EPIST. AD ROM.  128 }\pstart Id est, ipsa praedicatio est per Verbum Det, nihil enim debet praedicari, praeter Verbum Det. Vbi hie sunt Regulae Monachorum, et Canones Papistici?  \pend
\phantomsection
\addcontentsline{toc}{subsection}{\textit{Auditus autem. }}
\subsection*{\textit{Auditus autem. }}
\phantomsection
\addcontentsline{toc}{subsection}{\textit{Sed dico, an non audierunt ? }}
\subsection*{\textit{Sed dico, an non audierunt ? }}\pstart Est Occupatio: Nemo inquit dubitabit de Euan- gelii praedicatione, siquidem in omnem terram exiuit sonus, uel uox praedicantium Euangelium, et pa- cem ex Deo: uerum erunt multi, qui tamen Euan- gelio non credent. Vides et hodie, quod experientia docet, inuulgato Verbo non haberi fidem, tamen testimonium habet, etiam ab inimicis, quod sit Ver- bum Det? quare nulla reliqua eis erit excusatio, dan- nabuntur, quotquot non crediderint.   \pend
\phantomsection
\addcontentsline{toc}{subsection}{\textit{Sed dico, nunquid non cog- nouit lsrael. }}
\subsection*{\textit{Sed dico, nunquid non cog- nouit lsrael. }}\pstart Est ubique auditum, quomodo ergo Ifraël potuü ignorare? Sed uideamus, quid de hac re Scriptura le quatur.   \pend
\phantomsection
\addcontentsline{toc}{subsection}{\textit{Primus Moses dicit. }}
\subsection*{\textit{Primus Moses dicit. }}\pstart Id est, Moses, qui primus est inter eos, qui spiritu docente scripserunt, ille quoque significauit tudaeos noncredituros Verbo, ergo abiiciendos esse.  \pend\pstart Ego ad emulationem prouoca. bouos.  \pend
\section*{COMMENT. POMER.  }\pstart Sunt ex.32. Capite Deuter. Irritabo, inquit, animos uestros, ut inuideatis Gentibus gratiam Dei: uel ego efficiam, ut Gentes, quae hodie despectae sunt, fanctae sint, et hoc ubi uiderint ludaet, prae inuidia et cordis amaritudine, ipsi gratiam quoq conten- nent, et Christum abiicient: id quod et hodie ea gens fa- cit, quod pulchre in parabola Lucae expressum est de filio prodigo: frater enim potius uoluit carere pa- ternis bonis, quam istam gratiam uidere. Idem faci- unt hodie lusticiarii, qui quoque nobis inuident Dei gratiam, quam tamen ipsi nolunt complecti, imo per seuerant in sua caecitate. Stultam uocat gentem, quia Idola uenerabatur, et Deum ignorabat, que est uera stultitia.   \pend
\phantomsection
\addcontentsline{toc}{subsection}{\textit{Esaias autem post hunc audet. }}
\subsection*{\textit{Esaias autem post hunc audet. }}\pstart Id 'est, post Mosen, et Esaids dexit de Iadrorum abiectione et caecitate.  \pend
\phantomsection
\addcontentsline{toc}{subsection}{\textit{Inuentus fui iis qui me non quae rebant }}
\subsection*{\textit{Inuentus fui iis qui me non quae rebant }}\pstart Verba sunt Dei in Propheta. Est uero mira phra- sis, qua significat absque merito et studio nostro no- bis datam gratiam, quia non quaesitam inuenimus.  \pend
\phantomsection
\addcontentsline{toc}{subsection}{\textit{Conspicuus factus sum fjs etc. d est, illis adparui, et me agnouerunt per fidem, }}
\subsection*{\textit{Conspicuus factus sum fjs etc. d est, illis adparui, et me agnouerunt per fidem, }}\pstart Qui de me non interrogabant.  \pend
\section*{IN EPIST. AD ROM.  }
\marginpar{[ p.124 ]}\pstart Id est, tantum aberat, ut me quaererent, ut ne- fama quidem illis cognitus essem: ergo hic locus uel solus indicat nos saluatos, praeter omnem meritum-  \pend
\phantomsection
\addcontentsline{toc}{subsection}{\textit{Aduersus Israel autem dicit. }}
\subsection*{\textit{Aduersus Israel autem dicit. }}\pstart Id est, contra eos, qui uidebantur esse Dei popu- lus.   \pend
\phantomsection
\addcontentsline{toc}{subsection}{\textit{Toto die expandi. }}
\subsection*{\textit{Toto die expandi. }}\pstart Hunc locum multi de Christo cruci suffixum ex- ponunt, sed inepte: uox Det est in Propheta, di- centis: Ecce, ego tanquam pater filiis, porrexi popu- lo meo manus, ille uero auxilium meum, et mea be- neficia non agnouit. Commendat uero in primis hic locus nobis gratiam Dei: expandit ad nos manus, sponte offert beneficia, ipse solus dat nobis, nos illi uihil dare possumus.   \pend
\phantomsection
\addcontentsline{toc}{subsection}{\textit{Et contradicentem. }}
\subsection*{\textit{Et contradicentem. }}\pstart Nam quemadmodum hic populus, corde non credit, sic neque potest confiteri Deum, nam par ra- tio est impietatis, quae fidei: ubi fides est, ibi necesse confessionem quoque oris esse: ubi impietas in corde, ibi nihil aliud, quam blasphemiae nominis Dei. Sa- than enim agit corda, ut non praedicent gra- tiam, quae alians suis uiribus nituntur, et impossibile est arborem ma- lam, bonos ferre fru-  \pend
\vspace{0.5cm}\noindent
\vspace{0.2cm}\rule{1cm}{0.2pt}\\ 
\hspace{0.2cm}\textit{mg}
\footnotesize Gratiae com- nendatio.  
\normalsize| 
\section*{COMMENT. POMER.  }\pstart CCVPATIO est, diceret quis: Deus ele-  git populum hunc ludaeorum, et repulit eum, ergo electio et Dei promissio irrita facta est- Re- spondet Paulus: Non potest fieri, ut sit irrita-  \pend
\endnumbering\beginnumbering\section{CAP.XI.}
\phantomsection
\addcontentsline{toc}{subsection}{\textit{Nam et ego lsraelitasum. }}
\subsection*{\textit{Nam et ego lsraelitasum. }}\pstart Non omnes Iudaei sunt abiecti, nam et ego Iudaeus sum, attamen non abiectus. Magna haec Pault est fi- des de electione sua, tametsi unus ipse esset, audet constanter adfirmare, non omnes Iudaeos retectos. Sunt autem haec plena consolationis: duriora erant, quae dicebat in superiore Cap-  \pend\pstart Scilicet, etiam secundum carnem . TRIBVS BENIAMIN. Vt est in Philippen-  \pend
\phantomsection
\addcontentsline{toc}{subsection}{\textit{Exsemine. }}
\subsection*{\textit{Exsemine. }}
\phantomsection
\addcontentsline{toc}{subsection}{\textit{Non repulit Deus populum su- um, quem ante agnouerat. }}
\subsection*{\textit{Non repulit Deus populum su- um, quem ante agnouerat. }}\pstart Id est, Deus non abiecit eum populum, quem iam olim elegerat, uerum illum abiecit, qui Dei populi nomem habebat, uidebaturque esse semen Abrahae, ta- men non erat . Mirabiliter enim splendebat in oculis hominum, iusticia, vita sanctitas huius populi, at coram Deo nihil erat aliud, quam stercus : suam, non Dei quaerebat gloriam, ergo merito sunt abiecti. Ex hoc consequitur, non omnes item Christianos esse,  \pend
\vspace{0.5cm}\noindent
\vspace{0.2cm}\rule{1cm}{0.2pt}\\ 
\hspace{0.2cm}\textit{mg}
\footnotesize Non omnes Christiani, qui uidentur esse.  
\normalsize| 
\section*{IN EPIST. AD ROM.  125 }
\marginpar{[ p.]}\pstart qui Christiani uidentur aut uocantur, sed quos ille agnoscit esse Christianos. Non habet ecclesia Chri- sti externam laruam, qua agnoscas eam, quam tamen Papistae miris dolis confinxerunt. Solus uero Deus no- uit et qui sui sunt, multis enim haud dubie in die in dicii dicturus est: Non noui uos.  \pend
\phantomsection
\addcontentsline{toc}{subsection}{\textit{An nescitis de Elia. }}
\subsection*{\textit{An nescitis de Elia. }}\pstart Alterum est exemplum, quo confirmat non om- nes Iudaeos abiectos. Hunc locum, nisi teneas histo- riam, non intelliges: uidebantur sanctissimo utro Eliae declinasse ad impietatem omnes Iudaei, neminemque ; adeo superesse, qui de Deo recte sentiret, dicebat enim: O Domine, ecce ommes adorant Idola, abii- ciunt legem tuam, ego uero solus relictus sum ex omni- bus, et ex tanto numero, qui non declinarim ad Idola, et ad cultum Baal. Sed uide, quantum errarit hoc iudicio suo Propheta, nam dicitur ei, diuino respon- so: Ego Deus noul, non tu nosti, qui met sunt, erras, te solum superesse putas, qui me ex corde suspicias, falleris: habeo adhuc septem millia uirorum electorum, qui me uere colunt in fide, hos tu ignoras. Quare non est ut cogites, te unico occiso, neminem reliquum fore, qui pius sit. Ecce, tam sanctus Propheta, ex septem millibus uirorum, non unum nouit, non igi- tur statim est iudicandum de iis aut aliis: iudictum et sententia Deo relinquantur, qui uidet in ea tantum,  \pend
\section*{COMMENT. POMER.  }\pstart quae sunt in abscondito: sunt eius tudicia inscrutabilia, quos uult, seruat mirabiliter, ignorantibus etiam ip- sis sanctis. Stulte itaque faciunt hodie, qui temere pro nunciant de Euangello, dicentes: neminem uideo, qui melior fiat, quare non possum persuaderi, esse verum Evangelium . Non videbis o miser, quae Do- minus operetur in suis, putas ne tibi hoc mandatum a Deo, ut tantum numeres eius electos? Pudeat te tuae stulticiae et ignauiae.  \pend
\phantomsection
\addcontentsline{toc}{subsection}{\textit{bic igitur et in hoctempore. }}
\subsection*{\textit{bic igitur et in hoctempore. }}\pstart Accommodat historiam: et hodie, inquit, eadem manet ratio electionis, non omnes abiiciuntur, si - quidem et Petrus duabus contionibus, fere octo mil- lia hominum Iudaeorum ad Christum conuertit. Sal- ui fiunt omnes, quicunque eliguntur a Deo, quod fa- cit tantum per gratiam: si per gratiam, ubi sunt me- rita et opera nostra quae iactamus? Ideo addit  \pend
\phantomsection
\addcontentsline{toc}{subsection}{\textit{Quod si per gratiam, non iam ex opibus }}
\subsection*{\textit{Quod si per gratiam, non iam ex opibus }}\pstart Si reliquiae populi Dei saluantur tantum per gra- tiam, nihil est que gloriemur de nostra iusticia et san- ctttate: uerum tam sumus iniqui et impii, ut maxi- me hanc argumentationem, etiam rationis bumanae iudicio admittamus, tamen tantum honoris fidet ha- bere non possumus, ut palam et constanter dicamus: ex sola Dei gratia est salus, non ex meritis nostris. Sratiam quidem multi praedicant, quod ad nomen et  \pend
\vspace{0.5cm}\noindent
\vspace{0.2cm}\rule{1cm}{0.2pt}\\ 
\hspace{0.2cm}\textit{mg}
\footnotesize Iudicium illo- rum, qui di- cunt, nullos meliores rec di Euangelio. 
\normalsize| 
\section*{IN EPIST. AD ROM.  }
\marginpar{[ p.125 ]}\pstart tinet, tamen in corde gratiam non credunt, sed tan- tum sua opera meritoria, nihil uero conuenit gratiae cum meritis nostris: ubi gratia est, iam meritum lo- cum non potest habere. Aperta sunt haec, et lu- ce clariora.   \pend
\phantomsection
\addcontentsline{toc}{subsection}{\textit{bin ex operibus etc. }}
\subsection*{\textit{bin ex operibus etc. }}\pstart Vtcunque hoc argumentum uerteris, constabit: uix est clarior locus contra eos, qui nituntur iudicio- rationis, et nihil admittendum censent, quam quod ipsi possunt animo complecti. Conferant hic operis et gratiae rationem, et deprehendent, nisi fuerint trunci et lapides, gratia constare salutem, non operibus-  \pend
\phantomsection
\addcontentsline{toc}{subsection}{\textit{Quid igitur. }}
\subsection*{\textit{Quid igitur. }}\pstart Id est, quid plura dicam, haec est summa totius rei-  \pend
\phantomsection
\addcontentsline{toc}{subsection}{\textit{Quod quaerit lsrael, hoc non est adsecutus. }}
\subsection*{\textit{Quod quaerit lsrael, hoc non est adsecutus. }}\pstart Quaerit iusticiis carnis salutem, at non inuenit, nemo autem eam quaerere potest, nisi spiritus ipsum praeuenerit et ad quaerendum impulerit, quod uero Matth. 7. dicitur: Quaerite, et inuenietis, de Chri- stianis, id est, iis qui spiritu Det ducuntur, dicitur, ill inueniunt, alii iusticiarii non inuentunt.   \pend
\phantomsection
\addcontentsline{toc}{subsection}{\textit{Sed electio conlecuta est. }}
\subsection*{\textit{Sed electio conlecuta est. }}\pstart Id est, illi, qui electi sunt a Deo, saluantur, quod plene et eleganter dixit: sic et haereditas saepe pro ipsis haeredibus legitur.   \pend
\section*{COMMENT. POMER.  }
\phantomsection
\addcontentsline{toc}{subsection}{\textit{Reliqui uero excaecati sunt. }}
\subsection*{\textit{Reliqui uero excaecati sunt. }}\pstart Id est, quotquot sunt extra electionem, ii quan- tumuis uideantur esse iusti, caeci sunt, neque possunt oculos in Dei tusticiam, id est, gratiam intendere. Sunt autem iiomnes extra electionem, quicunque fidunt ope- ribus, quales sunt omnes, in regno Papae uiuentes, qui sunt inimici Dei, nam ducuntur carnali quadam sapientia, quae est stultitia coram Deo, qua et excae cantur. Hinc addit  \pend
\phantomsection
\addcontentsline{toc}{subsection}{\textit{Quemadmodum scriptum est, dedit eis spiritum compunctionis. }}
\subsection*{\textit{Quemadmodum scriptum est, dedit eis spiritum compunctionis. }}\pstart Vel exacerbationis, aut amarulentum spiritum, quo compungantur, seque ipsos mordeant, ira et in- uidia aestuent, quem spiritum uidemus habere omnes tusticiarios: hic spiritus alius non est, quam blasphemiae. Cum enim uident gloriam Det per Euangelii praedicationem manifestari, statim tanquam amentes prosiliunt, uo- ciferantes: Bona opera, bona opera, non Euange- lium me saluabit, quod quid aliud est, quam pessime- de Deo sentire, illiusque gratiam pedibus conculcare- Id faciunt hodie, quos Clericos uocamus, aut Sacer- dotes, et Venerabiles Altaristas, non reuerentur no- men magni Dei, statuunt sua, et defendunt tantum conuitiis et consuetudine, non Scripturis, siquidem eas nunquam attigerunt, cur autem attingerent, qui non iutelligant titulum unius libri Sacrorum Bi-  \pend
\vspace{0.5cm}\noindent
\vspace{0.2cm}\rule{1cm}{0.2pt}\\ 
\hspace{0.2cm}\textit{mg}
\footnotesize Spiritus blas- phemie 
\normalsize| 
\section*{IN EPIST. AD ROM.  }
\marginpar{[ p.117 ]}\pstart bliorum. Habuerunt olim Iudaei plus occasionis dan- nandt Euangelii, quam illi nostri, illorum enimce- remoniae ex Dei Verbo erant institutae, horum ex stul- tissimis et ineptissimis Monachis sunt profectae. Sed uide textum.   \pend
\phantomsection
\addcontentsline{toc}{subsection}{\textit{Oculos, ut non uideant. }}
\subsection*{\textit{Oculos, ut non uideant. }}\pstart Id est, excaecauit eos Deus, posteaq contempse. runt: Dei potentia factum est, ut excaecarentur, estque ́ caecitas eis poena contemptus Verbi: intellige autem oculos cordis. Item aures, quae quoque sant obturatae, ut nequeant audire quicquam. Et addidit non abs re  \pend
\phantomsection
\addcontentsline{toc}{subsection}{\textit{Vlque ad hodiernum diem. }}
\subsection*{\textit{Vlque ad hodiernum diem. }}\pstart Quasi diceret, et hodie illud uidetis fieri, adhuc perseuerant in infidelitate cordis, obturati contra Verbum Dei. Non longius petenda sunt exempla, sunt in omnium oculis: scripsit enim Paulus haec, quae se omnium temporum, in quibus praedicaretur Euan- gelium, rationibus accommodarent.  \pend
\phantomsection
\addcontentsline{toc}{subsection}{\textit{Et Dauid dicit, Vertatur mem- sa illorum in laqueum. }}
\subsection*{\textit{Et Dauid dicit, Vertatur mem- sa illorum in laqueum. }}\pstart Sed quae mensa in laqueum uerti potest? Ita intelli- ge: Scriptura Verbum Dei, nunc mensam aut panem, alians item, cibum, triticum, et aquam sapientiae sa- lutaris, uinumque optimum uocat. Dicit ergo Dauid: Scriptura, uel Verbum salutis, quod hactenus habue- runt Iudaei, uertetur eis in perditionem . Erant ipst  \pend
\vspace{0.5cm}\noindent
\vspace{0.2cm}\rule{1cm}{0.2pt}\\ 
\hspace{0.2cm}\textit{mg}
\footnotesize Caecitas con- tempti Verbi poena.  
\normalsize| 
\section*{COMMENT. POMER.  }\pstart Dei populus, solique ; habebant Verbum, sed non aliter quemadmodum asinus portans mysteria: quid eni attinent mysteria ad asinum, nihil horum intelligit, tamen por- tans ea ostentat? Ita Iusticiarii insigniter quoque glo- riantur se habere Scripturas, sed scire nolunt se eas ig- norare: non ergo bic de mensa Philosophorum, sed lu- daeorum dicit, nam ipsis solis erant commissa Dei ora- cula. Non caret Emphasi, ꝗ dicit, illorum. q. d. il- li ipsi, qui hactenus cibum habuerunt delicatum, nem pe, Verbum Dei, hi iam ueneno uescentur: uel, ipsis panis non in nutrimentum, sed perniciem mutabi- tur et interitum. Sic idem Verbum uidemus quosdam illuminare, alios autem excaecare, alio igitur docto- re opus est, quam homine, is enim tantum monet, hortatur, Deus uero aut illuminat, aut excaecat- Hinc in Esaia dicitur: Excaeca cor populi huius, scili- cet, tua praedicatione, fac ut Verbum meum sit eis la- queus, captio, offendiculum, et retaliatio, id est, nihil mali non accipiant, et extrahant e Verbo, ubi alii ex hoc spem in Deum concipiunt, ibi isti nihil ali- ud, quam desperationem imbibant, quod quid alind esse potest, quam aeternus cruciatus, horror et anxie- tas conscientiae? sed haec horrenda sunt, utinam uide- remus: non abs re his uocabulis, et Propheta et Da- uid usus, quo magis nos terreremur, et uide remus, plane nihil esse in nobis, quo Deum nobis concilie- mus. Laqueum dixit quiddam simplicius, sed ubi ad  \pend
\vspace{0.5cm}\noindent
\vspace{0.2cm}\rule{1cm}{0.2pt}\\ 
\hspace{0.2cm}\textit{mg}
\footnotesize Verbo dliii- luminantur, alii excaecan- tur.  
\normalsize| 
\section*{IN EPIST. AD ROM.  }
\marginpar{[ p.128 ]}\pstart dierint insuetum, et uident omnia consonare quae Iegunt praedicationt, mirantur uehementer, sed quia statim inueniunt quosdam locos, impingunt et ex- caecantur: qualis locus est in Primo ad Romanos: Non auditores, sed factores legis saluabuntur. In fine omnium illorum, retaliationem adposuit, sig- nificaturus caecitatem esse praemum illorum scande- lorum et poenan-  \pend
\phantomsection
\addcontentsline{toc}{subsection}{\textit{Obtenebrentur? }}
\subsection*{\textit{Obtenebrentur? }}\pstart Id est, excaecentur ocult eorum.  \pend
\phantomsection
\addcontentsline{toc}{subsection}{\textit{Et tergum illorum semper in- curua. }}
\subsection*{\textit{Et tergum illorum semper in- curua. }}\pstart Deprime eos oneribus conscientiae peccatorum, ne sursum suspicere in coelum possint, habeant ocu- los semper defixos in suam iusticiam, non in Dei gra- tiam. Ita profecto nunquam in libertatem, conscien- tiae iusticiariorum, et gaudium perduci possunt, las- santur et fatigantur onere iusticiae carnis, sed sine spe salutis aut consolationis.   \pend
\vspace{0.5cm}\noindent
\vspace{0.2cm}\rule{1cm}{0.2pt}\\ 
\hspace{0.2cm}\textit{mg}
\footnotesize Pax conscien- tiae non est Iu- sticiariis.  
\normalsize| 
\section*{COMMENT. POMER.  }
\phantomsection
\addcontentsline{toc}{subsection}{\textit{Dico igitur, Num ideo impege runt, ut conciderent ? }}
\subsection*{\textit{Dico igitur, Num ideo impege runt, ut conciderent ? }}\pstart Id est, ut caderent omnes israelitae secundum car- nem. Est uero sententia huius loci: Deus sic recepit Gentes, abiectis Iudaeis, non ut abiecti maneant, sed ut emulatione Gentium prouocentur, ut et ipsi gra- tiam sitiant et resipiscant. Vult itaque occupare omnia Paulus quae hic possent obiici, ut scilicet hoc Dei iudi- cium mitiget, ne desperemus de Dei bonitate, qui so- lus nouit, qua ratione quemuis saluari uelit : eius con- silia sunt abyssus multa, non exposita hominum inda- gini. Sed ut clarius cernas, quid hic locus uelit, assu- mito parabolam Lucae, de filio prodigo et fratre eius. Hic peccator est, ille Iusticiarius : hic relicto patre, omnia consumit et male perdit : ille rursus manet in agro paterno, et laboribus perpetuis distringitur, uideturque omnibus patri esse utilis, et rem facere, ob idque illi acceptum se fore magis sperat, et tandem consecuturum omnem haereditatem: nihil minus puta- uit futurum, quam rediturum alium fratrem in patris gratiam, tamen ubi eum rediisse intellexit, ac susce- ptum a patre, non uult ipse ingredi in domum pater- nam, etiam ipso rogante patre: abiicitur ergo et ex- cluditur, et qui uagus, et omnium decoctor fuerat, in locum eius subrogatur. Queri itaque posset, num ideo dle prodigus susceptus est, ut alter pertret? minime,  \pend
\vspace{0.5cm}\noindent
\vspace{0.2cm}\rule{1cm}{0.2pt}\\ 
\hspace{0.2cm}\textit{mg}
\footnotesize tudei, cur ab iecti- 
\normalsize| 
\hspace{0.2cm}\textit{mg}
\footnotesize Parabola de filio prodigo. 
\normalsize| 
\section*{IN EPIST. AD ROM.  E. }
\marginpar{[ p.]}\pstart sed ideo abiectus est, ut prouocaretur per prodigum fratrem, ad agnitionem liberalitatis, et charitatis paternae, et ne crederet se ex hoc haeredem fore bono- rum, quia diligenter in agro laborarat, quare etiam amice ipsum adloquitur pater: Fili, omnia mea tua sunt, ꝗ d. Non ob id quia laborasti in agro, sed quia te diligo, et es mihi natus filius: quid fecit bont ali- us frater tuus? certe nihil, tamen possidebit mea bo- na, saltem ob id, quia diligo eum: disce ergo ex hac mea beneficentia erga fratrem tuum, me non respice- re in opus aut laborem externum, sed quod facio, gra- tis facio, quia pater. Hac ratione uoluit Deus ludaeos Iusticiarios, per assumptionem gentium, ad agnitio- nem suae gratiae prouocare, quo discerent, Deum non respicere opera legis dum nos saluat, sed solum om- nia ipsum agere ex misericordia: sic debebant et no- stri hodie pudefieri, qui stupidi, tribuunt omnem iusti- ciam suis meritis, non diuinae gratiae: ergo cadunt quidam ex electis, non ut nunquam surgant, sed ul- discant, quomodo surgendum sit: non est ira haec Dei sine misericordia: ideo trascitur, ut tandem mise- reatur nostri. Abiecti sunt in ira Iudaei, sed ut olim per misericordiam restituerentur, quod factum est multis, qui cum aliquandiu uixissent in operibus le- gis, audientes Euangelium, dubitare coeperunt de illis, nam ne uerbum quidem inuenerunt in Scripturis, quod opponerent gratiae in defensionem operum. Sic  \pend
\vspace{0.5cm}\noindent
\vspace{0.2cm}\rule{1cm}{0.2pt}\\ 
\hspace{0.2cm}\textit{mg}
\footnotesize Electi quidam cur cadunt.  
\normalsize| 
\section*{COMMENT. POMER.  }\pstart ergo praeuenerunt in gratiae cognitionem, quam an- tea ignorabant. Idem hodie in quibusdam est cerne- re, qui relicta fiducia operum, se gratiae Dei tantum committunt: itaq mirabiliter Deus suas diuitias in no- bis manifestat, facit ex stercore aurum: per istas im- mundicias uitam aeternam nobis donat, ut negare non possimus, ipsum omnia nobiscum operari, secundum gratiam et bonitatem suam. Sed uideamus textum  \pend
\phantomsection
\addcontentsline{toc}{subsection}{\textit{Sed per lapsum. }}
\subsection*{\textit{Sed per lapsum. }}\pstart Id est, ꝗ ipsi Iudaei ceciderunt, Christumque non susceperunt, quod propter duriciam cordis factum est.  \pend
\phantomsection
\addcontentsline{toc}{subsection}{\textit{Salus contigit Gentibus. }}
\subsection*{\textit{Salus contigit Gentibus. }}\pstart Christus, salus Iudaeorum a patre ipsis missus, ut Paulus Acto.13. dicit: Vobis missum erat Ver- bum salutis, non est susceptus: ne irrita esset promssio, translata est haec salus ad Gentes, secundum hoc Esa- iae: Dedi te in lucem Gentium.  \pend
\phantomsection
\addcontentsline{toc}{subsection}{\textit{In hoc ut eos etc. }}
\subsection*{\textit{In hoc ut eos etc. }}\pstart Id est, ut extimularentur illa Det bonitate in Gen- tes Iudaei ad resipiscentiam.   \pend
\phantomsection
\addcontentsline{toc}{subsection}{\textit{Quod si lapsus illorum. }}
\subsection*{\textit{Quod si lapsus illorum. }}\pstart Multa hic argumenta congerit propter nos Gen- tes, ne condemnemus Iudaeos, ex quibus salus uenit, nam ex iis natus est Christus. Impia hactenus aestima tione factum est, ut Iudaeos tanq bestias haberemus,  \pend
\section*{IN EPIST . AD ROM.  }
\marginpar{[ p.n0 ]}\pstart ton uidebant quo interim honore digni erant, propter promissiones ipsis factas et Christum, quare nonest dicendum, ludaeus peior est cane, nam Christi est oc- cisor. Lauda potius arcanum Dei opus, quod ille te scire noluit, Iudaei seruierunt, tanq instrumenta Dei, nostrae redemptiont, tu quoque iam et ipsis seruito, ut aliquando tua pietate et fide prouocati, resipiscant.  Diuitiae Iudaeorum sunt, salus: nomine mundi Gen- tes intelliguntur, uel omnes homines.  \pend
\phantomsection
\addcontentsline{toc}{subsection}{\textit{Et diminutio illorum. }}
\subsection*{\textit{Et diminutio illorum. }}\pstart Si cum pauci Iudaeorum salui factisint, et illorum lapsus (qui malus est) nobis occasio fuit salutis,  \pend
\phantomsection
\addcontentsline{toc}{subsection}{\textit{Quanto magis et plenitudo illorum. }}
\subsection*{\textit{Quanto magis et plenitudo illorum. }}\pstart Id est, multo magis profutura est nobis ipsorum resi- piseentia: uel, si tam efficax fuit malum, multo effi- catius erit bonum. Estque hoc argumentum mire ar- gutum, quo persuadeat nobis Iudaeos redituros olim ad Messiam, quem nolebant recipere.  \pend
\phantomsection
\addcontentsline{toc}{subsection}{\textit{Vobis enim gentibus dico. }}
\subsection*{\textit{Vobis enim gentibus dico. }}\pstart Conuertit sermonem ad Gentes, ut familiarius ipsis loquatur: ne gloriamini, inquit, aduersus ludaeos, non auersemini eos, nam possunt conuerti: habent prae- terea magnam gloriam secundum carnem, quare et:am a nobis uenerandi sunt, tametsi ea gloria coram Deo nul- lasit, et sola ea quaee estinspiritu: terrcamur potius  \pend
\vspace{0.5cm}\noindent
\vspace{0.2cm}\rule{1cm}{0.2pt}\\ 
\hspace{0.2cm}\textit{mg}
\footnotesize Iudaei non sper- nendi.  
\normalsize| 
\section*{COMMENT. POMER.  }\pstart hoc exemplo, ꝗ ii ex quibus Christus natus est, prae- uaricati sint, et exciderint: nihil nostrae securitati permittamus, nam idem nobis, qui uidemur optimi Christiani, quod Iudaeis, accidere potest : hinc et em- phasi non caret ꝙ Gentibus dicit. q. d. Vos non estis secundum carnem ex Iudaeis, tamen recepti in gratiam, uidete ne superbiatis: beneficii diuint est, ꝗ et ego uobis Euangelium praedicem, inuulgemq, ut hoc Iu- daei uidentes, inuideant id uobis, et ex hoc tandem ad resipiscendum moueantur, ad Christumque ; redeant.  \pend
\phantomsection
\addcontentsline{toc}{subsection}{\textit{Ministerium meum illustro. }}
\subsection*{\textit{Ministerium meum illustro. }}\pstart Id est, cum magna gloria, scilicet nominis Chri- sti praedico Euangelium. Hoc est meum ministerium, non aliud mihi a Deo commissum, ideoque meum-  \pend
\phantomsection
\addcontentsline{toc}{subsection}{\textit{Si quo modo ad emulationem pro }}
\subsection*{\textit{Si quo modo ad emulationem pro }}
\phantomsection
\addcontentsline{toc}{subsection}{\textit{uocem carnem meam. }}
\subsection*{\textit{uocem carnem meam. }}\pstart Id est, Iudaeos, quia ipse sum Iudaeus natus. Vides Paulum integro et candido esse animo, ettam in ad- uersarios nominis Christi: id quod magnum est, nam natura nemo potest diligere inimicum.   \pend
\phantomsection
\addcontentsline{toc}{subsection}{\textit{Et saluos reddam nonnullos. }}
\subsection*{\textit{Et saluos reddam nonnullos. }}\pstart Quia non omnes saluantur, sed tantum reliquiae, id est, ii qui credunt praedicationi Euangelicae.  \pend
\phantomsection
\addcontentsline{toc}{subsection}{\textit{Namsi reiectio illorum. }}
\subsection*{\textit{Namsi reiectio illorum. }}\pstart Iisdem perpetuo fere hic argumentis utitur. Deus non uult esse sine populo, nam uno abiecto, statime  \pend
\vspace{0.5cm}\noindent
\vspace{0.2cm}\rule{1cm}{0.2pt}\\ 
\hspace{0.2cm}\textit{mg}
\footnotesize Pauli charitas erga Iudaeos.  
\normalsize| 
\section*{IN EPIST. AD ROM.  131 }
\marginpar{[ p.]}\pstart alium sibi assumit: nos Gentes, illis abiectis, Dei con- secutae sumus gratiam, et ipsi Deo reconciliatae. QVAE ERIT ASSVMPTIO. Scilicet Iudaeorum.   \pend
\phantomsection
\addcontentsline{toc}{subsection}{\textit{Nisi uita ex mortuis. }}
\subsection*{\textit{Nisi uita ex mortuis. }}\pstart Id est, post hanc reiectionem tanquam ex morte resurgent, nam dum sunt in infidelitate, in morte sunt, neque habentur populus Dei, quemadmodum pater de filio prodigo dicit: Filius meus erat mortuus etc-  \pend
\phantomsection
\addcontentsline{toc}{subsection}{\textit{Quod si primitiae sanctae. }}
\subsection*{\textit{Quod si primitiae sanctae. }}\pstart  \pend
\section*{COMMENT. POMER.  }\pstart contemnite Iudaeos, possunt et ipsi boni esse, si in fi- de suorum patrum, quos bonos habuerunt, perstiterint-  \pend\pstart Manet Paulus in sententia, tamen aliis atque aliis utens similibus. Radix est Abraham et patres, rami sunt ludaei credentes: uidete quaeso, inquit, sanctos natura, secundum carnem esse Iudaeos, uos uero im- pios, quia non ex bona radice profecti-  \pend
\phantomsection
\addcontentsline{toc}{subsection}{\textit{Et si radix sancta. }}
\subsection*{\textit{Et si radix sancta. }}
\phantomsection
\addcontentsline{toc}{subsection}{\textit{Quod si nomnulli rami defracti sunt. }}
\subsection*{\textit{Quod si nomnulli rami defracti sunt. }}\pstart Fieri potest, ut bona arbor, habeat poma aliqua verminosa et corrupta, intemperie aurae, tamen ni hilominus arbor bona est, et fert quaedam poma optima.  corrupta igitur abiiciuntur, integra seruantur. Rami hic quidam malt sunt, id est, Iudaei quidam non credunt, hi defringuntur: inseruntur autem alii ramt, tametsi na- tura non sint bont, in radicem aut arborem natura bonam: quibus significamini uos Gentes, estis enim rami syluestris arboris, id est, filii imptetatis: nunc autem per Dei mifericordiam abiectis ex bona arbore, id est, patribus sanctis Iudaeorum, ludaeis incredulis, uos insertae estis in locum eorum: Hinc addit  \pend
\phantomsection
\addcontentsline{toc}{subsection}{\textit{Iu uero cum esses oleaster. }}
\subsection*{\textit{Iu uero cum esses oleaster. }}\pstart Id est, uos Gentes quae non eratis uera olea. Ne- mo enim colebat hanc syluestrem oleam: Oleastrim, id est Gentes, errabant seductae carnis iudicio, Iudaei uero, patres et Mosen habuerunt.  \pend
\vspace{0.5cm}\noindent
\vspace{0.2cm}\rule{1cm}{0.2pt}\\ 
\hspace{0.2cm}\textit{mg}
\footnotesize Iudaei eiecti, Gentes inser- tae.  
\normalsize| 
\section*{IN EPIST. AD ROM.  132 }
\marginpar{[ p.]}\pstart Id est, in ludaeorum arborem bonam, uos Gentes estis insitae per fidem.  \pend
\phantomsection
\addcontentsline{toc}{subsection}{\textit{Insitus fuisti illis. }}
\subsection*{\textit{Insitus fuisti illis. }}
\phantomsection
\addcontentsline{toc}{subsection}{\textit{Et consors radicis et pinguedinis }}
\subsection*{\textit{Et consors radicis et pinguedinis }}\pstart Eam gratiam estis quoque consecutae, quae ludaeis data et missa erat. Oleaster nullam habet pinguetudi- nem, uel succum bonum ad ferendos fructus suaues,  \pend
\phantomsection
\addcontentsline{toc}{subsection}{\textit{Ne glorieris aduersus ramos. }}
\subsection*{\textit{Ne glorieris aduersus ramos. }}\pstart Id est, ideoque cum non sitis naturales rami oleae, ne- gloriemini aduersus Iudaeos, qui iam sunt rami defracti.  \pend
\phantomsection
\addcontentsline{toc}{subsection}{\textit{Quod si gloriaris, non tu radi- cem portas. }}
\subsection*{\textit{Quod si gloriaris, non tu radi- cem portas. }}\pstart Stultissime feceritis o Gentes, si gloriemini aduer- sus Iudaeos: radix est Iudaeorum, non Gentium : radix, id est, patres fidei uos portant, innitimini enim promii- sionibus ipsis factis. Vos non portatis radicem, id est non est profecta salus ex Gentibus, neque illis datae sunt promissiones, quare nolite superbire in Iudaeos.  DICES IGITVR DEFRACTI SVNT RAMI. Id est, abiecti Iudaei.  \pend
\phantomsection
\addcontentsline{toc}{subsection}{\textit{Vt ego insererer. }}
\subsection*{\textit{Vt ego insererer. }}\pstart Id est, nos Gentes reciperemur in gratiam. Recte, inquit Paulus. Sed uidete quaeso, cur abiecti sint Iudaei, et terreamnini: idem uobis euenire potest quod illis, iam uos in eam arborem per fidem estis iserte.  \pend
\vspace{0.5cm}\noindent
\vspace{0.2cm}\rule{1cm}{0.2pt}\\ 
\hspace{0.2cm}\textit{mg}
\footnotesize Non glorie- mur aduersus Iudaeos.  
\normalsize| 
\section*{COMMENT. POMER.  }\pstart quid si relaberemini in priorem infidelitatem, nonne et uos rursus abiiceremini.  \pend\pstart Nam si Deus naturalibus ramis. non pepercit.  \pend\pstart Est amplificatio: si Deus Iudaeos qui ex sanctis pa- tribus erant procreati, abiecit propter incredulitatem, quid futurum est uobis, qui estis malae arboris filii?  \pend
\phantomsection
\addcontentsline{toc}{subsection}{\textit{Ecce igitur bonitatem. }}
\subsection*{\textit{Ecce igitur bonitatem. }}\pstart Haec ex praedictis infert. q.d. In his omnibus ali- ud non est cernere, quam Dei bonitatem in credentes, seueritatem in impios et contemptores sui Verbi : non abiicit nisi malos, non seruat nisi bonos. Vos Gen- tes si manseritis in bonitate, in fide, erit uobis bonus: sin uero gloriari uolueritis de Iudaeorum calamitate, quid aliud expectabitis, quam ut rursus excindami- ni? Illi uero Iudaei si resipuerint, uobis abiectis, rursus adsumentur. Tam bonus, tanque ettam seuerus est Deus, ut non patiatur eos perire, quos videt in viam redire, neque etiam ferat eos qui cum boni fue- rint, degenerent. Vide igitur quam mire se Paulus transformet? nunc Iudaeos, nunc Gentes consolatur, ut discamus ad omnes pertinere De i gratiam, Deum non uelle abiicere ullum populum. Sed unde haec ha- bet Apostolus? ex charitatis aestu, hoc enim impelli- tur, ut simul consoletur et terreat, plus tamen in iis omnibus inest consolationis, quam terroris, si uideas.   \pend
\phantomsection
\addcontentsline{toc}{subsection}{\textit{}}
\subsection*{\textit{}}
\vspace{0.5cm}\noindent
\vspace{0.2cm}\rule{1cm}{0.2pt}\\ 
\hspace{0.2cm}\textit{mg}
\footnotesize Gratia ad om- nes pertinet.  
\normalsize| 
\section*{EPIST. AD ROM.  }
\marginpar{[ p. ]}
\phantomsection
\addcontentsline{toc}{subsection}{\textit{Non enim uolo uos ignorare fratres. }}
\subsection*{\textit{Non enim uolo uos ignorare fratres. }}\pstart Id est, uos Gentes, scitote non temere factum esse ꝙ ludaei sint abiecti, uos uero assumptae : Est admi- rando Dei et arcano consilio factum, qui nouit quo modo unumquemuis trahat ad se et allictat, quare no- lite esse insolentes: aut qui uideamini uobis esse me- liores et sanctiores ludaeis, cauete ne in hanc opi- nionem incidatis: Seruite Deo in timore potius, sic enim consistetis.  \pend
\phantomsection
\addcontentsline{toc}{subsection}{\textit{Exparte. }}
\subsection*{\textit{Exparte. }}\pstart Id est, aliquibus tantum Iudeis, nempe iis, qui erant increduli.   \pend
\phantomsection
\addcontentsline{toc}{subsection}{\textit{Donec plenitudo aduenerit. }}
\subsection*{\textit{Donec plenitudo aduenerit. }}\pstart Id est, omnes Gentes quae erant conuertendae. No- uit autem solus Deus hanc plenitudinem.  \pend
\phantomsection
\addcontentsline{toc}{subsection}{\textit{Et lictotus Israel saluus erit. }}
\subsection*{\textit{Et lictotus Israel saluus erit. }}\pstart Multi de ludaeis nouissimo tempore conuertendis intelligunt, ego uero, de toto Ifraële, id est, omni- bus electis, siue Iudaeis, siue Gentibus intelligo: nam et Paulus in ea est sententia, ut uocet eos tantum Ifraël qui credunt Verbo Dei: Vide Cap.9. Et alios locos  \pend
\phantomsection
\addcontentsline{toc}{subsection}{\textit{Sicut scriptum est , Adueniet. }}
\subsection*{\textit{Sicut scriptum est , Adueniet. }}\pstart Esaiae hic locus, aliis uerbis in Hebraeo scribitur, non est tamen alia sententia, quod ideo moneo, ut scia- mus, nihil retulisse apud Apostolos reddere uerba,  \pend
\section*{COMMENT. POMER.  }\pstart modo adsequerentur sententiam, usi enim sunt cita- tionibus Scripturae ex libris tum uulgatis: habebant Romani Graeca magis, quam Hebraica in manibus- Sic et nos hodie ueterem translationem citamus, in iis maxime quae ad fidei articulos pertinent, ne quid possint in nobis calumniari Papistae. ADVENIET ERGO QVI liberet, id est, Messias uel Christus-  \pend
\phantomsection
\addcontentsline{toc}{subsection}{\textit{Et auertet impietates a lacob. }}
\subsection*{\textit{Et auertet impietates a lacob. }}\pstart Nam neque lacob, aut ullus patrum, neque nos po- tuimus eas auertere: solius Christi id proprium est-  \pend
\phantomsection
\addcontentsline{toc}{subsection}{\textit{Et hoc illis a me testamentum. }}
\subsection*{\textit{Et hoc illis a me testamentum. }}\pstart Scilicet erit. Aliud erit quam hactenus fuit sub Mo- se, de eo Esa. 2. sic dixit. De Syon exibit lex. q. d.  Lex in monte olim data non ualebit, haec quae iam da- tur, est lex fpiritus uitae, testamentum pacis, remissio peccatorum.   \pend
\phantomsection
\addcontentsline{toc}{subsection}{\textit{Cum abstulero peccata eorum. }}
\subsection*{\textit{Cum abstulero peccata eorum. }}\pstart Ipsi non poterunt ea suis conatibus auferre, nec po- terunt mecum inire foedus, ideoque hoc erit, uel in hoc consistet meum testamentum, cum ipsi nihil fecerint ad remissionem peccatorum, ego uero solus abstulero peccata. Sed aduerte haec uerba in Esaia non inueniri qui ergo non sunt assueti Sacris literis, non deprehen- dent certo iudicio quid quandoque Apostoli, quid item Prophetae dixerint. Paulus sententiam Prophe- tae eandem expressit claruus, tametsi non uideas eadem  \pend
\vspace{0.5cm}\noindent
\vspace{0.2cm}\rule{1cm}{0.2pt}\\ 
\hspace{0.2cm}\textit{mg}
\footnotesize Vires huma- nae. 
\normalsize| 
\section*{IN EPIST . AD ROM.  }
\marginpar{[ p.134 ]}\pstart uerba. Hieronymus hunc locum Prophetae negligen- tius transtulit: sic igitur clarius, et aptius textum digeres: Et ueniet Syoni redemptor, et iis qui se- auertent ab impietate in lacob, dicit Dominus, et hoc testamentum meum ad illos: Spiritus meus qui est in te, et uerba mea, quae dedi in os tuum? non recedent de- ore tuo, et de ore seminis tui, et de ore seminis se- minis tui, amodo et usque in sempiternum etc. Cum ergo Paulus dixit: Cum abstulero peccata eorum, Esaiam expressit cum dicit: Spiritus meus qui est in te etc- Nam quid aliud est dare Spiritum sanctum, quam au- ferre peccata? Qui ergo hunc habet, et uerba Dei perpetuo in ore, huic remissa sunt omnia peccata, quia id Dei opus est. Vbi spiritus est Domini, ibi est libertas.   \pend
\phantomsection
\addcontentsline{toc}{subsection}{\textit{Secundum Euangelium quidem. }}
\subsection*{\textit{Secundum Euangelium quidem. }}\pstart Id est, propter praesentem Euangelii praedicatio- nem sunt inimici Dei, sunt enim excaecati, et inuident Gentibus gratiam, redibunt tamen aliqui olim ad Dominum, incipient cognoscere gratiam, eruntque ; amici Dei, quemadmodum uos Gentes nunc estis. Dei est consilium, qui uoluit, ut propter uos qui era- tis alieni ab ipso, abiicerentur, utque insereremini in oli- uam, qui eratis ex oleastro producti.  \pend
\phantomsection
\addcontentsline{toc}{subsection}{\textit{Secundum electionem autem, }}
\subsection*{\textit{Secundum electionem autem, }}\pstart Que est aeterna, et cuius Deum poenitere non potest  \pend
\vspace{0.5cm}\noindent
\vspace{0.2cm}\rule{1cm}{0.2pt}\\ 
\hspace{0.2cm}\textit{mg}
\footnotesize Spiritum san- ctum dare, est auferre pec- cata.  
\normalsize| 
\section*{COMMENT. POMER.  }\pstart dicitque de iis tantum, qui olim conuertendi sunt-  \pend
\phantomsection
\addcontentsline{toc}{subsection}{\textit{Dilecti sunt. }}
\subsection*{\textit{Dilecti sunt. }}\pstart Scilicet Deo, propter patres, id est, propter fidem quam olim secuturi sunt: implebuntur in eis promis- siones olim patribus eorum factae, nam non tantum iis factae sunt, sed etiam semini, ideoque clarius exponit, quod dixit: Electos non posse perire.   \pend
\phantomsection
\addcontentsline{toc}{subsection}{\textit{Namdona quidem. }}
\subsection*{\textit{Namdona quidem. }}\pstart Scilicet, quae sunt in tempore, sed ab aeterno prae- stituta: quod Deus donare uoluit, hoc nouit priusquam nasceremur, neque potest ille mutare consilium, im- muitabilis enim est: uidetur nobis quidem mutari in iis quae facit, tamen secus habet. Non ergo illum poenitet cuiusq sui instituti aut consilii, uerum, stul- tum fuerit humana ratione id uelle perquirere-  \pend
\phantomsection
\addcontentsline{toc}{subsection}{\textit{Quemadmodum et uos. }}
\subsection*{\textit{Quemadmodum et uos. }}\pstart Scilicet, Gentes quondam increduli fuistis Deo, eratis etiam pro tempore inimici Dei, nunc uero amici  \pend
\phantomsection
\addcontentsline{toc}{subsection}{\textit{Nunc autem misericordiam estis consecuti per incredulitatem illorum. }}
\subsection*{\textit{Nunc autem misericordiam estis consecuti per incredulitatem illorum. }}\pstart  \pend
\vspace{0.5cm}\noindent
\vspace{0.2cm}\rule{1cm}{0.2pt}\\ 
\hspace{0.2cm}\textit{mg}
\footnotesize Deus non mu- u suaconsila 
\normalsize| 
\section*{IN EPIST. AD ROM.  }
\marginpar{[ p.139 ]}
\phantomsection
\addcontentsline{toc}{subsection}{\textit{Ex eo quod uos misericordiam estis adepti. }}
\subsection*{\textit{Ex eo quod uos misericordiam estis adepti. }}\pstart Quia istam gratiam uidere non possunt, nec susti- nere commercia Gentium, quas uidebant Idelolatras, unde et magis inuidere illis coeperunt, quic quid con- tigit gratiae.   \pend
\phantomsection
\addcontentsline{toc}{subsection}{\textit{Vt et ipli misericordiam conse. }}
\subsection*{\textit{Vt et ipli misericordiam conse. }}\pstart Scilicet, ubi resipuerint, in hoc enim ceciderunt, ut per fidem Gentium ad Deum prouocarentur. Ita omnia Deus mirabiliter de nobis instituit, quos sal- uos uult, hos sinit ad tempus esse blasphemos, et ir- risores nominis sui: quemadmodum hodie in quibus-  \pend
\phantomsection
\addcontentsline{toc}{subsection}{\textit{}}
\subsection*{\textit{}}\pstart  \pend
\vspace{0.5cm}\noindent
\vspace{0.2cm}\rule{1cm}{0.2pt}\\ 
\hspace{0.2cm}\textit{mg}
\footnotesize Gratia.  
\normalsize| 
\section*{COMMENT. POMER.  }\pstart perspiciamus, sitiamusque uehementius, nam nisi hoc fieret, nunq quae uirtus eius, et quantus sit erga nos fauor illius, deprehenderemus. Sui itaque haec omnia sunt consilii, non nostri, quamuis multa curiose, vel inquiramus, et interim ratiocinemur, non est quod cum illo contendamus: melius omnia (licet stultissima quandoque nobis adpareant) instituit de nobis, quam un quam cogitare possimus: uult ergo ut ipsi accepta om- nia feramus, non abiiciamus quancunque occasionem, qua nos ad se allicere conatur. Sciamus nos eos esse, qui inuenimus, non quaerentes ipsum, dumque sumus blasphemi in nomem eius, tamen nobis nescientibus, nos ad eum properare. Quid igitur gloriabimur, quam uel simus, aut fuerimus sine peccato? aut ꝗ statim au- dita Euangelii praedicatione obsecuti sinus? grauissi- me impegimus in illum ante. Nolunt et hodte multi, ut maxime audiant praedicare, gratiam suscipere, ne- gant eam, et blasphemant: uerum iudicio Dei fit, quod eos agit, et in hoc impellit, quod quidem hor- rendum est auditu: et hic caro se continere non po- test, quin contra Deum expuat, et horribiliter in- saniat, cum non nouerit, quam bono animo, et pio- consilio id faciat, nescit ideo esse blasphemos, et con temptores Verbi quosdam, ut aliquando resipiscant, sentiantque maiori affectu, quantae sint diuitiae boni- tatis illius, qui haec agit sic et instituit, non ut per- dat suos, sed saluet : ergo addidit Paulus.  \pend
\vspace{0.5cm}\noindent
\vspace{0.2cm}\rule{1cm}{0.2pt}\\ 
\hspace{0.2cm}\textit{mg}
\footnotesize Gratia De- enmensa. 
\normalsize| 
\section*{IN EPIST. AD ROM.  }
\marginpar{[ p.136 ]}
\phantomsection
\addcontentsline{toc}{subsection}{\textit{Vt omnium milereretur. }}
\subsection*{\textit{Vt omnium milereretur. }}\pstart Ecce in ira uides misericordiam concludi: irae est, quam in peccato conclusit, gratiae est ꝗ misereatur, ergo qui conclusit, hic quoque solus miseretur nostri: non sa- mus in uinculis aut damnatione alicuius hominis iudi- cis, sed Dei, quare nemo nos extraxerit nisi ipse solus.  Qui ingemiscunt, et clamant, uidentes hanc damna tionem sui, illi exaudiuntur, et eripiuntur: qui uero non uident, quales sunt omnes Iusticiarii, ii in aeterno carcere detinebuntur: pater uult agnosct, non tyran- nus, si tantum in hoc intenderis animum, ꝗ conclu- sus sis sub peccatum, tyrannum habiturus eras perpetuum- Vbi autem hoc paternum cor quoque agnoueris, ut om- nium miseretur, patrem cernes ultro se tibi offerentem, si sitis misericordiam, misericors est, si iratum conci- pis, iratus est, itaque et Iudaei, et Gentes, hanc mise- ricordiam inuocantes saluantur. Neque uero etiam hic nos moretur praedestinationis quaestio, certa et rata haec sit sententia: quicunque inuocauerit nomen Do- mini saluus erit, non terreant peccata: et Paulus blas- phemus fuit, nunc tamen electum Dei organum habetur-  \pend
\phantomsection
\addcontentsline{toc}{subsection}{\textit{O profunditatem. }}
\subsection*{\textit{O profunditatem. }}\pstart Nolumus ut quispiam plura de praedestinatione a nobis expectet, q̃ Paulus dixit, breui enim senten- tia, omnia carnis mirabiliter excogitata problemata, soluit: si gratiam adsecutus es, fide, postq̃ peccatum tuum agnouisti, Deus non te deseret, siuero hanc non  \pend
\vspace{0.5cm}\noindent
\vspace{0.2cm}\rule{1cm}{0.2pt}\\ 
\hspace{0.2cm}\textit{mg}
\footnotesize Liber antur qui- se damnatos agnoscunt 
\normalsize| 
\section*{COMMENT. POMER.  }\pstart es adsecutus, nihil profuerit tibi ꝗ multa de praedesti natione disputes. Certe aliquid est gratiae, anxium esse te de praedestinatione, quia curae tibi est salus, sed eate nus tamen, ut omnia committas Deo, si hoc potueris facere, non anxium te reddet praedestinatio: at ubi ab ea salutis tuae rationem ordiri uolueris, in despera- tionem cades, nam negligis ea, quae te primum Deus scire uoluit. Ergo terret hac exclamatione superbos, qui praesumunt ea scire quae Dei sunt, metirique ; diui- na, humana ratione. q. d. Quid est q multa inquiri- tis, non potest horum reddi racio, in Det sunt manu et uoluntate, pudeat uos stultos uestrae curiositatis, quin potius in uos prinum respicite, qui ipsi sitis. Con solatur etiam interim adflictos. q. d.ut maxime Dei arcana ignoretis, tamen certi sitis Deum uos non deserturum, etiam si totum regnum Sathanae uos op- primat: uidetis hic abyssum ipsam, quam nondum pe netrare potestis, sed confidite tantum, ipse, quia pa terfamilians est, omnia secundum misericordiam suam. de uobis curabit.  \pend
\phantomsection
\addcontentsline{toc}{subsection}{\textit{Diuitiarum. }}
\subsection*{\textit{Diuitiarum. }}\pstart Id est abundantiae. ET SAPIENTIAE.  Quia Deus solus sapiens est, et cognoscit omnia, praeuidet et praescit omnia priusq fiant. Ad hanc sa- pientiam nullus mortalium adsurgere potest, quare ipsi gloria relinquenda est, et nos pace contenti simus iuxta Angelorum carmem, quod cecinerunt Christo nato  \pend
\vspace{0.5cm}\noindent
\vspace{0.2cm}\rule{1cm}{0.2pt}\\ 
\hspace{0.2cm}\textit{mg}
\footnotesize Praedestina tionis usus.  
\normalsize| 
\section*{IN EPIST. AD ROM.  }
\marginpar{[ p.u7 ]}\pstart Multa Deus fecit, quae ratio nostra non potest con cipere, tamen ausi sunt stulti ea tentare, sed oppres- si sunt a gloria, quemadmodum uidemus Philosophis accidisse: et quotus quisque est nostrum, qui intelligat quae sit substantia animae, cum tamen ex ea pendeat omnis nostra uita animalis: circunferimus eam, estque  ubique nobiscum, et ignoramus quid sit. Quam stul- tum igitur et ridiculum est, nos uelle diuina intellige re, cum haec uulgaria non capiamus: si in iis quae sunt naturae pugnant sapientum mundi sententiae, quanto magis pugnabunt in diuinarum rerum inda- gatione? Ergo consilia et iudicia sua, Deo relinqua mus, non sumamus in iis nobis quod supra captum nostrum est.   \pend
\phantomsection
\addcontentsline{toc}{subsection}{\textit{Quam inscrutabilia. }}
\subsection*{\textit{Quam inscrutabilia. }}
\phantomsection
\addcontentsline{toc}{subsection}{\textit{Et inperuestigabiles uiae eius. }}
\subsection*{\textit{Et inperuestigabiles uiae eius. }}\pstart Viae Det dicuntur ea, quaecunque Deus agit, quo- rum omnium nos nullam rationem reddere possumus, nam solus qui haec agit, rationem nouit.   \pend
\phantomsection
\addcontentsline{toc}{subsection}{\textit{Quis enim cognouit mentem Dñi? }}
\subsection*{\textit{Quis enim cognouit mentem Dñi? }}\pstart Id est, quis sententiam eius, aut quid ipse uelit, tenet? Haec in Esaia et alibi scripta legimus.  \pend
\phantomsection
\addcontentsline{toc}{subsection}{\textit{Aut quis illi fuit a consiliis ? }}
\subsection*{\textit{Aut quis illi fuit a consiliis ? }}\pstart Scilicet, ut faceret sic uel sic, quis ipsum edocuit? Qu'indo igitur fecerit omnia, antequa nasceremur, ut uoluit, cur illius iudicium et uoluntatem accusamus?  \pend
\vspace{0.5cm}\noindent
\vspace{0.2cm}\rule{1cm}{0.2pt}\\ 
\hspace{0.2cm}\textit{mg}
\footnotesize Non inucsti- ganda, quae supra captum nostrum.  
\normalsize| 
\section*{COMMENT. POMER.  }
\phantomsection
\addcontentsline{toc}{subsection}{\textit{Aut quis prior dedit illi? }}
\subsection*{\textit{Aut quis prior dedit illi? }}\pstart Et hoc contra nostram superbiam dictum est, qua fit, ut saepe irascamur Deo, quando ille nostras iusti- cias damnat. Commendatur uero gratia, cui soli acce- ptum ferre debemus, si saluamur, nam quaecumque nostrarum sunt uirium, per ignem purgabuntur, imo absumen- tur, ut id solum gloriam habeat, quod ex ipso est-  \pend
\phantomsection
\addcontentsline{toc}{subsection}{\textit{Quoniam ex illo etc. }}
\subsection*{\textit{Quoniam ex illo etc. }}\pstart Scilicet, sunt omnia, nihil est uspiam quod non ex ipso prodierit, pertinetque hoc ad creationem.   \pend
\phantomsection
\addcontentsline{toc}{subsection}{\textit{Et per illum. }}
\subsection*{\textit{Et per illum. }}\pstart Quia omnia conseruat, prouidet et gubernat, nam quemadmodum omnia ex ipso tantum processe- runt, sic quoque non possunt per alium administrari praeter ipsum.   \pend
\phantomsection
\addcontentsline{toc}{subsection}{\textit{Et in illum. }}
\subsection*{\textit{Et in illum. }}\pstart Scilicet redeunt et refluunt. Quid igitur ipsum ae- cusabimus, quod alios damnet, alios saluet?  \pend
\phantomsection
\addcontentsline{toc}{subsection}{\textit{Ipli gloria in secula. }}
\subsection*{\textit{Ipli gloria in secula. }}\pstart Ecce, nihil nostris uiribus tribuitur, nulla relin- quitur gloria, sed tantum confusio, quod etiam peti- mus in oratione: Sanctificetur nomen tuum, id est, nomen et gloria nostra pereat, et tu solus glorifi- ceris.   \pend
\marginpar{[ p.]}\pstart  \pend\pstart  \pend
\vspace{0.5cm}\noindent
\vspace{0.2cm}\rule{1cm}{0.2pt}\\ 
\hspace{0.2cm}\textit{mg}
\footnotesize 
\normalsize| 
\section*{COMMENT. POMER.  }\pstart Christum data fuerit, facies opera bona: alia quae ad mortificationem tuam: alia uero quae ad proximum pertine bunt, nulla autem ad Deum, cui tantum fide et spe, non operibus externis seruitur: nam est spiritus De- us, ergo tantum in spiritu uult coli, non externis ce remoniis, iohan.4. Bona opera, quae ad nostram per sonam pertinere dixi, ea omnia talia sint oportet, ut per ea eneces adfectus carnis tuae, quae semper resistit spiritui, ut in Septimo Capite dictum est, idque tam diu facies, donec tota caro moriatur, nam ipsa uiuente, semper uigent in ea adfectus mali. Reliqua omnia ad proximi seruitutem pertinent, quibus uel ipsius ege- statem leues, aut ipsum consoleris adflictum, aut ex errore reuoces, si uideas errantem: itaque sic orditur.  \pend
\phantomsection
\addcontentsline{toc}{subsection}{\textit{Obsecro uos per miserationes Dei. }}
\subsection*{\textit{Obsecro uos per miserationes Dei. }}\pstart Ecce et his uerbis gratiam et fidem eis inculcat, qui sunt Christiani-q. d. uos qui iam scitis et credi- tis Dei misericordia uos iustificatos, nihilque deesse uobis, quod ad salutem pertineat, moueamini quaeso, per illam misericordiam et beneficentiam Dei erga uos, ut omnia uestra opera ex fide instituatis: nam quicquid non est ex fide, peccatum est. Sed posset hic moueri quaestio, cur Paulus leges iis praescribit, qui- bus lex non est posita, quiq ab omni lege liberi sunt? Respondendum erit: leges has Christianis tantum scri- bi et propont, id est, iis qui credunt, et in quibus  \pend
\vspace{0.5cm}\noindent
\vspace{0.2cm}\rule{1cm}{0.2pt}\\ 
\hspace{0.2cm}\textit{mg}
\footnotesize Opera bona- 
\normalsize| 
\section*{IN EPIST. AD ROM.  }
\marginpar{[ p.139 ]}\pstart spiritus Dei sponte omnia facit, imo et secum adfert: quod si absit hic sptritus, Christianus dici nequit, neq potuerit quicquam facere, hilari et spontaneo ani- mo, quidquid faciet, id faciet tyranno et lege cogen- te. Vides uero diuersum hic in Paulo, non est legisla- tor, non minatur mortem, aut lapidationes transgres soribus, ut et nostri faciunt, cum dicunt: Manda- mus, Aggrauamus, Reaggrauamus, Anathemisa- mus etc. qui putant sibi potestatem tantum in destru- ctionem datam, ideoque suauissime inquit: Obfecro uos fratres, quid potuit esse lentus et mansuetius? non est haec uox Mosi, atque ita ne leges quidem istae sunt, nam sponte fiunt tantum per monitorem, non per exactorem. Sed oritur hic altera quaestio: quid opus est haec Christiano praescribi, qui per se non indigens legislatore, omnia faciat, Respondeo: bona opera hic non praescribuntur, ut doceantur, nam norunt ea omnia Christiani: praeterea, non ob eam caussam, quod debeant sic impelli et cogi Christiani, uerum cum uideamus nos esse carnem et sanguinem, quibus semper ad oscitantiam et pigriciam solicitamur, opus est piis hortationibus, quibus excitetur in nobis fides contra carnem, et eius inertiam, ut sic perpetuo exer- ceatur, ne unquam obdormiat, uacetque ; iis operibus quae testentur et coram Deo et proximo fidem no- stram esse integram et constantem. Itaque ; dicit Pau- lus, per eas miserationes, per quas redempti estis, uos  \pend
\vspace{0.5cm}\noindent
\vspace{0.2cm}\rule{1cm}{0.2pt}\\ 
\hspace{0.2cm}\textit{mg}
\footnotesize Euangelio, cur leges admisce- antur.  
\normalsize| 
\section*{COMMENT. POMER.  }\pstart obsecro, non autem est simplex obsecratio, sed talis quam Christiani negligere nequeant : estque ; potius ad iuratio, q̃ obsecratio, quam non uident leuiores spi- ritus, qui omnia et diuina et humana contemnunt-  \pend
\phantomsection
\addcontentsline{toc}{subsection}{\textit{Vt praebeatis corpora uestra. }}
\subsection*{\textit{Vt praebeatis corpora uestra. }}\pstart In quibus adhuc est peccatum, id est, uestrum car nem et sanguinem, uel ut uno uerbo dicam, uestrum totum Adamum, non dico de pede aut manu, sed de- omni quod uos estis, illud inquam praebete-  \pend
\phantomsection
\addcontentsline{toc}{subsection}{\textit{Hostiam uiuentem. }}
\subsection*{\textit{Hostiam uiuentem. }}\pstart Per appositionem haec dicuntur, estque ; sententia Apo- stoli. Si uero cultu Deum colere uolueritis, declina- bitis oblationes Iudaicas, in quibus tantum erant irra- tionales pecudes, offeretis autem uosmetipsos totos, id est, rationem vestram, viresivestras, in hanc enim unicam hostiam respexeunt legis illae victimae : quod si alium extra hunc cultum statvere volueritis, Mona chatum, scilicet, aut alia carnalia, erit quoque irratio- nalis, id est, absque Verbo Dei. Siquidem cum iam per fidem in eorde, toti sitis consecrati Deo, absurdum fuerit in externa respicere, quibus nihil salutis compa- ratur. Ergo, uestra corpora, quae uos impediunt, solicitantque ad peccatum, offerte Deo, curate ut mo- riantur, pereantque in uobis omnes adfectus, qui con- tra Deum pugnant, uiuum totum corpus non potest Deo offerri, quare mortificandum est per totam uitam,  \pend
\vspace{0.5cm}\noindent
\vspace{0.2cm}\rule{1cm}{0.2pt}\\ 
\hspace{0.2cm}\textit{mg}
\footnotesize Victima Chri- stianorum.  
\normalsize| 
\section*{IN EPIST. AD ROM.  }
\marginpar{[ p.140 ]}\pstart ut sic per mortem ad uitam properemus. Atque interim ra- tio ordinis eorum, quae hic tractat Paulus, semper tibi proposita sit, nam cum de hostia loquitur, de re quadam eximia, quae ad solos Sacerdotes pertineat, loquitur- Scit Christianos iam antea omnes, inunctos esse oleo spi- ritus laetitiae, consecratos soli Deo per fidem, ideo nunc iubet, ut quoq ea agant, quae sunt inunctorum et con secratoru, nempe, ut exhibeant oblationes eas, quas Deus offerri ipsis mandauit, quae aliud non sunt, quam ipsa eorum corpora, quae, quia omnes sumus Sacerdo tes, sunt talia sacrificia, ut omnino ea grata a nobis accipiat Deus, propter unicum et sacrosanctum sa- crificium Christi, qui semel oblatus est pro nobis omni- bus in cruce: in huius mortem sumus baptizati, est illa nostra facta mors, uictima huius, nostra est, nisi enim ea uera fuisset uictima, id est, nisi illa placuisset patri, non posset haec nostra uictima, id est, corporum no- strorum mortificatio, grata esse patri. Vocat uero hanc hostiam uiuentem, ut scias aliam hanc esse, quae olim in lege offerebantur: hostiae illae non erant uiuae, nostra autem uiua est, id est, talis, quae mortifi cando et moriendo uiuificetur. Dum enim offerimus et occidimus nostra corpora, interim semper magis uiuificantur, et abeunt in spiritum paulatim, saeui- unt adfectus carnis in spiritum, sed sensim tantum de- ficiunt, dum sic uerbo mortificantur. Ergo multo sunt digniora haec sacrificia Noui Testamenti, sacrificiis Ve- teris-Praeterea uult esse  \pend
\vspace{0.5cm}\noindent
\vspace{0.2cm}\rule{1cm}{0.2pt}\\ 
\hspace{0.2cm}\textit{mg}
\footnotesize Sacrificium Christi.  
\normalsize| 
\section*{COMMENT. POMER.  }\pstart Id est, talem quae non sit prophana, sed segregata per Verbum, ab iis quae sunt mundi. Hoc conuenit cum Veteri Testamento, animantia enim omnia, quae de- bebant immolari, erant prophana, ubt uero ad sacer dotes fuissent adducta, sancta habebantur, scilicet, Deo Sic item nostra saonficia, non debent pertinere ad mundum, sed solum Deum. Et addit  \pend
\phantomsection
\addcontentsline{toc}{subsection}{\textit{Sanctam hostiam. }}
\subsection*{\textit{Sanctam hostiam. }}\pstart Nisi hoc addidisset, nihil esset consolationis in iis uerbis, nam uehementer terretur et dolet caro, ubi audit quod debeat se totum resignare Deo, et abne- gare suas uires, utque cogatur dicere: Fiat uoluntas tua: adueniat regnum tuum. Sed tursum consolatur nos, cum dicit: Ista uestra hostia uiuens et sancta, placet Deo, uoluntatis est paternae, ut sic patiamini et mor- tificemini, in hoc estis eius filii. Talem igitur hostiam ubi feceritis uestra corpora, instituetis cultum Dei rationalem, id est, talem, in quo omnes uestrae uires illi offerantur secundum eius uoluntatem, quae ex ip- sius Verbo cognoscitur. Quomodo autem gratius quid quam offerretis?  \pend
\phantomsection
\addcontentsline{toc}{subsection}{\textit{Acceptam. }}
\subsection*{\textit{Acceptam. }}
\phantomsection
\addcontentsline{toc}{subsection}{\textit{Et ne accommodetis uos. }}
\subsection*{\textit{Et ne accommodetis uos. }}\pstart Exponit se Paulus. Ad hoc, inquit, mortifica- mur, ut uiuificemur, quod lex in suis sacrificiis non praestabat. offerri iubebat mortua, et mancbant mor-  \pend
\vspace{0.5cm}\noindent
\vspace{0.2cm}\rule{1cm}{0.2pt}\\ 
\hspace{0.2cm}\textit{mg}
\footnotesize Carni durum est se Deo re- signare.  
\normalsize| 
\section*{IN EPIST. AD ROM.  }
\marginpar{[ p.141 ]}\pstart tua: hic uero in nostro sacrificio, quia sumus sub gratia, non sub lege, offerimus nostra corpora et mor- tificamus ea, non ut mortua sint, sed ut uiuificentur. Ecce hic est mortificatio et uiuificatio simul: ibi sola erat mortificatio, quandoquidem nullae hostiae ad ui- tam redibant.   \pend
\phantomsection
\addcontentsline{toc}{subsection}{\textit{Ad figuram. }}
\subsection*{\textit{Ad figuram. }}\pstart Quae mundi sit figura notum est: Mundus quaeril honores, diuitias, uoluptates, et suam tranquillita- tem, ignorat hero quae Dei sunt, habetque ea pro nihi- lo, ut maxime multa de Deo loquatur: haec igitur est ipsissima mundi figura concupiscentia. 1. Iohan. 2.  ad hanc, inquit, ne accommodetis uos, id est, ne fa- ciatis uos similes mundo-  \pend
\phantomsection
\addcontentsline{toc}{subsection}{\textit{Sed transformemini. }}
\subsection*{\textit{Sed transformemini. }}\pstart Scilicet, a uetustate Adae, uel a figura seculi, id est, noui homines fiatis.  \pend
\phantomsection
\addcontentsline{toc}{subsection}{\textit{Per renouationem mentis. }}
\subsection*{\textit{Per renouationem mentis. }}\pstart Non per hypocrisin: non simulate aliam uitam fo- ris, ut faciunt hypocritae, qui Cappam aut Cilicium induunt, contegentes latronem intus latitantem, in quo putant se mundum deseruisse, ac se renouatos, cum aliud nihil sit, quam deceptio: ergo, inquit, mem- tibus renouemint, non per simulationem externam. Vide, quomodo et hic nostras uires damnet Aposto- lus, omnia sunt abolenda, nihil sit tam speciosum, aut  \pend
\vspace{0.5cm}\noindent
\vspace{0.2cm}\rule{1cm}{0.2pt}\\ 
\hspace{0.2cm}\textit{mg}
\footnotesize Mundus- 
\normalsize| 
\hspace{0.2cm}\textit{mg}
\footnotesize Hypocriste. 
\normalsize| 
\section*{COMMENT. POMER.  }\pstart egregium, quod non iudicio Dei prosternatur . Verum interim aduerte, quam diligenter semper loquatur de- sacrificio hoc corporum nostrorum, id scilicet nondum esse perfectum, uerum adhuc fiert, quia sacrificamus adhuc, necdum sacrificautmus plene: et quemadmo dum Christus obtulit se patri in cruce, sic nos iam cru- ci suffixi sumus et sacrificamur, destinati sumus toti morti, attamen adhuc uiuunt corpora nostra: haben- tur pro mortuis coram Deo, tametsi nondum prorsus mortua sunt. Sed haec uita eorum, quid aliud est, qui ta mortis? per eam enim in ipsam mortem transeunt, ut postea et ex illa morte, rursus in uitam aliam no- uam constituamur, uide Paulum.1. Corinih . 16.  Itaq sacrificium et holocaustoma Christianorum est iam in altare positum, sed adhuc consummandum est, ut in Isaac historia uidemus. Hinc est quod toties di- cat: Ne accommodetis: Transformemini et c. Perfi- ciendum est sacrificium, non somnianda ulla carnis li- bertas, quam hic uehementer damnari uidemus: sumus quidem filii Det, ut dicamus: Pater noster, sed dici- mus quoque : Remitte nobis debita nostra, quibus pre- cationibus conuincimur nos esse peccatores, si pecca- tores, quid aliud faciemus, q̃ ut peccatum aboleamus per fidem, quod uere Christianum est sacrificium? Sed quod addit, mirae est suauitatis: His quorum corpora- iaem sacrificantur.  \pend
\phantomsection
\addcontentsline{toc}{subsection}{\textit{Vt probetis quae sit uoluntas Dei. }}
\subsection*{\textit{Vt probetis quae sit uoluntas Dei. }}
\vspace{0.5cm}\noindent
\vspace{0.2cm}\rule{1cm}{0.2pt}\\ 
\hspace{0.2cm}\textit{mg}
\footnotesize Sacrificium Christianorum non semel per ficitur. 
\normalsize| 
\section*{IN EPIST. AD ROM.  }
\marginpar{[ p.142 ]}\pstart Probetis inquam, id est, fide cognoscatis, quod haec sit Det bona uoluntas, uos sic offerri, quemadmo dum Christus oblatus est. Gregorius et caetert, multa scripserunt de moribus Christianorum, quibus mise- ras uexarunt conscientias, sed haec alia sunt et multo digntora, quae hic legis. Hos mores, non format hu- mana prudentia uel industria, sed ipse spiritus eos se- cum adfert: ipsa ergo inquit experientia uos docebit num filii Dei sitis, num secundum ipsius uoluntatem ambuletis: carnales nunquam sunt certi, num eorum opera aut uita accepta sit Deo et perfecta, id est, ta- lis, quam Deus requirat, nam non sciunt in conscien- tiis suis, quae Dei sit uoluntas bona. Subiicit itaque for- mulas quasdam carnis mortificandae: quaedam eo per- tinent, ut adfectus cordis, quibus ad quaeque turpia et iniqua propensus es, eneces: aliae quomodo proxi- mum obserues. Non est uero hic multa disputatio, quomodo primum Venerem ieiuniis, aut aliis repri- mas, sed quomodo sublimem illum adfectum tuum per- das et destruas, ne scilicet infleris in animo tuo, sed humilieris, et abiiciaris coram teipso, et nisi hoc fa- ctum fuerit, fieri nequit, ut uitiis crassis, in tuam car nem et proximum, habenas imponas.  \pend
\phantomsection
\addcontentsline{toc}{subsection}{\textit{Dico igitur uobis per gratiam. }}
\subsection*{\textit{Dico igitur uobis per gratiam. }}\pstart Id est, per hoc Euangelium, quod mihi datum est, ut praedicem CVILIBET.  \pend
\vspace{0.5cm}\noindent
\vspace{0.2cm}\rule{1cm}{0.2pt}\\ 
\hspace{0.2cm}\textit{mg}
\footnotesize Formulae mor- tificandae car- nis 
\normalsize| 
\section*{COMMENT. POMER.  }\pstart Neminem excipio, hae leges omnibus praescribun tur, ut scias neminem esse, qui iis uitiis non inpingue tur, si enim abessent, cur haec moneret Paulus?  \pend
\phantomsection
\addcontentsline{toc}{subsection}{\textit{Ne quis arroganter de se sentiat. }}
\subsection*{\textit{Ne quis arroganter de se sentiat. }}\pstart Ne scilicet quis se aliquid esse putet, et alteri se- praeferat, quod natura nobis innatum est.   \pend
\phantomsection
\addcontentsline{toc}{subsection}{\textit{Supra quam oportet de se sentire- }}
\subsection*{\textit{Supra quam oportet de se sentire- }}\pstart Licet quidem, ut aliquid de nobis sentiamus, sed modeste et sobrie, id quod antea dixit: Ne quis arro- ganter sentiat . Sciat igitur Christianus, sibi gratiam datam, aut donum aliquod, non ut in hoc glorictur, et prae se contemnat fratrem, sed ut per illud fratri seruiat: qui donum habet Prophetiae, is sciat et sen- tiat se habere, interpretetur Scripturas, ut alii quoq inde suum peccatum agnoscant et aedificentur: ue- rum Sathanae astus non sinit nos esse modestos et so- brios in usu donorum, nam id agit, ut gloriemur in iis, putemusque multo nos esse meliores aliis. Si quid habemus, Det gratiae acceptum referamus, omnia enim ex illo habemus, nihil ex nobis: fur est et latro, qui sibi id arrogat, quod non suum est, sed Dei, ergo ad mensuram de nobis sentiamus, quod donis nostris uta- mur, non in nostram, fed illius gloriam faciamus, ag- noscamusque per haec nos non factos esse dominos, sed seruos aliorum. Siquidem in iis (si modo Christiani esse- mus) esset cernere, quomodo quisque alteri seruire de-  \pend
\vspace{0.5cm}\noindent
\vspace{0.2cm}\rule{1cm}{0.2pt}\\ 
\hspace{0.2cm}\textit{mg}
\footnotesize Donis quomo doutendum.  
\normalsize| 
\section*{IN EPIST. AD ROM.  143 }
\marginpar{[ p.]}\pstart beat, secundum rationem sui doni, id quod docet ip- sa fides: ergo adiecit  \pend
\phantomsection
\addcontentsline{toc}{subsection}{\textit{Vt cuique Deus partitus est mem- suram fidei. }}
\subsection*{\textit{Vt cuique Deus partitus est mem- suram fidei. }}\pstart Hic enim minus, hic plus credit: his haec, illis alia collata sunt dona, et quisque utitur pro mensura fi- det, id est, ut quenque impellit et hortatur sua fides, ad feruiendum proximo per ea dona, ita seruit. Qui- dam hodie se doctores esse putant, uoluntque omnes homi- nes docere, cum nondum possunt esse discipuli: quod quidem non est sentire secundum mensuram fidei, sed contra fidem. Sed uide, quomodo id similitudine quoq eleganti et satis manifesta nobis proponat.   \pend
\phantomsection
\addcontentsline{toc}{subsection}{\textit{ Quemadmodum enim in uno corpore membra. }}
\subsection*{\textit{ Quemadmodum enim in uno corpore membra. }}\pstart Oculus aliis membrum nobiltus uidetur, attamen, quo nobilior est, hoc magis omnibus aliis membris ser- uit . Item, manus praestat caeteris membris, sed non potest gloriari quod nobilius sit et dignius, imo ui- det se esse seruum omnium aliorum membrorum, par ratio est de omnibus membris, quae omnia unum cor- pus simul constituunt. Sic inter Christianos fieri quoq debet, inter hos non sit dominus aut tyrannus, ut ma- xime alius nobilius habeat donum altero, tamen non putet se esse in hoc dominum, sed seruum aliorum- Haec digna erant, quae in usum prouerbiorum trahe- rentur.   \pend
\section*{COMMENT. POMER.  }
\phantomsection
\addcontentsline{toc}{subsection}{\textit{Singulatim autem. }}
\subsection*{\textit{Singulatim autem. }}\pstart Id est, cum simus diuersa membra unius corporis Christi, non sumus nostra membra, sed aliorum: ut oculus non suum est membrum, sed aliorum membro- rum, et manus, totius corporis, non suum est mem- brum.Ita quoque et nos seruiamus aliis, non abuta- mur donis, quae nobis data sunt in proximi utilitatem, non ut nobis ea tantum sumamus, quod tam stultum est et ridiculum, quamut si manus dicat: manus sum, non propter alia membra, sed tantum propter me: Quae, dic quaeso, haec manus fuerit?  \pend
\phantomsection
\addcontentsline{toc}{subsection}{\textit{Sed tamen dona habentes. }}
\subsection*{\textit{Sed tamen dona habentes. }}\pstart Scilicet a Deo . IVXTA GRATIAM DA- TAM NOBIS VARIA. Non a nobis ea habe- mus, quo tollitur et meritum et studium nostrum, iuxta illud in Corinth. Quid habes quod non accepisti? Si respexeris simpliciter in corpus, nullam inuenies mem- brorum differentiam, uerum obseruatis officiis membro- rum, differre alia ab aliis uidebis, Ephes.4. Alios dedit Apostolos, alius docet, alius praeest Reip. hic literis intentus est etc. Vnum corpus est, sed diuer- sa officia.  \pend
\phantomsection
\addcontentsline{toc}{subsection}{\textit{Siue prophetiam. }}
\subsection*{\textit{Siue prophetiam. }}\pstart Si quis habet uirtutem interpretandi Scris aras, is ha- beat eam secundum fidei mensuram et proportionem: id est, utatur illa in fide, non arroganter, non doceat ea  \pend
\vspace{0.5cm}\noindent
\vspace{0.2cm}\rule{1cm}{0.2pt}\\ 
\hspace{0.2cm}\textit{mg}
\footnotesize Contra meri- unostra.  
\normalsize| 
\section*{IN EPIST. AD ROM.  }
\marginpar{[ p.144 ]}\pstart quae nescit, aut de quibus dubitat . Hic nullum illorum do- norum magnifacit Paulus, quae nostri Papistae in coelum ferebant, nempe, castitatem, edere semel in die, in ci- licio dormire etc. Dona quae hic numerantur, sunt dona spiritus, quibus proximo seruitur, tametsi uide- antur hypocritis esse minima et contemnenda, ha- bent enim sua, in quae respiciunt, nitentia multa et egregia hypocrisi. Primas itaque ; dat Apl's prophetiae, que eius usus in primis sit nobis necessarius, is eni pau- perrimus est (ut maxime opes possideat) qui Verbo Dei caret et non intelligit, quae sit uoluntas Dei bona: aperit autem nobis Prophetia, Verbi et uoluntatis diuinae arcana, nam uel Scripturas interpretatur, aut in san- ctis futura praedicit, ad cauendum scandala et Sathanae astus: in iis, cum fides omnia penetret et adsequatur, fieri nequit, ut quid inuenias, quod rationi fidei non sit consentaneum, tamen si audis quosdam, qui interpreten- tur Scripturas, et praedicant futura, contra fidei amus- sim, scias id non ex Dei, sed Sathanae spiritu profi- cisci, quales hactenus audiuimus in uniuersitatibus Papae, qui omnes Scripturas, ad Aristotelis somnia, exhibebant, et hodie audimus alios, qui subtilius, et alia impostura prophetant. Noli igitur credere pro- phetiae, quae fidei consimilis non est: qui prophetat, in primis certus sit, se ueritatem tenere, qui incertus est, aut qui non credit uera esse quae dicit, is mendax est, et falit dlios . Tales Prophetae fuerunt Ananiae  \pend
\vspace{0.5cm}\noindent
\vspace{0.2cm}\rule{1cm}{0.2pt}\\ 
\hspace{0.2cm}\textit{mg}
\footnotesize Prophetia.  
\normalsize| 
\section*{COMMENT. POMER.  }\pstart et Sedechias. Videant ergo nostri hodie, qui uolunt Scripturas interpretari, quomodo certi sint in sua con scientia de interpretationibus. Illi si in locum aliquem obscurum incidant, non uident, num eorum exposi tio fidei conueniat, sed hoc saltem, ut noui aliquid et arguti, quod tum spirituale dicunt, dicant: quae res facit, ut tandem uniuersa fere Scriptura in dubium uo cetur, nam uno loco perperam et contra simplicita- tem fidei exposito, necesse est, ut pugnantes multos locos deprehendas. Quare qui fidet regula exposue rit, difficiles locos, non errabit, neque alios quoque in errorem pertrahet. Idem est Dei spiritus, eadem sen tentia per totam Scripturam, non potest esse sui dissi milis: sed de iis satis.  \pend
\phantomsection
\addcontentsline{toc}{subsection}{\textit{Siue ministerium in administrando. }}
\subsection*{\textit{Siue ministerium in administrando. }}\pstart Id est, si quis minister est aliorum, pau perum aut infirmorum, diligenter illis administret, quales sunt Diaconi qui in ecclesiis haberi debent: sicut oculus non negligit suum officium, pes non quiescit perpetuo, ita nemo in ecclesia, segnius ea quae sui sunt muneris curet- Curabit autem quisque sua, si timor Domini (qui est ipsa fides) defuerit.  \pend
\phantomsection
\addcontentsline{toc}{subsection}{\textit{Siue qui docet in doctrina. }}
\subsection*{\textit{Siue qui docet in doctrina. }}\pstart Si cui datum est, ut docere possit, doceat, agat id in fide: et cui exhortatio credita est, is quoque exhor- tetur per fidem. Petrus: Qui loquitur quasi sermones  \pend
\vspace{0.5cm}\noindent
\vspace{0.2cm}\rule{1cm}{0.2pt}\\ 
\hspace{0.2cm}\textit{mg}
\footnotesize Scripturae in terpretatio fi deconuenat. 
\normalsize| 
\section*{145 IN EPIST. AD ROM.  }
\marginpar{[ p.]}\pstart Dei. Doctrina prior est, quia ignorantes docemus: exhortatio posterior, nam scientes exhortami. Mul- ti hodie sunt, qui uix doctrina nostra aegeant, tamen omnes exhortatione opus habent, cum simus adhue caro et sanguis, et quos Sathanas semper suis insi- diis obseruet.   \pend
\phantomsection
\addcontentsline{toc}{subsection}{\textit{Qui impertit, in simplicitate. }}
\subsection*{\textit{Qui impertit, in simplicitate. }}\pstart St quis donum a Deo habet, ut aliis largiri possit: is simplicissimo corde largiatur, non spectet homi- num plausum aut gloriam, non sciat eius sinistra quid dextra fecerit, Det gloriam tantum spectet, non sua me- rita: sciat se gratis accepisse a Deo, gratis igitur det, et semper omni aegenti.   \pend
\phantomsection
\addcontentsline{toc}{subsection}{\textit{Qui praeest, in diligentia. }}
\subsection*{\textit{Qui praeest, in diligentia. }}\pstart Si cut commissa est Respub-diligenter praesit. quod contra principes nostros dictum est, item Epi- scopos, qui non uocationem suam uident, sed tantum uentris sui curam, et mundi delitias.  \pend
\phantomsection
\addcontentsline{toc}{subsection}{\textit{Qui miseretur, in hilaritate. }}
\subsection*{\textit{Qui miseretur, in hilaritate. }}\pstart Si quis aegentibus subuenire uelit, faciat id hilari spiritu, non tristi, aut impedito, nam hilarem dato- rem diligit Deus. 2. Corinth.s.  \pend
\phantomsection
\addcontentsline{toc}{subsection}{\textit{Dilectio sit non simulata. }}
\subsection*{\textit{Dilectio sit non simulata. }}\pstart Vt diligas, scilicet, proximum ex corde sicut te ipsum, neque solum amicum, sed et inimicum, et si opus sit, etiam cum periculo uitae tuae ipsi seruias  \pend
\vspace{0.5cm}\noindent
\vspace{0.2cm}\rule{1cm}{0.2pt}\\ 
\hspace{0.2cm}\textit{mg}
\footnotesize Principum et Episcoporum officia.  
\normalsize| 
\section*{COMMENT. POMER.  }
\phantomsection
\addcontentsline{toc}{subsection}{\textit{Sitis odio persequentes quod malum est. }}
\subsection*{\textit{Sitis odio persequentes quod malum est. }}\pstart Id est, absque omnibus insidiis, dolo, hypocrisi et scandalo agite. Haec praescribuntur Christianis tantum, qui non sunt Christiani, ne umbram quidem horum norunt  \pend
\phantomsection
\addcontentsline{toc}{subsection}{\textit{Adhaerentes ei quod bonum est. }}
\subsection*{\textit{Adhaerentes ei quod bonum est. }}\pstart Id est, quod integrum, candidum et sanctum est, complectimini ex corde.  \pend
\phantomsection
\addcontentsline{toc}{subsection}{\textit{Per fraternam charitatem ad mu- tuo uos diligendos propensi. }}
\subsection*{\textit{Per fraternam charitatem ad mu- tuo uos diligendos propensi. }}\pstart Id est, ad hoc animi uestri sint inclinaum, quod uos fratres esse testetur. Vtinam sic adfecti essemus hodie, sciremus enim omnes homines esse fratres nostros, nul- la esset dissensio, nemo se altero meliorem putatret, nullas sectas excitarent somnia humana. Impie ergo docuerunt, qui dixerunt, ordinata charitas incipit a se: uerum est, charitatem ueram, ordinatam esse opor- tere a Deo, sed non confictam nostra opinione car- nis, sed ea non a se incipit, uerum in proximum pergit: Christianus homo omnia sua resignat aliis, tan- tum abest, ut a se incipiat: fides aliud non potest do- cere, q̃ ut teipsum odio habeas, et Christus non a se- incepit . Itaque impium est, ut praecipiatur quomodo nos amemus, siquidem plus satis huius ueneni habe- mus a natura. Vide ergoquomodo uocabulum chari- tatis in scriptura intelligas, non est alius affectus, q̃ quo  \pend
\vspace{0.5cm}\noindent
\vspace{0.2cm}\rule{1cm}{0.2pt}\\ 
\hspace{0.2cm}\textit{mg}
\footnotesize Charitas 
\normalsize| 
\section*{IN EPIST. AD ROM.  }
\marginpar{[ p.140 ]}\pstart rapimur ad seruiendum proximo, non ad gloriam, aut famam nobis conciliandam.  \pend
\phantomsection
\addcontentsline{toc}{subsection}{\textit{Fonore alius alium praecedentes. }}
\subsection*{\textit{Fonore alius alium praecedentes. }}\pstart Non sic simus adfecti, ut eum tantum honoremus, qui nos honorat, ut maxime maius donum habeamus, tamen alteri, qui minus habere uidetur, honor primum habea- tur a nobis: hoc est Charitatis, quae non uult, ut pro loco, aut honore contendamus. Qui ex his uerbis non se, sedalios putant admonitos, Christianum sentien- tes esse, ut ipsi neminem honorent, se uero honoran- dos putent, asini sunt, ferula digni, quemadmodum et illi, qui ex admonitione ut demus, censent sibi debere da- ri, non se debere dare, homines. s. egregie Christiani.  \pend
\phantomsection
\addcontentsline{toc}{subsection}{\textit{Studio non pigri. }}
\subsection*{\textit{Studio non pigri. }}\pstart Idest, alii pro aliis soliciti sitis, tanquam quodlibet membrum pro toto corpore solicitum est-  \pend
\phantomsection
\addcontentsline{toc}{subsection}{\textit{Spiritu feruentes. }}
\subsection*{\textit{Spiritu feruentes. }}\pstart Non sitis securi secundum carnem: egregia est Iro- nia. 1.Corinth.5. contra eos qui iam se omnia tenere putant: Iam saturi estis etc. Valt ergo nos strenue pergere in iis quae coepimus Christiani.  \pend
\phantomsection
\addcontentsline{toc}{subsection}{\textit{Tempori seruientes. }}
\subsection*{\textit{Tempori seruientes. }}\pstart Nam non simul, et ubique , et apud omnes omnia sunt facienda, quemadmodum oculi non semper faciunt suum officium, nonnunq dormiunt, et stomachus non semper estuat in alimentum, sic Christianus sciet  \pend
\section*{COMMENT. POMER.  }\pstart se attemperare omnibus horis, nunc enim consola- bitur, iam arguet proximum etc.   \pend
\phantomsection
\addcontentsline{toc}{subsection}{\textit{Spe gaudentes. }}
\subsection*{\textit{Spe gaudentes. }}\pstart Scilicet, in aduersis: Christianorum est, ut in sum- mis perturbationibus, meliora sperent, non diffidant, aut succumbant. IN ADFLICTIONE PA- TIENTES. Id est, perdurantes, ne excidatis a fide-  \pend
\phantomsection
\addcontentsline{toc}{subsection}{\textit{Precationi instantes. }}
\subsection*{\textit{Precationi instantes. }}\pstart Id est, non solum orate, uerum etiam instate, fit enim saepe ut non statim det Deus quod petitur: itaque non cel- sandum, sed semper pulsandum ostium, donec ape- riatur: nunquam autem desinit orare fides, nullam ha- bet horam aut diem in hoc praescriptam, siquidem sent per uidet sibi ab aduersario intentari malum.   \pend
\phantomsection
\addcontentsline{toc}{subsection}{\textit{Necessitatibus sanctorum communi. }}
\subsection*{\textit{Necessitatibus sanctorum communi. }}\pstart Id est, uestra impertientes: si quos uidetis Chri- stianos aegere, ilis quaecunque potestis communicate- Dicit uero haec de sanctis, et Christianis, non de ua- gis et hominibus nihili, qui tantum ex aliorum labo- ribus uiuere uolunt, ut ipsi non laborent, quales sunt multi hodie, quibus cum negatur a Christianis, cha- ritatis, ut dicunt, subsidium, statim obiiciunt, est hoc esse Christianos, nihil dare? quibus rursum oppona- tur: est hoc Christianos esse non laborare? et: Qui non laborat non manducet. Vide pri. Timoth. 5. et * Thessalo.3. tales esse excommunicandos.  \pend
\vspace{0.5cm}\noindent
\vspace{0.2cm}\rule{1cm}{0.2pt}\\ 
\hspace{0.2cm}\textit{mg}
\footnotesize Oratio non langueat. 
\normalsize| 
\hspace{0.2cm}\textit{mg}
\footnotesize Aegentibus, non ottosis dan- dum. 
\normalsize| 
\section*{IN EPIST. AD ROM.  }
\marginpar{[ p.140 ]}
\phantomsection
\addcontentsline{toc}{subsection}{\textit{Hospitalitatem sectantes. }}
\subsection*{\textit{Hospitalitatem sectantes. }}\pstart Id est, sitis paratt ad suscipiendum eos in hospitium quicunque ueniunt aduos Christiani. Olim enim cursi- tabant Apostoli, et alii Praedicatores de ciuitate in ciuitatem praedicantes, tales suscipiebantur a piis, et sanctis hospitio.  \pend
\phantomsection
\addcontentsline{toc}{subsection}{\textit{Bene loquamini de iis, qui uos etc. }}
\subsection*{\textit{Bene loquamini de iis, qui uos etc. }}\pstart Hoc est Christianorum, ut inquit Petrus, non so- lum de amicis, uerum et inimicis bene loqui, nan- si tantum de amicis bene loquamini, non pertinebit ad uos benedictio haereditatis, pri. Petri,  \pend
\phantomsection
\addcontentsline{toc}{subsection}{\textit{Gaudete cum gaudentibus. }}
\subsection*{\textit{Gaudete cum gaudentibus. }}\pstart Cum iis quibus laeta contingunt Haec ex similitu- dine membrorum corporis, ut habet Paulus in Co- rinth. intelliguntur, si unum membrum gaudet uel dolet, tum caetera omnia sic adficiuntur: ita commoda et incommoda fratrum ad nos quoque pertinere creda- mus: si igitur quisque personam alterius indueret, non esset opus toties, ac tanta de charitate precipere.  \pend
\phantomsection
\addcontentsline{toc}{subsection}{\textit{Flete cumflentibus. }}
\subsection*{\textit{Flete cumflentibus. }}\pstart Id est, quibus aduersa acciderunt.   \pend
\phantomsection
\addcontentsline{toc}{subsection}{\textit{Eodem animo alii in alios adfecti. }}
\subsection*{\textit{Eodem animo alii in alios adfecti. }}\pstart Id est, inuicem sitis concordes, etiamsi frater quii piam uos laedat: non commoueamini, tametsi tu ma- ior uel doctior sis, noli cum altero contendere, et ipsum afperneri.   \pend
\section*{COMMENT. POMER.  }
\phantomsection
\addcontentsline{toc}{subsection}{\textit{Non arroganter de uobis sentientes. }}
\subsection*{\textit{Non arroganter de uobis sentientes. }}\pstart Exponit quod iam dixit, non inflemini donis, unde fastidium fratrum nascitur. Haec uides subinde repeti, ut di- ligentius inculcentur : loquitur uero hic tantum de iis qui dona acceperunt, non de iis qui simulant: qui accepe- runt dona, primum utantur illis: secundo, non abutantur  \pend
\phantomsection
\addcontentsline{toc}{subsection}{\textit{Humilibus uos accommodantes. }}
\subsection*{\textit{Humilibus uos accommodantes. }}\pstart Id est, humilibus, et inferioribus uobis, uos at temperate, agnoscite eos fratres, quis non uideatis eos tantis donis insignitos, quantis uos, habetis enim unum patrem, et expectatis eandem haereditatem.   \pend
\phantomsection
\addcontentsline{toc}{subsection}{\textit{Ne sitis arrogantes apud uosmet. }}
\subsection*{\textit{Ne sitis arrogantes apud uosmet. }}\pstart Tertio hoc iam inculcat Paulus, nisi fecerit quis ea, de quibus praedixit in hoc Cap. non potest, non ar- roganter de se sentire: uult hanc radicem Apostolus pror- sus euelli, ex hac enim oia mala nascuntur, huius mali- tia fit, ut aliis, si quae ineis uidemus, sua dona inuidea mus, et nobis arrogemus supra dona nobis data, cum tamen ista amarulentia consequi nequeamus regnum Dei- Gaudendum est mihi, si uideam in alio dona, quae ego non habeo, nam scio ipsum ea habere, et in meam utilitatem  \pend
\phantomsection
\addcontentsline{toc}{subsection}{\textit{Neque cuiquam malum etc. }}
\subsection*{\textit{Neque cuiquam malum etc. }}\pstart Ex naturae adfectu, nihil possumus ferre in fratri- bus, quod ex diametro cum charitate pugnat :ficri non potest, ut in quotidiano conuictu non fiamus dis- cordes, igitur reconciliemur, non malum pro male  \pend
\vspace{0.5cm}\noindent
\vspace{0.2cm}\rule{1cm}{0.2pt}\\ 
\hspace{0.2cm}\textit{mg}
\footnotesize Arrogan- tia 
\normalsize| 
\section*{IN EPIST. AD ROM.  }
\marginpar{[ p.149 ]}\pstart reddamus, quod Ethnici faciunt . PROVIDE PARANTES. Id est, facientes.  \pend
\phantomsection
\addcontentsline{toc}{subsection}{\textit{In conspectu omnium hominum. }}
\subsection*{\textit{In conspectu omnium hominum. }}\pstart Ne male audiat Euangelium, ne dicant: sunt isti Euangelici? Tota uestra uita in oculis etiam hominum ta- lis sit, quam etiam aduersarii non possint calumniari, ut ora eorum obstruantur, qui male uobis loquuntur.  \pend
\phantomsection
\addcontentsline{toc}{subsection}{\textit{Si fieri potest, quantum in uobis est. }}
\subsection*{\textit{Si fieri potest, quantum in uobis est. }}\pstart Id est, etiam cum infidelibus, et inimicis quantum licebit, habete pacem. NON VOSMETIPSOS VLCISCENTES. Reconciliemini, non alite odia in corde uestro.  \pend
\phantomsection
\addcontentsline{toc}{subsection}{\textit{Quin potius date locumirae. }}
\subsection*{\textit{Quin potius date locumirae. }}\pstart Sinite, ut omnino obliuiscamini irae, sic alibi dicit: Sol non occidat super iracundiam uestram, id est, amprimum reconciliemini fratri, si offensi fueritis.   \pend
\phantomsection
\addcontentsline{toc}{subsection}{\textit{Mihi ultio ego rependam. }}
\subsection*{\textit{Mihi ultio ego rependam. }}\pstart Ad me solum, dicit Dominus in Mose, pertinet uin- dicta: si qui faciunt iniurias ego rependam, non uos. Vbi igitur quispiam ultionem sibi uindicarit, sciat se hoc sibi sumpsisse quod Dei est, huic nihil prospere succes- surum est: committenda oia sunt Deo, quem nihil latet-  \pend
\phantomsection
\addcontentsline{toc}{subsection}{\textit{Si igitur esurit inimicus tuus. }}
\subsection*{\textit{Si igitur esurit inimicus tuus. }}\pstart Nos oibus debemus benefacere, tantum abest, ut uin dicemus, non solum uerbo dilige, sed et re ipsa que sententia ex Prouerbiis desumpta est:  \pend
\vspace{0.5cm}\noindent
\vspace{0.2cm}\rule{1cm}{0.2pt}\\ 
\hspace{0.2cm}\textit{mg}
\footnotesize Vindicta.  
\normalsize| 
\section*{COMMENT. POMER.  }
\phantomsection
\addcontentsline{toc}{subsection}{\textit{Hoc enim si feceris, carbones ignis. }}
\subsection*{\textit{Hoc enim si feceris, carbones ignis. }}\pstart Haebraismus est, id est, efficies, ut ipse sibi irasca- tur, agnoscens tuam bonitatem, et suam maliciam, seipsum arguet dicens: ah quid feci? cur eum per secutus sum, qui mihi bene uult? Sic fit, ut multipli- ci dilectione, quam ille cogetur sentire, eum cruciabis aut ures: hoc est nobilissimum uictoriae genus, benefi- ciis uincere, de qua re quidem multa prodiderunt Phi- losophi, at nihil praestare potuerunt, scribere potue- runt, persuadere minime.  \pend
\phantomsection
\addcontentsline{toc}{subsection}{\textit{Ne uincaris a malo. }}
\subsection*{\textit{Ne uincaris a malo. }}\pstart Id est, ab iniuria, uel persecutore:uel, ne sis impo- tenti animo, sed forti, et constanti.  \pend
\phantomsection
\addcontentsline{toc}{subsection}{\textit{Imo uince bono malum. }}
\subsection*{\textit{Imo uince bono malum. }}\pstart Id est, iniuriam benefactis repende . Hactenus itaque de moribus Christianis dixit multa Paulus, quae ad nos et proximum pertinent, nunc alias paraeneses proponit.   \pend\pstart CAP. XIII.  \pend\pstart ARNALES, audita Buangelii praedicatione, A statim abutuntur Christiana libertate, nolentes cuiquam obedire: dicunt enim, Christianus est dominus omnium, ergo non ero cuiquam̃ subditus: quasi uero Chri- stianus sit, qui ita sentit, id quod nostro tempore ui- demus frequens esse, ideo Paulus nos admonet, ue  \pend
\vspace{0.5cm}\noindent
\vspace{0.2cm}\rule{1cm}{0.2pt}\\ 
\hspace{0.2cm}\textit{mg}
\footnotesize Beneficiis uin cendi mali. 
\normalsize| 
\hspace{0.2cm}\textit{mg}
\footnotesize Obediendum magistratui 
\normalsize| 
\section*{IN EPIST. AD ROM.  }
\marginpar{[ p.149 ]}\pstart dominis huius temporis, subditi simus, non tantum propter publicam pacem, sed etiam propter conscientiam  \pend
\phantomsection
\addcontentsline{toc}{subsection}{\textit{Omnis anima potestatibus. }}
\subsection*{\textit{Omnis anima potestatibus. }}\pstart Hebraismus est: de Christianis haec dicit, quid ad Apostolos non pertinent ii, qui foris sunt, ut in Corin. dicit. Nostri principes deberent multum gaudere, etiam secundum carnem de talt praedicatione: at quia excaecati sunt, aliud non possunt, quam blasphema- re, et impudentissime mentiri, dicentes, nos tantum praedicare, ne obediatur, sed breui mercedem suam accepturt sunt, nisi meliorem gloriam dent Verbo Dei, iam reuelato. Nos pro illis docemus: Papistae se dominos eorum faciebant, tamen propensiores semper in illos quam in nos adhuc hodie sunt. At haec sunt Dei iudicia.  \pend
\phantomsection
\addcontentsline{toc}{subsection}{\textit{Non enim est potestas, nisi a Deo. }}
\subsection*{\textit{Non enim est potestas, nisi a Deo. }}\pstart Est hoc labefactare, aut infringere, autoritatem Magistratuum, et potestatum mundi? Vix potest in tota Scriptura tam insigne Principum mundi Enco- mium inuenti, quam habet, uel hic unicus locus, neque tam stultus quispiam est, ut intelligat, hic Paulum de alia, quam mundi potestate loqui, a Deo ordinata, hoc est, de illa, quae portat gladium in defenssionem bo- norunt, cui propter damus honorem, tributum etc- An haec obscura sunt? non loquitur de potestate Sa- thanae, aut inferorum.   \pend
\phantomsection
\addcontentsline{toc}{subsection}{\textit{Quae uero sunt potestates? }}
\subsection*{\textit{Quae uero sunt potestates? }}
\vspace{0.5cm}\noindent
\vspace{0.2cm}\rule{1cm}{0.2pt}\\ 
\hspace{0.2cm}\textit{mg}
\footnotesize Magistratus autoritas non labefactatur Euangelio.  
\normalsize| 
\section*{COMMENT. POMER.  }\pstart Ecce de iis potestatibus loquitur, quae a Deo sunt in- stitutae et ordinatae.   \pend
\phantomsection
\addcontentsline{toc}{subsection}{\textit{Itaque quisquis resistit potestati, Dei ordinationi resistit. }}
\subsection*{\textit{Itaque quisquis resistit potestati, Dei ordinationi resistit. }}\pstart Quia Deus eam ordinauit. IVDICIVM. Con demnationem, sententiam mortis. Alians multa de pote- state gladii scripta sunt, non sunt aliae potestates quam Principes, Iudices, Magistratus, hi soli gladium gerunt: Papistae qui quoque sibi potestatem somniant, non ge- runt gladium, sed Plattas, et infulas, de quibus hic nihil dicitur, neque Paulus de aliis loqui hic potest, quam de iis qui suo tempore gladium habebant, in hoc, ut publicam pacem conseruarent. Est praeterea et haec potestas gladii, in Decalogo confirmata, in quo non solum praecipitur, ut honores parentes (nam id statim calumniarentur aduersarii, et quique audien- tes, hie intelligi etiam potestates) uerum, a diecta sunt et illa: Non occides, Non furaberis, Non moecha- beris, ex quibus consequitur, iudices esse oportere, quemnadmodum in legibus Mosi clarius uides. Occi- debatur enim qui alium occiderat: Moechus item da- bat poenas, sed sub iudicibus . Si gitur non volumus esse Fures. Moechos, Homicidas, necesse est esse Ma- gistr atusquos  timeant mali . Qui itaq se spirituales vo- cant, nihil ex hoc loco habent, quo suam fictam po- testatem defendant, nam non habent cam ordinatam  \pend
\vspace{0.5cm}\noindent
\vspace{0.2cm}\rule{1cm}{0.2pt}\\ 
\hspace{0.2cm}\textit{mg}
\footnotesize Gladii pote- stas.  
\normalsize| 
\hspace{0.2cm}\textit{mg}
\footnotesize Potestas eccle- siastica, ut uo- cant.  
\normalsize| 
\section*{EPIST. AD ROM.  }
\marginpar{[ p.150 ]}\pstart a Deo: sed aiunt: habemus potestatem ligandi et soluendi, Recte: sed haec non est potestas, uerum ministerium praedicandi Euangelii, quo annunciatur peccatorum remissio, quodque ; ad Christianos, non ad malos pertinet. Proinde, si quis dicat non esse malis Principibus obediendum, is longe fallitur, ut maxime regna ui inuaserint, et tyrannide occuparint, Dei uoluntate factum est, nihil ad nos, modo nihilcon- tra Deum praecipiant. Certum est, omnes fere Prin- cipes nostros, a principio primum occupationis ter- rarum (quemadmodum hodie Turca multa occupat) fraude, tyrannide, ambitione, cupiditate dominandi sua regna occupasse, tamen non sunt abiiciendi, sed eis obediendum, quia ordinati sunt a Deo, in hoc, ut nobis praesint et prosint, nam utcunque primi occu- patores non recte senserint, tamen Deus recte sensit, qui tales uoluit esse contra legum transgressores: hinc et Dii sunt in mundo, secundum Scripturam, in loco Dei, quod ad officium eorum attinet.  \pend
\phantomsection
\addcontentsline{toc}{subsection}{\textit{Nam principes non terrori sunt. }}
\subsection*{\textit{Nam principes non terrori sunt. }}\pstart Sed quia hodie abutuntur sua potestate, om- nia inuertunt, sunt enim terrori bene agentibus, et iis maextme qui pietatem Euangelii sectantur, impune agunt, quiduis sub ipsis impii, cogunt expilationibus, terroribus, et tormentis, ut abnegent boni Christum, et eius Euangelium: Rursus Papae scelera, quae mundus  \pend
\vspace{0.5cm}\noindent
\vspace{0.2cm}\rule{1cm}{0.2pt}\\ 
\hspace{0.2cm}\textit{mg}
\footnotesize Obediendum etiam malis. 
\normalsize| 
\section*{COMMENT. POMER.  }\pstart uix potest ferre, defendunt. Itaq quicunque sub iis he- die securus uiuere uolet, omnia faciat scelera et ma- la, quemadmodum ille dixit: Aude aliquid breuibus gyaris et carcere dignum. Si uis esse aliquid, pro bitas laudatur et alget.  \pend
\phantomsection
\addcontentsline{toc}{subsection}{\textit{Vis aut non timere potestatem. Scilicet, quae portat gladium, }}
\subsection*{\textit{Vis aut non timere potestatem. Scilicet, quae portat gladium, }}
\phantomsection
\addcontentsline{toc}{subsection}{\textit{Quod bonum est facito. }}
\subsection*{\textit{Quod bonum est facito. }}\pstart Non loquitur de bono coram Deo, sed ciuili bo- nitate, quam mundus iudicat, et Iudex uel Princeps commendat et laudat.  \pend
\phantomsection
\addcontentsline{toc}{subsection}{\textit{Dei enim minister est in bonum }}
\subsection*{\textit{Dei enim minister est in bonum }}\pstart Id est, is qui gerit gladium, a Deo ordinatus est, ut faciat in suo magistratu quod Deus uult-Si bene- feceris, dicet te bonum ciuem: sin minus, occidet- TIME. Scilicet gladii potestatem. FRVSTRA.  Id est, uindicabit gladio.  \pend
\phantomsection
\addcontentsline{toc}{subsection}{\textit{Vltor ad iram. }}
\subsection*{\textit{Vltor ad iram. }}\pstart Scilicet contra te- Totus hic locus contra eos est, qui Christiana libertate abutuntur, uultque Apostolus, ut omnes simus subditi potestatibus, et fratribums ser- uiamus, non enim Christiana libertas, est libertas car- nis: sed conscientiae, quae est per Christum, in hoc, ut nihil peccatorum uideamus, quod nos possit dan- nare, si ad alia traxeris hanc libertatem, non es Chri- stianus, non ideo te liberauit Christus, ut perturbes  \pend
\vspace{0.5cm}\noindent
\vspace{0.2cm}\rule{1cm}{0.2pt}\\ 
\hspace{0.2cm}\textit{mg}
\footnotesize Christiana li- bertas quae.  
\normalsize| 
\endnumbering
\end{pages}
\end{document}
        